\documentclass[12pt]{article}

\usepackage[utf8]{inputenc}
\usepackage{CJKutf8}
\usepackage{amsmath,amssymb}
\usepackage{amsthm}
\usepackage{geometry}
\geometry{
  a4paper,
  top=25mm,
  bottom=25mm,
  left=20mm,
  right=20mm
}
\usepackage{xcolor}
\usepackage{url}
\usepackage[unicode,colorlinks=true,linkcolor=blue,urlcolor=blue]{hyperref}
\usepackage{xurl}
\Urlmuskip=0mu plus 2mu
\sloppy
\usepackage{tocloft}
\usepackage{array}
\usepackage{longtable}

% 目录中 section 条目使用点线填充到页码
\renewcommand{\cftsecleader}{\cftdotfill{\cftdotsep}}
\renewcommand{\refname}{参考文献}

% 基本定理环境（工作笔记：形式系统风格；参考 tex21）
\newtheorem{definition}{定义}[section]
\newtheorem{axiom}{公理}[section]
\newtheorem{assumption}{假设}[section]
\newtheorem{theorem}{定理}[section]
\newtheorem{proposition}{命题}[section]
\newtheorem{corollary}{推论}[section]
\newtheorem{lemma}{引理}[section]
\newtheorem{remark}{备注}[section]
\newtheorem{certificate}{证书}[section]

% 记号（仅用于本笔记自洽引用，避免依赖主稿宏定义）
\newcommand{\code}[1]{\path{#1}}
\newcommand{\GFramework}{\textsf{G-Framework}}

% 核心符号（本章自洽）
\newcommand{\GZLS}{\ensuremath{\mathsf{GZLS}}}
\newcommand{\GZLOC}{\ensuremath{\mathsf{GZLOC}}}
\newcommand{\GZLCurv}{\ensuremath{\mathsf{GZLCurv}}}
\newcommand{\GZTLC}{\ensuremath{\mathsf{GZTLC}}}
\newcommand{\JacGap}{\ensuremath{\mathsf{JacGap}}}
\newcommand{\LBOPB}{\ensuremath{\mathrm{LBOPB}}}
\newcommand{\PTOPB}{\ensuremath{\mathrm{PTOPB}}}
\newcommand{\ETOPB}{\ensuremath{\mathrm{ETOPB}}}
\newcommand{\OPB}{\ensuremath{\mathrm{OPB}}}
\newcommand{\RTSC}{\textsf{RTSC}}
\newcommand{\Conn}{\ensuremath{\mathsf{Conn}}}
\newcommand{\Curv}{\ensuremath{\mathsf{Curv}}}
\newcommand{\EDOM}{\ensuremath{\mathcal{M}^{\mathrm{EDOM}}}}
\newcommand{\PGOM}{\ensuremath{\mathcal{M}^{\mathrm{PGOM}}}}

% 避免少量不可控的换行溢出（尤其是混排 CJK 与等宽字体时）

\hfuzz=700pt
\hbadness=10000

\begin{document}
\begin{CJK*}{UTF8}{gkai}

\title{Applied Math Vol II：从 \LBOPB{} 到 \PTOPB{} 的可迁移模板\\（工作笔记：形式系统版）}
\author{Gao Zheng（高政）}
\date{2026-01-16}
\maketitle

\section*{前言}
\addcontentsline{toc}{section}{前言}

本文的目标不是单独整理 \LBOPB{}、\PTOPB{}、RTSC 虚拟实验或逆向构造，而是在
\GFramework{} 的应用数学链条中，把总算子主纤维丛、算子幺半群、离散 GRL-SAC、
环境同伦复用、路线插件化与控制一等化组织成一个可迁移模板。若这些对象分散出现，
它们容易只表现为工程组件；本文将其压合为从低层算子基线到物理目标逆向设计的
统一形式系统。

本文的第一层立场是可迁移模板的总算子化。\LBOPB{} 到 \PTOPB{} 的迁移并不只是
更换问题名称，而是把候选路线、控制节点、环境适配器、局部优化器与路径积分稳定性
界统一放入总算子主纤维丛中。由此，不同应用域可以共享同一套算子分解、纤维扩张
与同伦复用接口。

本文的第二层立场是 RTSC 逆向构造的分层化。室温超导路线在本文中被处理为虚拟
实验与逆向构造问题，而不是无条件物理结论。可行性需要按层划分：哪些节点可被
结构控制，哪些节点需要插件路线，哪些节点只能通过环境适配器间接调节，哪些节点
需要承认为暂不可控。总雅可比间隙给出这种可控性的审计判据。

本文的第三层立场是控制一等化与路线坍缩防护。逆向设计若只追求单一路线，容易把
复杂搜索压缩为路线坍缩幻觉；本文通过子丛、子幺半群、联合共设计与导航证书，
把路线选择提升为可比较、可复验、可迁移的控制对象。控制不再只是工程参数，而是
与路线结构同等重要的一等对象。

本文的第四层立场是虚拟实验的接口化。虚拟实验不是替代真实实验的结论机器，而是
用于组织候选路线、环境适配器、反馈指标和失败证书的工程接口。它的价值在于把
不可直接穷举的实验空间转化为可插件化、可重放、可审计的路径集合，为后续真实
验证保留边界。

本文的第五层立场是环境同伦复用。不同问题域中的环境对象虽然物理语义不同，但其
在算法层的作用可以通过适配器、路径积分稳定性界和同伦复用接口统一表达。这样，
从 \LBOPB{} 到 \PTOPB{} 的迁移不再只是经验类比，而是一个可定义的环境对象变换。

本文的第六层立场是路线插件化。路线不是写死的单条流程，而是可以作为子丛、子幺
半群或插件族被加载、比较、复验和替换。插件化允许系统在失败时换路，在证书不足
时降级，在新信息出现时扩展，从而避免把逆向设计锁死在早期猜测中。

本文的第七层立场是总 JacGap 的可控性判据。RTSC 或其他复杂逆向设计任务是否可控，
不能只看某个局部指标，而要看总雅可比间隙能否被可作用的格点、结构控制方向和
证书链共同压低。总 JacGap 将路线、控制、环境和结构缺口合成为一个可审计目标。

因此，本文在整体 \GFramework{} 中承担的是“应用数学可迁移逆向设计模板”的作用。
它向前连接 PTOPB/ETOPB 二分与 law-space 算子体系，向中间连接虚拟环境、路线插件、
GRL 路径积分与总 JacGap，向后为开放系统治理、同伦命中和 PDE 残差范式提供可迁移
工程模板。概括地说，本文把“如何把一个应用问题迁移到另一个应用问题”推进为
总算子、纤维丛、环境同伦与可控性证书共同约束的形式系统。

更具体地说，本文把逆向设计从“寻找一个答案”改写为“组织一个可迁移的路线族”。
路线族可以被插件化、被虚拟环境测试、被总 JacGap 审计，也可以在失败时替换或
降级。这样的逆向设计不再依赖单次灵感，而依赖可维护的路线生态。

从方法论上看，本文的关键是把控制、路线和环境放到同一个主纤维丛模板里。控制
不再只是后端调参，路线不再只是前端选择，环境也不再只是外部条件；三者共同决定
问题是否可迁移、可测试、可审计和可继续扩展。

从整体谱系看，本文为 \GFramework{} 提供应用数学方向的“工程母版”。后续任何需要
从抽象形式系统落到具体复杂设计任务的稿件，都可以借用这里的总算子、插件路线、
环境同伦和可控性证书接口。它不是 RTSC 单一案例，而是可迁移逆向设计的模板。

\begin{center}
\begingroup
\setlength{\fboxsep}{10pt}
\setlength{\fboxrule}{0.8pt}
\fcolorbox{black}{gray!6}{\begin{minipage}{0.94\linewidth}
\textbf{开篇重点：以总雅可比间隙（Total JacGap）作为 \PTOPB$_{\mathrm{RTSC}}$ 的可控性判据。}
\[
\boxed{
\begin{aligned}
&\mathrm{RTSC}\ \text{在给定算子集合下可控}
\ \Longleftrightarrow\
\inf_{w\in\mathcal{W}_{\mathrm{adm}}}\ \mathcal{S}(w)\le \varepsilon
\\
&\Longleftrightarrow\
\text{每个子节点的 }g_i\text{均存在可作用的格点/结构控制方向，}
\\
&\hspace{4em}\text{使 }\mathbf{g}\to 0.
\end{aligned}
}
\]
形式化定义与证明见第 \ref{sec:ptopb-rtsc-total-jacgap} 节的 Proposition~\ref{prop:rtsc-jacgap-controllability-main}。
\end{minipage}}
\endgroup
\end{center}

\setcounter{tocdepth}{3}
\setcounter{secnumdepth}{3}
\tableofcontents
\clearpage

\section*{形式系统写作约定（公理降级、证书与谓词因果链）}
\addcontentsline{toc}{section}{形式系统写作约定（公理降级、证书与谓词因果链）}
本文采用“形式逻辑 + 证书”的工作笔记体例，并统一使用\textbf{基于数学符号推演逻辑的谓词因果链}来组织证明要点。
为避免口头叙述与工程接口脱钩，正文默认遵循如下约定：
\begin{itemize}
  \item \textbf{公理降级}：不单列“公理”环境；凡需作为基础假设固定者，均以“公理（降级为定理）”写成\textbf{定理}，
  并配套可复算的\textbf{证书}与\textbf{证明}。若断言包含按步/按层/按回合/按长度的全称迭代结构，默认给出\textbf{数学归纳法证书}。
  \item \textbf{定理/引理/命题/推论}：每条编号断言均同时给出（i）证书（可审计见证/日志/不变量/基步-归纳步等），
  （ii）证明（以谓词链闭合方式把推理写成可复核的符号推演步骤）。
  \item \textbf{备注}：只用于动机、解释、示例与工程语境对齐，不承担证明义务。
  \item \textbf{证书调用}：对外复用时建议成对引用“断言 + 证书”（如 \texttt{命题~\ref{prop:rtsc-jacgap-controllability-main}} 与
    \texttt{证书~\ref{cert:rtsc-jacgap-controllability-main}}），
  以确保他处调用时具备可追溯的接口；证书标签命名约定与索引见附录B。
\end{itemize}
\clearpage

%%%%%%%%%%%%%%%%%%%%%%%%%%%%%%%%%%%%%%%%%%%%%%%%%%%%%%%%%%%%
\section{可迁移模板：总算子主纤维丛、算子幺半群与离散 GRL-SAC}
%%%%%%%%%%%%%%%%%%%%%%%%%%%%%%%%%%%%%%%%%%%%%%%%%%%%%%%%%%%%

\subsection{定义：\LBOPB{} 的结构骨架（形式来源）}

\begin{definition}[\LBOPB{}（生命总算子主纤维丛）的结构骨架]
\label{def:lbopb-skeleton}
在 \code{docs/AppMathVol2/Pub_GFramework_AppMath_Vol2.tex} 中，\LBOPB{} 被定义为元组
\[
  \LBOPB
  :=
  \bigl(
    M^{\mathrm{life}},
    P^{\mathrm{life}},
    G^{\mathrm{life}},
    \Conn^{\mathrm{life}};
    \{\mathcal{M}^{D}\}_{D \in \mathcal{D}_{\mathrm{life}}}
  \bigr),
\]
其中 $M^{\mathrm{life}}\subset \GZLS$ 为生命态底空间，$P^{\mathrm{life}}\to M^{\mathrm{life}}$ 为主纤维丛，$G^{\mathrm{life}}$ 为结构群，
  $\Conn^{\mathrm{life}}$ 为法联络，其曲率 $\Curv^{\mathrm{life}}$ 记录非交换/路径依赖信息；$\{\mathcal{M}^{D}\}$ 为分学科域算子幺半群族，要求与
  $G^{\mathrm{life}}$ 的主作用及 $\Conn^{\mathrm{life}}$ 的可允许变化方向兼容。
\end{definition}

\begin{certificate}[文本证书：\LBOPB{} 的结构定义位置]
\label{cert:lbopb-definition-source}
在 \code{docs/AppMathVol2/Pub_GFramework_AppMath_Vol2.tex} 中，\LBOPB{} 的结构性定义以“Life Bio-Operator Principal Bundle (LBOPB)
  ”给出（参见该文件中 \texttt{definition} 环境附近的定义段落），并明确列出
$(M^{\mathrm{life}},P^{\mathrm{life}},G^{\mathrm{life}},\Conn^{\mathrm{life}};\{\mathcal{M}^{D}\})$ 的元组口径与“闭系统约束、兼容性、
  曲率记录 Jacobian-gap”的解释。
\end{certificate}

\begin{remark}[工程可迁移性：抽象骨架与域语义解耦]
Definition~\ref{def:lbopb-skeleton} 的关键点是：其结构骨架不依赖“生命科学”特有对象才能成立；其依赖的是
（i）法空间子域底空间 $M\subset\GZLS$，（ii）主纤维丛与结构群处理表述冗余，（iii）联络/曲率刻画允许变化与路径依赖，
（iv）可组合的域算子幺半群刻画可执行干预/操作序列。
因此该骨架可以作为“可替换域名的模板”，用于物理/工程逆向设计的形式化迁移。
\end{remark}

\subsection{定理（公理降级）：\LBOPB{} 结构的域替换模板性}

\begin{theorem}[模板性（公理降级）：由 \LBOPB{} 诱导 \PTOPB/\ETOPB{} 的结构性构造]
\label{thm:lbopb-template}
设给定一个域替换映射（语义替换）$\sigma:\mathcal{D}_{\mathrm{life}}\to \mathcal{D}_{\mathrm{phys}}$，
以及与之相容的结构替换
\[
(M^{\mathrm{life}},P^{\mathrm{life}},G^{\mathrm{life}},\Conn^{\mathrm{life}},\{\mathcal{M}^{D}\}_{D\in\mathcal{D}
  _{\mathrm{life}}})
\ \longmapsto\
(M^{\mathrm{phys}},P^{\mathrm{phys}},G^{\mathrm{phys}},\Conn^{\mathrm{phys}},\{\mathcal{M}^{D}\}_{D\in\mathcal{D}
  _{\mathrm{phys}}}),
\]
其中 $M^{\mathrm{phys}}\subset\GZLS$、$P^{\mathrm{phys}}\to M^{\mathrm{phys}}$、$G^{\mathrm{phys}}$、$\Conn^{\mathrm{phys}}$
  与 $\{\mathcal{M}^{D}\}$ 满足与 Definition~\ref{def:lbopb-skeleton} 同型的“闭系统/兼容性/曲率记录非交换”约束。
则可定义物理总算子主纤维丛
\[
  \PTOPB
  :=
  \bigl(
    M^{\mathrm{phys}},
    P^{\mathrm{phys}},
    G^{\mathrm{phys}},
    \Conn^{\mathrm{phys}};
    \{\mathcal{M}^{D}\}_{D \in \mathcal{D}_{\mathrm{phys}}}
  \bigr),
\]
并可将其作为物理/工程逆向设计的连续层形式骨架（\ETOPB{} 作为面向工程对象的同型记号亦可按同一方式给出）。
\end{theorem}

\begin{certificate}[数学归纳法证书：按“结构元组替换步数”验证同型性]
\label{cert:lbopb-template-induction}
将 \LBOPB{} 的结构元组拆为五块并按顺序替换：
\[
\begin{aligned}
T_1:&\ M^{\mathrm{life}}\mapsto M^{\mathrm{phys}},\\
T_2:&\ P^{\mathrm{life}}\mapsto P^{\mathrm{phys}},\\
T_3:&\ G^{\mathrm{life}}\mapsto G^{\mathrm{phys}},\\
T_4:&\ \Conn^{\mathrm{life}}\mapsto \Conn^{\mathrm{phys}},\\
T_5:&\ \{\mathcal{M}^{D}\}_{D\in\mathcal{D}_{\mathrm{life}}}\mapsto
\{\mathcal{M}^{D}\}_{D\in\mathcal{D}_{\mathrm{phys}}}.
\end{aligned}
\]
定义归纳谓词 $P(k)$（$k\in\{0,1,2,3,4,5\}$）为：
\[
P(k):\ \text{前 }k\text{ 步替换后所得对象与 Definition~\ref{def:lbopb-skeleton} 的结构约束同型成立。}
\]
证书化要求为：
\[
P(0)\ \text{成立},\qquad
(\forall k\in\{0,1,2,3,4\})\ \bigl(P(k)\Rightarrow P(k+1)\bigr).
\]
\end{certificate}

\begin{proof}[证明（数学归纳法证书；谓词因果链）]
由 Certificate~\ref{cert:lbopb-template-induction}：
\begin{itemize}
  \item 基步：$P(0)$ 成立（未替换时即为 \LBOPB{} 的结构定义）；
  \item 归纳步：对 $k\in\{0,1,2,3,4\}$，若 $P(k)$ 成立，则第 $k{+}1$ 步替换保持“闭系统/兼容性/曲率记录非交换”的同型约束，从而 $P(k{+}1)$ 成立。
\end{itemize}
由数学归纳法得 $P(5)$ 成立，即替换后所得结构与 \LBOPB{} 的形式骨架同型。将替换后的五元组按 Definition~\ref{def:lbopb-skeleton} 的元组口径组装，即得 \PTOPB{} 的定义与结论。证毕。

\end{proof}

\subsection{定义：\PTOPB/\ETOPB{} 的工程语义口径}

\begin{definition}[\PTOPB{}（物理总算子主纤维丛）]
\label{def:ptopb}
称元组
\[
  \PTOPB
  :=
  \bigl(
    M^{\mathrm{phys}},
    P^{\mathrm{phys}},
    G^{\mathrm{phys}},
    \Conn^{\mathrm{phys}};
    \{\mathcal{M}^{D}\}_{D \in \mathcal{D}_{\mathrm{phys}}}
  \bigr)
\]
为物理总算子主纤维丛，若满足：
\begin{itemize}
  \item $M^{\mathrm{phys}}\subset\GZLS$ 表示“物理系统/工程构件”的法级状态空间（包含结构/材料/边界条件/制造与法规等闭系统约束）；
  \item $P^{\mathrm{phys}}\to M^{\mathrm{phys}}$ 为主纤维丛，其纤维承载内部自由度与表述冗余（如坐标系/网格/参数化冗余/多物理场离散表示选择）；
  \item $G^{\mathrm{phys}}$ 表示不改变物理对象本体但改变表示的等价变换群；
  \item $\Conn^{\mathrm{phys}}$ 给出可允许的“法级变化方向”，其曲率 $\Curv^{\mathrm{phys}}$ 记录不可交换与路径依赖；
  \item $\{\mathcal{M}^{D}\}_{D \in \mathcal{D}_{\mathrm{phys}}}$ 为多域算子幺半群族，用于生成可组合的工程操作序列，并与 $G^{\mathrm{phys}}$ 作用与
    $\Conn^{\mathrm{phys}}$ 兼容。
\end{itemize}
\end{definition}

\begin{remark}[域集合示例]
可取
\[
\mathcal{D}_{\mathrm{phys}}
=
\{\text{几何/拓扑},\ \text{材料},\ \text{制造},\ \text{控制},\ \text{载荷/环境},\ \text{可靠性/安全},\ \text{成本/供应链}\},
\]
并按工程对象族的实际需求调整域集合与域内算子字典。
\end{remark}

\subsection{定理（公理降级）：联络/曲率的路径依赖语义}

\begin{theorem}[路径依赖语义（公理降级）：曲率记录允许变化的非交换性]
\label{thm:connection-curvature-semantics}
在 \PTOPB{} 的结构中，若 $\Conn^{\mathrm{phys}}$ 为主丛 $P^{\mathrm{phys}}\to M^{\mathrm{phys}}$ 上的法联络，并以其曲率 $\Curv^{\mathrm{phys}}
  $ 刻画允许变化的非交换性，
则“先执行变化 $a$ 再执行变化 $b$”与“先执行变化 $b$ 再执行变化 $a$”的差异可由 $\Curv^{\mathrm{phys}}$ 的非平凡性刻画；从而 $\Curv^{\mathrm{phys}}\neq 0$
  表示工程逆向设计中的顺序依赖/路径依赖结构不可忽略。
\end{theorem}

\begin{certificate}[文本证书：联络/曲率与非交换性的语义来源]
\label{cert:connection-curvature-source}
\code{docs/AppMathVol2/Pub_GFramework_AppMath_Vol2.tex} 在 \LBOPB{} 的结构性定义段明确指出：$\Conn^{\mathrm{life}}$ 指定允许的无穷小变化方向，其曲率
  $\Curv^{\mathrm{life}}$ “measures the failure of admissible changes to commute”（记录 Jacobian-gap 数据），并用于刻画路径依赖。
  \cite{GaoZhengAppVol2}
\end{certificate}

\begin{proof}[证明（谓词因果链）]
由 Certificate~\ref{cert:connection-curvature-source}，对联络/曲率语义取谓词链：
\[
\begin{aligned}
Q_1:&\ \Curv\ \text{刻画允许变化的不交换性},\\
Q_2:&\ \Curv\neq 0\Rightarrow \exists(a,b)\ \text{使得先后顺序导致结果差异}.
\end{aligned}
\]
由 $Q_1$ 的语义定义可推出 $Q_2$；因此 $\Curv^{\mathrm{phys}}\neq 0$ 即对应非交换性与路径依赖不可忽略，得到结论。证毕。
\end{proof}

\subsection{定义：刚性干预层与工程设计算子幺半群}

\begin{definition}[工程设计算子幺半群 \EDOM{}（刚性干预层口径）]
\label{def:edom}
设 \EDOM{} 为工程离散设计原语的幺半群，其元素为可执行的设计原语（例如加孔、改厚度、换材料牌号、调整热处理工艺参数、变更控制律模块、器件选型替换等），乘法为可执行序列的连接（concatenation）。
一个完整设计方案以算子词表示：
\[
w = a_1 a_2 \cdots a_T,\qquad a_t\in \EDOM.
\]
\end{definition}

\begin{definition}[域幺半群到 \EDOM{} 的同态投影]
\label{def:domain-hom-to-edom}
对每个物理/工程域 $D\in\mathcal{D}_{\mathrm{phys}}$，设其域算子幺半群为 $\mathcal{M}^{D}$。
称映射
\[
\Phi^{D}:\mathcal{M}^{D}\to \EDOM
\]
为域算子到统一执行层的同态投影，若 $\Phi^{D}$ 为幺半群同态并满足：域内可执行操作在统一执行层中可用 \EDOM{} 的算子词表示，从而跨域组合可在同一“可执行设计程序语言”中完成一致性检查。
\end{definition}

\subsection{定理（公理降级）：刚性干预层使逆向设计可离散可编程}

\begin{theorem}[离散可编程性（公理降级）：干预方案可表示为幺半群词并进入离散 GRL-SAC]
\label{thm:rigid-layer-discrete}
若存在从连续层（联络/曲率驱动的允许变化）到刚性干预层（\EDOM{}）的投影，使得每个干预方案对应于 \EDOM{} 中的有限词 $w$，
且每个词的语义在满足闭系统约束与兼容性约束的条件下可由仿真器/模型复核，
则逆向设计问题可等价地表述为“在 \EDOM{} 的词空间上定义代价/回报泛函，并用离散 GRL-SAC 对算子序列进行搜索与优化”。
\end{theorem}

\begin{certificate}[数学归纳法证书：按词长度 $T$ 归纳的“可执行计划=词”口径]
\label{cert:rigid-layer-word-induction}
定义谓词 $R(T)$ 为：“任意长度不超过 $T$ 的干预计划可表示为 \EDOM{} 中的词，并可由闭系统约束下的可复核语义解释。”
证书化要求为：
\[
R(1)\ \text{成立},\qquad
(\forall T\in\mathbb{N}_{\ge 1})\ \bigl(R(T)\Rightarrow R(T+1)\bigr),
\]
其中递推步使用幺半群的连接封闭性（将长度 $T$ 的词与一个原语 $a_{T+1}$ 连接得到长度 $T{+}1$ 的词）。
\end{certificate}

\begin{proof}[证明（数学归纳法证书；谓词因果链）]
由 Certificate~\ref{cert:rigid-layer-word-induction}：
\begin{itemize}
  \item 基步：$R(1)$ 成立（单步干预即为一个设计原语 $a_1\in\EDOM$）；
  \item 归纳步：若 $R(T)$ 成立，则任意长度 $T{+}1$ 的计划可写为 $w a_{T+1}$（其中 $w$ 长度为 $T$），由 \EDOM{} 的封闭性与语义复核假设得 $R(T{+}1)$ 成立。
\end{itemize}
由归纳法得对任意有限长度计划均可表示为词；在此基础上，对词空间定义代价/回报并调用离散 GRL-SAC 优化即得到定理结论。证毕。
\end{proof}

\subsection{定义：离散 GRL-SAC 逆向搜索问题（工程版）}

\begin{definition}[离散 GRL-SAC（算子幺半群动作空间口径）]
\label{def:grl-sac-engineering}
设 $\mathcal{S}$ 为离散/离散化状态空间，$\mathcal{A}$ 为有限动作空间。工程逆向设计中可取：
\[
\mathcal{A}\subseteq \EDOM\ \ \text{或}\ \ \mathcal{A}\subseteq \mathcal{W}_{\mathrm{act}}\subseteq \mathcal{W}(\EDOM),
\]
其中 $\mathcal{W}(\EDOM)$ 为 \EDOM{} 生成的词空间，$\mathcal{W}_{\mathrm{act}}$ 由长度/相邻约束/域使用约束等裁剪得到。
给定转移核 $P(s'\mid s,a)$（由仿真器/模型近似）与回报函数 $r(s,a)$（由语义指标与路径积分泛函诱导），定义熵正则化目标：
\[
  J(\pi_{\theta})
  :=
  \mathbb{E}\Biggl[
    \sum_{t=0}^{\infty}
    \gamma^{t} \bigl(
      r(s_{t},a_{t}) + \alpha \mathcal{H}(\pi_{\theta}(\cdot \mid s_{t}))
    \bigr)
  \Biggr],
\]
其中 $\gamma\in(0,1)$、$\alpha>0$，$\mathcal{H}$ 为 Shannon 熵。
\end{definition}

\begin{certificate}[文本证书：离散 GRL-SAC 的动作空间与目标形式]
\label{cert:grl-sac-source}
\code{docs/AppMathVol2/Pub_GFramework_AppMath_Vol2.tex} 在 “Discrete GRL-SAC on operator monoids” 的定义段给出：状态空间 $\mathcal{S}
  $、动作空间 $\mathcal{A}$（基本算子或算子词）、转移核 $P(s'\mid s,a)$、回报 $r(s,a)$，以及熵正则化目标 $J(\pi_\theta)$ 的形式。\cite{GaoZhengAppVol2}
\end{certificate}

\begin{proof}[证明（谓词因果链）]
由 Certificate~\ref{cert:grl-sac-source}，离散 GRL-SAC 的抽象构件为谓词链：
\[
S_1:\ (\mathcal{S},\mathcal{A},P,r)\ \text{给定},\qquad
S_2:\ J(\pi_\theta)\ \text{按熵正则形式定义},\qquad
S_3:\ \text{优化 }J\ \text{得到策略搜索问题}.
\]
由 $S_1\Rightarrow S_2\Rightarrow S_3$ 可得 Definition~\ref{def:grl-sac-engineering} 的形式与 Vol II 抽象同型，因此可作为工程逆向设计的统一优化口径。证毕。

\end{proof}

\subsection{命题：工程逆向设计与 Vol II 模板的结构同型}

\begin{proposition}[结构同型：从生命域到物理/工程域的同型迁移]
\label{prop:isomorphic-migration}
若以 $\sigma$ 将生命域集合 $\mathcal{D}_{\mathrm{life}}$ 替换为物理域集合 $\mathcal{D}_{\mathrm{phys}}$，
并以 \PTOPB{} 替换 \LBOPB{}（Theorem~\ref{thm:lbopb-template}），同时以 \EDOM{} 替换 \PGOM{} 作为刚性干预层（Theorem~\ref{thm:rigid-layer-discrete}），
则 Vol II 中“总算子主纤维丛 + 多域算子幺半群 + 刚性干预层 + 离散 GRL-SAC”的结构性模板可在物理/工程逆向设计中同型复用。
\end{proposition}

\begin{certificate}[证书：同型迁移的结构对应表]
\label{cert:isomorphic-migration}
同型迁移采用如下对应：
\[
\begin{aligned}
&M^{\mathrm{life}}\mapsto M^{\mathrm{phys}},\quad
P^{\mathrm{life}}\mapsto P^{\mathrm{phys}},\quad
G^{\mathrm{life}}\mapsto G^{\mathrm{phys}},\quad
\Conn^{\mathrm{life}}\mapsto \Conn^{\mathrm{phys}},\\
&\{\mathcal{M}^{D}\}_{D\in\mathcal{D}_{\mathrm{life}}}\mapsto \{\mathcal{M}^{D}\}_{D\in\mathcal{D}_{\mathrm{phys}}},
  \quad
\PGOM\mapsto \EDOM,\quad
(\mathcal{S},\mathcal{A},P,r,J)\ \text{保持为同型抽象构件}.
\end{aligned}
\]
\end{certificate}

\begin{proof}[证明（谓词因果链）]
由 Certificate~\ref{cert:isomorphic-migration} 给出的结构对应，结合 Theorem~\ref{thm:lbopb-template} 与 Theorem~\ref{thm:rigid-layer-discrete}，
可得连续层骨架（主丛/结构群/联络/曲率）与离散层骨架（动作=算子或算子词、GRL-SAC 目标）均按同型替换保持定义结构；
因此整体模板可在物理/工程域中复用，得到结论。证毕。
\end{proof}

\subsection{推论：内环 RL + 外环规则/证书配置的可控性收益}

\begin{corollary}[两层结构推论：外环裁剪与版本化提升可控性]
\label{cor:two-loop-control}
在 Definition~\ref{def:grl-sac-engineering} 的离散 GRL-SAC 结构下，若将安全、法规、制造、材料适用性、供应链等硬约束以外环规则/证书形式实现：
\begin{enumerate}
  \item 对动作空间 $\mathcal{A}$ 作裁剪（禁止某些算子或相邻组合）；
  \item 对词空间 $\mathcal{W}_{\mathrm{act}}$ 施加序列约束与版本化指标配置；
\end{enumerate}
则可在不改变内环推断算法结构的前提下显著降低搜索空间并提高可解释性与可部署性。
\end{corollary}

\begin{certificate}[证书：外环约束作为动作空间与词空间的可计算裁剪]
\label{cert:outer-loop-pruning}
外环约束以集合与谓词给出：
\[
\mathcal{A}_{\mathrm{adm}}:=\{a\in\mathcal{A}\mid \mathrm{Adm}(a)\},\qquad
\mathcal{W}_{\mathrm{adm}}:=\{w\in\mathcal{W}_{\mathrm{act}}\mid \mathrm{SeqAdm}(w)\},
\]
其中 $\mathrm{Adm}$ 与 $\mathrm{SeqAdm}$ 为可计算判定谓词（由规则库/证书库实现）。
\end{certificate}

\begin{proof}[证明（谓词因果链）]
由 Certificate~\ref{cert:outer-loop-pruning}，外环约束将动作/词空间替换为其可计算可接受子集：
\[
\mathcal{A}\mapsto \mathcal{A}_{\mathrm{adm}},\qquad
\mathcal{W}_{\mathrm{act}}\mapsto \mathcal{W}_{\mathrm{adm}}.
\]
该替换不改变 $P(s'\mid s,a)$、$r(s,a)$ 与 $J(\pi_\theta)$ 的形式，只改变可选动作集合，从而不改变内环推断算法的结构定义；同时由于集合裁剪，搜索空间规模下降且约束语义可被版本化记录。故得到结论。证毕。
\end{proof}

\subsection{命题：可行性判据与三项主要难点}

\begin{proposition}[可行性判据：定义可行不等于实现可行]
\label{prop:feasibility-vs-implementation}
在 \PTOPB{} 与 \EDOM{} 已按 Definition~\ref{def:ptopb} 与 Definition~\ref{def:edom} 给定时，框架在数学定义层面成立。
在工程实现层面，其可行性主要取决于以下三项判据是否可被证书化满足：
\begin{enumerate}
  \item \textbf{算子词表粒度：}动作空间规模与可控性之间的折中可被外环规则稳定约束；
  \item \textbf{闭系统边界：}哪些外部因素应内化进 $M^{\mathrm{phys}}$（闭系统）与哪些应作为扰动进入 $P(s'\mid s,a)$ 的划分具有可复核口径；
  \item \textbf{仿真一致性区间：}当 $P$ 由仿真器近似时，外环证书能给出“仿真--现实”一致性的保守区间与失配惩罚，使内环优化不至于收敛到仿真投机解。
\end{enumerate}
\end{proposition}

\begin{certificate}[证书：三项判据的可检验形式]
\label{cert:feasibility-criteria}
三项判据对应的可检验形式为：
\begin{enumerate}
  \item 词表粒度：存在上界 $|\mathcal{A}_{\mathrm{adm}}|\le N_{\max}$ 与最大词长 $T\le T_{\max}$ 的规则证书；
  \item 闭系统边界：存在对外部变量集合 $\mathcal{E}_{\mathrm{in}}$（内化）与 $\mathcal{E}_{\mathrm{out}}$（扰动）的一致分类证书；
  \item 仿真一致性：存在误差界证书 $\mathrm{Err}(P,P_{\mathrm{real}})\le \varepsilon$ 或相应的风险惩罚下界形式。
\end{enumerate}
\end{certificate}

\begin{proof}[证明（谓词因果链）]
由 Certificate~\ref{cert:feasibility-criteria}，三项判据均可被表述为可检验约束（集合上界、分类一致性、误差界/惩罚界）。
若三类证书均成立，则外环可将动作空间与仿真一致性约束纳入系统定义，使 Definition~\ref{def:grl-sac-engineering} 的 $P$ 与 $r$ 具备可复核边界，从而实现层面的不确定性被纳入可控范围；
  反之若任一证书缺失，则实现层面可能出现动作空间不可控、闭系统边界漂移或仿真投机解等风险。证毕。
\end{proof}

\subsection{备注：最小切入路径（可部署优先）}

\begin{remark}[最小可行路径]
为避免“全物理宇宙”式的大规模不可控建模，可采用如下最小切入路径：
\begin{enumerate}
  \item 固定工程对象族与明确目标（例如热管理部件、功率变换器、结构件、天线、微电网控制器等），从而固定 $M^{\mathrm{phys}}$ 的闭系统边界口径；
  \item 先选取少量域幺半群（例如几何/材料/制造，或几何/控制/成本），构造小规模 \EDOM{} 与 $\mathcal{W}_{\mathrm{act}}$；
  \item 先在外环固化制造可行性、安全下界与法规约束等规则/证书，以裁剪动作空间并为 $r(s,a)$ 提供可解释惩罚项；
  \item 再将离散 GRL-SAC 接入仿真器完成“逆向达标”闭环，以目标达标率、约束满足率与样本效率作为验证指标；
  \item 该版本稳定后再逐步并入更多域（可靠性/法规/供应链/运维），并在每次扩展时同步补齐证书化边界。
\end{enumerate}
\end{remark}

\clearpage
%%%%%%%%%%%%%%%%%%%%%%%%%%%%%%%%%%%%%%%%%%%%%%%%%%%%%%%%%%%%
\section{\PTOPB{} 框架下 RTSC 的虚拟实验与逆向构造：可行性分层}
\label{sec:ptopb-rtsc}
%%%%%%%%%%%%%%%%%%%%%%%%%%%%%%%%%%%%%%%%%%%%%%%%%%%%%%%%%%%%

\subsection{问题设定与评估对象}

\begin{remark}[评估对象]
在 \PTOPB{} 框架下讨论“室温超导（RTSC）”时，本章对两类能力进行区分评估：
\begin{enumerate}
  \item \textbf{虚拟实验筛选能力（in-silico screening）：}仅依赖前向模型/仿真器，能将候选材料、结构、外场与工艺在多接近目标规格的意义下排序与筛选；
  \item \textbf{按需逆向构造能力（inverse construction）：}给定目标规格（例如 $T_c\ge 300\,\mathrm{K}$、$P\approx 1\,\mathrm{atm}$、可承受磁场与电流密度、
    可制造与可复现等），算法输出材料与工艺路径（算子词），并满足证书化可检验约束。
\end{enumerate}
两者的差异在于：前者只需要“可计算的前向映射”，而后者还需要“目标集合非空”与“模型误差可控到足以支持逆向优化”。
\end{remark}

\subsection{定义：\PTOPB{} 的超导特化对象与三域幺半群}

\begin{definition}[\PTOPB$_{\mathrm{SC}}$：超导逆向设计的总算子主纤维丛]
\label{def:ptopb-sc}
定义超导逆向设计的 \PTOPB{} 特化对象为
\[
  \PTOPB_{\mathrm{SC}}
  :=
  \bigl(
    M^{\mathrm{SC}},
    P_{\mathrm{tot}},
    G_{\mathrm{tot}},
    \Conn;\,
    \mathcal{M}_{\mathrm{env}},\,
    \mathcal{M}_{\mathrm{mat}},\,
    \mathcal{M}_{\mathrm{th}}
  \bigr),
\]
其中：
\begin{itemize}
  \item $M^{\mathrm{SC}}\subset \GZLS$ 为“超导系统态”的闭系统法级状态空间；
  \item $\mathcal{M}_{\mathrm{env}}$ 为环境操作幺半群（压力/温度/应力应变/磁电场/泵浦驱动/边界浴等）；
  \item $\mathcal{M}_{\mathrm{mat}}$ 为材料操作幺半群（化学计量、掺杂、缺陷、层状/界面、晶格畸变、纳米结构等）；
  \item $\mathcal{M}_{\mathrm{th}}$ 为理论操作幺半群（理论与数值近似选择，例如
  BCS/\allowbreak Eliashberg、 DFT+\allowbreak phonon、 DFT+\allowbreak DMFT、
  spin-\allowbreak fluctuation、 界面模型、 非平衡 Floquet 等）。
\end{itemize}
\end{definition}

\begin{definition}[总操作幺半群与算子词]
\label{def:ptopb-sc-word}
定义总操作幺半群为半直积
\[
\mathcal{M}_{\mathrm{tot}}
:=
\mathcal{M}_{\mathrm{env}}\rtimes(\mathcal{M}_{\mathrm{mat}}\rtimes \mathcal{M}_{\mathrm{th}}).
\]
一次“设计/合成/调控路径”表示为算子词
\[
w = a_1 a_2 \cdots a_T,\qquad a_t\in \mathcal{M}_{\mathrm{tot}},
\]
其中 $T\in\mathbb{N}$ 为词长。
\end{definition}

\subsection{定理（公理降级）：算子词表示完备性（按词长归纳）}

\begin{theorem}[算子词表示完备性（公理降级）]
\label{thm:ptopb-sc-word-completeness}
若 $\mathcal{M}_{\mathrm{tot}}$ 的元素均可执行且其连接运算对应可执行序列的连接，
则任何有限步的环境/材料/理论组合操作均可表示为一个有限词 $w\in \mathcal{W}(\mathcal{M}_{\mathrm{tot}})$。
\end{theorem}

\begin{certificate}[数学归纳法证书：按词长 $T$ 的表示归纳]
\label{cert:ptopb-sc-word-completeness-induction}
定义谓词 $C(T)$ 为：“任意 $T$ 步组合操作可表示为长度为 $T$ 的算子词”。证书化要求为：
\[
C(1)\ \text{成立},\qquad
(\forall T\in\mathbb{N}_{\ge 1})\ \bigl(C(T)\Rightarrow C(T+1)\bigr),
\]
其中递推步使用幺半群封闭性：将长度 $T$ 的词与第 $T{+}1$ 步操作连接得到长度 $T{+}1$ 的词。
\end{certificate}

\begin{proof}[证明（数学归纳法证书；谓词因果链）]
由 Certificate~\ref{cert:ptopb-sc-word-completeness-induction}：
\begin{itemize}
  \item 基步：$C(1)$ 成立（一步操作即为一个元素 $a_1\in\mathcal{M}_{\mathrm{tot}}$）；
  \item 归纳步：若 $C(T)$ 成立，则任意 $T{+}1$ 步操作可写为 $w a_{T+1}$，由幺半群封闭性得到长度 $T{+}1$ 的词表示，因此 $C(T{+}1)$ 成立。
\end{itemize}
由归纳法得任意有限步组合操作均可表示为算子词，证毕。
\end{proof}

\subsection{定义：前向映射、观测向量与目标规格}

\begin{definition}[前向映射与观测向量]
\label{def:forward-map-y}
给定算子词 $w$，定义前向映射输出为观测向量
\[
y(w):=
\bigl(
T_c(w),\,P(w),\,J_c(w),\,H_{c2}(w),\,\mathrm{Stab}(w),\,\mathrm{Proc}(w),\,\mathrm{Unc}(w)
\bigr),
\]
其中分别表示临界温度、所需压力、临界电流密度、上临界场、（热/动力学）稳定性指标、工艺可实现性指标与模型不确定度指标。
\end{definition}

\begin{definition}[目标规格与可达目标集合]
\label{def:target-set}
给定目标规格 $y^\star$ 与约束集合 $\mathcal{C}$，定义可达目标集合为
\[
\mathcal{W}_{\mathrm{goal}}(y^\star,\mathcal{C})
:=
\{\,w \mid y(w)\ \text{满足目标规格 }y^\star\ \text{且}\ w\ \text{满足约束}\ \mathcal{C}\,\}.
\]
\end{definition}

\subsection{定义：逆向构造的代价泛函与 GRL 路径积分口径}

\begin{definition}[代价泛函 $\mathcal{S}(w)$（证书与不确定度作为一等项）]
\label{def:cost-functional}
设 $[x]_+:=\max(x,0)$。给定目标 $T_c^\star,P^\star$ 以及权重 $\alpha,\beta,\gamma,\delta,\eta>0$，定义常用代价泛函为
\[
\mathcal{S}(w)=
\alpha\,[T_c^\star - T_c(w)]_+
\;+\;
\beta\,[P(w)-P^\star]_+
\;+\;
\gamma\,\mathrm{Instab}(w)
\;+\;
\delta\,\mathrm{Improc}(w)
\;+\;
\eta\,\mathrm{Unc}(w).
\]
其中 $\mathrm{Instab}$ 为稳定性罚项，$\mathrm{Improc}$ 为工艺不可实现罚项，$\mathrm{Unc}$ 为模型不确定导致的风险罚项。
\end{definition}

\begin{definition}[GRL 路径积分/分配函数口径]
\label{def:grl-partition}
设 $\pi$ 为在词空间上的策略分布，$\beta_{\mathrm{inv}}>0$ 为探索--收敛强度参数，定义
\[
\mathcal{Z}_{\mathrm{GRL}}(\beta_{\mathrm{inv}};\pi)
:=
\sum_{w}
\exp\!\bigl(-\beta_{\mathrm{inv}}\,\mathcal{S}(w)\bigr)\,\mathbb{P}_{\pi}(w).
\]
\end{definition}

\subsection{证书库：RTSC 评估的现状锚点（外部材料）}

\begin{certificate}[证书：常压常规（electron--phonon）机制下室温 $T_c$ 的证据性上限倾向]
\label{cert:natcomm-2025-max-tc}
公开研究对常规（electron--phonon / Eliashberg）超导在常压实现室温 $T_c$ 的可行性给出不利结论倾向：在大规模计算数据与统计分析下，
常压实现室温“常规超导”极不可能；且随着预测 $T_c$ 提升，热力学稳定性变差将显著抬高合成难度。\cite{NatComm2025MaxTcConvAmbient}
\end{certificate}

\begin{certificate}[证书：高压氢化物（superhydrides）高 $T_c$ 的理论--实验一致性样本]
\label{cert:nsr-superhydrides-review}
高压氢化物在中高压条件下出现高 $T_c$ 的路线具有较强的理论--实验匹配度，典型案例包括 LaH$_{10}$ 在高压下报告的高 $T_c$ 及其独立确认；同时该路线的关键张力在于“高压稳定氢笼结构”与“常压使用”的冲突。
  \cite{NSR2024Superhydrides}
\end{certificate}

\begin{certificate}[证书：近常压/室温个别报道的可重复性风险样本]
\label{cert:pubmed-retraction}
近年关于“近常压/室温超导”的个别报道出现撤稿与可重复性危机，属于必须显式纳入“证书/可复现约束”的风险样本。\cite{PubMedRetractionLutetiumHydride}
\end{certificate}

\begin{certificate}[证书：常压附近 $\sim 100\,\mathrm{K}$ 量级理论候选与非简谐/量子核效应]
\label{cert:rbph3-ambient-100k}
公开理论工作提出：在常压附近可达到约 $100\,\mathrm{K}$ 量级 $T_c$ 的候选体系（如 RbPH$_3$）需要将离子量子非简谐效应等稳定化机制纳入模型；该结论为理论预测口径，需与实验复核分离管理。
  \cite{ScienceDirectRbPH3}
\end{certificate}

\begin{certificate}[证书：常压高温超导数据集与基准回归口径]
\label{cert:htsc-2025-dataset}
公开基准数据集与任务定义（如 HTSC-2025）可用于常压高温超导候选的 $T_c$ 预测回归与模型对比，但不替代实验事实；其价值在于提供可复核的训练/验证分割与基线指标。\cite{arXivHTSC2025}
\end{certificate}

\subsection{定理：\PTOPB{} 将可行性转写为“目标集合是否为空”与“到目标集合的距离”}

\begin{theorem}[可行性等价：目标集合非空与代价下确界]
\label{thm:feasibility-as-emptiness}
在 Definition~\ref{def:target-set} 与 Definition~\ref{def:cost-functional} 下：
\begin{enumerate}
  \item 若 $\mathcal{W}_{\mathrm{goal}}(y^\star,\mathcal{C})\neq\varnothing$，则存在 $w$ 使 $\mathcal{S}(w)
    =0$（在罚项与阈值按目标规格定义的口径下）；
  \item 若 $\mathcal{W}_{\mathrm{goal}}(y^\star,\mathcal{C})=\varnothing$，则 $\inf_{w}\mathcal{S}(w)>0$ 或者
    $\inf_w\mathcal{S}(w)$ 由稳定性/压力/不确定度罚项主导，从而给出“目标不可达或近似不可达”的可计算证据链条入口。
\end{enumerate}
\end{theorem}

\begin{certificate}[证书：代价泛函的阈值化构造口径]
\label{cert:thresholded-cost}
在 Definition~\ref{def:cost-functional} 中，若将稳定性与工艺罚项定义为阈值化形式（例如 $\mathrm{Instab}(w)=0$ 当且仅当稳定性满足要求，$\mathrm{Improc}(w)=0$
  当且仅当工艺可实现），并将 $T_c^\star,P^\star$ 与其他约束统一纳入目标规格，则 $\mathcal{S}(w)=0$ 等价于“全部目标与约束均满足”。
\end{certificate}

\begin{proof}[证明（谓词因果链）]
由 Certificate~\ref{cert:thresholded-cost}，可将谓词链写为：
\[
A(w):=\bigl(w\in\mathcal{W}_{\mathrm{goal}}(y^\star,\mathcal{C})\bigr),\qquad
B(w):=\bigl(\mathcal{S}(w)=0\bigr).
\]
证书给出 $A(w)\Leftrightarrow B(w)$。因此：
\begin{itemize}
  \item 若 $\mathcal{W}_{\mathrm{goal}}\neq\varnothing$，则存在 $w$ 使 $A(w)$ 成立，从而存在 $w$ 使 $B(w)$ 成立；
  \item 若 $\mathcal{W}_{\mathrm{goal}}=\varnothing$，则对任意 $w$ 有 $\neg A(w)$，从而对任意 $w$ 有 $\neg B(w)$，即不存在 $\mathcal{S}(w)
    =0$；此时 $\inf_w\mathcal{S}(w)$ 的正下界或主导项即给出“不可达/近似不可达”的可计算入口。
\end{itemize}
证毕。
\end{proof}

\subsection{定义：三类 RTSC 目标（A/B/C）}

\begin{definition}[目标 A/B/C：按压力与温度规格分层]
\label{def:rtsc-targets}
以 $(T_c^\star,P^\star)$ 与稳定性/工艺/可复现约束为核心，定义三类目标：
\begin{enumerate}
  \item \textbf{目标 A（常压稳态室温可工程化）：} $T_c^\star\ge 300\,\mathrm{K}$ 且 $P^\star\approx 1\,\mathrm{atm}$，
    并要求稳定性与工艺可实现性满足工程化约束；
  \item \textbf{目标 B（中高压下接近室温的高 $T_c$）：} 在允许中高压操作的条件下追求高 $T_c$（例如 $T_c^\star\ge 200\text{--}250\,\mathrm{K}$ 且
    $P^\star\gg 1\,\mathrm{atm}$），并以压力操作作为显式可执行环境算子；
  \item \textbf{目标 C（常压但非室温或亚稳态维持）：} 目标常压或近常压下的中高 $T_c$（例如 $50\text{--}120\,\mathrm{K}$ 量级），
    或在高压合成后以亚稳态淬火/维持到常压的可实现性为附加目标。
\end{enumerate}
\end{definition}

\subsection{定理：三类目标在 \PTOPB$_{\mathrm{SC}}$ 中的可行性分层（证书化口径）}

\begin{theorem}[可行性分层：目标 A/B/C]
\label{thm:rtsc-feasibility-tier}
在 \PTOPB$_{\mathrm{SC}}$（Definition~\ref{def:ptopb-sc}）下，若将理论操作幺半群 $\mathcal{M}_{\mathrm{th}}$
分解为机制族子幺半群（例如常规 electron--phonon 子族与高压氢化物相关子族等），则：
\begin{enumerate}
  \item \textbf{目标 A：}当 $\mathcal{M}_{\mathrm{th}}$ 限制在常规 electron--phonon / Eliashberg 机制族时，
    依据现有证据倾向（Certificate~\ref{cert:natcomm-2025-max-tc}），目标集合的非空性难以被证书化建立；在该机制族子空间上更适合作为“可证伪/可量化接近程度”的问题（即
    $\inf_w\mathcal{S}(w)$ 的收敛行为）；
  \item \textbf{目标 B：}当 $\mathcal{M}_{\mathrm{env}}$ 允许中高压操作且 $\mathcal{M}_{\mathrm{th}}$ 覆盖高压氢化物相关模型链时，
    存在更强的理论--实验一致性样本（Certificate~\ref{cert:nsr-superhydrides-review}），从而更接近“可操作的逆向设计算法框架”；
  \item \textbf{目标 C：}在常压附近取得 $\sim 100\,\mathrm{K}$ 量级 $T_c$ 的理论候选存在公开预测（Certificate~\ref{cert:rbph3-ambient-100k}），
    但其可行性依赖于将关键稳定化机制显式纳入 $\mathcal{M}_{\mathrm{th}}$，并对 $\mathrm{Unc}(w)$ 进行证书化管理。
\end{enumerate}
\end{theorem}

\begin{certificate}[证书：三类目标的外部证据对应]
\label{cert:rtsc-tier-evidence}
目标分层对应证据入口为：
\[
\text{目标 A}\Rightarrow \text{Certificate~\ref{cert:natcomm-2025-max-tc}},\quad
\text{目标 B}\Rightarrow \text{Certificate~\ref{cert:nsr-superhydrides-review}},\quad
\text{目标 C}\Rightarrow \text{Certificate~\ref{cert:rbph3-ambient-100k}}.
\]
其中目标 A 的证据为“不利倾向”，目标 B 为“较强一致性样本”，目标 C 为“理论预测且需机制补全”。
\end{certificate}

\begin{proof}[证明（谓词因果链）]
由 Certificate~\ref{cert:rtsc-tier-evidence}，分别取谓词链：
\[
\begin{aligned}
P_A:&\ \mathcal{M}_{\mathrm{th}}\subseteq \mathcal{M}_{\mathrm{th}}^{\mathrm{conv}}\wedge \text{证据倾向不利},\\
P_B:&\ \mathcal{M}_{\mathrm{env}}\ \text{允许高压}\wedge \text{存在一致性样本},\\
P_C:&\ \text{存在理论候选}\wedge \text{需显式纳入机制项}.
\end{aligned}
\]
则由各证书内容可推出：
\[
\begin{aligned}
P_A&\Rightarrow \text{目标 A 的非空性难以证书化建立},\\
P_B&\Rightarrow \text{目标 B 更可操作},\\
P_C&\Rightarrow \text{目标 C 依赖机制补全与不确定度管理}.
\end{aligned}
\]
三者合取即得结论。证毕。
\end{proof}

\subsection{引理：虚拟实验的适用域与不确定度项的必要性}

\begin{lemma}[虚拟实验可用域：机制族决定前向误差的可控性]
\label{lem:insilico-validity}
在 \PTOPB$_{\mathrm{SC}}$ 中：
\begin{enumerate}
  \item 对常规 electron--phonon 机制链与高压氢化物相关模型链，前向模拟链（如
  DFT+\allowbreak phonon+\allowbreak Eliashberg/\allowbreak Allen--Dynes）相对成熟，
  $\mathrm{Unc}(w)$ 更可能被证书化上界约束；
  \item 对强关联/非常规配对/界面诱导/非平衡驱动等模型族，近似敏感性强，若不将 $\mathrm{Unc}(w)$ 作为代价一等项，则逆向优化存在“模型投机解”风险。
\end{enumerate}
\end{lemma}

\begin{certificate}[证书：前向模型成熟度与不确定度管理的必要性]
\label{cert:uncertainty-necessity}
对不同机制族，前向模型链条的成熟度与误差结构差异显著；在不确定度较大的模型族中，将 $\mathrm{Unc}(w)$ 纳入 $\mathcal{S}(w)$ 是防止优化过程利用模型漏洞的必要手段。该证书为工程方法学口径，
  可由实际对比试验（仿真--实验偏差统计）逐步量化与替换。
\end{certificate}

\begin{proof}[证明（谓词因果链）]
由 Certificate~\ref{cert:uncertainty-necessity}，取谓词链：
\[
\begin{aligned}
Q_1:&\ \text{机制族不同}\Rightarrow \text{前向误差结构不同},\\
Q_2:&\ \text{误差结构不可控}\wedge \mathrm{Unc}\ \text{未入代价}\Rightarrow \text{优化可能产生投机解}.
\end{aligned}
\]
由 $Q_1$ 与 $Q_2$ 得到结论（1）（2）。证毕。
\end{proof}

\subsection{命题：可复现证书约束作为硬门槛并入代价泛函}

\begin{proposition}[可复现门槛：以证书谓词裁剪/惩罚非复现区域]
\label{prop:cert-gating}
定义可复现证书谓词 $\mathrm{Cert}(w)\in\{0,1\}$，并对代价作门槛化修正：
\[
\mathcal{S}_{\mathrm{cert}}(w):=\mathcal{S}(w)+M\cdot(1-\mathrm{Cert}(w)),
\qquad M\gg 0.
\]
则在 $M$ 足够大时，任意 $\mathrm{Cert}(w)=0$ 的候选将被优化过程在代价意义下显式排除或强烈惩罚，从而避免资源集中于不可复现区域。
\end{proposition}

\begin{certificate}[证书：不可复现风险样本与门槛化必要性]
\label{cert:reproducibility-risk}
近常压/室温个别报道存在撤稿与可重复性危机（Certificate~\ref{cert:pubmed-retraction}），因此将“可复现性”作为硬约束或强惩罚项并入代价泛函属于必要的风险控制口径。
\end{certificate}

\begin{proof}[证明（谓词因果链）]
由 Certificate~\ref{cert:reproducibility-risk} 取风险控制谓词 $R(w):=\neg \mathrm{Cert}(w)$。
若 $R(w)$ 成立，则 $\mathcal{S}_{\mathrm{cert}}(w)=\mathcal{S}(w)+M\ge M$。取任意满足 $\mathrm{Cert}(w)=1$ 的候选 $w'$，有 $\mathcal{S}
  _{\mathrm{cert}}(w')=\mathcal{S}(w')$。
当 $M\gg \sup_w \mathcal{S}(w)$（或在可比集合内 $M$ 大于可接受代价上界）时，优化 $\min_w \mathcal{S}_{\mathrm{cert}}(w)$ 将优先在 $\mathrm{Cert}(w)
  =1$ 的子域内搜索；因此不可复现候选被排除/强惩罚。证毕。
\end{proof}

\subsection{推论：按需逆向构造成立的两条充分条件}

\begin{corollary}[按需逆向构造的充分条件（非空性 + 误差可控）]
\label{cor:inverse-design-conditions}
在 \PTOPB$_{\mathrm{SC}}$ 下，若同时满足：
\begin{enumerate}
  \item 目标集合非空：$\mathcal{W}_{\mathrm{goal}}(y^\star,\mathcal{C})\neq\varnothing$；
  \item 前向误差可控：存在证书化上界使 $\mathrm{Unc}(w)$ 在搜索域内被控制到足以支持优化收敛的尺度；
\end{enumerate}
则“按需逆向构造”可被表述为可计算的最优化问题（例如 $\min_w \mathcal{S}_{\mathrm{cert}}(w)$ 或等价 GRL 形式），并具备可收敛的工程解释。
\end{corollary}

\begin{certificate}[证书：两条件合取的逻辑充分性]
\label{cert:inverse-design-sufficient}
条件（1）保证存在可达零代价点（Theorem~\ref{thm:feasibility-as-emptiness}），条件（2）保证代价景观不会被不可控模型误差主导（Lemma~\ref{lem:insilico-validity}
  的口径）。两者合取给出“逆向构造可计算且可解释”的充分条件。
\end{certificate}

\begin{proof}[证明（谓词因果链）]
由 Certificate~\ref{cert:inverse-design-sufficient}，取谓词链：
\[
S_1:\ \mathcal{W}_{\mathrm{goal}}\neq\varnothing \Rightarrow \exists w:\mathcal{S}(w)=0,\quad
S_2:\ \mathrm{Unc}\ \text{可控}\Rightarrow \text{优化过程不被误差主导}.
\]
由 $S_1\wedge S_2$ 可得逆向构造问题在形式上等价为可计算最优化，并具备可收敛的工程解释。证毕。
\end{proof}

\subsection{备注：最小闭环形态（虚拟--实验主动学习）}

\begin{remark}[最小闭环形态]
对目标 A 或对高不确定度机制族的目标 C，仅依赖纯虚拟逆向构造往往不足以形成证书化闭环。较稳妥的最小闭环形态为：
\begin{enumerate}
  \item 以 $\mathcal{S}_{\mathrm{cert}}(w)$ 作为统一目标，显式包含 $T_c$、压力、稳定性、工艺、证书与不确定度项；
  \item 将 $\mathcal{M}_{\mathrm{th}}$ 分层为“快速近似（筛选）$\to$ 中等精度（排序）$\to$ 高精度（少量验证）”的理论算子链；
  \item 每轮由内环（GRL-SAC/等价优化）产生少量“高期望改进 + 低不确定度 + 可证书化”的候选，进入真实实验；
  \item 实验结果反哺更新 $\mathcal{M}_{\mathrm{th}}$ 的参数与 $\mathrm{Unc}$ 评估，使证书化边界逐步收敛。
\end{enumerate}
\end{remark}

\clearpage
%%%%%%%%%%%%%%%%%%%%%%%%%%%%%%%%%%%%%%%%%%%%%%%%%%%%%%%%%%%%
\section{\PTOPB$_{\mathrm{SC}}$ 的路线插件化：子丛/子幺半群与统一引擎}
\label{sec:ptopb-plugin}
%%%%%%%%%%%%%%%%%%%%%%%%%%%%%%%%%%%%%%%%%%%%%%%%%%%%%%%%%%%%

\subsection{定义：路线算子丛与母体 \PTOPB$_{\mathrm{SC}}$ 的包含关系}

\begin{definition}[路线算子丛 $\OPB^{(k)}_{\RTSC}$]
\label{def:opb-k}
设 $\mathcal{K}$ 为机制路线索引集合（例如高压氢化物、界面/层状增强、非平衡驱动、强关联非常规等的分类索引）。
对任意 $k\in\mathcal{K}$，定义第 $k$ 条路线对应的路线算子丛为
\[
  \OPB^{(k)}_{\RTSC}
  :=
  \bigl(
    M^{(k)},
    P^{(k)},
    G^{(k)},
    \Conn^{(k)};\,
    \mathcal{M}^{(k)}_{\mathrm{env}},\,
    \mathcal{M}^{(k)}_{\mathrm{mat}},\,
    \mathcal{M}^{(k)}_{\mathrm{th}}
  \bigr),
\]
其中 $M^{(k)}\subseteq M^{\mathrm{SC}}$ 为该机制路线下的可行子域，$\mathcal{M}^{(k)}_{\mathrm{env/mat/th}}$ 为该路线允许的环境/材料/理论操作字典（子幺半群）。
\end{definition}

\begin{definition}[母体包含与路线约束投影]
\label{def:mother-inclusion-projection}
称 \PTOPB$_{\mathrm{SC}}$（Definition~\ref{def:ptopb-sc}）为路线族 $\{\OPB^{(k)}_{\RTSC}\}_{k\in\mathcal{K}}$ 的\textbf{母体}，
  若满足包含关系
\[
M^{\mathrm{SC}} \supseteq \bigcup_{k\in\mathcal{K}} M^{(k)},
\qquad
\mathcal{M}_{\bullet} \supseteq \bigcup_{k\in\mathcal{K}} \mathcal{M}^{(k)}_{\bullet},
\quad \bullet\in\{\mathrm{env},\mathrm{mat},\mathrm{th}\}.
\]
并定义路线约束投影（或称切片）为一族映射
\[
\Pi_k:\PTOPB_{\mathrm{SC}}\to \OPB^{(k)}_{\RTSC},
\]
其作用是：在母体中施加“机制选择 + 结构约束 + 动作字典限制”，并将对象/态射限制到路线 $k$ 的子域。
\end{definition}

\begin{certificate}[证书：路线插件化的结构约束口径]
\label{cert:plugin-structure}
路线插件化以三类可核验输入给出：
\begin{enumerate}
  \item \textbf{动作字典：}给出 $\mathcal{M}^{(k)}_{\mathrm{env/mat/th}}$ 作为 $\mathcal{M}_{\mathrm{env/mat/th}}$ 的子幺半群；
  \item \textbf{前向评估器分层：}给出前向评估器族 $F^{(k)}_{\mathrm{fast/mid/hi}}$，用于实现 $w\mapsto y(w)$ 的多保真近似；
  \item \textbf{证书谓词：}给出可复现与记录字段的证书谓词 $\mathrm{Cert}^{(k)}(w)$（或其等价的证书化约束集合）。
\end{enumerate}
上述三项共同确定 $\Pi_k$ 的限制口径与路线 $k$ 的可执行搜索域。
\end{certificate}

\subsection{定理（公理降级）：母体统一引擎与插件接口的可组合性}

\begin{theorem}[统一引擎（公理降级）：固定目标函数与策略更新，插件仅改变可行域与评估器]
\label{thm:unified-engine}
在 Definition~\ref{def:ptopb-sc-word}--Definition~\ref{def:grl-partition} 的口径下，设母体引擎由下列三项组成：
\begin{enumerate}
  \item 统一目标与代价：$\mathcal{S}_{\mathrm{cert}}(w)$（由 Definition~\ref{def:cost-functional} 与 Proposition~\ref{prop:cert-gating} 给出）；
  \item 统一策略更新：以 Definition~\ref{def:grl-partition} 的分配函数口径（或等价的离散 GRL-SAC 目标）更新策略 $\pi$；
  \item 统一空集/近似空集判别入口：以 Theorem~\ref{thm:feasibility-as-emptiness} 的“目标集合非空 $\Leftrightarrow$ 零代价点存在”作为形式判据入口。
\end{enumerate}
若对每个 $k\in\mathcal{K}$ 给定 Certificate~\ref{cert:plugin-structure} 的插件三元组，则路线 $k$ 的逆向构造问题可表述为同一引擎在受限可行域上的实例：
\[
\min_{w\in \mathcal{W}^{(k)}_{\mathrm{adm}}}\ \mathcal{S}_{\mathrm{cert}}^{(k)}(w),
\]
其中 $\mathcal{W}^{(k)}_{\mathrm{adm}}$ 由 $\mathcal{M}^{(k)}_{\mathrm{env/mat/th}}$ 与序列约束诱导，$\mathcal{S}_{\mathrm{cert}}
  ^{(k)}$ 由 $\mathrm{Cert}^{(k)}$ 与不确定度评估器替换得到；引擎骨架（目标形式与更新机制）保持不变。
\end{theorem}

\begin{certificate}[证书：引擎不变性与插件替换点]
\label{cert:engine-invariance}
引擎骨架只依赖以下接口输出：
\[
y(w),\quad \mathrm{Unc}(w),\quad \mathrm{Cert}(w),\quad \text{以及可计算的序列约束谓词}.
\]
插件仅替换：
\[
\mathcal{M}_{\mathrm{env/mat/th}}\mapsto \mathcal{M}^{(k)}_{\mathrm{env/mat/th}},\quad
F_{\mathrm{fast/mid/hi}}\mapsto F^{(k)}_{\mathrm{fast/mid/hi}},\quad
\mathrm{Cert}\mapsto \mathrm{Cert}^{(k)}.
\]
\end{certificate}

\begin{proof}[证明（谓词因果链）]
由 Certificate~\ref{cert:engine-invariance}，引擎对插件的依赖被限制在“可行域（动作字典/序列约束）+ 前向评估器 + 证书谓词”三类替换点。
代价形式 $\mathcal{S}_{\mathrm{cert}}$ 与分配函数口径 $\mathcal{Z}_{\mathrm{GRL}}$ 均保持同型，仅将 $y,\mathrm{Unc},\mathrm{Cert}$ 的实现替换为路线
  $k$ 的对应版本；
因此路线 $k$ 的问题即为母体引擎在受限域上的实例，结论成立。证毕。
\end{proof}

\subsection{引理：机制选择作为理论算子的一阶内化（更快的形式来源）}

\begin{lemma}[机制内化：路线索引 $k$ 可作为理论操作字典的元素或标签]
\label{lem:k-in-mth}
在 Definition~\ref{def:ptopb-sc} 的口径下，若将“机制路线选择”编码为理论操作幺半群 $\mathcal{M}_{\mathrm{th}}$ 的元素或标签（例如以算子 $k\in \mathcal{M}
  _{\mathrm{th}}$ 表示进入某一理论链与近似层级），
则母体引擎可在同一算子词空间上先进行广域搜索，再通过路线投影 $\Pi_k$ 收缩到具体路线子域进行细化。
\end{lemma}

\begin{certificate}[证书：机制标签与理论算子链的同构口径]
\label{cert:k-as-theory-operator}
对任意路线 $k$，其核心差异之一是理论链与近似层级的选择（对应 $\mathcal{M}^{(k)}_{\mathrm{th}}$ 的取值与序列约束）。因此可将“路线选择”视为在理论操作字典中的可执行选择，并由 $\mathcal{M}
  _{\mathrm{th}}$ 的子幺半群限制实现。
\end{certificate}

\begin{proof}[证明（谓词因果链）]
由 Certificate~\ref{cert:k-as-theory-operator}，路线差异可由理论操作字典的子幺半群限制刻画。将该限制编码为词空间中的元素/标签选择，即可在母体词空间先行搜索；
随后通过 $\Pi_k$ 将候选限制到 $k$ 的可行子域完成细化。证毕。
\end{proof}

\subsection{命题：三项优势的形式化解释（更快/更全面/更可操作）}

\begin{proposition}[优势命题：统一引擎下的速度、完整性与可操作性]
\label{prop:three-advantages}
在 Theorem~\ref{thm:unified-engine} 与 Lemma~\ref{lem:k-in-mth} 的前提下：
\begin{enumerate}
  \item \textbf{更快：}路线 $k$ 的引入在形式系统层面表现为对可行域与评估器的替换（Certificate~\ref{cert:plugin-structure}），而不需要重写目标函数与策略更新。
  并且机制选择可在词空间中以内化方式处理（Lemma~\ref{lem:k-in-mth}），从而减少“先定路线后建模”的早期锁定风险。
  \item \textbf{更全面：}统一观测向量与代价泛函将稳定性、工艺可实现性、不确定度与可复现证书作为一等项（见 Definition~\ref{def:forward-map-y}、Definition~\ref{def:cost-functional} 与 Proposition~\ref{prop:cert-gating}）。
  因此“理论可行但不稳定、数值好看但不可制造、结果不可复现”等失败模式可转写为同一优化问题中的罚项或硬门槛。
  \item \textbf{更可操作：}算子词 $w$ 作为可执行序列（Theorem~\ref{thm:ptopb-sc-word-completeness}），与证书谓词 $\mathrm{Cert}^{(k)}$
    的门槛化约束共同保证输出可被记录、复核与复现实验流程对齐。
\end{enumerate}
\end{proposition}

\begin{certificate}[证书：三项优势的依赖关系]
\label{cert:three-advantages-deps}
（1）依赖 Theorem~\ref{thm:unified-engine} 与 Lemma~\ref{lem:k-in-mth}；（2）依赖 Definition~\ref{def:cost-functional} 与
  Proposition~\ref{prop:cert-gating}；（3）依赖 Theorem~\ref{thm:ptopb-sc-word-completeness} 与路线证书谓词的可执行定义。
\end{certificate}

\begin{proof}[证明（谓词因果链）]
由 Certificate~\ref{cert:three-advantages-deps}：
\begin{itemize}
  \item 对（1）：由 Theorem~\ref{thm:unified-engine} 得“插件仅替换可行域与评估器”，由 Lemma~\ref{lem:k-in-mth} 得“路线选择可内化进理论操作字典”，
    从而形式上减少早期锁定并允许统一引擎复用；
  \item 对（2）：由代价定义与证书门槛化得失败模式在同一 $\mathcal{S}_{\mathrm{cert}}$ 中可比较；
  \item 对（3）：由词表示完备性得输出为可执行序列，叠加证书谓词得输出可复核。
\end{itemize}
三者分别推出命题三项结论，证毕。
\end{proof}

\subsection{推论：Vol I 路线的模块化表达（插件目录口径）}

\begin{corollary}[模块化推论：路线以插件目录形式并入母体而不改变引擎]
\label{cor:plugin-catalog}
将 Vol I 中的每条 \RTSC{} 路线表示为插件
\[
\mathrm{Plugin}(k):=\bigl(\mathcal{M}^{(k)}_{\mathrm{env/mat/th}},\ F^{(k)}_{\mathrm{fast/mid/hi}},\ \mathrm{Cert}^{(k)}
  \bigr),
\]
并令母体问题为
\[
\mathrm{Solve}:\ \min_{w}\mathcal{S}_{\mathrm{cert}}(w),
\]
则可在不改变 $\mathrm{Solve}$ 的形式与策略更新机制的前提下，通过更换 $\mathrm{Plugin}(k)$ 实现路线切换；换言之，“更换插件而不重写引擎”在形式系统层面成立。
\end{corollary}

\begin{certificate}[证书：Vol I 路线作为插件的可核验输入项]
\label{cert:vol1-plugin}
Vol I 路线的最小插件输入可归约为三项：动作字典、前向评估器分层与证书字段集合（与 Certificate~\ref{cert:plugin-structure} 同型）。该归约为工程接口口径，不依赖对路线叙事的额外假设。
\end{certificate}

\begin{proof}[证明（谓词因果链）]
由 Theorem~\ref{thm:unified-engine}，插件替换不改变引擎骨架；由 Certificate~\ref{cert:vol1-plugin}，Vol I 的路线可按同型三元组给出。
因此路线切换等价于插件切换，而 $\mathrm{Solve}$ 的形式保持不变。证毕。
\end{proof}

\clearpage
%%%%%%%%%%%%%%%%%%%%%%%%%%%%%%%%%%%%%%%%%%%%%%%%%%%%%%%%%%%%
\section{环境同伦复用：虚拟环境对象、适配器与 GRL 路径积分稳定性界}
\label{sec:env-homotopy-reuse}
%%%%%%%%%%%%%%%%%%%%%%%%%%%%%%%%%%%%%%%%%%%%%%%%%%%%%%%%%%%%

\subsection{定义：虚拟环境作为可比较对象}

\begin{definition}[虚拟环境对象 $\mathcal{E}$]
\label{def:env-object}
称一个“可用于逆向设计/虚拟实验的环境对象”为
\[
\mathcal{E}
:=
\Bigl(
M;\,P\to M;\,G;\,\Conn;\,
\mathcal{M}_{\mathrm{env}},\mathcal{M}_{\mathrm{mat}},\mathcal{M}_{\mathrm{th}};\,
F;\,\mathrm{Cert}
\Bigr),
\]
其中：
\begin{itemize}
  \item $M\subseteq \GZLS$ 为底空间（法空间中的可观测态/可控态集合，并含闭系统约束）；
  \item $P\to M$、$G$、$\Conn$ 用于编码等价表述、隐变量、路径依赖与不可交换性（由联络/曲率语义刻画）；
  \item $\mathcal{M}_{\mathrm{env}}$、$\mathcal{M}_{\mathrm{mat}}$、$\mathcal{M}_{\mathrm{th}}$ 为环境/材料/理论近似三类算子幺半群；
  \item $F$ 为前向评估器族（仿真器/代理模型/多保真链），用于实现 $w\mapsto y(w)$ 与相关代价项；
  \item $\mathrm{Cert}$ 为证书接口（可复现、可核验、可记录的最低约束），用于定义 $\mathrm{Cert}(w)\in\{0,1\}$ 或等价的证书化判据集合。
\end{itemize}
\end{definition}

\begin{definition}[总算子词空间]
\label{def:env-wordspace}
定义总算子幺半群为
\[
\mathcal{M}_{\mathrm{tot}}
:=
\mathcal{M}_{\mathrm{env}}\rtimes(\mathcal{M}_{\mathrm{mat}}\rtimes \mathcal{M}_{\mathrm{th}}),
\]
并记其词空间为 $\mathcal{M}_{\mathrm{tot}}^{\ast}$。对 $w\in \mathcal{M}_{\mathrm{tot}}^{\ast}$，$F$ 给出该词的前向输出（观测向量、稳定性/工艺/不确定度指标等）
  ，并诱导代价泛函 $\mathcal{S}_{\mathcal{E}}(w)$。
\end{definition}

\begin{remark}[PTOPB/LBOPB 与环境对象类]
在该口径下，\LBOPB{} 与 \PTOPB{} 的差异主要体现在：底空间 $M$ 的语义域、以及可用算子字典与评估器族 $(\mathcal{M}_{\bullet},F)$ 的取值差异。
因此二者可被视为同一“环境对象类”中的不同实例：它们在结构层面可比较，而在语义层面可替换。
\end{remark}

\subsection{定义：同伦适配器与三项兼容性条件}

\begin{definition}[同伦适配器（三元组）]
\label{def:homotopy-adapter}
给定两个环境对象 $\mathcal{E}_1,\mathcal{E}_2$，称三元组
\[
(\Phi,\Psi,\phi)
\]
为从 $\mathcal{E}_1$ 到 $\mathcal{E}_2$ 的同伦适配器，若：
\begin{enumerate}
  \item $\Phi:\mathcal{E}_1\to\mathcal{E}_2$、$\Psi:\mathcal{E}_2\to\mathcal{E}_1$ 是结构保持映射（环境适配器），在底空间、联络结构、
    算子字典与证书接口上保持可解释性；
  \item $\phi:\mathcal{M}_{\mathrm{tot}}^{(1)\ast}\to \mathcal{M}_{\mathrm{tot}}^{(2)\ast}$ 为词编译映射，将 $\mathcal{E}_1$
    的算子词编译为 $\mathcal{E}_2$ 的算子词；
  \item 存在“可观测意义下的同伦”关系，使得
  \[
  \Psi\circ\Phi \simeq \mathrm{Id}_{\mathcal{E}_1},
  \qquad
  \Phi\circ\Psi \simeq \mathrm{Id}_{\mathcal{E}_2},
  \]
  其中 $\simeq$ 的判定以可观测输出、证书接口与代价残差界为口径（见下文定义）。
\end{enumerate}
\end{definition}

\begin{definition}[三项兼容性条件：词编译、代价兼容、证书可迁移]
\label{def:compatibility-three}
称 $(\Phi,\Psi,\phi)$ 满足三项兼容性条件，若存在常数 $\varepsilon\ge 0$ 使：
\begin{enumerate}
  \item \textbf{词编译兼容：}对任意词 $w\in \mathcal{M}_{\mathrm{tot}}^{(1)\ast}$，其结构（连接/生成元约束）在编译后被保持，即 $\phi$ 与词连接相容（见
    Theorem~\ref{thm:word-compilation}）。
  \item \textbf{代价同伦兼容：}对任意词 $w$，存在残差 $r(w)$ 满足
  \[
  \mathcal{S}_{\mathcal{E}_1}(w)=\mathcal{S}_{\mathcal{E}_2}(\phi(w)) + r(w),
  \qquad |r(w)|\le \varepsilon.
  \]
  \item \textbf{证书接口可迁移：}对任意词 $w$，有
  \[
  \mathrm{Cert}_1(w)=1\ \Longrightarrow\ \mathrm{Cert}_2(\phi(w))=1,
  \]
  并且对称地可在另一方向给出对应条件（用于双向复用时）。
\end{enumerate}
\end{definition}

\subsection{定理（公理降级）：词编译兼容性可由词长归纳验证}

\begin{theorem}[词编译兼容性（公理降级）：生成元一致 $\Rightarrow$ 词级一致]
\label{thm:word-compilation}
设 $\phi$ 在生成元集合上给定，并满足对任意生成元 $a$ 有 $\phi(a)\in \mathcal{M}_{\mathrm{tot}}^{(2)}$；且满足连接相容性：对任意词 $w$ 与生成元 $a$，
\[
\phi(wa)=\phi(w)\,\phi(a).
\]
则对任意 $T\in\mathbb{N}$ 与任意长度为 $T$ 的词 $w=a_1a_2\cdots a_T$，有
\[
\phi(w)=\phi(a_1)\phi(a_2)\cdots \phi(a_T),
\]
从而 $\phi$ 将词结构一致地编译到目标词空间。
\end{theorem}

\begin{certificate}[数学归纳法证书：按词长 $T$ 的编译一致性]
\label{cert:word-compilation-induction}
定义谓词 $P(T)$ 为：“对任意长度为 $T$ 的词 $w=a_1\cdots a_T$，有 $\phi(w)=\phi(a_1)\cdots\phi(a_T)$”。证书化要求为：
\[
P(1)\ \text{成立},\qquad
(\forall T\in\mathbb{N}_{\ge 1})\ \bigl(P(T)\Rightarrow P(T+1)\bigr),
\]
其中递推步使用连接相容性 $\phi(wa)=\phi(w)\phi(a)$。
\end{certificate}

\begin{proof}[证明（数学归纳法证书；谓词因果链）]
由 Certificate~\ref{cert:word-compilation-induction}：
\begin{itemize}
  \item 基步：$P(1)$ 成立，因 $\phi(a_1)=\phi(a_1)$；
  \item 归纳步：若 $P(T)$ 成立，则任意长度 $T{+}1$ 的词可写为 $w'a_{T+1}$，由连接相容性与归纳假设得
  \[
  \phi(w'a_{T+1})=\phi(w')\phi(a_{T+1})=\phi(a_1)\cdots\phi(a_T)\phi(a_{T+1}),
  \]
  即 $P(T{+}1)$ 成立。
\end{itemize}
由归纳法得对任意 $T$ 成立，证毕。
\end{proof}

\subsection{定理：证书边界与可行域的可迁移裁剪}

\begin{theorem}[证书可迁移性：失败域与可行域的编译裁剪]
\label{thm:cert-reuse}
设 $\mathcal{E}_1,\mathcal{E}_2$ 与 $\phi$ 满足 Definition~\ref{def:compatibility-three} 的证书可迁移条件。令
\[
\mathcal{B}_1^{\mathrm{fail}}:=\{w\mid \mathrm{Cert}_1(w)=0\},\qquad
\mathcal{B}_2^{\mathrm{fail}}:=\{u\mid \mathrm{Cert}_2(u)=0\}.
\]
则对任意 $w\in\mathcal{M}_{\mathrm{tot}}^{(1)\ast}$，若 $w\in \mathcal{B}_1^{\mathrm{fail}}$，则 $\phi(w)\in \mathcal{B}
  _2^{\mathrm{fail}}$；因此可将 $\phi(\mathcal{B}_1^{\mathrm{fail}})$ 作为在 $\mathcal{E}_2$ 上的可复用裁剪证书，用于直接裁剪搜索空间。
\end{theorem}

\begin{certificate}[证书：证书谓词的单调迁移口径]
\label{cert:cert-monotone-transfer}
证书可迁移条件 $\mathrm{Cert}_1(w)=1\Rightarrow \mathrm{Cert}_2(\phi(w))=1$ 等价于其逆否命题 $\mathrm{Cert}_2(\phi(w))=0\Rightarrow
  \mathrm{Cert}_1(w)=0$ 在定义域上的成立；因此失败域在编译映射下可被一致传播。
\end{certificate}

\begin{proof}[证明（谓词因果链）]
由 Certificate~\ref{cert:cert-monotone-transfer}，对任意 $w$ 若 $\mathrm{Cert}_1(w)=0$，若反设 $\mathrm{Cert}_2(\phi(w))=1$，
  则由可迁移条件推出 $\mathrm{Cert}_1(w)=1$，矛盾。故 $\mathrm{Cert}_2(\phi(w))=0$，即 $\phi(w)\in \mathcal{B}_2^{\mathrm{fail}}$。证毕。
\end{proof}

\subsection{定义：GRL 路径积分与自由能}

\begin{definition}[环境分配函数与自由能]
\label{def:env-free-energy}
给定环境 $\mathcal{E}$、策略分布 $\pi$ 与温度参数 $\beta>0$，定义分配函数
\[
\mathcal{Z}_{\mathcal{E}}(\beta;\pi)
:=
\sum_{w\in \mathcal{M}_{\mathrm{tot}}^{\ast}}
\exp\!\bigl(-\beta\,\mathcal{S}_{\mathcal{E}}(w)\bigr)\,\mathbb{P}_{\pi}(w),
\]
并定义自由能（软化的最优可达代价）为
\[
\mathcal{F}_{\mathcal{E}}(\beta;\pi)
:=
-\frac{1}{\beta}\log \mathcal{Z}_{\mathcal{E}}(\beta;\pi).
\]
\end{definition}

\subsection{定理：代价残差界蕴含自由能稳定性界}

\begin{theorem}[自由能稳定性：残差界 $\Rightarrow$ 自由能差界]
\label{thm:free-energy-stability}
设 $\mathcal{E}_1,\mathcal{E}_2$ 之间存在词编译映射 $\phi$，且 $\phi$ 为词空间之间的双射。令 $\pi_1$ 为 $\mathcal{E}_1$ 上的词分布，
并令 $\pi_2:=\phi_{\#}\pi_1$ 为其在 $\mathcal{E}_2$ 上的推前分布，即对任意 $w$ 有 $\mathbb{P}_{\pi_2}(\phi(w))=\mathbb{P}_{\pi_1}(w)$。
若存在 $\varepsilon\ge 0$ 使对任意词 $w$ 有
\[
\bigl|\mathcal{S}_{\mathcal{E}_1}(w)-\mathcal{S}_{\mathcal{E}_2}(\phi(w))\bigr|\le \varepsilon,
\]
则对任意 $\beta>0$ 有自由能差界
\[
\bigl|\mathcal{F}_{\mathcal{E}_1}(\beta;\pi_1)-\mathcal{F}_{\mathcal{E}_2}(\beta;\pi_2)\bigr|\le \varepsilon.
\]
\end{theorem}

\begin{certificate}[证书：分配函数夹逼不等式]
\label{cert:partition-sandwich}
若对任意 $w$ 有 $-\varepsilon\le \mathcal{S}_{\mathcal{E}_1}(w)-\mathcal{S}_{\mathcal{E}_2}(\phi(w))\le \varepsilon$，
则对任意 $\beta>0$ 有逐项夹逼：
\[
e^{-\beta\varepsilon}\exp\!\bigl(-\beta\,\mathcal{S}_{\mathcal{E}_2}(\phi(w))\bigr)
\le
\exp\!\bigl(-\beta\,\mathcal{S}_{\mathcal{E}_1}(w)\bigr)
\le
e^{\beta\varepsilon}\exp\!\bigl(-\beta\,\mathcal{S}_{\mathcal{E}_2}(\phi(w))\bigr),
\]
并在与推前分布配对求和后得到
\[
e^{-\beta\varepsilon}\mathcal{Z}_{\mathcal{E}_2}(\beta;\pi_2)
\le
\mathcal{Z}_{\mathcal{E}_1}(\beta;\pi_1)
\le
e^{\beta\varepsilon}\mathcal{Z}_{\mathcal{E}_2}(\beta;\pi_2).
\]
\end{certificate}

\begin{proof}[证明（谓词因果链）]
由 Certificate~\ref{cert:partition-sandwich} 得
\[
e^{-\beta\varepsilon}\le \frac{\mathcal{Z}_{\mathcal{E}_1}(\beta;\pi_1)}{\mathcal{Z}_{\mathcal{E}_2}(\beta;\pi_2)}\le
  e^{\beta\varepsilon}.
\]
取对数并乘以 $-1/\beta$，得到
\[
-\varepsilon\le \mathcal{F}_{\mathcal{E}_1}(\beta;\pi_1)-\mathcal{F}_{\mathcal{E}_2}(\beta;\pi_2)\le \varepsilon,
\]
即结论。证毕。
\end{proof}

\subsection{推论：最优路径族的权重扰动界与策略复用}

\begin{corollary}[路径族稳定性：低代价路径权重的有界重加权]
\label{cor:path-family-stability}
在 Theorem~\ref{thm:free-energy-stability} 的条件下，定义环境诱导的路径分布
\[
\rho_{\mathcal{E}}(w)
\,\propto\,
\exp\!\bigl(-\beta\,\mathcal{S}_{\mathcal{E}}(w)\bigr)\,\mathbb{P}_{\pi}(w).
\]
则对任意词 $w$，$\rho_{\mathcal{E}_1}(w)$ 与 $\rho_{\mathcal{E}_2}(\phi(w))$ 之间的相对权重差异受 $e^{\beta\varepsilon}$ 控制，
从而“低代价路径族”的分布结构在残差界下稳定；据此可将 $\mathcal{E}_1$ 上学得的策略 $\pi_1$ 通过推前分布初始化 $\pi_2$，以实现策略复用。
\end{corollary}

\begin{certificate}[证书：Boltzmann 权重比的上界]
\label{cert:boltzmann-ratio}
若对任意 $w$ 有 $|\mathcal{S}_{\mathcal{E}_1}(w)-\mathcal{S}_{\mathcal{E}_2}(\phi(w))|\le \varepsilon$，则
\[
e^{-\beta\varepsilon}
\le
\frac{\exp(-\beta\mathcal{S}_{\mathcal{E}_1}(w))}{\exp(-\beta\mathcal{S}_{\mathcal{E}_2}(\phi(w)))}
\le
e^{\beta\varepsilon}.
\]
归一化常数的差异由 Theorem~\ref{thm:free-energy-stability} 控制，因此整体路径分布的重加权因子受同一指数界控制。
\end{certificate}

\begin{proof}[证明（谓词因果链）]
由 Certificate~\ref{cert:boltzmann-ratio} 得到逐项权重比的指数界；再由 Theorem~\ref{thm:free-energy-stability} 得到归一化常数（分配函数）比的指数界。
两者合并即得路径分布的相对权重扰动界；推前分布 $\pi_2=\phi_{\#}\pi_1$ 作为初始化由定义直接给出。证毕。
\end{proof}

\subsection{定义：残差类与不确定度的一等化管理}

\begin{definition}[残差 $\mathrm{Res}(w)$ 与代价的一等化并入]
\label{def:residual-class}
在同伦适配器存在但代价兼容残差无法忽略的情形下，可定义残差度量
\[
\mathrm{Res}(w):=\bigl|\mathcal{S}_{\mathcal{E}_1}(w)-\mathcal{S}_{\mathcal{E}_2}(\phi(w))\bigr|,
\]
并将其作为一等项并入代价：
\[
\mathcal{S}(w)\leftarrow \mathcal{S}(w)+\eta\,\mathrm{Unc}(w)+\kappa\,\mathrm{Res}(w),
\qquad \eta,\kappa>0.
\]
\end{definition}

\begin{remark}[残差管理的作用边界]
残差的一等化管理用于显式标注“同伦复用失效或不稳定”的区域：当 $\mathrm{Res}(w)$ 在某些词上显著增大时，
表示编译适配器与前向评估器在该区域出现结构性失配；将其纳入代价可促使优化过程把算力集中到“残差大/不确定度大”的区域，
以实现适配器与环境定义的迭代修复（例如重新定义子幺半群、补齐证书字段、或提升评估器保真度）。
\end{remark}

\subsection{推论：复用能力的三元输出（策略/证书/评估器）}

\begin{corollary}[复用三元组：策略复用、证书复用与多保真评估器复用]
\label{cor:triple-reuse}
若存在同伦适配器满足 Definition~\ref{def:compatibility-three}，则可得到三类可复用能力：
\begin{enumerate}
  \item \textbf{策略复用：}由 Corollary~\ref{cor:path-family-stability}，可用 $\pi_2=\phi_{\#}\pi_1$ 作为新环境的初始化或约束生成器；
  \item \textbf{证书复用：}由 Theorem~\ref{thm:cert-reuse}，可将已证伪区域与已验证边界作为可调用证书在新环境中裁剪搜索空间；
  \item \textbf{评估器复用：}在同伦残差较小的区域，可共享粗粒度评估器（fast tier）；仅对残差大或不确定度大的区域调用更高保真评估器（hi tier），从而将算力集中在“结构未对齐”的部分。
\end{enumerate}
\end{corollary}

\begin{certificate}[证书：复用三元组的依赖链]
\label{cert:triple-reuse-deps}
（1）由 Theorem~\ref{thm:free-energy-stability} 与 Corollary~\ref{cor:path-family-stability}；
（2）由 Theorem~\ref{thm:cert-reuse}；
（3）由 Definition~\ref{def:residual-class} 的残差度量与多保真评估器的分层实现口径。
\end{certificate}

\begin{proof}[证明（谓词因果链）]
由 Certificate~\ref{cert:triple-reuse-deps}：
\begin{itemize}
  \item（1）路径族稳定性给出策略推前初始化的可解释性；
  \item（2）证书失败域的编译传播给出可复用裁剪；
  \item（3）残差度量给出“何处可共享 fast tier、何处必须调用 hi tier”的可计算判据。
\end{itemize}
三者分别推出推论三项结论，证毕。
\end{proof}

\clearpage
%%%%%%%%%%%%%%%%%%%%%%%%%%%%%%%%%%%%%%%%%%%%%%%%%%%%%%%%%%%%
\section{控制一等化与路线坍缩幻觉：联合共设计与导航证书}
\label{sec:control-route-hallucination}
%%%%%%%%%%%%%%%%%%%%%%%%%%%%%%%%%%%%%%%%%%%%%%%%%%%%%%%%%%%%

\subsection{定义：将控制算子提升为一等生成元}

\begin{definition}[控制幺半群与控制一等化的总算子幺半群]
\label{def:ctrl-first-class}
在 \PTOPB$_{\mathrm{SC}}$ 的三域分解（Definition~\ref{def:ptopb-sc}）之外，引入控制/外场/非平衡驱动的算子幺半群
\[
\mathcal{M}_{\mathrm{ctrl}},
\]
其元素为可执行的闭环/时变控制词（例如时间序列外场、反馈律、泵浦协议、界面调制序列等）。
定义控制一等化后的总算子幺半群为
\[
\mathcal{M}_{\mathrm{tot}}^{\mathrm{ctrl}}
:=
\mathcal{M}_{\mathrm{mat}}
\rtimes
\mathcal{M}_{\mathrm{ctrl}}
\rtimes
\mathcal{M}_{\mathrm{env}}
\rtimes
\mathcal{M}_{\mathrm{th}},
\]
并以 $\mathcal{M}_{\mathrm{tot}}^{\mathrm{ctrl}\ast}$ 表示其词空间。
\end{definition}

\begin{definition}[域投影与联合决定口径]
\label{def:domain-projection}
对任意词 $w\in \mathcal{M}_{\mathrm{tot}}^{\mathrm{ctrl}\ast}$，定义域投影 $w_{\mathrm{mat}},w_{\mathrm{ctrl}},w_{\mathrm{env}},
  w_{\mathrm{th}}$ 为：在词 $w$ 中仅保留属于相应域的生成元（丢弃其它域生成元）所得到的子词。
在该口径下，前向输出（Definition~\ref{def:forward-map-y}）写为
\[
y(w)=y\bigl(w_{\mathrm{mat}},w_{\mathrm{ctrl}},w_{\mathrm{env}},w_{\mathrm{th}}\bigr),
\]
特别地，$T_c$ 等关键指标由材料与控制共同决定：
\[
T_c=T_c\bigl(w_{\mathrm{mat}},w_{\mathrm{ctrl}},w_{\mathrm{env}},w_{\mathrm{th}}\bigr).
\]
\end{definition}

\begin{remark}[可达集合的投影稀薄性]
\label{rem:reachable-projection-sparse}
令可达目标集合（Definition~\ref{def:target-set}）为
\[
\mathcal{W}^\star:=\mathcal{W}_{\mathrm{goal}}(y^\star,\mathcal{C}).
\]
“材料维度的投影稀薄/近似为空”可表述为：在仅考虑材料子词的投影下，
\[
\pi_{\mathrm{mat}}(\mathcal{W}^\star):=\{w_{\mathrm{mat}}\mid \exists w\in\mathcal{W}^\star\},
\]
该集合在候选材料词空间中占比极小；而在联合空间 $(w_{\mathrm{mat}},w_{\mathrm{ctrl}})$ 中，可达集合可能非空且连通。
该结构对应“目标由材料与控制共同约束”的情形：材料投影可能稀薄，但联合投影可形成可搜索的连通块。
\end{remark}

\subsection{定理（公理降级）：控制一等化的表示完备性（按词长归纳）}

\begin{theorem}[表示完备性（公理降级）：四域总词空间覆盖任意有限步组合操作]
\label{thm:ctrl-first-class-completeness}
若 $\mathcal{M}_{\mathrm{mat}}$、$\mathcal{M}_{\mathrm{ctrl}}$、$\mathcal{M}_{\mathrm{env}}$、$\mathcal{M}_{\mathrm{th}}$
  的生成元均可执行，且连接运算对应可执行序列的连接，则任何有限步的“材料+控制+环境+理论”组合操作均可表示为某个词 $w\in\mathcal{M}_{\mathrm{tot}}^{\mathrm{ctrl}\ast}$。
\end{theorem}

\begin{certificate}[数学归纳法证书：按词长 $T$ 的总词表示归纳]
\label{cert:ctrl-first-class-completeness-induction}
定义谓词 $P(T)$ 为：“任意由 $T$ 步域生成元组成的可执行序列，可表示为长度为 $T$ 的词 $w\in\mathcal{M}_{\mathrm{tot}}^{\mathrm{ctrl}\ast}$”。证书化要求为：
\[
P(1)\ \text{成立},\qquad
(\forall T\in\mathbb{N}_{\ge 1})\ \bigl(P(T)\Rightarrow P(T+1)\bigr),
\]
其中递推步使用幺半群封闭性：将长度 $T$ 的词与第 $T{+}1$ 步生成元连接得到长度 $T{+}1$ 的词。
\end{certificate}

\begin{proof}[证明（数学归纳法证书；谓词因果链）]
由 Certificate~\ref{cert:ctrl-first-class-completeness-induction}：
\begin{itemize}
  \item 基步：$P(1)$ 成立（一步操作即为一个生成元，亦即长度为 1 的词）；
  \item 归纳步：若 $P(T)$ 成立，则任意 $T{+}1$ 步操作可写为 $w a_{T+1}$，由封闭性得 $w a_{T+1}\in\mathcal{M}_{\mathrm{tot}}^{\mathrm{ctrl}\ast}
    $，从而 $P(T{+}1)$ 成立。
\end{itemize}
由数学归纳法得结论。证毕。
\end{proof}

\subsection{定义：路线索引与“路线坍缩”现象}

\begin{definition}[路线索引 $k$ 与策略坍缩]
\label{def:route-collapse}
记路线/机制索引为离散变量 $k\in\mathcal{K}$（可理解为理论链、稳定化假设、近似层级与评估器族的选择）。
在 GRL/策略优化中，若联合策略分布 $\pi(w,k)$ 满足
\[
\pi(k)\to \delta_{k_0},
\qquad
\pi(w\mid k_0)\ \text{在词空间上趋于近确定性},
\]
则称发生\textbf{路线坍缩}：策略在早期就锁定到单一路线 $k_0$，并在该子空间内强收敛。
\end{definition}

\subsection{命题：路线坍缩幻觉（模型自洽但现实不可达）}

\begin{proposition}[路线坍缩幻觉：模型代价小不等价于现实代价小]
\label{prop:route-hallucination}
设对每条路线 $k$ 存在前向评估器 $\widehat{F}_k$ 并诱导模型代价 $\mathcal{S}_k(w)$；设真实系统对应的代价为 $\mathcal{S}_{\mathrm{real}}(w)$。
若存在 $k_0$ 与词 $w$ 使
\[
\mathcal{S}_{k_0}(w)\ \text{很小},\qquad
\mathcal{S}_{\mathrm{real}}(w)\ \text{不小},
\]
且在优化过程中未将不确定度/残差/证书门槛作为一等项并入代价。
（参见 Definition~\ref{def:cost-functional}、Definition~\ref{def:residual-class} 与 Proposition~\ref{prop:cert-gating} 的结构。）
则路线坍缩（Definition~\ref{def:route-collapse}）可能导致策略对 $k_0$ 的置信度迅速上升并在 $k_0$ 的模型子空间内自洽化收敛，从而产生“模型上越来越好、现实上并未接近目标”的路线坍缩幻觉。
\end{proposition}

\begin{certificate}[证书：幻觉成立的最小条件（失配 + 缺少校准项 + 坍缩）]
\label{cert:route-hallucination}
本命题使用三条证书化前提：
\begin{enumerate}
  \item \textbf{失配存在：}存在 $w$ 使 $\mathcal{S}_{k_0}(w)$ 与 $\mathcal{S}_{\mathrm{real}}(w)$ 显著不一致；
  \item \textbf{缺少校准项：}优化目标未显式包含 $\mathrm{Unc}(w)$、$\mathrm{Res}(w)$ 与 $\mathrm{Cert}(w)$ 的强约束/强惩罚；
  \item \textbf{策略坍缩：}发生 Definition~\ref{def:route-collapse} 所述坍缩。
\end{enumerate}
\end{certificate}

\begin{proof}[证明（谓词因果链）]
由 Certificate~\ref{cert:route-hallucination}，取谓词链：
\[
A:\ \exists w\ \bigl(\mathcal{S}_{k_0}(w)\ll 1 \wedge \mathcal{S}_{\mathrm{real}}(w)\not\ll 1\bigr),\quad
B:\ \text{缺少校准项},\quad
C:\ \pi(k)\to\delta_{k_0}.
\]
在 $A$ 成立且 $B$ 成立时，模型优化可以在 $k_0$ 的评估器漏洞/系统性偏差方向上持续降低 $\mathcal{S}_{k_0}$ 而不受惩罚；
在 $C$ 成立时，策略无法通过再探索回到其它路线或通过控制一等化扩展可行域，因此“模型自洽收敛”与“现实不可达”可并存。
故得到命题结论。证毕。
\end{proof}

\subsection{定理：防止路线幻觉的可计算导航机制（熵底座 + 校准项）}

\begin{theorem}[导航定理：在路线维度保留熵底座并加入校准项可抑制坍缩幻觉]
\label{thm:navigation-against-hallucination}
设联合分配函数扩展为同时对路线与词求和：
\[
\mathcal{Z}(\beta;\pi)
:=
\sum_{k\in\mathcal{K}}
\sum_{w\in\mathcal{M}_{\mathrm{tot}}^{\mathrm{ctrl}\ast}}
\exp\!\bigl(-\beta\,\mathcal{S}_k(w)\bigr)\,\pi(w,k),
\]
并对路线边缘分布施加熵底座约束 $H(\pi(k))\ge h_{\min}>0$，同时将 $\mathrm{Unc}(w)$、$\mathrm{Res}(w)$ 与 $\mathrm{Cert}(w)$
  以一等项并入代价（Definition~\ref{def:residual-class}、Proposition~\ref{prop:cert-gating}）。
则在该约束下，策略在早期难以发生不可逆的单一路线坍缩，且模型失配造成的“虚假低代价路径”将更容易在 $\mathrm{Unc}/\mathrm{Res}/\mathrm{Cert}$ 项上被识别并被路径积分权重压制，从而抑制路线坍缩幻觉。
\end{theorem}

\begin{certificate}[证书：熵底座与校准项对坍缩的抑制口径]
\label{cert:entropy-calibration}
（i）熵底座 $H(\pi(k))\ge h_{\min}$ 防止 $\pi(k)$ 在早期坍缩为 $\delta_{k_0}$；
（ii）将 $\mathrm{Unc}/\mathrm{Res}/\mathrm{Cert}$ 并入代价会使“仅在单一路线评估器中自洽”的候选在跨评估器一致性与可复现约束上付出代价，从而在路径积分权重下难以主导概率质量。
\end{certificate}

\begin{proof}[证明（谓词因果链）]
由 Certificate~\ref{cert:entropy-calibration}：
\begin{itemize}
  \item（熵底座）$H(\pi(k))\ge h_{\min}$ 排除 $\pi(k)$ 过早趋于点质量，保持跨路线探索；
  \item（校准项）若某词 $w$ 仅在单一路线评估器下代价偏低，则其跨评估器残差/不确定度或证书门槛更可能暴露不一致，从而提升有效代价并降低权重 $\exp(-\beta\mathcal{S}_k(w))$。
\end{itemize}
二者合取得到“保持导航、抑制幻觉”的结论。证毕。
\end{proof}

\subsection{定义：材料--控制联合共设计（双层可计算形式）}

\begin{definition}[联合共设计：材料词与控制词的双层最优化]
\label{def:bilevel-codesign}
将总词分为材料词 $w_m$ 与控制词 $w_c$（通过域投影得到），定义双层优化形式：
\[
w_c^\ast(w_m):=\arg\min_{w_c}\ \mathcal{S}(w_m,w_c),
\qquad
w_m^\ast:=\arg\min_{w_m}\ \mathcal{S}\bigl(w_m,w_c^\ast(w_m)\bigr).
\]
其中 $\mathcal{S}$ 可取 $\mathcal{S}_{\mathrm{cert}}$ 或包含残差校准项的扩展代价。
\end{definition}

\begin{proposition}[联合共设计的保守性：避免过早淘汰“可被控制补齐”的材料候选]
\label{prop:codesign-benefit}
在 Definition~\ref{def:bilevel-codesign} 的双层形式下，外层对材料候选的比较基于其在最优控制词下的可达上限（最小代价）。
因此当目标可达集合在材料投影下稀薄（Remark~\ref{rem:reachable-projection-sparse}）时，该形式可避免因“控制未充分探索”而过早淘汰材料候选，从而提高命中“材料+控制窄组合”的概率。
\end{proposition}

\begin{certificate}[证书：以 $w_c^\ast(w_m)$ 表示材料候选的可达上限]
\label{cert:codesign-upper-envelope}
对任意材料候选 $w_m$，量 $J(w_m):=\min_{w_c}\mathcal{S}(w_m,w_c)$ 给出在允许控制共设计时该材料的最优可达代价（上限包络）。以 $J(w_m)$
  比较材料可避免把控制不足导致的劣化误判为材料本质不可达。
\end{certificate}

\begin{proof}[证明（谓词因果链）]
由 Certificate~\ref{cert:codesign-upper-envelope}，外层比较的是 $J(w_m)$ 而不是固定控制下的代价。
因此当某材料需要特定控制词才能进入可达域时，双层形式允许控制在内层补齐；从而材料不会因控制缺失而被误淘汰。证毕。
\end{proof}

\subsection{推论：路线幻觉判别证书与触发机制（可计算口径）}

\begin{corollary}[判别证书：代价下降主要来自单一路线且残差不降 \& 证书难以满足]
\label{cor:hallucination-certificate}
设优化迭代产生序列 $\{(\pi_t,\mathcal{S}_t)\}$。若在一段窗口内满足：
\begin{enumerate}
  \item 路线边缘分布快速坍缩：$\pi_t(k)$ 近似集中在单一路线 $k_0$；
  \item 模型代价持续下降，但残差 $\mathrm{Res}(w)$ 的统计量不降反升（或维持高位），且证书谓词 $\mathrm{Cert}(w)$ 难以达成；
\end{enumerate}
则可在证书口径下判定发生“路线坍缩幻觉”风险，并触发强制的再探索机制：提高路线熵底座、扩展 $\mathcal{M}_{\mathrm{ctrl}}$ 的动作字典或提升评估器保真度。
\end{corollary}

\begin{certificate}[证书：触发条件的可观测统计量]
\label{cert:hallucination-trigger}
触发条件可由可观测统计量实现：路线熵 $H(\pi_t(k))$、残差均值/分位数 $\mathbb{E}[\mathrm{Res}]$、证书通过率 $\mathbb{P}(\mathrm{Cert}=1)$、以及代价序列
  $\mathbb{E}[\mathcal{S}]$ 的下降趋势。
\end{certificate}

\begin{proof}[证明（谓词因果链）]
由 Certificate~\ref{cert:hallucination-trigger}，上述三类统计量分别对应：
（i）是否坍缩，（ii）是否跨评估器一致，（iii）是否满足可复现门槛。
当（i）成立且（ii）（iii）不改善时，命题 \ref{prop:route-hallucination} 的三条前提在统计意义下被激活，故可判定幻觉风险并触发 Theorem~\ref{thm:navigation-against-hallucination} 所述导航修正机制。证毕。
\end{proof}

\begin{remark}[执行性边界]
上述判别证书不声称“定位真实机制”，其功能是：在优化过程早期阻止不可逆坍缩，并将算力重新分配到“路线不确定/控制未充分/评估器失配”的区域。
在实践中，$\mathcal{M}_{\mathrm{ctrl}}$ 的定义与 $\mathrm{Res}$ 的构造应优先采用可复核的观测量与评估器对照，以确保触发机制本身可被审计与版本化。
\end{remark}

\clearpage
%%%%%%%%%%%%%%%%%%%%%%%%%%%%%%%%%%%%%%%%%%%%%%%%%%%%%%%%%%%%
\section{\LBOPB/\PTOPB{} 的广度动机：可达性诊断优先于类内构造}
\label{sec:breadth-as-diagnosis}
%%%%%%%%%%%%%%%%%%%%%%%%%%%%%%%%%%%%%%%%%%%%%%%%%%%%%%%%%%%%

\subsection{备注：同伦视角下的“诊断环境”与“治疗工具”}

\begin{remark}[同伦等价：先把方向做成可比较诊断，再在类内做构造]
在 \LBOPB/\PTOPB{} 的语言中，可将路线/机制/近似层级索引 $k$ 理解为对“机制子丛/评估器族/理论链”的离散选择；其效果是将可行域与代价结构限制在某个同伦意义下可对比的类内。
因此“方向”本质上是选同伦类（选路线/机制子丛），而“构造”是在该类内对可执行算子词 $w$ 做局部最优。
这解释了为什么 \LBOPB/\PTOPB{} 首先应被构造为\textbf{诊断环境}（把问题放进可比较、可复用、可证书化的虚拟宇宙里），其次才是\textbf{治疗工具}（在选定类内做精细构造与优化）。
\end{remark}

\subsection{定义：证书化动作集、诊断泛函与治疗解}

\begin{definition}[路线索引 $k$、可执行算子词 $w$ 与证书化动作集 $\mathcal{A}_k(\varepsilon)$]
\label{def:certified-action-set}
记路线/机制/近似层级索引为离散变量 $k\in\mathcal{K}$（参见 Definition~\ref{def:route-collapse}；并可视为理论域操作字典的子幺半群选择，参见 Lemma~\ref{lem:k-in-mth}）。
记可执行算子词 $w\in \mathcal{M}_{\mathrm{tot}}^{\mathrm{ctrl}\ast}$（忽略控制域时可退化为 $w\in\mathcal{M}_{\mathrm{tot}}^{\ast}$）。
对每条路线 $k$ 给定：
\[
\mathcal{S}_k:\mathcal{M}_{\mathrm{tot}}^{\mathrm{ctrl}\ast}\to\mathbb{R}_{\ge 0},\qquad
\mathrm{Cert}^{(k)}:\mathcal{M}_{\mathrm{tot}}^{\mathrm{ctrl}\ast}\to\{0,1\},\qquad
\mathrm{Res}^{(k)}:\mathcal{M}_{\mathrm{tot}}^{\mathrm{ctrl}\ast}\to\mathbb{R}_{\ge 0}.
\]
给定容差 $\varepsilon\ge 0$，定义路线 $k$ 下可被证书化/可复现/残差可控的动作集合（不是全空间）为
\[
\mathcal{A}_k(\varepsilon)
:=
\{\,w:\ \mathrm{Cert}^{(k)}(w)=1,\ \mathrm{Res}^{(k)}(w)\le \varepsilon\,\}.
\]
\end{definition}

\begin{definition}[诊断（方向选择）与治疗（类内构造）]
\label{def:diagnosis-and-treatment}
定义诊断泛函
\[
\mathcal{D}(k;\varepsilon)
:=
\inf_{w\in \mathcal{A}_k(\varepsilon)}\ \mathcal{S}_k(w),
\]
并将固定路线 $k$ 的类内最优化（治疗）定义为
\[
w^\ast(k;\varepsilon)\in \arg\min_{w\in \mathcal{A}_k(\varepsilon)}\ \mathcal{S}_k(w).
\]
若 $\mathcal{A}_k(\varepsilon)=\emptyset$，可约定 $\mathcal{D}(k;\varepsilon)=+\infty$，表示“在该路线下不可证书化地可达”。
\end{definition}

\begin{definition}[有限词长诊断下界]
\label{def:truncated-diagnosis}
记词长为 $|w|$。对 $T\in\mathbb{N}_{\ge 1}$ 定义截断动作集与对应诊断下界
\[
\mathcal{A}_{k,T}(\varepsilon):=\{w\in\mathcal{A}_k(\varepsilon)\mid |w|\le T\},
\qquad
\mathcal{D}_T(k;\varepsilon):=\inf_{w\in \mathcal{A}_{k,T}(\varepsilon)} \mathcal{S}_k(w).
\]
\end{definition}

\subsection{定理（公理降级）：把“方向选择”压成可计算的等式链}

\begin{theorem}[可计算等式链（公理降级）：诊断泛函由有限词长下界逼近]
\label{thm:computable-diagnosis-chain}
对任意路线 $k$ 与容差 $\varepsilon\ge 0$，序列 $\{\mathcal{D}_T(k;\varepsilon)\}_{T\ge 1}$ 单调不增，且满足
\[
\mathcal{D}(k;\varepsilon)
=
\inf_{T\in\mathbb{N}_{\ge 1}}\mathcal{D}_T(k;\varepsilon)
=
\inf_{T\in\mathbb{N}_{\ge 1}}\ \inf_{w\in \mathcal{A}_{k,T}(\varepsilon)} \mathcal{S}_k(w).
\]
因此“诊断”可被实现为：在增加词长预算 $T$ 的同时，对每个 $k$ 持续更新其可达下界估计并进行跨 $k$ 的排序与比较。
\end{theorem}

\begin{certificate}[数学归纳法证书：按词长 $T$ 的下界单调性]
\label{cert:computable-diagnosis-induction}
定义谓词 $P(T)$ 为：
\[
P(T):\ \mathcal{A}_{k,T}(\varepsilon)\subseteq \mathcal{A}_{k,T+1}(\varepsilon)\ \Rightarrow\ \mathcal{D}_{T+1}(k;
  \varepsilon)\le \mathcal{D}_T(k;\varepsilon).
\]
证书化要求为：
\[
P(1)\ \text{成立},\qquad
(\forall T\in\mathbb{N}_{\ge 1})\ \bigl(P(T)\Rightarrow P(T+1)\bigr).
\]
\end{certificate}

\begin{proof}[证明（数学归纳法证书；谓词因果链）]
由 Definition~\ref{def:truncated-diagnosis}，对任意 $T$ 有蕴含
\[
|w|\le T\ \Rightarrow\ |w|\le T+1,
\]
从而得到集合包含 $\mathcal{A}_{k,T}(\varepsilon)\subseteq \mathcal{A}_{k,T+1}(\varepsilon)$，故对任意 $w\in \mathcal{A}_{k,T}
  (\varepsilon)$ 有
\[
\mathcal{D}_{T+1}(k;\varepsilon)\le \mathcal{S}_k(w).
\]
取 $w$ 的下确界即得 $\mathcal{D}_{T+1}(k;\varepsilon)\le \mathcal{D}_T(k;\varepsilon)$，从而由 Certificate~\ref{cert:computable-diagnosis-induction} 可视作按词长 $T$ 的归纳单调性证书成立。
进一步地，任取 $w\in\mathcal{A}_k(\varepsilon)$，令 $T:=|w|$ 则 $w\in \mathcal{A}_{k,T}(\varepsilon)$，故
\[
\inf_{T}\mathcal{D}_T(k;\varepsilon)\le \mathcal{S}_k(w).
\]
对 $w\in\mathcal{A}_k(\varepsilon)$ 取下确界得 $\inf_T\mathcal{D}_T(k;\varepsilon)\le \mathcal{D}(k;\varepsilon)$；反向不等式由
  $\mathcal{A}_{k,T}(\varepsilon)\subseteq \mathcal{A}_k(\varepsilon)$ 直接给出 $\mathcal{D}(k;\varepsilon)\le \mathcal{D}
  _T(k;\varepsilon)$，再对 $T$ 取下确界即得 $\mathcal{D}(k;\varepsilon)\le \inf_T\mathcal{D}_T(k;\varepsilon)$。
两侧合并得到定理结论。证毕。
\end{proof}

\subsection{定理：诊断优先（方向优先）比类内构造更决定成败}

\begin{theorem}[诊断优先：全局最优由 $\mathcal{D}(k)$ 的比较决定]
\label{thm:diagnosis-priority}
给定候选路线集合 $\mathcal{K}_0\subseteq\mathcal{K}$，定义
\[
\mathcal{D}^\star(\mathcal{K}_0;\varepsilon):=\inf_{k\in\mathcal{K}_0}\mathcal{D}(k;\varepsilon).
\]
则在“先选路线再在类内优化”的范式下：
\[
\inf_{k\in\mathcal{K}_0}\ \inf_{w\in\mathcal{A}_k(\varepsilon)} \mathcal{S}_k(w)
\,=\,
\mathcal{D}^\star(\mathcal{K}_0;\varepsilon),
\]
而任意固定路线 $k$ 的最优治疗值为 $\mathcal{D}(k;\varepsilon)$。
因此\textbf{诊断阶段对 $k$ 的判别与比较}（排序 $\mathcal{D}(k)$ 并识别主导障碍）比在固定 $k$ 内把 $w$ 做到极致更决定能否逼近目标。
\end{theorem}

\begin{certificate}[证书：全局最小值的嵌套下确界展开]
\label{cert:diagnosis-priority}
由 Definition~\ref{def:diagnosis-and-treatment}，对每个 $k$ 有
\[
\mathcal{D}(k;\varepsilon)=\inf_{w\in\mathcal{A}_k(\varepsilon)}\mathcal{S}_k(w).
\]
对候选集 $\mathcal{K}_0$ 取下确界即得
\[
\inf_{k\in\mathcal{K}_0}\mathcal{D}(k;\varepsilon)=\inf_{k\in\mathcal{K}_0}\ \inf_{w\in\mathcal{A}_k(\varepsilon)}
  \mathcal{S}_k(w).
\]
\end{certificate}

\begin{proof}[证明（谓词因果链）]
由 Certificate~\ref{cert:diagnosis-priority}，记目标谓词为
\[
Q:\ \inf_{k\in\mathcal{K}_0}\ \inf_{w\in\mathcal{A}_k(\varepsilon)}\mathcal{S}_k(w)=\mathcal{D}^\star(\mathcal{K}_0;
  \varepsilon).
\]
由定义 $\mathcal{D}^\star(\mathcal{K}_0;\varepsilon):=\inf_{k\in\mathcal{K}_0}\mathcal{D}(k;\varepsilon)$ 与 $\mathcal{D}(k;
  \varepsilon)=\inf_{w\in\mathcal{A}_k(\varepsilon)}\mathcal{S}_k(w)$，逐步代入即可得到 $Q$ 成立。
同时，对固定 $k$ 的最优治疗值恰为 $\inf_{w\in\mathcal{A}_k(\varepsilon)}\mathcal{S}_k(w)=\mathcal{D}(k;\varepsilon)$。
因此若选择了不良路线 $k$，其类内优化再“极致”也无法跨越 $\mathcal{D}(k;\varepsilon)$ 的下界；而跨 $k$ 的诊断排序直接决定能否进入更低下界的同伦类。证毕。
\end{proof}

\subsection{引理：选错同伦类将导出近似不可达下界}

\begin{lemma}[类内下界：$\mathcal{D}(k;\varepsilon)>0 \Rightarrow$ 类内近似不可达]
\label{lem:wrong-class}
若在某路线 $k$ 下有 $\mathcal{D}(k;\varepsilon)>0$，则对任意 $w\in\mathcal{A}_k(\varepsilon)$ 都有
\[
\mathcal{S}_k(w)\ge \mathcal{D}(k;\varepsilon)>0.
\]
特别地，当 $\mathcal{S}_k(w)=0$ 被用作“达成目标”的判据时，$\mathcal{D}(k;\varepsilon)>0$ 意味着目标在该同伦类/机制子丛内近似不可达；继续在固定 $k$
  内做深度构造只会把算力推向一个“看似自洽的类内局部最优”，即路线偏差的幻觉。
\end{lemma}

\begin{certificate}[证书：下确界是全体取值的下界]
\label{cert:inf-lower-bound}
对任意非空集合 $A$ 与函数 $f$，量 $\inf_{x\in A}f(x)$ 是 $f(A)$ 的下界，即
\[
(\forall x\in A)\quad f(x)\ge \inf_{y\in A}f(y).
\]
\end{certificate}

\begin{proof}[证明（谓词因果链）]
由 Certificate~\ref{cert:inf-lower-bound}，取 $A:=\mathcal{A}_k(\varepsilon)$ 与 $f:=\mathcal{S}_k$，则对任意 $w\in\mathcal{A}
  _k(\varepsilon)$ 有
\[
\mathcal{S}_k(w)\ge \inf_{u\in\mathcal{A}_k(\varepsilon)}\mathcal{S}_k(u)=\mathcal{D}(k;\varepsilon).
\]
若 $\mathcal{D}(k;\varepsilon)>0$，则上式立即给出 $\mathcal{S}_k(w)>0$，从而得到引理结论。证毕。
\end{proof}

\subsection{命题：广度的单调性与“建模宇宙排除真实解”风险}

\begin{proposition}[广度单调性：扩大候选路线集合不会更差]
\label{prop:breadth-monotonicity}
令 $\mathcal{K}_1\subseteq \mathcal{K}_2 \subseteq \mathcal{K}$，则
\[
\mathcal{D}^\star(\mathcal{K}_2;\varepsilon) \le \mathcal{D}^\star(\mathcal{K}_1;\varepsilon).
\]
若存在 $k^\star\in\mathcal{K}$ 使 $\mathcal{D}(k^\star;\varepsilon)=0$（或足够小），但 $k^\star\notin\mathcal{K}_1$，则可能出现
  $\mathcal{D}^\star(\mathcal{K}_1;\varepsilon)>0$：
即在诊断阶段就把正确同伦类/机制子丛排除在候选空间之外——此后再强的类内优化也无法“救回”被排除的路线。
\end{proposition}

\begin{certificate}[证书：下确界对候选集的集合单调性]
\label{cert:inf-set-monotone}
若 $B_1\subseteq B_2$，则对任意函数 $g$ 有
\[
\inf_{x\in B_2}g(x)\le \inf_{x\in B_1}g(x).
\]
\end{certificate}

\begin{proof}[证明（谓词因果链）]
取 $g(k):=\mathcal{D}(k;\varepsilon)$，由 Certificate~\ref{cert:inf-set-monotone} 与 $\mathcal{K}_1\subseteq\mathcal{K}_2$ 得
\[
\inf_{k\in\mathcal{K}_2}\mathcal{D}(k;\varepsilon)\le \inf_{k\in\mathcal{K}_1}\mathcal{D}(k;\varepsilon),
\]
即 $\mathcal{D}^\star(\mathcal{K}_2;\varepsilon)\le \mathcal{D}^\star(\mathcal{K}_1;\varepsilon)$。
若存在 $k^\star\notin\mathcal{K}_1$ 且 $\mathcal{D}(k^\star;\varepsilon)=0$，则 $\inf_{k\in\mathcal{K}_1}\mathcal{D}(k;
  \varepsilon)$ 不会“看到”该零值项，因而可能严格大于 0。
结合 Lemma~\ref{lem:wrong-class}，这意味着后续在 $\mathcal{K}_1$ 内做任意类内优化都无法达到 0 代价基准。证毕。
\end{proof}

\subsection{推论：算力分配原则与“幻觉预警”触发口径}

\begin{corollary}[执行原则：先做广度诊断，再做深度治疗]
\label{cor:diagnosis-before-treatment}
在预算受限的逆向设计任务（尤其是 \RTSC{} 这类“目标极窄、模型易自洽化”的任务）中，应优先将算力用于：
\begin{enumerate}
  \item \textbf{诊断：}扩大候选路线集合 $\mathcal{K}_0$ 与动作字典（含 $\mathcal{M}_{\mathrm{ctrl}}$ 的控制协议覆盖），并估计/排序 $\mathcal{D}_T(k;
    \varepsilon)$，强化 $\mathrm{Res}/\mathrm{Cert}$ 的交叉检验，使 $\mathcal{D}(k;\varepsilon)$ 的比较可信；
  \item \textbf{治疗：}仅对少数被诊断为“可达或近可达”的路线 $k$ 做深度类内优化以逼近 $w^\ast(k;\varepsilon)$。
\end{enumerate}
并且：若在固定路线 $k$ 内出现“$\mathcal{S}_k(w)$ 持续下降，但 $\mathrm{Res}^{(k)}(w)$ 不降反升或 $\mathrm{Cert}^{(k)}(w)$ 长期难以满足”，则应立即视为方向偏航的预警，
  触发对 $\mathcal{K}_0$ 与动作字典的扩展与再探索（参见 Corollary~\ref{cor:hallucination-certificate}）。
\end{corollary}

\begin{certificate}[证书：诊断优先的可计算收益来自下界排序与集合扩张单调性]
\label{cert:diagnosis-before-treatment}
（1）由 Theorem~\ref{thm:computable-diagnosis-chain}，有限词长下界 $\mathcal{D}_T(k;\varepsilon)$ 可随 $T$ 递增逼近诊断值，并用于跨 $k$ 排序；
（2）由 Proposition~\ref{prop:breadth-monotonicity}，扩大 $\mathcal{K}_0$ 不会使 $\mathcal{D}^\star$ 变差，反而可避免排除真实路线；
（3）由 Corollary~\ref{cor:hallucination-certificate}，在固定 $k$ 内代价下降但残差/证书不改善时存在路线坍缩幻觉风险，应触发再诊断。
\end{certificate}

\begin{proof}[证明（谓词因果链）]
由 Certificate~\ref{cert:diagnosis-before-treatment}：
\begin{itemize}
  \item（1）给出“先诊断”的可计算实现：用 $\mathcal{D}_T$ 的序列下界对路线排序；
  \item（2）给出“做广度”的形式收益：扩大候选集只会降低或保持全局可达下界，且避免把零值路线排除在外；
  \item（3）给出“幻觉预警”的触发口径：当类内优化仅降低模型代价而不改善残差/证书时，应回到诊断层面扩张候选。
\end{itemize}
三者合取即可推出推论的执行原则与触发机制。证毕。
\end{proof}

\subsection{备注：三类广度收益（候选空间/诊断语言/结构障碍可观测性）}

\begin{remark}[广度 $\approx$ 可复用的同伦分类器]
将 \LBOPB/\PTOPB{} 视为虚拟环境对象 $\mathcal{E}$（Definition~\ref{def:env-object}），其“广度”带来的核心收益可整理为三类：
\begin{enumerate}
  \item \textbf{候选空间足够大：}让路线索引 $k$ 的候选集合覆盖材料、控制协议、非平衡驱动与关键理论算子（如非简谐、量子核效应等），避免在诊断阶段把正确同伦类从建模宇宙中删除；
  \item \textbf{诊断语言可跨域复用：}将 $\mathrm{Cert}(w)$、$\mathrm{Res}(w)$、$\mathrm{Unc}(w)$ 做成统一接口，
    使算力沉淀为“已证伪区域/已验证边界/不确定区轮廓”的可迁移资产，而不是一次性的仿真日志；
  \item \textbf{结构障碍暴露为可观测特征：}当真正阻断可达性的因素来自同伦类不对、障碍类不为零或曲率导致不可交换路径依赖时，宽覆盖的任务族有助于这些结构性信号反复出现、可比对、可证书化，从而减少“在错误地图上置信度趋近
    1”的情况。
\end{enumerate}
因此可将本节主旨压缩为：
\[
\begin{aligned}
\text{\LBOPB/\PTOPB{} 的广度}
&\approx
\text{把可达性诊断做成可复用的同伦分类器}\\
&\Rightarrow
\text{再在正确类内做构造。}
\end{aligned}
\]
\end{remark}

\clearpage
%%%%%%%%%%%%%%%%%%%%%%%%%%%%%%%%%%%%%%%%%%%%%%%%%%%%%%%%%%%%
\section{对 \RTSC{} 尾部风险内参稿的形式增益：同伦复用、诊断优先与反偏航约束}
\label{sec:tex14-rtsc-tail-risk-gains}
%%%%%%%%%%%%%%%%%%%%%%%%%%%%%%%%%%%%%%%%%%%%%%%%%%%%%%%%%%%%

\subsection{备注：从“多补一条路线”到“可迁移诊断资产”}

\begin{remark}[增益主旨：把算力投入转写为可迁移资产]
对“尾部风险 + \RTSC{} + 绝对锚点/\textsf{MetaSol}”这一形式口径而言，本次讨论的增益不在于“多补了一条路线叙事”，而在于：
将 \LBOPB/\PTOPB{} 视为虚拟环境对象（Definition~\ref{def:env-object}）的同伦类，从而把\textbf{算力投入}转写为可迁移的\textbf{诊断边界}、\textbf{证书接口}
  与\textbf{探索分布}资产。
在该口径下，文档中已有的主线（\GZLS / Jacobian-gap（\textsf{JacGap}）/ \textsf{PFB-GNLA} / GRL path-integral）可被推进到更可计算、可复用、
  且更能防止“路线自洽偏航”的层面：算力不再仅产出一次性的仿真日志，而被沉淀为可迁移的“已证伪区域/已验证边界/不确定区轮廓”。
\end{remark}

\subsection{定义：\textsf{MetaSol} 口径下的可复用资产}

\begin{definition}[可复用资产三元组：诊断边界、证书接口与探索分布]
\label{def:reusable-asset-triple}
给定虚拟环境对象 $\mathcal{E}$（Definition~\ref{def:env-object}），其证书谓词记为 $\mathrm{Cert}_{\mathcal{E}}(w)\in\{0,1\}$。
定义证书失败域（可复用诊断边界）为
\[
\mathcal{B}_{\mathcal{E}}^{\mathrm{fail}}
:=
\{w\mid \mathrm{Cert}_{\mathcal{E}}(w)=0\}.
\]
并将三元组
\[
\mathfrak{A}(\mathcal{E})
:=
\bigl(\mathcal{B}_{\mathcal{E}}^{\mathrm{fail}},\ \mathrm{Cert}_{\mathcal{E}},\ \pi_{\mathcal{E}}\bigr)
\]
称为环境 $\mathcal{E}$ 上的可复用资产，其中 $\pi_{\mathcal{E}}$ 为 GRL/path-integral 中的探索分布（Definition~\ref{def:env-free-energy} 的
  $\pi$ 角色）。
\end{definition}

\begin{definition}[同伦编译与资产推前]
\label{def:asset-pushforward}
设 $\mathcal{E}_1,\mathcal{E}_2$ 之间存在同伦适配器/编译映射 $\phi$（Definition~\ref{def:compatibility-three}），并令推前分布 $\pi_2:=\phi_{\#}
  \pi_1$（Theorem~\ref{thm:free-energy-stability} 的口径）。
则称资产 $\mathfrak{A}(\mathcal{E}_1)$ 可\textbf{推前复用}到 $\mathcal{E}_2$：失败域通过 $\phi$ 传播为 $\phi(\mathcal{B}_{\mathcal{E}_1}
  ^{\mathrm{fail}})$，证书接口通过 $\mathrm{Cert}_1\Rightarrow \mathrm{Cert}_2\circ\phi$ 的单调条件迁移，探索分布通过推前分布初始化。
\end{definition}

\begin{certificate}[文本证书：\textsf{MetaSol}/绝对锚点口径的定义位置]
\label{cert:tex14-metasol-source}
在“绝对锚点”与 “MetaSol 的范畴化定义”口径中，统一母体、
  证书接口与统一动力学（path-integral/GRL）共同构成元方案接口。
本章仅抽取其中的“统一接口 + 可复用资产”视角，用于刻画“最优元方案”在算力经济学意义上的增益。
\end{certificate}

\subsection{定理（公理降级）：把“同伦复用”写成可计算的链式结论}

\begin{theorem}[M6（公理降级）：虚拟环境同伦复用条款的链式闭包]
\label{thm:m6-homotopy-reuse-chain}
给定环境链 $\mathcal{E}_1,\mathcal{E}_2,\ldots,\mathcal{E}_{N+1}$ 及对应的词编译双射
\[
\phi_i:\mathcal{M}_{\mathrm{tot}}^{(i)\ast}\to \mathcal{M}_{\mathrm{tot}}^{(i+1)\ast}\qquad (i=1,\ldots,N),
\]
使得对每一步 $i$ 都满足：
\begin{enumerate}
  \item \textbf{代价同伦残差界：}存在 $\varepsilon_i\ge 0$ 使对任意词 $w$ 有
  \[
  \bigl|\mathcal{S}_{\mathcal{E}_i}(w)-\mathcal{S}_{\mathcal{E}_{i+1}}(\phi_i(w))\bigr|\le \varepsilon_i;
  \]
  \item \textbf{证书单调迁移：}对任意词 $w$ 有
  \[
  \mathrm{Cert}_{\mathcal{E}_i}(w)=1\ \Longrightarrow\ \mathrm{Cert}_{\mathcal{E}_{i+1}}(\phi_i(w))=1;
  \]
  \item \textbf{探索分布推前：}定义 $\pi_{i+1}:=\phi_{i\#}\pi_i$。
\end{enumerate}
令复合编译 $\Phi_{N:1}:=\phi_N\circ\cdots\circ\phi_1$。则：
\begin{enumerate}
  \item \textbf{诊断边界可迁移：}
  \[
  \Phi_{N:1}\bigl(\mathcal{B}^{\mathrm{fail}}_{\mathcal{E}_1}\bigr)\subseteq \mathcal{B}^{\mathrm{fail}}_{\mathcal{E}
    _{N+1}};
  \]
  \item \textbf{自由能可控漂移：}对任意 $\beta>0$ 有
  \[
  \bigl|\mathcal{F}_{\mathcal{E}_1}(\beta;\pi_1)-\mathcal{F}_{\mathcal{E}_{N+1}}(\beta;\pi_{N+1})\bigr|
  \le
  \sum_{i=1}^{N}\varepsilon_i.
  \]
\end{enumerate}
因此在环境链上，算力投入可被转写为可复用资产：诊断边界与探索分布可随编译映射推前传播，且其在自由能意义下的漂移由残差和控制。
\end{theorem}

\begin{certificate}[数学归纳法证书：按复用步数 $N$ 的链式闭包]
\label{cert:m6-homotopy-reuse-induction}
定义谓词 $P(N)$ 为：“在 $N$ 步环境链上，Theorem~\ref{thm:m6-homotopy-reuse-chain} 的两条结论同时成立”。
证书化要求为：
\[
P(1)\ \text{成立},\qquad
(\forall N\in\mathbb{N}_{\ge 1})\ \bigl(P(N)\Rightarrow P(N+1)\bigr).
\]
其中：
\begin{itemize}
  \item $P(1)$ 的两条结论分别由 Theorem~\ref{thm:cert-reuse} 与 Theorem~\ref{thm:free-energy-stability} 给出；
  \item 归纳步对失败域使用集合包含的传递性，对自由能差界使用三角不等式并叠加一步残差界。
\end{itemize}
\end{certificate}

\begin{proof}[证明（数学归纳法证书；谓词因果链）]
由 Certificate~\ref{cert:m6-homotopy-reuse-induction}：
\begin{itemize}
  \item 基步 $N=1$：失败域迁移由 Theorem~\ref{thm:cert-reuse} 得
  \[
  \phi_1(\mathcal{B}^{\mathrm{fail}}_{\mathcal{E}_1})\subseteq \mathcal{B}^{\mathrm{fail}}_{\mathcal{E}_2};
  \]
  自由能差界由 Theorem~\ref{thm:free-energy-stability} 得
  \[
  \bigl|\mathcal{F}_{\mathcal{E}_1}(\beta;\pi_1)-\mathcal{F}_{\mathcal{E}_2}(\beta;\pi_2)\bigr|\le \varepsilon_1.
  \]
  \item 归纳步：假设 $P(N)$ 成立。记 $\Phi_{N:1}:=\phi_N\circ\cdots\circ\phi_1$，并令 $\Phi_{N+1:1}:=\phi_{N+1}\circ \Phi_{N:1}$。
  由归纳假设有 $\Phi_{N:1}(\mathcal{B}^{\mathrm{fail}}_{\mathcal{E}_1})\subseteq \mathcal{B}^{\mathrm{fail}}_{\mathcal{E}_{N+1}
    }$，再由两环境版本的失败域迁移（Theorem~\ref{thm:cert-reuse} 应用于 $(\mathcal{E}_{N+1},\mathcal{E}_{N+2},\phi_{N+1})$）得
  \[
  \phi_{N+1}\bigl(\mathcal{B}^{\mathrm{fail}}_{\mathcal{E}_{N+1}}\bigr)\subseteq \mathcal{B}^{\mathrm{fail}}_{\mathcal{E}
    _{N+2}}.
  \]
  合并两条包含关系得到 $\Phi_{N+1:1}(\mathcal{B}^{\mathrm{fail}}_{\mathcal{E}_1})\subseteq \mathcal{B}^{\mathrm{fail}}_{\mathcal{E}
    _{N+2}}$。
  对自由能差界，由三角不等式与归纳假设得
  \[
  \begin{aligned}
  &\bigl|\mathcal{F}_{\mathcal{E}_1}(\beta;\pi_1)-\mathcal{F}_{\mathcal{E}_{N+2}}(\beta;\pi_{N+2})\bigr|\\
  &\le
  \bigl|\mathcal{F}_{\mathcal{E}_1}(\beta;\pi_1)-\mathcal{F}_{\mathcal{E}_{N+1}}(\beta;\pi_{N+1})\bigr|
  +
  \bigl|\mathcal{F}_{\mathcal{E}_{N+1}}(\beta;\pi_{N+1})-\mathcal{F}_{\mathcal{E}_{N+2}}(\beta;\pi_{N+2})\bigr|.
  \end{aligned}
  \]
  第一项由 $P(N)$ 上界为 $\sum_{i=1}^{N}\varepsilon_i$；第二项由 Theorem~\ref{thm:free-energy-stability} 上界为 $\varepsilon_{N+1}$。
  因此得到 $\sum_{i=1}^{N+1}\varepsilon_i$ 的上界，归纳步成立。
\end{itemize}
由数学归纳法得结论。证毕。
\end{proof}

\subsection{定理：算力经济学口径下的“近最优推前复用”}

\begin{theorem}[近最优推前复用：同伦残差界蕴含最优性缺口上界]
\label{thm:near-optimal-pushforward}
设两环境 $\mathcal{E}_1,\mathcal{E}_2$ 之间存在词编译双射 $\phi$，且满足残差界
\[
\bigl|\mathcal{S}_{\mathcal{E}_1}(w)-\mathcal{S}_{\mathcal{E}_2}(\phi(w))\bigr|\le \varepsilon.
\]
定义最优自由能值
\[
\mathcal{F}_{\mathcal{E}}^\star(\beta):=\inf_{\pi}\ \mathcal{F}_{\mathcal{E}}(\beta;\pi).
\]
则对任意 $\beta>0$ 有：
\begin{enumerate}
  \item \textbf{最优值稳定：}$\bigl|\mathcal{F}_{\mathcal{E}_1}^\star(\beta)-\mathcal{F}_{\mathcal{E}_2}^\star(\beta)\bigr|\le
    \varepsilon$；
  \item \textbf{近最优可迁移：}若 $\pi_1$ 在 $\mathcal{E}_1$ 上为 $\delta$-近最优，即
  \[
  \mathcal{F}_{\mathcal{E}_1}(\beta;\pi_1)\le \mathcal{F}_{\mathcal{E}_1}^\star(\beta)+\delta,
  \]
  且令 $\pi_2:=\phi_{\#}\pi_1$，则 $\pi_2$ 在 $\mathcal{E}_2$ 上为 $(\delta+2\varepsilon)$-近最优：
  \[
  \mathcal{F}_{\mathcal{E}_2}(\beta;\pi_2)\le \mathcal{F}_{\mathcal{E}_2}^\star(\beta)+\delta+2\varepsilon.
  \]
\end{enumerate}
该结论给出“算力复用”的可计算口径：在残差有界的同伦类内，策略/探索分布可通过推前初始化获得可控的最优性缺口，从而减少重复试错成本。
\end{theorem}

\begin{certificate}[证书：由自由能稳定性界推得最优值与近最优差界]
\label{cert:near-optimal-pushforward}
（1）由 Theorem~\ref{thm:free-energy-stability}，对任意 $\pi_1$ 与 $\pi_2=\phi_{\#}\pi_1$ 有
\[
\bigl|\mathcal{F}_{\mathcal{E}_1}(\beta;\pi_1)-\mathcal{F}_{\mathcal{E}_2}(\beta;\pi_2)\bigr|\le \varepsilon;
\]
（2）对任意函数族 $f(\pi)$，有下确界不等式：$\inf_{\pi} f(\pi)\le f(\pi_0)$。
\end{certificate}

\begin{proof}[证明（谓词因果链）]
由 Certificate~\ref{cert:near-optimal-pushforward}：
\begin{itemize}
  \item 最优值稳定：任取 $\pi_1$，令 $\pi_2=\phi_{\#}\pi_1$，则
  \[
  \mathcal{F}_{\mathcal{E}_2}^\star(\beta)\le \mathcal{F}_{\mathcal{E}_2}(\beta;\pi_2)\le \mathcal{F}_{\mathcal{E}_1}
    (\beta;\pi_1)+\varepsilon.
  \]
  对 $\pi_1$ 取下确界得到 $\mathcal{F}_{\mathcal{E}_2}^\star(\beta)\le \mathcal{F}_{\mathcal{E}_1}^\star(\beta)+\varepsilon$。
  交换 $\mathcal{E}_1,\mathcal{E}_2$（使用 $\phi^{-1}$）同理得 $\mathcal{F}_{\mathcal{E}_1}^\star(\beta)\le \mathcal{F}
    _{\mathcal{E}_2}^\star(\beta)+\varepsilon$，合并得结论（1）。
  \item 近最优可迁移：由 Theorem~\ref{thm:free-energy-stability} 得
  \[
  \mathcal{F}_{\mathcal{E}_2}(\beta;\pi_2)\le \mathcal{F}_{\mathcal{E}_1}(\beta;\pi_1)+\varepsilon
  \le \mathcal{F}_{\mathcal{E}_1}^\star(\beta)+\delta+\varepsilon.
  \]
  再由（1）得 $\mathcal{F}_{\mathcal{E}_1}^\star(\beta)\le \mathcal{F}_{\mathcal{E}_2}^\star(\beta)+\varepsilon$，代入即得
  \[
  \mathcal{F}_{\mathcal{E}_2}(\beta;\pi_2)\le \mathcal{F}_{\mathcal{E}_2}^\star(\beta)+\delta+2\varepsilon.
  \]
\end{itemize}
综上得到定理结论。证毕。
\end{proof}

\subsection{引理：\RTSC{} 的“材料$\times$控制窄组合”要求控制一等化}

\begin{lemma}[控制一等化的必要性：忽略控制将产生伪不可达诊断]
\label{lem:ctrl-first-class-necessary}
固定路线 $k$ 与容差 $\varepsilon\ge 0$，令 $\mathcal{A}_k(\varepsilon)$ 与 $\mathcal{D}(k;\varepsilon)$ 如 Definition~\ref{def:certified-action-set}--Definition~\ref{def:diagnosis-and-treatment}。
定义“控制忽略”的子动作集为
\[
\mathcal{A}_k^{\mathrm{mat}}(\varepsilon)
:=
\{w\in \mathcal{A}_k(\varepsilon)\mid |w_{\mathrm{ctrl}}|=0\},
\qquad
\mathcal{D}^{\mathrm{mat}}(k;\varepsilon):=\inf_{w\in \mathcal{A}_k^{\mathrm{mat}}(\varepsilon)} \mathcal{S}_k(w).
\]
若存在可达词 $w^\star\in\mathcal{A}_k(\varepsilon)$ 使 $\mathcal{S}_k(w^\star)=0$ 且 $|w^\star_{\mathrm{ctrl}}|>0$，并且
  $\mathcal{A}_k^{\mathrm{mat}}(\varepsilon)=\emptyset$（或其下确界严格大于 $0$），则
\[
\mathcal{D}(k;\varepsilon)=0,\qquad \mathcal{D}^{\mathrm{mat}}(k;\varepsilon)=+\infty\ \ (\text{或 } >0).
\]
即：在 \RTSC{} 这类“目标可能是材料$\times$控制小众组合”的任务中，若未将控制算子提升为一等生成元（Definition~\ref{def:ctrl-first-class}），
  则诊断阶段可能直接把真实可达解排除在建模宇宙之外。
\end{lemma}

\begin{certificate}[证书：缩小动作集只会抬高诊断下界]
\label{cert:ctrl-first-class-necessary}
若 $B_1\subseteq B_2$，则对任意函数 $g$ 有 $\inf_{x\in B_2}g(x)\le \inf_{x\in B_1}g(x)$。
特别地，若 $B_1=\emptyset$，则可约定 $\inf_{x\in B_1}g(x)=+\infty$。
\end{certificate}

\begin{proof}[证明（谓词因果链）]
由 Definition~\ref{def:certified-action-set} 与 Definition~\ref{def:domain-projection}，$\mathcal{A}_k^{\mathrm{mat}}
  (\varepsilon)\subseteq \mathcal{A}_k(\varepsilon)$。
因此由 Certificate~\ref{cert:ctrl-first-class-necessary} 得
\[
\mathcal{D}(k;\varepsilon)=\inf_{w\in\mathcal{A}_k(\varepsilon)}\mathcal{S}_k(w)\le \inf_{w\in\mathcal{A}_k^{\mathrm{mat}
  }(\varepsilon)}\mathcal{S}_k(w)=\mathcal{D}^{\mathrm{mat}}(k;\varepsilon).
\]
若存在 $w^\star\in\mathcal{A}_k(\varepsilon)$ 使 $\mathcal{S}_k(w^\star)=0$，则 $\mathcal{D}(k;\varepsilon)\le 0$；结合
  $\mathcal{S}_k\ge 0$ 得 $\mathcal{D}(k;\varepsilon)=0$。
若同时 $\mathcal{A}_k^{\mathrm{mat}}(\varepsilon)=\emptyset$，则按约定 $\mathcal{D}^{\mathrm{mat}}(k;\varepsilon)=+\infty$，
  从而得到“忽略控制导致伪不可达”的结论。证毕。
\end{proof}

\subsection{命题：反路线偏差自洽偏航的可审计机制}

\begin{proposition}[反偏航约束：将 $\mathrm{Unc}/\mathrm{Res}/\mathrm{Cert}$ 并入代价可压制自洽偏航]
\label{prop:anti-drift-constraint}
定义反偏航代价为
\[
\mathcal{S}_{\mathrm{anti}}(w)
:=
\mathcal{S}(w)+\eta\,\mathrm{Unc}(w)+\kappa\,\mathrm{Res}(w)+M\cdot(1-\mathrm{Cert}(w)),
\qquad \eta,\kappa\ge 0,\ M\gg 0.
\]
则：
\begin{enumerate}
  \item 任意 $\mathrm{Cert}(w)=0$ 的候选满足 $\mathcal{S}_{\mathrm{anti}}(w)\ge M$，将被路径积分权重 $\exp(-\beta \mathcal{S}
    _{\mathrm{anti}}(w))$ 强抑制；
  \item 当 $\mathrm{Res}(w)$ 或 $\mathrm{Unc}(w)$ 升高时，候选的权重将被显式压制，从而减少“只在单一路线评估器中自洽”的投机解主导概率质量；
  \item 叠加路线熵底座约束（Theorem~\ref{thm:navigation-against-hallucination}），可形成可审计的反偏航机制：在错误地图上置信度难以不可逆地趋近 1。
\end{enumerate}
\end{proposition}

\begin{certificate}[证书：反偏航代价由残差一等化与证书门槛化拼接得到]
\label{cert:anti-drift-constraint}
（1）由 Definition~\ref{def:residual-class}，$\mathrm{Unc}(w)$ 与 $\mathrm{Res}(w)$ 可作为代价一等项并入；
（2）由 Proposition~\ref{prop:cert-gating}，$M\cdot(1-\mathrm{Cert}(w))$ 给出硬门槛化惩罚；
（3）由 Theorem~\ref{thm:navigation-against-hallucination}，在加入上述校准项且保持路线熵底座时可抑制坍缩幻觉。
\end{certificate}

\begin{proof}[证明（谓词因果链）]
由 Certificate~\ref{cert:anti-drift-constraint}：
\begin{itemize}
  \item 对（1）：若 $\mathrm{Cert}(w)=0$，则 $\mathcal{S}_{\mathrm{anti}}(w)\ge M$，故其路径积分权重满足 $\exp(-\beta \mathcal{S}
    _{\mathrm{anti}}(w))\le e^{-\beta M}$，在 $M\gg 0$ 时被强抑制；
  \item 对（2）：$\mathrm{Res}/\mathrm{Unc}$ 的上升直接抬高 $\mathcal{S}_{\mathrm{anti}}$，从而压制“跨评估器不一致/不确定度高”的候选主导概率质量；
  \item 对（3）：将上述校准项与路线熵底座同时施加，满足 Theorem~\ref{thm:navigation-against-hallucination} 的结构，从而在形式系统层面抑制路线坍缩幻觉。
\end{itemize}
三者合取即得命题三条结论。证毕。
\end{proof}

\subsection{推论：研发—治理联动链（尾部风险止损 + 诊断优先 + 再治疗下注）}

\begin{corollary}[两阶段联动：先止损治理，再做诊断探索，后做类内治疗]
\label{cor:tail-risk-rd-chain}
在同时面对“尾部风险必须抑制”与“\RTSC{} 机制路线高度不确定”的任务中，可将可执行策略组织为两阶段联动链：
\begin{enumerate}
  \item \textbf{止损治理：}以证书门槛化代价 $\mathcal{S}_{\mathrm{cert}}$（Proposition~\ref{prop:cert-gating}）强制排除不可复现/不可审计区域，
    保证底线风险约束不被牺牲；
  \item \textbf{诊断探索：}按 Theorem~\ref{thm:diagnosis-priority} 与 Corollary~\ref{cor:diagnosis-before-treatment} 先比较路线诊断值
    $\mathcal{D}(k;\varepsilon)$，并在控制一等化（Lemma~\ref{lem:ctrl-first-class-necessary}）下扩展候选空间；
  \item \textbf{类内治疗：}仅对少数被诊断为“可达或近可达”的路线做深度类内优化，并用反偏航约束（Proposition~\ref{prop:anti-drift-constraint}）持续监控 $\mathrm{Res}
    /\mathrm{Unc}/\mathrm{Cert}$，一旦触发偏航预警（Corollary~\ref{cor:hallucination-certificate}）则返回诊断阶段扩张候选。
\end{enumerate}
该链条为内参稿提供可执行的“研发—治理联动”决策口径：在不牺牲止损底线的前提下，将算力优先配置给可迁移诊断资产，再在正确同伦类内下注深度构造。
\end{corollary}

\begin{certificate}[证书：两阶段联动链由门槛化、诊断优先与反偏航机制合取得到]
\label{cert:tail-risk-rd-chain}
（1）由 Proposition~\ref{prop:cert-gating}，证书门槛化将不可复现区域转写为强惩罚；
（2）由 Theorem~\ref{thm:diagnosis-priority} 与 Corollary~\ref{cor:diagnosis-before-treatment}，路线比较应优先于类内深度优化；
（3）由 Lemma~\ref{lem:ctrl-first-class-necessary}，控制一等化可避免在诊断阶段排除真实可达解；
（4）由 Proposition~\ref{prop:anti-drift-constraint} 与 Corollary~\ref{cor:hallucination-certificate}，可形成可审计的偏航预警与回退机制。
\end{certificate}

\begin{proof}[证明（谓词因果链）]
由 Certificate~\ref{cert:tail-risk-rd-chain}：
\begin{itemize}
  \item（1）确保“尾部风险必须抑制”的底线约束以硬门槛形式进入优化；
  \item（2）（3）共同确保诊断阶段的候选空间不被过早缩窄，并以 $\mathcal{D}(k;\varepsilon)$ 的比较驱动路线选择；
  \item（4）确保在类内治疗过程中，若出现“模型自洽但跨评估器不一致/不可复现”的信号，则触发回退到诊断阶段而非继续加码局部拟合。
\end{itemize}
因此两阶段联动链在形式系统层面成立。证毕。
\end{proof}

\subsection{备注：落点建议（面向内参稿的结构插入位置）}

\begin{remark}[落点建议：把 M6 写成 \textsf{MetaSol} 的新增条款]
若将本节内容整理为 \textsf{MetaSol} 口径下的新增条款：
\begin{enumerate}
  \item Theorem~\ref{thm:m6-homotopy-reuse-chain} 可作为 \textsf{MetaSol} 的新增条款（示意名：M6 Virtual-Environment Homotopy Reuse），
    用于补强“最优元方案”的算力经济学理由；
  \item Theorem~\ref{thm:near-optimal-pushforward} 与 Corollary~\ref{cor:path-family-stability} 可作为“统一动力学/路径积分接口”的可计算补充：
    明确给出策略/探索分布推前复用的最优性缺口上界；
  \item Proposition~\ref{prop:anti-drift-constraint} 可放入文档的“GRL/path-integral 权重设计”附近，作为反路线偏航的可审计约束口径；
  \item Corollary~\ref{cor:tail-risk-rd-chain} 可放入“尾部风险必须抑制”与“\RTSC{} 路线下注”之间，作为资源配置的决策链条。
\end{enumerate}
\end{remark}

\clearpage
%%%%%%%%%%%%%%%%%%%%%%%%%%%%%%%%%%%%%%%%%%%%%%%%%%%%%%%%%%%%
\section{\PTOPB$_{\mathrm{RTSC}}$ 的基准主丛候选：基底/纤维划分与截面族}
\label{sec:rtsc-benchmark-ptopb-candidates}
%%%%%%%%%%%%%%%%%%%%%%%%%%%%%%%%%%%%%%%%%%%%%%%%%%%%%%%%%%%%

\subsection{备注：基准需要同时支撑虚拟实验、逆向构造与证书化复用}

\begin{remark}[基准目标：把“路线/管线/控制协议”内化为截面选择]
\RTSC{} 基准（benchmark）的目标不仅是“收集数据”，而是要同时支撑：
（i）可比较的虚拟实验（跨路线/跨评估器可对齐），（ii）可执行的逆向构造（从规格到可执行方案），（iii）可证书化的复用（$\mathrm{Cert}/\mathrm{Res}/\mathrm{Unc}$ 作为一等接口沉淀为可迁移资产）。
本节以“主丛/广义主丛”的口径给出一组 \PTOPB$_{\mathrm{RTSC}}$ 的主丛候选组合：每个组合都明确
（1）哪些变量被提升为\textbf{基底坐标}（适合做诊断与跨路线比较），（2）哪些变量被放入\textbf{纤维/截面}（表述冗余、近似层、离散化、隐变量与控制协议实现细节等）。
当纤维中仅含表示结构群 $G_{\mathrm{rep}}$ 时它退化为真正的主纤维丛；当纤维还包含控制/工艺等“实现空间”因子时，可将其视为带 $G_{\mathrm{rep}}$ 作用的广义主丛/关联丛候选。
\end{remark}

\subsection{定义：RTSC 基准母体变量分块（统一复用口径）}

\begin{definition}[RTSC 基准母体状态分块]
\label{def:rtsc-benchmark-blocks}
将一个“可被虚拟实验/证书化记录”的 \RTSC{} 研发状态写为五块：
\[
x
\equiv
\bigl(
x_{\mathrm{mat}},\ x_{\mathrm{env}},\ x_{\mathrm{ctrl}},\ x_{\mathrm{proc}},\ x_{\mathrm{obs}}
\bigr).
\]
其中：
\begin{enumerate}
  \item \textbf{材料块：}
  \[
  x_{\mathrm{mat}}=(x_{\mathrm{comp}},x_{\mathrm{struct}},x_{\mathrm{micro}}),
  \]
  例如化学计量/掺杂向量、晶体/层状/界面描述、缺陷/应变/颗粒度/相分离等微结构指标（用可证书化特征替代全原子坐标）；
  \item \textbf{环境块：}
  \[
  x_{\mathrm{env}}=(T,P,B,\epsilon_{\mathrm{strain}},\ldots),
  \]
  例如温度、压力、磁场、应变、边界浴与散热条件等；
  \item \textbf{控制块：}
  \[
  x_{\mathrm{ctrl}}=u(\cdot)\ \text{或其有限维参数化 }\theta_u,
  \]
  例如泵浦波形、门控序列、压力/温度扫描路径、反馈律等；
  \item \textbf{工艺块：}
  \[
  x_{\mathrm{proc}}=\text{工艺路径（序列）及阶段态},
  \]
  例如烧结/退火/淬火/薄膜生长参数、界面制备路径等；
  \item \textbf{可观测/证书块：}
  \[
  x_{\mathrm{obs}}=\bigl(T_c,J_c,H_{c2},\mathrm{Stab},\mathrm{ProcFeas},\mathrm{Cert},\mathrm{Res},\mathrm{Unc}\bigr),
  \]
  其中 $(T_c,J_c,H_{c2})$ 为性能指标，$\mathrm{Stab}/\mathrm{ProcFeas}$ 为稳定性/可制造性指标，$\mathrm{Cert}/\mathrm{Res}/\mathrm{Unc}$
    为证书、残差与不确定度接口（参见 Definition~\ref{def:residual-class} 与 Proposition~\ref{prop:cert-gating} 的口径）。
\end{enumerate}
\end{definition}

\subsection{定义：表示冗余结构群（或群胚）与路线选择的内化}

\begin{definition}[表示结构群（或群胚） $G_{\mathrm{rep}}$]
\label{def:g-rep}
引入表示冗余的结构对象 $G_{\mathrm{rep}}$（可取为群或广义群胚），其元素 $g\in G_{\mathrm{rep}}$ 用于吸收同一物理对象的不同表示与建模选择，包括但不限于：
\begin{enumerate}
  \item 坐标/网格/离散化/求解器参数等数值表示选择；
  \item 理论近似层与评估器族的路线选择（如 electron--phonon/Eliashberg、非简谐/量子核效应、高压氢化物链、强关联链、非平衡 Keldysh/Floquet 链等）；
  \item 规范/参数化冗余与等价变换。
\end{enumerate}
并假设 $G_{\mathrm{rep}}$ 作用于“实现空间”，使得等价表示位于同一轨道；该作用用于把“路线/管线”内化为局部截面选择。
\end{definition}

\begin{remark}[路线作为截面族]
在本节的候选组合中，“局部截面”统一解释为：对给定基底点，选定一条路线/一个求解管线/一个控制协议族，并把抽象点落成可执行对象。
这与前文“路线作为插件”的口径一致（参见 Corollary~\ref{cor:plugin-catalog}）。
\end{remark}

\subsection{定义：基准主丛候选模板（基底--纤维划分）}

\begin{definition}[主丛候选模板：由变量块划分诱导的平凡束]
\label{def:rtsc-benchmark-bundle-template}
令四个“配置块”的坐标空间分别为 $M_{\mathrm{mat}},M_{\mathrm{env}},M_{\mathrm{ctrl}},M_{\mathrm{proc}}$，
并记
\[
\mathcal{B}:=\{\mathrm{mat},\mathrm{env},\mathrm{ctrl},\mathrm{proc}\}.
\]
对任意子集 $I\subseteq\mathcal{B}$，定义：
\[
M_I:=\prod_{b\in I} M_b,
\qquad
F_I:=\Bigl(\prod_{b\in \mathcal{B}\setminus I} M_b\Bigr)\times G_{\mathrm{rep}},
\qquad
P_I:=M_I\times F_I,
\]
并令投影 $\pi_I:P_I\to M_I$ 为 $\pi_I(x,f)=x$。
对任意映射 $f(\cdot):M_I\to F_I$，定义截面
\[
s(x):=(x,f(x)).
\]
该截面满足 $\pi_I\circ s=\mathrm{id}_{M_I}$，可解释为“在基底坐标给定时，选择纤维实现与表示冗余（路线/管线/控制协议）”。
\end{definition}

\subsection{定理（公理降级）：分块提升生成候选主丛（按块数归纳）}

\begin{theorem}[生成定理（公理降级）：逐块提升到基底即可生成任意划分主丛候选]
\label{thm:block-lift-generates-bundles}
取任意排序 $(b_1,b_2,b_3,b_4)$ 为 $\mathcal{B}$ 的一个排列，并对 $k\in\{0,1,2,3,4\}$ 定义
\[
I_k:=\{b_1,\ldots,b_k\}.
\]
则对每个 $k$，由 Definition~\ref{def:rtsc-benchmark-bundle-template} 可构造候选束 $\pi_{I_k}:P_{I_k}\to M_{I_k}$，
且任意“可执行路线/管线/控制协议”的实现选择都可等价表述为某个截面 $s:M_{I_k}\to P_{I_k}$。
因此对任意子集 $I\subseteq\mathcal{B}$，选择一个使 $I=I_{|I|}$ 的排序即可得到相应的主丛候选；换言之，
“谁在基底、谁在纤维”的任意划分均可由“逐块提升到基底”的有限步过程生成。
\end{theorem}

\begin{certificate}[数学归纳法证书：按“已提升到基底的块数 $k$”归纳]
\label{cert:block-lift-induction}
定义谓词 $P(k)$ 为：“对集合 $I_k$，束 $\pi_{I_k}:P_{I_k}\to M_{I_k}$ 存在且任意映射 $f:M_{I_k}\to F_{I_k}$ 都诱导一个截面 $s(x)=(x,f(x))$”。
证书化要求为：
\[
P(0)\ \text{成立},\qquad
(\forall k\in\{0,1,2,3\})\ \bigl(P(k)\Rightarrow P(k+1)\bigr).
\]
\end{certificate}

\begin{proof}[证明（数学归纳法证书；谓词因果链）]
由 Definition~\ref{def:rtsc-benchmark-bundle-template}：
\begin{itemize}
  \item 基步 $k=0$：$M_{I_0}$ 为平凡积（可视为一点或空积），$P_{I_0}=M_{I_0}\times F_{I_0}$，投影与截面按定义成立，故 $P(0)$ 成立；
  \item 归纳步：若 $P(k)$ 成立，则 $M_{I_{k+1}}=M_{I_k}\times M_{b_{k+1}}$，$F_{I_{k+1}}=(\prod_{b\in\mathcal{B}\setminus I_{k+1}}
    M_b)\times G_{\mathrm{rep}}$，从而 $P_{I_{k+1}}=M_{I_{k+1}}\times F_{I_{k+1}}$ 与投影 $\pi_{I_{k+1}}$ 均按同一模板定义；任意映射 $f:
    M_{I_{k+1}}\to F_{I_{k+1}}$ 仍可写成截面 $s(x)=(x,f(x))$，故 $P(k+1)$ 成立。
\end{itemize}
由数学归纳法得对 $k=0,1,2,3,4$ 均成立。对任意 $I\subseteq\mathcal{B}$，选取将 $I$ 排到前 $|I|$ 个位置的排序即可得到 $I=I_{|I|}$，从而得到定理结论。证毕。
\end{proof}

\subsection{组合 A：物理配置为基底、表示/近似为纤维（推荐基准主丛）}

\begin{definition}[组合 A：配置--表示主丛]
\label{def:rtsc-bundle-A}
定义
\[
M_A:=M_{\mathrm{mat}}\times M_{\mathrm{env}}\times M_{\mathrm{ctrl}}\times M_{\mathrm{proc}},
\qquad
P_A:=M_A\times G_{\mathrm{rep}},
\qquad
\pi_A(x,g)=x.
\]
将路线族索引为 $k\in\mathcal{K}$。称 $\{s_k\}_{k\in\mathcal{K}}$ 为一组\textbf{路线截面族}，若每个 $s_k:M_A\to P_A$ 满足 $\pi_A\circ
  s_k=\mathrm{id}$，
并可写为
\[
s_k(x)=(x,\,g_k(x)),
\qquad g_k:M_A\to G_{\mathrm{rep}}.
\]
\end{definition}

\begin{lemma}[路线选择等价于局部截面选择]
\label{lem:route-as-section}
在 Definition~\ref{def:rtsc-bundle-A} 的口径下，给定路线 $k$ 等价于给定一个局部截面 $s_k$（即给定表示选择函数 $g_k$）。
\end{lemma}

\begin{certificate}[证书：在平凡束 $P_A=M_A\times G_{\mathrm{rep}}$ 上截面与函数一一对应]
\label{cert:section-function-bijection}
对平凡束 $\pi_A:M_A\times G_{\mathrm{rep}}\to M_A$，任意截面 $s$ 必存在唯一函数 $g(\cdot):M_A\to G_{\mathrm{rep}}$ 使 $s(x)=(x,g(x))$；
  反之任意 $g(\cdot)$ 都诱导一个截面。
\end{certificate}

\begin{proof}[证明（谓词因果链）]
由 Certificate~\ref{cert:section-function-bijection}，取谓词链：
\[
Q_1:\ s\ \text{是截面}\ \Rightarrow\ \exists! g(\cdot)\ \text{使 } s(x)=(x,g(x)),\qquad
Q_2:\ g(\cdot)\ \Rightarrow\ s(x)=(x,g(x))\ \text{且 }\pi_A\circ s=\mathrm{id}.
\]
因此“选路线 $k$（即选 $g_k$）”与“选截面 $s_k$”等价。证毕。
\end{proof}

\begin{remark}[组合 A 的工程含义]
组合 A 将“材料+环境+控制+工艺”作为可比较基底坐标，将“近似链/离散化/求解器/规范选择”压入纤维并通过路线截面族表达；
因此它天然支持跨路线的诊断比较与策略复用，并与“方向优先于构造”的诊断范式一致（参见 Theorem~\ref{thm:diagnosis-priority}）。
\end{remark}

\subsection{组合 B：材料为基底、控制为纤维（材料$\to$最优控制截面）}

\begin{definition}[组合 B：材料--环境--工艺基底 + 控制纤维]
\label{def:rtsc-bundle-B}
定义
\[
M_B:=M_{\mathrm{mat}}\times M_{\mathrm{env}}\times M_{\mathrm{proc}},
\qquad
P_B:=M_B\times M_{\mathrm{ctrl}}\times G_{\mathrm{rep}},
\qquad
\pi_B(x,u,g)=x.
\]
称截面 $s_B:M_B\to P_B$ 为\textbf{最优控制截面}，若它可写为
\[
s_B(x)=\bigl(x,\,u^\ast(x),\,g^\ast(x)\bigr),
\]
其中 $u^\ast(x)$ 表示在给定材料/环境/工艺下的（近）最优控制协议。
\end{definition}

\begin{proposition}[双层结构：截面选择等价于“固定材料后先控后治”的纤维最小化]
\label{prop:bundle-B-bilevel}
给定代价泛函 $\mathcal{S}:P_B\to\mathbb{R}_{\ge 0}$，定义材料基底点 $x\in M_B$ 的纤维最优包络
\[
\mathcal{V}_B(x):=\inf_{(u,g)}\ \mathcal{S}(x,u,g),
\qquad (u,g)\in M_{\mathrm{ctrl}}\times G_{\mathrm{rep}}.
\]
若存在最优控制截面 $s_B$ 使对任意 $x$ 有
\[
\mathcal{S}(s_B(x))=\mathcal{V}_B(x),
\]
则“组合 B 的截面选择”在几何上等价于对控制纤维做最小化；
该口径对应前文双层共设计形式（Definition~\ref{def:bilevel-codesign}）中“对每个材料候选，先求最优控制词”的结构。
\end{proposition}

\begin{certificate}[证书：纤维下确界与最优截面的等价展开]
\label{cert:bundle-B-bilevel}
对任意函数 $f$ 与集合 $A$，若存在 $a^\ast\in A$ 使 $f(a^\ast)=\inf_{a\in A}f(a)$，则 $a^\ast$ 实现该下确界，且对任意 $a$ 有 $f(a^\ast)\le f(a)$。
\end{certificate}

\begin{proof}[证明（谓词因果链）]
由 Definition~\ref{def:rtsc-bundle-B}，截面 $s_B$ 给出 $x\mapsto (u^\ast(x),g^\ast(x))$。
由 Certificate~\ref{cert:bundle-B-bilevel}，若对任意 $x$ 有 $\mathcal{S}(x,u^\ast(x),g^\ast(x))=\inf_{(u,g)}\mathcal{S}(x,u,g)$，
  则 $\mathcal{S}(s_B(x))=\mathcal{V}_B(x)$，且对任意控制/表示选择 $(u,g)$ 有
\[
\mathcal{S}(s_B(x))\le \mathcal{S}(x,u,g).
\]
因此“截面选择”与“控制纤维最小化”在形式上等价，并与双层共设计口径一致。证毕。
\end{proof}

\begin{remark}[组合 B 的诊断意义]
组合 B 直接编码“材料$\times$控制窄组合”的可能性：材料在基底上不可达不等价于“材料+最优控制”不可达。
因此它适合作为“材料候选筛选”基准：以 $\mathcal{V}_B(x)$（最优控制包络）而非固定控制下的代价排序材料候选。
\end{remark}

\subsection{组合 C：控制为基底、材料为纤维（控制优先诊断）}

\begin{definition}[组合 C：控制--环境基底 + 材料/工艺纤维]
\label{def:rtsc-bundle-C}
定义
\[
M_C:=M_{\mathrm{ctrl}}\times M_{\mathrm{env}},
\qquad
P_C:=M_C\times M_{\mathrm{mat}}\times M_{\mathrm{proc}}\times G_{\mathrm{rep}},
\qquad
\pi_C(u,e,\cdot)= (u,e).
\]
称截面 $s_C:M_C\to P_C$ 为“控制优先截面”，若它将给定控制协议族与环境点映射为某个材料/工艺实现：
\[
s_C(u,e)=\bigl((u,e),\,x_{\mathrm{mat}}^\ast(u,e),\,x_{\mathrm{proc}}^\ast(u,e),\,g^\ast(u,e)\bigr).
\]
\end{definition}

\begin{lemma}[控制优先诊断：将控制提升为基底可避免早期把控制降格为扰动]
\label{lem:ctrl-as-base-diagnosis}
在 \RTSC{} 这类“真实可达解可能依赖控制协议”的任务中，若以组合 C 的方式将 $M_{\mathrm{ctrl}}$ 提升为基底坐标，
则“控制同伦类/协议族”的可达性差异可被直接诊断与排序；从而降低“先固定材料路线、后把控制当扰动”导致的伪不可达风险（参见 Lemma~\ref{lem:ctrl-first-class-necessary}）。
\end{lemma}

\begin{certificate}[证书：把控制放入基底等价于强制在控制维度做诊断排序]
\label{cert:ctrl-as-base-diagnosis}
由 Definition~\ref{def:rtsc-bundle-C}，控制协议 $u$ 不再属于纤维最小化的“细节变量”，而成为基底坐标；因此任何基于基底排序/比较的诊断泛函都会显式依赖 $u$，
  从而迫使探索先覆盖控制协议族再进入类内优化。
\end{certificate}

\begin{proof}[证明（谓词因果链）]
由 Certificate~\ref{cert:ctrl-as-base-diagnosis} 取谓词链：
\[
\begin{aligned}
R_1:&\ u\in M_C\ \Rightarrow\ \text{诊断比较发生在 }(u,e)\text{ 层},\\
R_2:&\ \text{诊断发生在 }(u,e)\text{ 层}\Rightarrow \text{控制不会被默认降格为扰动}.
\end{aligned}
\]
由 $R_1,R_2$ 得结论；结合 Lemma~\ref{lem:ctrl-first-class-necessary} 的“忽略控制将产生伪不可达诊断”即可得到本引理主旨。证毕。
\end{proof}

\subsection{组合 D：机制不变量为基底、具体实现为纤维（诊断优先的机制坐标化）}

\begin{definition}[组合 D：机制不变量基底 $M_{\mathrm{inv}}$ 与实现纤维]
\label{def:rtsc-bundle-D}
令 $M_{\mathrm{inv}}$ 为一组粗粒度机制坐标的空间（示例：$(\lambda,\omega_{\log},N(0),\mu^\ast,U/W,\dim,\ldots)$；仅作示意，可按体系扩展）。
给定一个“机制抽取/粗粒化”映射
\[
\Psi:\ M_{\mathrm{mat}}\times M_{\mathrm{ctrl}}\times M_{\mathrm{proc}}\times G_{\mathrm{rep}}\to M_{\mathrm{inv}},
\qquad z=\Psi(x_{\mathrm{mat}},x_{\mathrm{ctrl}},x_{\mathrm{proc}},g).
\]
定义集合形式的候选束
\[
P_D:=\bigl\{(z,x_{\mathrm{mat}},x_{\mathrm{ctrl}},x_{\mathrm{proc}},g)\mid z=\Psi(x_{\mathrm{mat}},x_{\mathrm{ctrl}},
  x_{\mathrm{proc}},g)\bigr\},
\qquad
\pi_D(z,\cdot)=z.
\]
其纤维 $\pi_D^{-1}(z)$ 为“实现该机制不变量 $z$ 的所有具体实现”。
\end{definition}

\begin{theorem}[机制诊断：在不变量坐标上定义可达性下界并先于治疗排序]
\label{thm:diagnosis-on-invariants}
给定代价泛函 $\mathcal{S}:P_D\to\mathbb{R}_{\ge 0}$，定义机制坐标上的诊断下界
\[
\mathcal{D}_{\mathrm{inv}}(z):=\inf_{p\in \pi_D^{-1}(z)}\ \mathcal{S}(p).
\]
则对任意截面 $s_D:M_{\mathrm{inv}}\to P_D$ 与任意 $z$ 有
\[
\mathcal{S}\bigl(s_D(z)\bigr)\ge \mathcal{D}_{\mathrm{inv}}(z),
\]
从而比较 $\mathcal{D}_{\mathrm{inv}}(z)$ 可在“机制层”先行诊断可达性，再下沉到具体实现做类内优化（治疗）。
\end{theorem}

\begin{certificate}[证书：下确界是函数取值的下界]
\label{cert:diagnosis-on-invariants}
对任意集合 $A$ 与函数 $f$，量 $\inf_{x\in A}f(x)$ 是 $f(A)$ 的下界，即 $(\forall x\in A)\ f(x)\ge \inf_{y\in A}f(y)$（参见
  Certificate~\ref{cert:inf-lower-bound} 的同型口径）。
\end{certificate}

\begin{proof}[证明（谓词因果链）]
由 Certificate~\ref{cert:diagnosis-on-invariants}，对 $A:=\pi_D^{-1}(z)$ 与 $f:=\mathcal{S}$，任取 $p\in \pi_D^{-1}(z)$ 有
\[
\mathcal{S}(p)\ge \inf_{q\in \pi_D^{-1}(z)}\mathcal{S}(q)=\mathcal{D}_{\mathrm{inv}}(z).
\]
又因 $s_D(z)\in \pi_D^{-1}(z)$，代入即得 $\mathcal{S}(s_D(z))\ge \mathcal{D}_{\mathrm{inv}}(z)$。证毕。
\end{proof}

\begin{remark}[组合 D 的使用方式]
组合 D 把“方向/诊断”提升为机制不变量坐标上的可计算对象：先比较 $\mathcal{D}_{\mathrm{inv}}(z)$ 的可达性下界，再选择若干 $z$ 区域进入材料$\times$控制$\times$工艺的具体构造。
这与第 6 章的“诊断优先于治疗”同型（Definition~\ref{def:diagnosis-and-treatment}）。
\end{remark}

\subsection{组合 E：性能目标为基底、设计实现为纤维（逆向构造接口）}

\begin{definition}[组合 E：目标规格基底 $M_{\mathrm{target}}$ 与逆向设计纤维]
\label{def:rtsc-bundle-E}
令目标规格空间为
\[
M_{\mathrm{target}}
:=
\bigl(T_c^\star,P^\star,J_c^\star,H_{c2}^\star,\mathrm{Stab}^\star,\mathrm{Proc}^\star\bigr),
\]
并令可执行方案的词空间为 $\mathcal{W}:=\mathcal{M}_{\mathrm{tot}}^{\mathrm{ctrl}\ast}$（Definition~\ref{def:ctrl-first-class}）。
设前向输出/观测映射 $y:\mathcal{W}\to M_{\mathrm{obs}}$（参见 Definition~\ref{def:forward-map-y} 的同型口径），并以 $[w]$ 表示在表示冗余
  $G_{\mathrm{rep}}$ 诱导的等价关系下的等价类。
定义候选束
\[
P_E:=\bigl\{(y^\star,[w])\mid y([w])\succeq y^\star,\ w\in\mathcal{W}\bigr\},
\qquad
\pi_E(y^\star,[w])=y^\star.
\]
其纤维 $\pi_E^{-1}(y^\star)$ 即“满足目标规格的可执行方案集合”（与 Definition~\ref{def:target-set} 的目标集口径同型）。
\end{definition}

\begin{theorem}[逆向构造 = 截面：策略分布可视为随机截面]
\label{thm:inverse-design-as-section}
若对某个目标子域 $U\subseteq M_{\mathrm{target}}$，对任意 $y^\star\in U$ 都有 $\pi_E^{-1}(y^\star)\neq\emptyset$，
则存在（至少一个）局部截面 $s_E:U\to P_E$，其形式为
\[
s_E(y^\star)=\bigl(y^\star,[w^\ast(y^\star)]\bigr).
\]
进一步地，若以 $\pi(\cdot\mid y^\star)$ 表示在给定目标下对方案词 $w$ 的策略分布，则 $\pi(\cdot\mid y^\star)$ 可被解释为 $U$ 上的\textbf{随机截面}：为每个
  $y^\star$ 给出一个分布并在 $\pi_E^{-1}(y^\star)$ 上采样方案；
在 GRL/path-integral 口径下，该随机截面可通过权重 $\exp(-\beta\mathcal{S}(w))$ 学习与更新（Definition~\ref{def:env-free-energy}）。
\end{theorem}

\begin{certificate}[证书：非空纤维 $\Rightarrow$ 可选代表元形成截面]
\label{cert:inverse-design-as-section}
若对每个 $y^\star\in U$，集合 $\pi_E^{-1}(y^\star)$ 非空，则可为每个 $y^\star$ 选择一个代表元 $[w^\ast(y^\star)]$；将该选择函数与 $\pi_E$ 组装即可得到截面
  $s_E$，并满足 $\pi_E\circ s_E=\mathrm{id}_U$。
\end{certificate}

\begin{proof}[证明（谓词因果链）]
由 Certificate~\ref{cert:inverse-design-as-section}，对每个 $y^\star\in U$ 选取一个 $[w^\ast(y^\star)]\in\pi_E^{-1}(y^\star)$，并定义
  $s_E(y^\star)=(y^\star,[w^\ast(y^\star)])$。
则由 Definition~\ref{def:rtsc-bundle-E} 有 $\pi_E(s_E(y^\star))=y^\star$，即 $\pi_E\circ s_E=\mathrm{id}_U$，故截面存在。
随机截面的解释来自于：对每个 $y^\star$ 用策略分布 $\pi(\cdot\mid y^\star)$ 在 $\pi_E^{-1}(y^\star)$ 上采样代表元；而 GRL/path-integral 的权重更新由
  Definition~\ref{def:env-free-energy} 给出。证毕。
\end{proof}

\begin{remark}[组合 E 的定位]
组合 E 是“按需逆向构造”的自然接口：基底即规格，截面即方案；当与证书门槛化与残差/不确定度校准项结合时，可将“治疗（构造）”直接写成可审计的截面/随机截面选择过程（参见 Proposition~\ref{prop:anti-drift-constraint}）。
\end{remark}

\subsection{组合 F：证书声明为基底、证据链为纤维（抑制路线幻觉）}

\begin{definition}[组合 F：声明基底 $M_{\mathrm{claim}}$ 与证据链纤维]
\label{def:rtsc-bundle-F}
令声明空间为
\[
M_{\mathrm{claim}}
:=
\bigl(x_{\mathrm{obs}}^{\mathrm{claim}},\ \text{测量口径/协议}\bigr),
\]
并令 $\mathcal{E}_{\mathrm{evi}}$ 表示证据链空间（原始数据、校准、样品表征、重复实验、对照组、误差模型等的可审计组合）。
定义候选束
\[
P_F:=M_{\mathrm{claim}}\times \mathcal{E}_{\mathrm{evi}},
\qquad
\pi_F(c,e)=c.
\]
称截面 $s_F$ 为“证书截面”，若对任意声明 $c$，其证据链 $e=s_F(c)$ 满足 $\mathrm{Cert}(e)=1$。
\end{definition}

\begin{proposition}[证书主丛：把可复现提升为一等坐标可压制路线自洽偏航]
\label{prop:certificate-bundle-anti-hallucination}
在 Definition~\ref{def:rtsc-bundle-F} 的口径下，将 $\mathrm{Cert}=1$ 的证书截面作为硬约束（或用门槛化代价 $M\cdot(1-\mathrm{Cert})$ 强惩罚，见
  Proposition~\ref{prop:cert-gating}），
则任何缺乏证据链支撑的“自洽但不可复现”路线都难以在优化/路径积分权重中获得主导概率质量，从而抑制路线偏差幻觉（参见 Corollary~\ref{cor:hallucination-certificate}）。
\end{proposition}

\begin{certificate}[证书：证书门槛化会显式排除不可复现区域]
\label{cert:certificate-bundle-anti-hallucination}
由 Proposition~\ref{prop:cert-gating}，若采用 $\mathcal{S}_{\mathrm{cert}}=\mathcal{S}+M\cdot(1-\mathrm{Cert})$ 且 $M\gg 0$，则
  $\mathrm{Cert}=0$ 的候选在代价意义下被显式排除或强惩罚；结合 Corollary~\ref{cor:hallucination-certificate} 的触发口径，可将“缺证据链导致的自洽偏航”转写为可计算的预警信号。

\end{certificate}

\begin{proof}[证明（谓词因果链）]
由 Certificate~\ref{cert:certificate-bundle-anti-hallucination}，取谓词链：
\[
\begin{aligned}
S_1:&\ \mathrm{Cert}=0\Rightarrow \mathcal{S}_{\mathrm{cert}}\ge M,\\
S_2:&\ M\gg 0\Rightarrow \exp(-\beta \mathcal{S}_{\mathrm{cert}})\ \text{被强抑制},\\
S_3:&\ \text{被强抑制}\Rightarrow \text{不可复现路线难以主导概率质量}.
\end{aligned}
\]
由 $S_1,S_2,S_3$ 得结论；再由 Corollary~\ref{cor:hallucination-certificate} 可将“代价下降但残差/证书不改善”解释为偏航预警并触发再诊断。证毕。
\end{proof}

\subsection{推论：如何把 A--F 当作 \RTSC{} 基准主丛候选组合来用}

\begin{corollary}[选型推论：按目标选择组合 A--F]
\label{cor:rtsc-bundle-selection}
在 \PTOPB$_{\mathrm{RTSC}}$ 的基准设计中：
\begin{enumerate}
  \item 若目标是“统一诊断 + 防偏航 + 多路线可复用”，优先采用组合 A，并将路线管线实现为截面族 $\{s_k\}$（Definition~\ref{def:rtsc-bundle-A}）；
  \item 若目标是刻意捕捉“材料$\times$控制窄组合”，优先采用组合 B 或组合 C：分别对应“材料$\to$最优控制”与“控制优先诊断”（Definition~\ref{def:rtsc-bundle-B}、
    Definition~\ref{def:rtsc-bundle-C}）；
  \item 若目标是把“诊断优先于治疗”落成可计算事实（比较哪类机制更可能可达），优先采用组合 D（Definition~\ref{def:rtsc-bundle-D}）；
  \item 若目标是“按需逆向构造”并与 GRL/path-integral 直接对接，优先采用组合 E（Definition~\ref{def:rtsc-bundle-E}、Theorem~\ref{thm:inverse-design-as-section}）；
  \item 若目标是把“证书链”作为硬约束并压制路线幻觉，优先采用组合 F。
  （见 Definition~\ref{def:rtsc-bundle-F} 与 Proposition~\ref{prop:certificate-bundle-anti-hallucination}。）
\end{enumerate}
\end{corollary}

\begin{certificate}[证书：选型依据来自“基底坐标决定诊断粒度”与“纤维内化实现细节”]
\label{cert:rtsc-bundle-selection}
（1）由组合 A--F 的定义，基底选择决定“诊断比较发生在哪些坐标上”，纤维选择决定“哪些变量通过截面/最小化实现”；（2）由 Lemma~\ref{lem:ctrl-first-class-necessary} 与
  Lemma~\ref{lem:ctrl-as-base-diagnosis}，控制在 \RTSC{} 中不宜默认降格；（3）由 Theorem~\ref{thm:diagnosis-priority} 与
  Corollary~\ref{cor:hallucination-certificate}，方向诊断与反偏航约束决定是否会在错误地图上自洽化。
\end{certificate}

\begin{proof}[证明（谓词因果链）]
由 Certificate~\ref{cert:rtsc-bundle-selection}：
\begin{itemize}
  \item（1）给出“为何不同组合对应不同基准任务”的结构原因；
  \item（2）给出“为何 B/C 对材料$\times$控制窄组合重要”的逻辑原因；
  \item（3）给出“为何 A/D/E/F 分别对应诊断、机制层诊断、治疗接口与证书抑幻觉”的形式原因。
\end{itemize}
因此推论各条选型建议成立。证毕。
\end{proof}

\begin{remark}[进一步细化：最小基底、最小生成元与插件接口]
若需把某一组合（例如组合 A 作为 \RTSC{} 基准主丛）细化为“可直接落地的接口规范”，可继续给出：
（i）基底坐标块的最小必需集合（哪些必须可证书化），（ii）$G_{\mathrm{rep}}$ 的最小生成元（哪些表示冗余必须被吸收），（iii）截面族 $\{s_k\}$ 的插件接口（每个截面必须交付哪些评估器分层、
  残差/不确定度构造与证书字段）。
该细化可与前文插件目录口径（Corollary~\ref{cor:plugin-catalog}）一致化，以便在同一引擎下复用策略与诊断资产。
\end{remark}

%%%%%%%%%%%%%%%%%%%%%%%%%%%%%%%%%%%%%%%%%%%%%%%%%%%%%%%%%%%%
\section{\RTSC{} 基准主丛的基底切换：子丛限制与纤维包络}
\label{sec:rtsc-benchmark-base-switching}
%%%%%%%%%%%%%%%%%%%%%%%%%%%%%%%%%%%%%%%%%%%%%%%%%%%%%%%%%%%%

\subsection{备注：A--F 不是多套主丛，而是同一母体上的拆分切换}

\begin{remark}[统一图景：组合 A--F 作为基底/纤维拆分的不同选择]
组合 A--F 并不是互不相干的多套主丛，而是在同一个 \RTSC{} 基准母体（benchmark）下，对“基底/纤维拆分”做不同选择。
这些基底切换通过两类标准操作彼此连通：
\begin{enumerate}
  \item \textbf{子丛限制（sub-bundle restriction）：}在最细基底上取子域 $M'\subset M_\star$ 并限制得到 $P_\star|_{M'}$；
  \item \textbf{纤维包络（bundle envelope / fiber thickening）：}用压缩/遗忘/聚合映射 $p_i:M_\star\twoheadrightarrow M_i$ 将被遗忘自由度并入纤维，
    得到新投影 $\pi_i:=p_i\circ\pi_\star$，其纤维“变厚”。
\end{enumerate}
因此“子纤维丛（或纤维包络）构成新基底”的严格含义是：同一总空间在不同基底上表现为不同投影/商/限制，而不需要引入互相冲突的多套对象。
\end{remark}

\subsection{定义：最细基准母体 $M_\star$ 与基准主丛 $\pi_\star$}

\begin{definition}[最细基准母体与基准主丛]
\label{def:rtsc-finest-benchmark-bundle}
取最细基准母体基底为
\[
M_\star
:=
M_{\mathrm{mat}}\times M_{\mathrm{env}}\times M_{\mathrm{ctrl}}\times M_{\mathrm{proc}},
\]
必要时可再乘上可观测空间因子（例如 $M_{\mathrm{obs}}$ 或其子域）以便把测量口径也纳入基底。
令表示冗余结构群（或群胚）为 $G_{\mathrm{rep}}$（Definition~\ref{def:g-rep}）。
定义基准主丛
\[
\pi_\star:P_\star\to M_\star
\]
为主 $G_{\mathrm{rep}}$-丛；在“基准主丛取平凡化模型”的工作笔记口径下，可取
\[
P_\star:=M_\star\times G_{\mathrm{rep}},
\qquad
\pi_\star(x,g)=x,
\]
这与组合 A（Definition~\ref{def:rtsc-bundle-A}）及模板构造（Definition~\ref{def:rtsc-benchmark-bundle-template} 取 $I=\mathcal{B}$）同型。

\end{definition}

\subsection{定义：基底切换映射 $p_i$（压缩/遗忘/聚合）}

\begin{definition}[基底切换映射]
\label{def:base-switch-map}
称映射
\[
p_i: M_\star \twoheadrightarrow M_i
\]
为一次\textbf{基底切换}（压缩/遗忘/聚合）操作。
典型例子包括：
\begin{enumerate}
  \item \textbf{组合 B（遗忘控制）：}
  \[
  p_B(x_{\mathrm{mat}},x_{\mathrm{env}},x_{\mathrm{ctrl}},x_{\mathrm{proc}})
  :=
  (x_{\mathrm{mat}},x_{\mathrm{env}},x_{\mathrm{proc}}),
  \]
  得到基底 $M_B=M_{\mathrm{mat}}\times M_{\mathrm{env}}\times M_{\mathrm{proc}}$（Definition~\ref{def:rtsc-bundle-B}）；
  \item \textbf{组合 C（控制优先）：}
  \[
  p_C(x_{\mathrm{mat}},x_{\mathrm{env}},x_{\mathrm{ctrl}},x_{\mathrm{proc}})
  :=
  (x_{\mathrm{ctrl}},x_{\mathrm{env}}),
  \]
  得到基底 $M_C=M_{\mathrm{ctrl}}\times M_{\mathrm{env}}$（Definition~\ref{def:rtsc-bundle-C}）；
  \item \textbf{组合 D（机制不变量聚合）：}取某个粗粒化/不变量抽取映射 $p_D:M_\star\to M_{\mathrm{inv}}$（Definition~\ref{def:rtsc-bundle-D}）；
  \item \textbf{组合 E（目标规格聚合）：}取由前向输出与规格映射诱导的 $p_E:M_\star\to M_{\mathrm{target}}$（Definition~\ref{def:rtsc-bundle-E}）；
  \item \textbf{组合 F（声明口径聚合）：}取由“性能声明 + 测量口径”诱导的 $p_F:M_\star\to M_{\mathrm{claim}}$（Definition~\ref{def:rtsc-bundle-F}）。

\end{enumerate}
\end{definition}

\subsection{引理：对子域限制得到子主丛（sub-bundle restriction）}

\begin{lemma}[子丛限制：限制到 $M'\subset M_\star$ 仍为主丛]
\label{lem:subbundle-restriction}
若 $\pi_\star:P_\star\to M_\star$ 为主 $G_{\mathrm{rep}}$-丛，且 $M'\subset M_\star$ 为任意子域，则限制丛
\[
P_\star|_{M'}:=\pi_\star^{-1}(M')\to M'
\]
仍为主 $G_{\mathrm{rep}}$-丛（其结构群与右作用由 $P_\star$ 继承）。
\end{lemma}

\begin{certificate}[证书：主丛在子基底上的限制保持主丛结构]
\label{cert:subbundle-restriction}
若 $P\to M$ 为主 $G$-丛，则对任意子集 $M'\subset M$，限制丛 $P|_{M'}\to M'$ 仍为主 $G$-丛：自由性与传递性在每条纤维上逐点继承。
\end{certificate}

\begin{proof}[证明（谓词因果链）]
由 Certificate~\ref{cert:subbundle-restriction}，取谓词链：
\[
A:\ P_\star\to M_\star\ \text{为主 }G_{\mathrm{rep}}\text{-丛},\qquad
B:\ M'\subset M_\star.
\]
由 $A$ 与 $B$ 得到：对任意 $m'\in M'$，限制纤维满足 $P_\star|_{M'}{}^{-1}(m')=P_\star^{-1}(m')$；
因此 $G_{\mathrm{rep}}$ 在该纤维上的自由与传递作用保持不变，从而限制丛仍为主丛。证毕。
\end{proof}

\subsection{命题：纤维包络（fiber thickening）与纤维公式}

\begin{definition}[纤维包络投影]
\label{def:bundle-envelope}
给定基底切换映射 $p_i:M_\star\twoheadrightarrow M_i$（Definition~\ref{def:base-switch-map}），定义包络投影为
\[
\pi_i:=p_i\circ \pi_\star:\ P_\star\to M_i.
\]
称 $\pi_i$ 为 $\pi_\star$ 在 $p_i$ 下诱导的\textbf{纤维包络}（fiber thickening / bundle envelope）。
\end{definition}

\begin{proposition}[纤维包络公式：遗忘自由度并入纤维]
\label{prop:envelope-fiber-formula}
在 Definition~\ref{def:bundle-envelope} 下，对任意 $m_i\in M_i$，包络纤维满足集合恒等式
\[
\pi_i^{-1}(m_i)
=
\bigcup_{m_\star\in p_i^{-1}(m_i)}\ \pi_\star^{-1}(m_\star).
\]
因此被 $p_i$ 遗忘/聚合的自由度（即 $p_i^{-1}(m_i)$ 的变化）整体并入纤维，形成“纤维变厚”的包络效应。
\end{proposition}

\begin{certificate}[证书：复合映射的原像展开]
\label{cert:envelope-fiber-formula}
对映射 $f:X\to Y$、$g:Y\to Z$ 与元素 $z\in Z$，有集合恒等式
\[
(g\circ f)^{-1}(z)=\bigcup_{y\in g^{-1}(z)} f^{-1}(y).
\]
\end{certificate}

\begin{proof}[证明（谓词因果链）]
取 $f=\pi_\star$、$g=p_i$，由 Certificate~\ref{cert:envelope-fiber-formula} 得
\[
(p_i\circ\pi_\star)^{-1}(m_i)
=
\bigcup_{m_\star\in p_i^{-1}(m_i)}\ \pi_\star^{-1}(m_\star).
\]
又由 Definition~\ref{def:bundle-envelope}，左侧即 $\pi_i^{-1}(m_i)$，故命题成立。证毕。
\end{proof}

\subsection{定理（公理降级）：包络算子在有限次基底切换下封闭（按步数归纳）}

\begin{theorem}[包络闭包（公理降级）：有限次压缩等价于一次合成压缩]
\label{thm:envelope-closure}
设给定一列基底切换映射
\[
p_1:M_\star\twoheadrightarrow M_1,\quad
p_2:M_1\twoheadrightarrow M_2,\quad \ldots,\quad
p_N:M_{N-1}\twoheadrightarrow M_N,
\]
并定义合成映射 $p_{N:1}:=p_N\circ\cdots\circ p_1$。
令 $\pi_{N:1}:=p_{N:1}\circ\pi_\star$。
则：
\begin{enumerate}
  \item \textbf{投影合成性：}$\pi_{N:1}$ 与按序逐次定义的包络投影一致；
  \item \textbf{纤维公式：}对任意 $m_N\in M_N$，有
  \[
  \pi_{N:1}^{-1}(m_N)
  =
  \bigcup_{m_\star\in p_{N:1}^{-1}(m_N)}\ \pi_\star^{-1}(m_\star).
  \]
\end{enumerate}
因此对任意有限步的“基底压缩/遗忘/聚合”，都可等价转写为一次合成压缩，并在集合意义下给出显式可计算的纤维膨胀公式。
\end{theorem}

\begin{certificate}[数学归纳法证书：按压缩步数 $N$ 归纳]
\label{cert:envelope-closure-induction}
定义谓词 $P(N)$ 为：“对 $N$ 步压缩的合成映射 $p_{N:1}$，Theorem~\ref{thm:envelope-closure} 的两条结论同时成立”。
证书化要求为：
\[
P(1)\ \text{成立},\qquad
(\forall N\in\mathbb{N}_{\ge 1})\ \bigl(P(N)\Rightarrow P(N+1)\bigr).
\]
\end{certificate}

\begin{proof}[证明（数学归纳法证书；谓词因果链）]
由 Certificate~\ref{cert:envelope-closure-induction}：
\begin{itemize}
  \item 基步 $N=1$：$p_{1:1}=p_1$，$\pi_{1:1}=p_1\circ\pi_\star$；纤维公式由 Proposition~\ref{prop:envelope-fiber-formula} 立即给出，故
    $P(1)$ 成立。
  \item 归纳步：假设 $P(N)$ 成立。令 $p_{N+1:1}:=p_{N+1}\circ p_{N:1}$，则
  \[
  \pi_{N+1:1}
  =
  p_{N+1:1}\circ\pi_\star
  =
  p_{N+1}\circ(p_{N:1}\circ\pi_\star),
  \]
  即投影合成性成立。
  对纤维公式，令 $m_{N+1}\in M_{N+1}$，由 Certificate~\ref{cert:envelope-fiber-formula} 得
  \[
  \pi_{N+1:1}^{-1}(m_{N+1})
  =
  \bigcup_{m_N\in p_{N+1}^{-1}(m_{N+1})}\ (p_{N:1}\circ\pi_\star)^{-1}(m_N).
  \]
  再由归纳假设 $P(N)$，对每个 $m_N$ 有
  \[
  (p_{N:1}\circ\pi_\star)^{-1}(m_N)
  =
  \bigcup_{m_\star\in p_{N:1}^{-1}(m_N)}\ \pi_\star^{-1}(m_\star).
  \]
  将其代回并合并并集，得到
  \[
  \pi_{N+1:1}^{-1}(m_{N+1})
  =
  \bigcup_{m_\star\in p_{N+1:1}^{-1}(m_{N+1})}\ \pi_\star^{-1}(m_\star),
  \]
  即 $P(N+1)$ 成立。
\end{itemize}
由数学归纳法得对任意 $N$ 成立，证毕。
\end{proof}

\subsection{命题：基底切换下的交换图与“投影/商/限制”统一表述}

\begin{proposition}[交换图命题：基底压缩可与商化 $q_i$ 组成交换方块]
\label{prop:commuting-square}
给定基底切换映射 $p_i:M_\star\twoheadrightarrow M_i$，并给定一个等价关系 $\sim_i$ 作用于 $P_\star$，满足
\[
u\sim_i v\ \Longrightarrow\ p_i(\pi_\star(u))=p_i(\pi_\star(v)).
\]
定义商空间 $P_i:=P_\star/\!\sim_i$ 与商映射 $q_i:P_\star\to P_i$。
则映射
\[
\pi_i:P_i\to M_i,\qquad
\pi_i\bigl(q_i(u)\bigr):=p_i(\pi_\star(u))
\]
良定义，并给出交换图
\[
\begin{array}{ccc}
P_\star & \xrightarrow{\ q_i\ } & P_i \\
\downarrow \pi_\star & & \downarrow \pi_i \\
M_\star & \xrightarrow{\ p_i\ } & M_i
\end{array}
\qquad
\text{满足}\quad
\pi_i\circ q_i=p_i\circ\pi_\star.
\]
当 $\sim_i$ 取平凡等价关系时，$P_i=P_\star$，该交换图退化为包络投影（Definition~\ref{def:bundle-envelope}）；
当 $P_\star$ 限制到子域 $M'\subset M_\star$ 时，可取 $q_i$ 为包含映射以得到子丛限制的交换表述（Lemma~\ref{lem:subbundle-restriction}）。
\end{proposition}

\begin{certificate}[证书：商映射下的良定义判据]
\label{cert:commuting-square}
若函数 $f:P_\star\to M_i$ 在等价类上常值，即 $u\sim_i v\Rightarrow f(u)=f(v)$，则可定义商函数 $\overline{f}:P_i\to M_i$ 使 $\overline{f}\circ
  q_i=f$，且 $\overline{f}$ 良定义。
\end{certificate}

\begin{proof}[证明（谓词因果链）]
令 $f(u):=p_i(\pi_\star(u))$。由命题前提 $u\sim_i v\Rightarrow p_i(\pi_\star(u))=p_i(\pi_\star(v))$，即 $f$ 在等价类上常值。
因此由 Certificate~\ref{cert:commuting-square}，存在良定义函数 $\pi_i$ 满足 $\pi_i\circ q_i=f=p_i\circ\pi_\star$，即交换图成立。证毕。
\end{proof}

\subsection{命题：纤维包络的“主丛化”版本（结构群扩展的充分条件）}

\begin{proposition}[主丛化命题：若遗忘自由度可被群作用参数化，则可扩展结构群]
\label{prop:principalization}
在包络投影 $\pi_i=p_i\circ\pi_\star$ 的情形下，若满足下列可验证条件：
\begin{enumerate}
  \item 存在群（或群胚）$H_i$ 作用于 $M_\star$，且对每个 $m_i\in M_i$，$H_i$ 在 $p_i^{-1}(m_i)$ 上自由且传递；
  \item 该作用可提升到 $P_\star$ 上，并与 $G_{\mathrm{rep}}$ 的右作用满足半直积兼容（即存在同态 $H_i\to \mathrm{Aut}(G_{\mathrm{rep}})$ 使得两者可组装为
    $G_i:=G_{\mathrm{rep}}\rtimes H_i$ 的作用）；
\end{enumerate}
则 $P_\star$ 可被视为主 $G_i$-丛在基底 $M_i$ 上的一个实现：包络纤维中的“被遗忘自由度”由 $H_i$ 方向参数化，表示冗余由 $G_{\mathrm{rep}}$ 方向参数化，从而保持“可操作的纤维方向”语义。
\end{proposition}

\begin{certificate}[证书：自由传递作用 + 兼容提升 $\Rightarrow$ 主丛化]
\label{cert:principalization}
若群 $H$ 在每个 $p^{-1}(m)$ 上自由且传递，并可提升到总空间与既有结构群作用形成半直积，则扩展后的群在包络纤维上可实现自由传递作用，从而给出主丛结构。
\end{certificate}

\begin{proof}[证明（谓词因果链）]
由 Certificate~\ref{cert:principalization}，取谓词链：
\[
A:\ H_i\ \text{在 }p_i^{-1}(m_i)\text{ 上自由传递},\quad
B:\ H_i\ \text{作用提升到 }P_\star,\quad
C:\ \text{与 }G_{\mathrm{rep}}\text{ 组装为 }G_i=G_{\mathrm{rep}}\rtimes H_i.
\]
由 $A$ 得到：任意两点 $m_\star,m_\star'\in p_i^{-1}(m_i)$ 可由唯一 $h\in H_i$ 连接；由 $B$ 得到：该 $h$ 可在总空间把 $\pi_\star^{-1}(m_\star)$ 运送到
  $\pi_\star^{-1}(m_\star')$；
由 $C$ 得到：在每条 $\pi_\star$-纤维内仍可用 $G_{\mathrm{rep}}$ 处理表示冗余，而跨 $p_i^{-1}(m_i)$ 的“遗忘方向”由 $H_i$ 处理。
因此包络纤维上的自由度可被 $G_i$ 的作用统一参数化，得到主丛化语义。证毕。
\end{proof}

\subsection{推论：组合 A--F 的统一解释与截面互化问题}

\begin{corollary}[统一推论：A--F 皆为同一 $P_\star$ 的投影/商/限制]
\label{cor:af-unified}
在 Definition~\ref{def:rtsc-finest-benchmark-bundle} 的基准主丛 $\pi_\star:P_\star\to M_\star$ 下：
\begin{enumerate}
  \item 组合 A 对应“保持最细基底”，即 $(P_\star,M_\star,\pi_\star)$；
  \item 组合 B/C/D/E/F 可分别由某个基底切换映射 $p_i:M_\star\to M_i$（Definition~\ref{def:base-switch-map}）诱导包络投影
    $\pi_i=p_i\circ\pi_\star$（Definition~\ref{def:bundle-envelope}），并在需要时再通过商化 $q_i$（Proposition~\ref{prop:commuting-square}）消除等价表述；
  \item 若进一步对某类材料家族/压力区间/控制协议族取子域 $M'\subset M_\star$，则通过 Lemma~\ref{lem:subbundle-restriction} 得到子丛限制版本。
\end{enumerate}
因此 A--F 并非冲突的多方案，而是同一基准主丛在不同基底上的投影/商/限制；这与“诊断优先、可复用检验、同伦类复用”的动机一致。
同时，截面在不同基底下的关系可被解释为一个提升问题：在包络基底 $M_i$ 上选截面等价于在每个 $p_i^{-1}(m_i)$ 上选择代表性的细基底点并给出其表示/管线（即在 $P_\star$ 上选择一条落点）。
\end{corollary}

\begin{certificate}[证书：包络投影与交换图给出“同一总空间的不同投影”]
\label{cert:af-unified}
（1）由 Proposition~\ref{prop:envelope-fiber-formula}，包络投影将被遗忘自由度并入纤维；（2）由 Proposition~\ref{prop:commuting-square}，可通过
  $q_i$ 对总空间做商化并保持交换性；（3）由 Lemma~\ref{lem:subbundle-restriction}，对子域限制得到子丛。
\end{certificate}

\begin{proof}[证明（谓词因果链）]
由 Certificate~\ref{cert:af-unified}：
\begin{itemize}
  \item（1）（2）给出“压缩基底 $\Rightarrow$ 纤维膨胀 + 可选商化”的统一形式；
  \item（3）给出“加条件/取子域 $\Rightarrow$ 子丛限制”的统一形式。
\end{itemize}
将三者分别对应到组合 A（不切换）、组合 B/C/D/E/F（基底压缩与必要的商化）、以及路线化子域（子丛限制），即可得到推论结论。证毕。
\end{proof}

\begin{remark}[与第 6--8 章的接口关系]
第 8 章给出了 \PTOPB$_{\mathrm{RTSC}}$ 的多种候选基底/纤维拆分；本章说明这些拆分可在同一基准主丛上通过“限制/包络/商化”互相连通。
结合第 6 章的诊断泛函 $\mathcal{D}(\cdot)$（Definition~\ref{def:diagnosis-and-treatment}）与第 7 章的同伦复用条款（Theorem~\ref{thm:m6-homotopy-reuse-chain}），可将“基底切换”解释为：先在更粗基底上做可迁移诊断与比较，再在子域限制下进入类内治疗与证书化落地。
\end{remark}

%%%%%%%%%%%%%%%%%%%%%%%%%%%%%%%%%%%%%%%%%%%%%%%%%%%%%%%%%%%%
\section{丛塔与切面包络：基底切换的多级化记号}
\label{sec:bundle-tower-section-envelope}
%%%%%%%%%%%%%%%%%%%%%%%%%%%%%%%%%%%%%%%%%%%%%%%%%%%%%%%%%%%%

\subsection{备注：将组合 A--F 压缩为“丛塔/基底切换 + 切面包络”}

\begin{remark}[统一化目标：同一总空间的多级视图]
本章将第 8 章的组合 A--F 进一步统一为一个更干净的框架：
\textbf{丛塔（bundle tower）/基底切换（base switching）+ 切面包络（section envelope）}。
其核心点是：
\begin{enumerate}
  \item 固定一个“全量基准母体”作为最细基底 $M_\star$，并在其上给定基准主丛 $\pi_\star:P_\star\to M_\star$（Definition~\ref{def:rtsc-finest-benchmark-bundle}）；
  \item 一阶基底切换由满射 $p_i:M_\star\twoheadrightarrow M_i$ 给出，并诱导同一总空间上的新投影 $\pi_i:=p_i\circ\pi_\star:P_\star\to
    M_i$（Definition~\ref{def:bundle-envelope}）；
  \item 在固定基底 $M_i$ 下，不同路线/管线/控制协议族对应局部截面（切面）$s_i^{(j)}$；
  \item 二阶“切面包络”将截面从“零厚度代表”扩展为可控邻域 $E_{ij}$，并用细粒度自由度 $m_{ij}$（参数空间 $\Xi_{ij}$）显式参数化，从而得到二阶丛（或等价的细化基底形式）。
\end{enumerate}
该统一化的直接用途是：把“方向=诊断”写成对不同 $p_i$ 与不同 $s_i^{(j)}$ 的可计算比较，把“构造=治疗”写成对包络邻域 $E_{ij}$（或提升后的细化基底 $U_i^{(j)}\times\Xi_{ij}$）的局部深挖；
  并且 $m_{ij}$ 提供抑制“路线自洽偏航”的可审计抓手。
\end{remark}

\subsection{定义：全量基底 $M_\star=M_{(1,2,3,\dots)}(m_{(1,2,3,\dots)})$}

\begin{definition}[全量坐标块与全量基底]
\label{def:full-base-blocks}
为避免与“基底切换索引 $i$”混淆，本章用圆括号下标表示全量坐标块：
\[
m_{(1,2,3,\dots)}:=\bigl(m_{(1)},m_{(2)},m_{(3)},\dots\bigr),
\]
其中 $m_{(a)}$ 表示第 $a$ 个“诊断层面不可再丢”的坐标块（可为材料/环境/控制/工艺/观测/证书等任意分块；分块方式不影响本章的抽象结构）。
定义全量基底为
\[
M_\star
:=
M_{(1,2,3,\dots)}\bigl(m_{(1,2,3,\dots)}\bigr).
\]
在 \RTSC{} 语境中，$m_{(1,2,3,\dots)}$ 可取为 Definition~\ref{def:rtsc-benchmark-blocks} 的五块拼接（必要时将测量口径也并入），并与
  Definition~\ref{def:rtsc-finest-benchmark-bundle} 的 $M_\star$ 识别。
\end{definition}

\subsection{定义：基准一阶主丛与表示结构对象 $G_{\mathrm{rep}}$}

\begin{definition}[基准一阶主丛]
\label{def:first-order-benchmark-bundle}
固定表示冗余结构对象 $G_{\mathrm{rep}}$（Definition~\ref{def:g-rep}），并令
\[
\pi_\star:P_\star\to M_\star
\]
为主 $G_{\mathrm{rep}}$-丛（Definition~\ref{def:rtsc-finest-benchmark-bundle}）。
在工作笔记的平凡化模型中，可取 $P_\star=M_\star\times G_{\mathrm{rep}}$ 与 $\pi_\star(x,g)=x$，但本章的“基底切换/包络”结论并不依赖于平凡化。
\end{definition}

\subsection{定义：一阶基底切换与包络投影}

\begin{definition}[一阶基底切换与包络投影]
\label{def:first-order-base-switching}
对任意 $i$，取基底切换（压缩/遗忘/聚合）映射
\[
p_i: M_\star \twoheadrightarrow M_i,
\qquad
M_i:=M_i\bigl(m_{(1,2,3,\dots)}\bigr),
\]
并在同一总空间 $P_\star$ 上定义诱导投影
\[
\pi_i:=p_i\circ\pi_\star:\ P_\star\to M_i.
\]
该构造与第 9 章的包络投影定义（Definition~\ref{def:bundle-envelope}）同型。
\end{definition}

\begin{lemma}[纤维包络并公式：遗忘自由度整体并入纤维]
\label{lem:first-order-fiber-union}
在 Definition~\ref{def:first-order-base-switching} 下，对任意 $u\in M_i$ 有
\[
\pi_i^{-1}(u)=\bigcup_{x\in p_i^{-1}(u)}\ \pi_\star^{-1}(x).
\]
\end{lemma}

\begin{certificate}[证书：包络纤维公式由复合原像展开得到]
\label{cert:first-order-fiber-union}
由 Proposition~\ref{prop:envelope-fiber-formula}（或其证书 Certificate~\ref{cert:envelope-fiber-formula}），对复合映射
  $\pi_i=p_i\circ\pi_\star$ 的原像可展开为对 $p_i^{-1}(u)$ 的并集，从而得到纤维包络公式。
\end{certificate}

\begin{proof}[证明（谓词因果链）]
由 Certificate~\ref{cert:first-order-fiber-union}，将 $p_i$ 与 $\pi_\star$ 代入 Proposition~\ref{prop:envelope-fiber-formula}
  即得结论。证毕。
\end{proof}

\subsection{定义：一阶切面（局部截面）与路线索引 $j$}

\begin{definition}[一阶切面：局部截面 $s_i^{(j)}$]
\label{def:first-order-sections}
在基底 $M_i$ 上，称映射
\[
s_i^{(j)}:\ U_i^{(j)}\subseteq M_i\to P_\star
\]
为一阶切面（局部截面），若满足
\[
\pi_i\circ s_i^{(j)}=\mathrm{Id}_{U_i^{(j)}}.
\]
其中索引 $j$ 用于标记“同一基底 $M_i$ 下的不同路线/管线/规范选择/控制协议族”，并可与“路线插件”口径一致化（参见 Corollary~\ref{cor:plugin-catalog}）。
\end{definition}

\begin{remark}[切面与路线的等价语义]
在平凡化模型 $P_\star=M_\star\times G_{\mathrm{rep}}$ 下，切面 $s_i^{(j)}$ 等价于一组“表示/管线选择函数”，即对每个粗粒度基底点选取一个可执行的表示与求解管线；
因此“路线选择=切面选择”可被严格表达（参见 Lemma~\ref{lem:route-as-section} 的同型口径）。
\end{remark}

\subsection{定义：二阶切面包络 $E_{ij}$ 与细粒度自由度 $m_{ij}$}

\begin{definition}[切面包络：邻域 $E_{ij}$ 与参数空间 $\Xi_{ij}$]
\label{def:section-envelope}
给定一阶切面 $s_i^{(j)}:U_i^{(j)}\to P_\star$（Definition~\ref{def:first-order-sections}）。
称 $E_{ij}\subseteq P_\star$ 为该切面的\textbf{包络邻域}，若
\[
s_i^{(j)}\bigl(U_i^{(j)}\bigr)\subseteq E_{ij}.
\]
引入细粒度自由度参数空间 $\Xi_{ij}$，并记其坐标为
\[
m_{ij}\in \Xi_{ij}.
\]
若存在映射（可视作“管状邻域/参数化”）
\[
\Phi_{ij}:\ U_i^{(j)}\times \Xi_{ij}\to E_{ij}\subseteq P_\star
\]
满足
\[
\pi_i\bigl(\Phi_{ij}(u,m_{ij})\bigr)=u,
\qquad
\Phi_{ij}(u,0)=s_i^{(j)}(u),
\]
则称 $(E_{ij},\Xi_{ij},\Phi_{ij})$ 为一组\textbf{二阶切面包络数据}：它将切面从“零厚度代表”扩展为在纤维方向上的可控扰动族，并把隐含细节自由度 $m_{ij}$ 显式化。
\end{definition}

\subsection{命题：切面包络导出二阶丛（两种等价形态）}

\begin{proposition}[二阶丛：同基底纤维细化 vs.\ 提升细粒度自由度为基底]
\label{prop:second-order-envelope}
在 Definition~\ref{def:section-envelope} 下，二阶包络可用两种等价形态表达：
\begin{enumerate}
  \item \textbf{形态 1（同一基底上的纤维细化）：}定义限制投影
  \[
  \pi_{ij}:=\pi_i|_{E_{ij}}:\ E_{ij}\to U_i^{(j)}.
  \]
  该形式用于“在固定路线附近做鲁棒性/不确定度分析”。
  \item \textbf{形态 2（提升 $m_{ij}$ 为显式二阶基底）：}定义二阶基底
  \[
  M_{ij}^{(2)}:=U_i^{(j)}\times \Xi_{ij},
  \]
  并以 $\tilde P_{ij}:=M_{ij}^{(2)}$ 作为带坐标的重参数化总空间，定义投影
  \[
  \tilde\pi_{ij}:\tilde P_{ij}\to M_{ij}^{(2)},\qquad \tilde\pi_{ij}=\mathrm{Id}_{M_{ij}^{(2)}},
  \]
  同时用 $\Phi_{ij}$ 将其嵌入原总空间：$\Phi_{ij}:\tilde P_{ij}\to P_\star$。
  该形式用于“把细节自由度提升为可被 GRL 直接优化的状态变量”。
\end{enumerate}
并且：若 $\Phi_{ij}$ 对每个 $u$ 在纤维方向上为单射（或在所关心的邻域内可逆），则两种形态在“可计算接口”意义下等价：形态 2 是形态 1 的坐标化版本。
\end{proposition}

\begin{certificate}[证书：限制投影与坐标化嵌入的等价口径]
\label{cert:second-order-envelope}
（1）限制映射 $E_{ij}\subseteq P_\star$ 上的投影 $\pi_i$ 自然给出 $\pi_{ij}=\pi_i|_{E_{ij}}$；（2）若存在参数化 $\Phi_{ij}:U_i^{(j)}
  \times\Xi_{ij}\to E_{ij}$ 且满足 $\pi_i\circ\Phi_{ij}=\mathrm{pr}_{U}$，则 $(u,m_{ij})$ 可作为 $E_{ij}$ 的局部坐标；在该坐标下，
  “总空间到基底”的投影可取恒等映射 $\tilde\pi_{ij}=\mathrm{Id}$，而 $\Phi_{ij}$ 负责把坐标化对象回嵌到原总空间。
\end{certificate}

\begin{proof}[证明（谓词因果链）]
由 Certificate~\ref{cert:second-order-envelope}：
\begin{itemize}
  \item 形态 1：$E_{ij}\subseteq P_\star$ 且 $\pi_i:P_\star\to M_i$ 已定义，因此限制得到 $\pi_{ij}=\pi_i|_{E_{ij}}$；
  \item 形态 2：由 $\pi_i\circ\Phi_{ij}=\mathrm{pr}_{U}$，对任意 $(u,m_{ij})$，$\Phi_{ij}(u,m_{ij})$ 位于 $\pi_i^{-1}(u)$，故 $(u,
    m_{ij})$ 可作为该邻域的显式坐标；令 $\tilde P_{ij}=U_i^{(j)}\times\Xi_{ij}$ 与 $\tilde\pi_{ij}=\mathrm{Id}$ 即得到坐标化的“二阶环境空间”。
\end{itemize}
若进一步假设 $\Phi_{ij}$ 在纤维方向上（局部）单射，则不同 $m_{ij}$ 对应不同的邻域点，从而形态 2 与形态 1 在该邻域内一一对应，得到“坐标化等价”。证毕。
\end{proof}

\subsection{定理（公理降级）：丛塔构造在有限层级下封闭（按层数归纳）}

\begin{theorem}[丛塔封闭性（公理降级）：基底切换--切面--包络可递归生成有限层级丛塔]
\label{thm:bundle-tower-closure}
从基准一阶主丛 $\pi_\star:P_\star\to M_\star$ 出发，对任意有限层数 $L\in\mathbb{N}$，若给定一组分层数据：
\begin{enumerate}
  \item \textbf{第 1 层（基底切换）：}选择某个 $i$ 并给定 $p_i:M_\star\twoheadrightarrow M_i$；
  \item \textbf{第 2 层（切面选择）：}在 $M_i$ 上给定局部切面 $s_i^{(j)}:U_i^{(j)}\to P_\star$；
  \item \textbf{第 3 层及以后（包络细化）：}对每一层 $\ell\in\{1,\ldots,L\}$，给定参数空间 $\Xi_{ij}^{[\ell]}$ 与映射
  \[
  \Phi_{ij}^{[\ell]}:\ U_i^{(j)}\times \Xi_{ij}^{[\ell]}\to P_\star
  \]
  满足 $\pi_i\circ \Phi_{ij}^{[\ell]}=\mathrm{pr}_{U}$，并在 $\ell=1$ 时满足 $\Phi_{ij}^{[1]}(u,0)=s_i^{(j)}(u)$；
\end{enumerate}
则可递归定义一个多指标丛塔（以“细化层数 $\ell$”为塔高）：
\[
M_{ij}^{(0)}:=U_i^{(j)},\qquad
M_{ij}^{(\ell)}:=U_i^{(j)}\times \prod_{t=1}^{\ell}\Xi_{ij}^{[t]}\quad(\ell\ge 1),
\]
并令 $\tilde P_{ij}^{(\ell)}:=M_{ij}^{(\ell)}$、$\tilde\pi_{ij}^{(\ell)}=\mathrm{Id}_{M_{ij}^{(\ell)}}$，以及嵌入映射
\[
\Phi_{ij}^{(\ell)}:\tilde P_{ij}^{(\ell)}\to P_\star
\]
由逐层拼接构造（将各层参数映射串接到同一 $P_\star$ 的邻域中）。
因此：有限层级的“基底切换--切面--包络细化”可被统一编码为一个可计算的丛塔对象，且每一层的细粒度自由度都被显式提升为基底坐标。
\end{theorem}

\begin{certificate}[数学归纳法证书：按包络细化层数 $\ell$ 归纳]
\label{cert:bundle-tower-closure-induction}
定义谓词 $P(\ell)$ 为：“$\ell$ 阶细化基底 $M_{ij}^{(\ell)}$ 与坐标化总空间 $\tilde P_{ij}^{(\ell)}$ 可按 Theorem~\ref{thm:bundle-tower-closure} 构造，且存在嵌入 $\Phi_{ij}^{(\ell)}$ 使 $\pi_i\circ \Phi_{ij}^{(\ell)}=\mathrm{pr}_{U}$（对 $U_i^{(j)}$ 分量）
  ”。
证书化要求为：
\[
P(0)\ \text{成立},\qquad
(\forall \ell\in\mathbb{N})\ \bigl(P(\ell)\Rightarrow P(\ell+1)\bigr).
\]
\end{certificate}

\begin{proof}[证明（数学归纳法证书；谓词因果链）]
由 Certificate~\ref{cert:bundle-tower-closure-induction}：
\begin{itemize}
  \item 基步 $\ell=0$：取 $M_{ij}^{(0)}=U_i^{(j)}$，并由切面条件 $\pi_i\circ s_i^{(j)}=\mathrm{Id}$ 得到嵌入 $\Phi_{ij}^{(0)}:=s_i^{(j)}
    $ 满足 $\pi_i\circ \Phi_{ij}^{(0)}=\mathrm{pr}_{U}$，故 $P(0)$ 成立；
  \item 归纳步：假设 $P(\ell)$ 成立。定义
  \[
  M_{ij}^{(\ell+1)}:=M_{ij}^{(\ell)}\times \Xi_{ij}^{[\ell+1]}=U_i^{(j)}\times \prod_{t=1}^{\ell+1}\Xi_{ij}^{[t]}.
  \]
  取 $\tilde P_{ij}^{(\ell+1)}:=M_{ij}^{(\ell+1)}$ 与 $\tilde\pi_{ij}^{(\ell+1)}=\mathrm{Id}$。
  由给定的参数映射 $\Phi_{ij}^{[\ell+1]}$（满足 $\pi_i\circ \Phi_{ij}^{[\ell+1]}=\mathrm{pr}_{U}$）以及归纳假设中已有的 $\Phi_{ij}^{(\ell)}$，
    可将“已有参数元组 + 新增参数”组合为一个映射 $\Phi_{ij}^{(\ell+1)}:\tilde P_{ij}^{(\ell+1)}\to P_\star$，并保持 $\pi_i\circ\Phi_{ij}
    ^{(\ell+1)}=\mathrm{pr}_{U}$（因为每一层映射对 $U$ 分量都保持不变）。
  因此 $P(\ell+1)$ 成立。
\end{itemize}
由数学归纳法得对任意 $\ell\le L$ 成立，从而 Theorem~\ref{thm:bundle-tower-closure} 成立。证毕。
\end{proof}

\subsection{推论：二阶/多阶包络自由度提供“反偏航”的可审计接口}

\begin{corollary}[反偏航推论：细粒度自由度显式化 $\Rightarrow$ 可证书化抑制自洽偏航]
\label{cor:envelope-anti-drift}
在 Theorem~\ref{thm:bundle-tower-closure} 的丛塔口径下，包络细化引入的参数块 $m_{ij}$（及其多阶扩展）可被视为“最易导致路线自洽偏航的隐藏自由度”的显式坐标化版本。
若将这些自由度对应的不确定度/残差/证书谓词一并纳入代价（Proposition~\ref{prop:anti-drift-constraint}），则：
\begin{enumerate}
  \item 在同一切面邻域内，对 $m_{ij}$ 的扫描可把“看似自洽的局部最优”转写为 $\mathrm{Res}/\mathrm{Unc}/\mathrm{Cert}$ 的可观测偏航信号；
  \item 一旦触发“代价下降但残差不降 \& 证书难以满足”的预警（Corollary~\ref{cor:hallucination-certificate}），即可回退到更粗基底层（切换 $p_i$ 或扩张 $\Xi_{ij}$）
    进行再诊断，而不是继续在单一路线内加码治疗。
\end{enumerate}
\end{corollary}

\begin{certificate}[证书：反偏航代价与偏航预警口径已在前文给出]
\label{cert:envelope-anti-drift}
\begin{enumerate}
  \item 由 Proposition~\ref{prop:anti-drift-constraint}，\\
  $\mathrm{Unc}/\mathrm{Res}/\mathrm{Cert}$ 可作为一等项并入代价并压制投机解；
  \item 由 Corollary~\ref{cor:hallucination-certificate}，\\
  固定路线内“代价下降但残差/证书不改善”构成可计算预警；
  \item 由本章的包络塔定义，$m_{ij}$ 将隐藏自由度显式化，从而可被纳入上述三类统计量的计算与审计。
\end{enumerate}
\end{certificate}

\begin{proof}[证明（谓词因果链）]
由 Certificate~\ref{cert:envelope-anti-drift}：
\begin{itemize}
  \item（1）给出“将隐藏自由度显式化并纳入代价”的反偏航机制；
  \item（2）给出“何时应停止类内加码并回退诊断”的可计算触发口径；
  \item（3）说明 $m_{ij}$ 的显式化使上述机制可被直接实现于包络邻域的扫描与统计上。
\end{itemize}
三者合取即可推出推论两条结论。证毕。
\end{proof}

\subsection{备注：多指标索引丛塔记号（面向后续递归展开）}

\begin{remark}[多指标索引：用 $\alpha\in\mathbb{N}^k$ 统一“切换—切面—包络—再细化”]
为避免后续在 \RTSC{} 基准章节中把“切换—切面—包络—再细化”写散成叙事，可引入多指标索引丛塔记号：
\[
(M_\alpha,\ P_\alpha,\ \pi_\alpha),\qquad \alpha\in \mathbb{N}^k,
\]
其中 $\alpha$ 的不同维分别对应：基底切换层级（$p_i$）、路线切面层级（$s_i^{(j)}$）、包络细化层级（$m_{ij}$）以及其它需要显式化的自由度层级。
本章的 Theorem~\ref{thm:bundle-tower-closure} 给出“有限塔高可被归纳验证”的形式系统口径，可作为后续递归展开的统一模板。
\end{remark}

\clearpage

%%%%%%%%%%%%%%%%%%%%%%%%%%%%%%%%%%%%%%%%%%%%%%%%%%%%%%%%%%%%
\section{总雅可比间隙（Total JacGap）可控性接口}
\label{sec:ptopb-rtsc-total-jacgap}
%%%%%%%%%%%%%%%%%%%%%%%%%%%%%%%%%%%%%%%%%%%%%%%%%%%%%%%%%%%%

\subsection{形式逻辑表达约定}

\begin{remark}[形式逻辑：公理降级、证书与谓词因果链]
本章延续第 1 章的统一表达口径：将机制性断言以“公理降级为定理”的形式给出，并显式附带数学归纳法证书；证明采用基于数学符号推演的谓词因果链。
在此口径下，“总雅可比间隙（Total JacGap）决定可控性”不再是修辞，而被编码为可计算泛函项 $G_{\mathrm{tot}}$ 与可控性谓词 $\inf_w\mathcal{S}[w]\le \varepsilon$
  之间的可验证逻辑关系。
\end{remark}

\subsection{定义：基准母体 $M_\ast$、子节点坐标块 $m_i$ 与子节点法对象 $\mathcal{L}_i$}

\begin{definition}[基准母体与子节点坐标块]
\label{def:rtsc-jacgap-base-subnodes}
令全量基准母体基底为
\[
M_\ast:=M_\star=M_{(1,2,3,\dots)}\bigl(m_{(1,2,3,\dots)}\bigr)
\]
（见 Definition~\ref{def:full-base-blocks}），并记
\[
m_{(1,2,3,\dots)}=(m_1,m_2,\dots,m_N)
\]
为其子节点坐标块拼接：每个 $m_i$ 对应一类可分离的机制/结构/约束块（例如材料晶格块、缺陷块、界面块、控制协议块、理论近似块、测量口径与证书块等）。
\end{definition}

\begin{definition}[子节点法对象族（局部约束族）]
\label{def:rtsc-subnode-law}
对每个子节点 $i\in\{1,\dots,N\}$，定义一个\textbf{子节点法对象}（局部约束族）
\[
\mathcal{L}_i=\mathcal{L}_i\bigl(m_i;\ m_{\neq i}\bigr),
\qquad
m_{\neq i}:=(m_1,\dots,m_{i-1},m_{i+1},\dots,m_N),
\]
其语义为：在给定全局状态 $m\in M_\ast$ 时，第 $i$ 个子节点应满足的局部结构/闭合/一致性约束。
并令总法对象（总约束族）为合取：
\[
\mathcal{L}_{\mathrm{tot}}(m):=\bigwedge_{i=1}^N \mathcal{L}_i(m).
\]
\end{definition}

\subsection{定义：子节点 Jacobiator-gap 与雅可比间隙向量 $\mathbf{g}(m)$}

\begin{definition}[子节点 Jacobiator-gap（JacGap）与间隙字量]
\label{def:rtsc-jacgap-subnode}
对每个子节点 $i$，取一个“障碍向量空间” $V_i$（可为抽象线性空间或可测特征空间），并定义其 Jacobiator/闭合失配对象
\[
\mathbf{J}_i(m)\in V_i.
\]
定义对应的\textbf{雅可比间隙字量}为
\[
g_i(m):=\|\mathbf{J}_i(m)\|_{V_i}\ \ge 0,
\]
并汇总得到雅可比间隙向量
\[
\mathbf{g}(m):=(g_1(m),g_2(m),\dots,g_N(m))\in \mathbb{R}_{\ge 0}^N.
\]
\end{definition}

\subsection{定义：总雅可比间隙（Total JacGap）$G_{\mathrm{tot}}(m)$}

\begin{definition}[总雅可比间隙：加权和与向量化版本]
\label{def:rtsc-total-jacgap}
给定权重 $w_i>0$，定义总雅可比间隙（Total JacGap）为加权和：
\[
G_{\mathrm{tot}}(m):=\sum_{i=1}^N w_i\,g_i(m).
\]
亦可取某个“总障碍空间” $V$ 与线性映射 $W_i:V_i\to V$，定义向量化总 JacGap：
\[
\mathbf{J}_{\mathrm{tot}}(m):=\sum_{i=1}^N W_i\,\mathbf{J}_i(m),
\qquad
G_{\mathrm{tot}}(m):=\|\mathbf{J}_{\mathrm{tot}}(m)\|_{V}.
\]
在形式系统中，两种定义的差异可通过证书化等价（范数等价常数）吸收；工程实现上更常用加权和形式以便直接进入优化器。
\end{definition}

\subsection{定义：目标法对象与可满足性谓词 $\mathrm{RTSC}(m)$}

\begin{definition}[目标谓词：总 JacGap 零点 + 证书门槛]
\label{def:rtsc-predicate-jacgap}
定义“室温超导”（\RTSC{}）在形式系统中的可满足性谓词为：
\[
\mathrm{RTSC}(m)\ \Longleftrightarrow\ \bigl(G_{\mathrm{tot}}(m)=0\bigr)\ \wedge\ \mathrm{Cert}(m)=1.
\]
其中 $\mathrm{Cert}(m)\in\{0,1\}$ 为可复现证书谓词（Proposition~\ref{prop:cert-gating} 的口径，可包含测量口径、重复性、对照组、磁响应与误差模型等证据链约束）。
\end{definition}

\subsection{定义：算子词路径 $w$、轨线 $m(t)$ 与变分泛函 $\mathcal{S}[w]$}

\begin{definition}[路径与变分泛函：Total JacGap 作为主导项进入 GRL 路径积分]
\label{def:rtsc-jacgap-action}
令 $w\in\mathcal{W}_{\mathrm{adm}}$ 表示允许的“设计/工艺/控制/测量”算子词（可视为 $\mathcal{M}_{\mathrm{tot}}^\ast$ 的一个可行动作子集），其在基底上诱导一条轨线
  $m(t)\in M_\ast$（$t\in[0,T]$）。
定义变分泛函（Action）为
\[
\begin{aligned}
\mathcal{S}[w]
:=
\int_{0}^{T}\Bigl(
&G_{\mathrm{tot}}(m(t))
\,+\,
\lambda\,\mathrm{Res}(m(t))
\,+\,
\eta\,\mathrm{Unc}(m(t))
\\
&+\,
M\cdot\bigl(1-\mathrm{Cert}(m(t))\bigr)
\Bigr)\,dt,
\qquad \lambda,\eta,M\ge 0,
\end{aligned}
\]
其中 $\mathrm{Res}/\mathrm{Unc}$ 的一等化并入沿用 Definition~\ref{def:residual-class} 的口径（此处将其视为随轨线可计算的统计量），$M\gg 0$
  为证书门槛惩罚（Proposition~\ref{prop:cert-gating}）。
由此得到路径积分权重
\[
\mathbb{P}(w)\ \propto\ \exp\!\bigl(-\beta\,\mathcal{S}[w]\bigr),
\qquad \beta>0.
\]
\end{definition}

\begin{remark}[接口解释：总 JacGap 定义可优化空间的几何]
在 Definition~\ref{def:rtsc-jacgap-action} 下，$G_{\mathrm{tot}}$ 不仅是“评价指标”，而是直接进入代价并决定权重的主导项：
优化/采样过程会把概率质量压到那些能降低 $G_{\mathrm{tot}}$ 的“可控路径族”上。
因此“总 JacGap 决定可控性”可以被直接理解为：$G_{\mathrm{tot}}$ 的可达下界决定 $\inf_w\mathcal{S}[w]$ 的可达下界。
\end{remark}

\subsection{定义：可控性谓词（阈值化）}

\begin{definition}[可控性谓词：以 $\inf_w\mathcal{S}(w)\le \varepsilon$ 表达]
\label{def:rtsc-jacgap-controllable}
给定阈值 $\varepsilon\ge 0$，称 “\RTSC{} 在算子集合 $\mathcal{W}_{\mathrm{adm}}$ 下（$\varepsilon$-）可控” 当且仅当
\[
\inf_{w\in\mathcal{W}_{\mathrm{adm}}}\ \mathcal{S}[w]\ \le\ \varepsilon.
\]
当 $\varepsilon=0$ 且 $\mathrm{Cert}$ 可满足时，对应严格零点可达；当 $\varepsilon>0$ 时对应“近似可达/近似可证书化”的工程口径。
\end{definition}

\subsection{定理（公理降级）：不可作用子节点 $\Rightarrow$ 总 JacGap 下界锁死可控性}

\begin{theorem}[不可作用下界（公理降级）：单一障碍子节点锁死可控性]
\label{thm:rtsc-jacgap-unactuated-lower-bound}
设存在某个子节点索引 $i_0$ 与常数 $c_{i_0}>0$，使得对任意允许路径 $w\in\mathcal{W}_{\mathrm{adm}}$ 与任意 $t\in[0,T]$ 都有
\[
g_{i_0}\bigl(m(t)\bigr)\ge c_{i_0}.
\]
则对任意 $w\in\mathcal{W}_{\mathrm{adm}}$ 都有
\[
\mathcal{S}[w]\ \ge\ \int_0^T G_{\mathrm{tot}}(m(t))\,dt\ \ge\ w_{i_0}\,c_{i_0}\,T,
\]
从而
\[
\inf_{w\in\mathcal{W}_{\mathrm{adm}}}\ \mathcal{S}[w]\ \ge\ w_{i_0}\,c_{i_0}\,T.
\]
因此：若 $\varepsilon < w_{i_0}c_{i_0}T$，则 \RTSC{} 在该算子集合下不可控；更强地，$G_{\mathrm{tot}}(m)=0$ 的零点不可能沿允许路径被达到。
\end{theorem}

\begin{certificate}[数学归纳法证书：按子节点数 $N$ 归纳（加权和下界）]
\label{cert:rtsc-jacgap-unactuated-induction}
定义谓词 $P(N)$ 为：“对任意 $g_1,\dots,g_N\ge 0$，若存在 $i_0\le N$ 使 $g_{i_0}\ge c_{i_0}>0$，则
\[
\sum_{i=1}^N w_i g_i\ \ge\ w_{i_0}\,c_{i_0}.
\]
” 证书化要求为：
\[
P(1)\ \text{成立},\qquad
(\forall N\in\mathbb{N})\ \bigl(P(N)\Rightarrow P(N+1)\bigr).
\]
\end{certificate}

\begin{proof}[证明（数学归纳法证书；谓词因果链）]
由 Certificate~\ref{cert:rtsc-jacgap-unactuated-induction}：
\begin{itemize}
  \item 基步 $N=1$：若 $g_1\ge c_1$，则 $w_1 g_1\ge w_1 c_1$，故 $P(1)$ 成立；
  \item 归纳步：假设 $P(N)$ 成立。对 $N+1$ 维向量 $(g_1,\dots,g_{N+1})$，若 $i_0\le N$，则由归纳假设
  \[
  \sum_{i=1}^{N} w_i g_i\ge w_{i_0}c_{i_0}\ \Rightarrow\ \sum_{i=1}^{N+1} w_i g_i=\sum_{i=1}^{N} w_i g_i+w_{N+1}g_{N+1}
    \ge w_{i_0}c_{i_0};
  \]
  若 $i_0=N+1$，则 $\sum_{i=1}^{N+1} w_i g_i\ge w_{N+1}g_{N+1}\ge w_{N+1}c_{N+1}$。故 $P(N+1)$ 成立。
\end{itemize}
由数学归纳法，$P(N)$ 对任意 $N$ 成立，从而对任意轨线 $m(t)$ 有 $G_{\mathrm{tot}}(m(t))\ge w_{i_0}c_{i_0}$。再由 Definition~\ref{def:rtsc-jacgap-action} 的各项非负可推出
\[
\mathcal{S}[w]\ge \int_0^T G_{\mathrm{tot}}(m(t))\,dt\ge w_{i_0}c_{i_0}T.
\]
因此得到 Theorem~\ref{thm:rtsc-jacgap-unactuated-lower-bound}。证毕。
\end{proof}

\subsection{命题：局部可控性充分条件（秩/灵敏度条件）}

\begin{proposition}[局部可控性充分条件：$\partial\mathbf{g}/\partial u$ 满秩 $\Rightarrow$ 可把 $\mathbf{g}$ 压到任意小]
\label{prop:rtsc-jacgap-rank-sufficient}
设存在一组可控参数 $u\in U\subset\mathbb{R}^d$（来自 $\mathcal{M}_{\mathrm{ctrl}}$ 的有效作用方向），使得在某个邻域内可将轨线写为 $m=m(u)$，从而 $\mathbf{g}$
  诱导为 $\mathbf{g}(u)$。
若在某点 $u_0$ 满足
\[
\mathrm{rank}\Bigl(\frac{\partial \mathbf{g}}{\partial u}(u_0)\Bigr)=N,
\]
则存在 $u_0$ 的邻域使得：对任意 $\delta>0$，存在 $u$ 使 $\|\mathbf{g}(u)\|_\infty\le \delta$，从而 $G_{\mathrm{tot}}(m(u))$
  可被压到任意小（再叠加证书门槛即可形成“可证书化近可达”的局部可控口径）。
\end{proposition}

\begin{certificate}[证书：满秩 $\Rightarrow$ 局部可达（隐函数/正则值口径）]
\label{cert:rtsc-jacgap-rank}
若映射 $\mathbf{g}:U\to\mathbb{R}^N$ 在 $u_0$ 处雅可比满秩，则 $u_0$ 为正则点；在局部意义下 $\mathbf{g}(U)$ 含有 $0$ 的一个邻域（或至少在各分量上可独立调节），因此可把
  $\mathbf{g}$ 压到任意小。
\end{certificate}

\begin{proof}[证明（谓词因果链）]
由 Certificate~\ref{cert:rtsc-jacgap-rank}，取谓词链：
\[
\begin{aligned}
A:&\ \mathrm{rank}\Bigl(\frac{\partial \mathbf{g}}{\partial u}(u_0)\Bigr)=N,\\
B:&\ \mathbf{g}\ \text{在 }u_0\ \text{附近局部可达（分量可独立调节）},\\
C:&\ (\forall \delta>0)\ \exists u:\ \|\mathbf{g}(u)\|_\infty\le \delta.
\end{aligned}
\]
由 $A\Rightarrow B$（满秩给出局部可达性），由 $B\Rightarrow C$（局部像含小邻域），从而得到结论。证毕。
\end{proof}

\subsection{定义：格点可变调节与晶格控制子幺半群 $\mathcal{M}_{\mathrm{lat}}$}

\begin{definition}[格点可变调节（可作用自由度）]
\label{def:rtsc-lattice-control}
称子节点 $i$ 含有“格点可变调节”自由度，若其坐标块显式包含某类晶格/结构变量（如晶格常数、应变张量、局部畸变场、界面错配等）：
\[
m_i\supset a_i.
\]
并称控制集合包含“晶格控制”，若存在控制子幺半群（或可行动作子集）
\[
\mathcal{M}_{\mathrm{lat}}\subset \mathcal{M}_{\mathrm{ctrl}}
\]
使其在该自由度上产生非平凡作用：
\[
a_i \xrightarrow{\ \mathcal{M}_{\mathrm{lat}}\ } a_i'\quad\text{且}\quad a_i'\neq a_i\ \text{在一般位置成立}.
\]
\end{definition}

\subsection{命题：格点可调是 Total JacGap 可控性的必要接口}

\begin{proposition}[必要接口：晶格不可作用 $\Rightarrow$ JacGap 下界 $\Rightarrow$ 不可控]
\label{prop:rtsc-lattice-necessary}
设某子节点 $i_0$ 的间隙主要由格点/结构变量 $a_{i_0}$ 主导，且在允许控制集合下 $a_{i_0}$ 不可被作用（即 Definition~\ref{def:rtsc-lattice-control}
  的“晶格控制”条件不成立）。
若进一步存在常数 $c_{i_0}>0$ 使得在该不可作用子域内对任意 $m$ 都有 $g_{i_0}(m)\ge c_{i_0}$，则由 Theorem~\ref{thm:rtsc-jacgap-unactuated-lower-bound} 得 \RTSC{} 在该算子集合下不可控。
\end{proposition}

\begin{certificate}[证书：不可作用自由度导致子节点间隙有界下不去]
\label{cert:rtsc-lattice-necessary}
若 $g_{i_0}$ 对 $a_{i_0}$ 具有非平凡灵敏度且真实可达的 $a_{i_0}$ 被固定在某子域内，则 $g_{i_0}$ 的可达下界可为正；此时该下界作为“结构性障碍”可锁死总 JacGap 的下界。
\end{certificate}

\begin{proof}[证明（谓词因果链）]
由 Certificate~\ref{cert:rtsc-lattice-necessary} 得到谓词
\[
A:\ a_{i_0}\ \text{不可作用}\Rightarrow (\exists c_{i_0}>0)\ \forall m:\ g_{i_0}(m)\ge c_{i_0}.
\]
由 $A$ 触发 Theorem~\ref{thm:rtsc-jacgap-unactuated-lower-bound} 的前提，从而推出不可控结论。证毕。
\end{proof}

\subsection{命题：将“格点可调”落到切面包络 $m_{ij}$ 的硬条件}

\begin{proposition}[包络合格性：格点自由度应放入 $m_{ij}$ 并被 $\mathcal{M}_{\mathrm{lat}}$ 有效作用]
\label{prop:rtsc-envelope-lattice-qualified}
沿用第 10 章的切面包络口径：一阶切面 $s_i^{(j)}$ 表示路线/管线选择，二阶包络自由度 $m_{ij}$（Definition~\ref{def:section-envelope}、Proposition~\ref{prop:second-order-envelope}）表征该路线邻域的细粒度可控扰动族。
称包络 $m_{ij}$ \textbf{合格}，若同时满足：
\[
\text{(i) }a\in m_{ij},\qquad
\text{(ii) }\mathcal{M}_{\mathrm{lat}}\ \text{在该包络上有效作用}.
\]
若包络不合格，则该切面邻域的搜索等价于在“固定晶格的子空间”内做局部优化，易产生“路线偏差自洽”的幻觉，并可能由 Proposition~\ref{prop:rtsc-lattice-necessary} 锁死 JacGap 下界。
\end{proposition}

\begin{certificate}[证书：包络自由度显式化 + 格点可作用 $\Rightarrow$ 避免固定晶格子空间偏航]
\label{cert:rtsc-envelope-lattice-qualified}
（1）由第 10 章的包络塔口径（Theorem~\ref{thm:bundle-tower-closure}），$m_{ij}$ 用于显式化容易被忽略的细粒度自由度；
（2）由 Proposition~\ref{prop:rtsc-lattice-necessary}，若晶格自由度不可作用，则可能出现子节点 JacGap 的结构性下界；
两者合取给出：将晶格自由度显式放入包络并纳入可控作用，是抑制固定晶格子空间偏航的必要工程接口。
\end{certificate}

\begin{proof}[证明（谓词因果链）]
由 Certificate~\ref{cert:rtsc-envelope-lattice-qualified}，取谓词链：
\[
\begin{aligned}
A:&\ m_{ij}\ \text{显式化细粒度自由度},\\
B:&\ a\ \text{若不可作用则 }g_{i_0}\text{可能有正下界},\\
C:&\ a\in m_{ij}\wedge \mathcal{M}_{\mathrm{lat}}\ \text{有效}\Rightarrow \text{避免“固定晶格子空间”搜索}.
\end{aligned}
\]
由 $A$ 与 $B$ 得到：若不将 $a$ 纳入 $m_{ij}$ 且不可作用，则可能锁死 JacGap 下界并诱发偏航；由 $C$ 给出合格性条件的工程解释。证毕。
\end{proof}

\subsection{命题：总 JacGap 可控性主命题（开篇重点的形式化版本）}

\begin{proposition}[主命题：总 JacGap 可控性判据（\PTOPB$_{\mathrm{RTSC}}$）]
\label{prop:rtsc-jacgap-controllability-main}
在 Definition~\ref{def:rtsc-jacgap-action}--Definition~\ref{def:rtsc-jacgap-controllable} 的口径下，
\[
\boxed{
\begin{aligned}
&\mathrm{RTSC}\ \text{在给定算子集合下可控}
\ \Longleftrightarrow\
\inf_{w\in\mathcal{W}_{\mathrm{adm}}}\ \mathcal{S}(w)\le \varepsilon
\\
&\Longleftrightarrow\
\text{每个子节点的 }g_i\text{均存在可作用的格点/结构控制方向，}
\\
&\hspace{4em}\text{使 }\mathbf{g}\to 0.
\end{aligned}
}
\]
\end{proposition}

\begin{certificate}[证书：主命题由“定义等价 + 不可作用下界 + 局部可达条件”拼接得到]
\label{cert:rtsc-jacgap-controllability-main}
（1）由 Definition~\ref{def:rtsc-jacgap-controllable}，可控性以 $\inf_w\mathcal{S}[w]\le \varepsilon$ 定义；
（2）由 Theorem~\ref{thm:rtsc-jacgap-unactuated-lower-bound} 与 Proposition~\ref{prop:rtsc-lattice-necessary}，
  若存在不可作用子节点（尤其缺少格点/结构控制自由度）导致 $g_i$ 有正下界，则 $\inf_w\mathcal{S}[w]$ 被锁死在正下界之上；
（3）由 Proposition~\ref{prop:rtsc-jacgap-rank-sufficient}，若每个子节点的间隙向量对可控参数具有足够灵敏度（满秩），则 $\mathbf{g}$ 可在局部被压到任意小，从而使
  $G_{\mathrm{tot}}$ 与 $\mathcal{S}[w]$ 可被压到阈值以内（再叠加证书门槛得到可证书化口径）。
\end{certificate}

\begin{proof}[证明（谓词因果链）]
由 Certificate~\ref{cert:rtsc-jacgap-controllability-main}：
\begin{itemize}
  \item（定义等价）由（1）可得：$\mathrm{Controllable}\ \Leftrightarrow\ \inf_w\mathcal{S}[w]\le \varepsilon$；
  \item（必要性）由（2）可得：若缺少可作用的格点/结构控制方向使某些 $g_i$ 形成正下界，则 $\inf_w\mathcal{S}[w]$ 被下界锁死，故不可控；
  \item（充分性）由（3）可得：若每个子节点均存在足够灵敏度的可控方向使 $\mathbf{g}\to 0$，则 $G_{\mathrm{tot}}$ 可被压到任意小，从而 $\mathcal{S}[w]$ 可被压到阈值以内，
    得到可控性。
\end{itemize}
三者合取即得到命题中的等价链。证毕。
\end{proof}

\begin{remark}[与丛塔/切面包络的接口：把“格点可调”固化为插件/包络门槛]
Proposition~\ref{prop:rtsc-envelope-lattice-qualified} 给出一个直接可执行的工程口径：对任意路线切面 $s_i^{(j)}$，
其包络自由度 $m_{ij}$ 必须显式包含晶格/畸变/应变/界面错配等变量，并且存在可作用的 $\mathcal{M}_{\mathrm{lat}}$。
否则该切面附近的搜索只是在“固定晶格子空间”内做治疗，极易触发“路线偏差自洽”的幻觉预警（Corollary~\ref{cor:hallucination-certificate}），应回退到更粗层做诊断扩展，而不是继续加码局部优化。
\end{remark}

%%%%%%%%%%%%%%%%%%%%%%%%%%%%%%%%%%%%%%%%%%%%%%%%%%%%%%%%%%%%
\section{扰动抑制与回归控制：同伦扭曲的快慢层机制}
\label{sec:rtsc-disturbance-regression}
%%%%%%%%%%%%%%%%%%%%%%%%%%%%%%%%%%%%%%%%%%%%%%%%%%%%%%%%%%%%

\subsection{形式逻辑表达约定}

\begin{remark}[控制论口径：把 \RTSC{} 写成最小总 JacGap 的扰动抑制问题]
本章将第 11 章的“总雅可比间隙 $G_{\mathrm{tot}}$ 可控性判据”进一步组织为一个可用于 \PTOPB/GRL 的控制论接口：
\begin{enumerate}
  \item \textbf{目标对象：}将“室温超导”写成最小化势能 $V(m)$ 的控制目标，其中 $V$ 的主导项为 $G_{\mathrm{tot}}$，并把 $\mathrm{Res}/\mathrm{Unc}
    /\mathrm{Cert}$ 作为硬约束/强惩罚项；
  \item \textbf{动力学对象：}把材料状态、环境与控制协议都提升为控制变量；外部扰动（例如极端电磁/供电/散热/噪声冲击）作为不可控输入；
  \item \textbf{机制对象：}用“同伦扭曲”（主丛纤维方向的快速规范变换/切面切换）+“基底回归”（材料/环境的慢层调节）组成扰动后的回归机制。
\end{enumerate}
\end{remark}

\subsection{定义：最小总 JacGap 的势能 $V(m)$（把“最小电阻”严格化）}

\begin{definition}[势能：总 JacGap + 残差/不确定度/证书门槛]
\label{def:rtsc-potential-V}
令全量状态为 $m(t)\in M_\ast$（Definition~\ref{def:rtsc-jacgap-base-subnodes}），并沿用总雅可比间隙 $G_{\mathrm{tot}}(m)
  $（Definition~\ref{def:rtsc-total-jacgap}）。
定义 Lyapunov 型势能（可接受度势能）为
\[
V(m)
:=
G_{\mathrm{tot}}(m)
\,+\,
\lambda\,\mathrm{Res}(m)
\,+\,
\eta\,\mathrm{Unc}(m)
\,+\,
M\cdot\bigl(1-\mathrm{Cert}(m)\bigr),
\qquad \lambda,\eta,M\ge 0.
\]
势能越小，表示越接近“低总 JacGap + 可复现 + 跨评估器一致”的可接受状态。
\end{definition}

\begin{remark}[与第 11 章的接口]
Definition~\ref{def:rtsc-potential-V} 中的 $V$ 可视为第 11 章 Action 主导项的瞬时代价：
将 $V(m(t))$ 沿轨线积分即可得到 “总 JacGap + 证书门槛” 对路径权重的主导影响（Definition~\ref{def:rtsc-jacgap-action}）。
本章进一步把“可控材料/可控环境”的带宽与成本写成额外正则项，从而得到更标准的扰动抑制—回归控制形式。
\end{remark}

\subsection{定义：受扰动控制系统与可控集合逼近项}

\begin{definition}[受扰动动力学：$\dot m=f(m)+B(m)u+E(m)d$]
\label{def:rtsc-disturbed-dynamics}
称 \RTSC{} 的 \PTOPB{} 环境在抽象层满足受扰动控制系统形式，若存在向量场 $f$ 与算子 $B,E$ 使得
\[
\dot m(t)=f(m(t)) + B(m(t))\,u(t) + E(m(t))\,d(t),
\]
其中
\[
u(t)=\bigl(u_{\mathrm{mat}}(t),u_{\mathrm{env}}(t),u_{\mathrm{ctrl}}(t)\bigr)\in\mathcal{U}_{\mathrm{adm}}(m(t))
\]
为可行控制（材料/环境/协议的联合可控变量），$d(t)$ 为外部扰动（不可控输入）。
\end{definition}

\begin{definition}[扩展 Action：控制代价与可控集合逼近]
\label{def:rtsc-control-action}
给定可控集合 $\mathcal{C}\subset M_\ast$（表示“可制造/可运行/可测量”的可行域），定义距离逼近项
\[
\mathrm{Dist}(m,\mathcal{C}):=\inf_{c\in\mathcal{C}}\ \|m-c\|.
\]
给定正定矩阵 $R\succ 0$，定义控制代价 $\|u\|_R^2:=u^\top R u$。
定义扩展 Action（用于 GRL/path-integral 的扰动抑制—回归优化）为
\[
\mathcal{S}_{\mathrm{ctrl}}[w]
:=
\int_{0}^{T}\Bigl(
V(m(t))
\,+\,
\|u(t)\|_R^2
\,+\,
\mathrm{Dist}(m(t),\mathcal{C})
\Bigr)\,dt.
\]
其中 $w$ 表示由算子词与控制信号共同编码的可执行路径（与第 11 章的 $w\in\mathcal{W}_{\mathrm{adm}}$ 口径一致化）。
\end{definition}

\subsection{命题：扰动尖峰导致权重重分配（$V$ 的尖峰 $\Rightarrow$ 路径积分再加权）}

\begin{proposition}[扰动尖峰：外部扰动可导致 $V$（尤其 $G_{\mathrm{tot}}$）突增]
\label{prop:rtsc-disturbance-spike}
在 Definition~\ref{def:rtsc-disturbed-dynamics} 下，若扰动 $d(t)$ 在某时段显著增大并使得
\[
\langle \nabla V(m(t)),\,E(m(t))\,d(t)\rangle\ \text{在该时段主导},
\]
则 $V(m(t))$ 可出现尖峰（spike），特别地可表现为某些子节点间隙 $g_i(m(t))$ 的突增（环境块、控制执行块、证书口径块等被冲击）。
在 GRL/path-integral 口径下，这将导致 $\exp\!\bigl(-\beta\,\mathcal{S}_{\mathrm{ctrl}}[w]\bigr)$ 的权重结构发生重分配，从而触发在线重规划。
\end{proposition}

\begin{certificate}[证书：链式法则给出 $\dot V$ 的扰动项]
\label{cert:rtsc-disturbance-spike}
若 $V$ 可微，则沿动力学有
\[
\dot V(m(t))=\langle \nabla V(m(t)),\,f(m(t))+B(m(t))u(t)\rangle+\langle \nabla V(m(t)),\,E(m(t))d(t)\rangle.
\]
当第二项在某时段占主导且为正时，$V$ 将在该时段上升并可形成尖峰；由于扩展 Action 将 $V$ 沿时间积分并进入路径积分权重，因此尖峰将直接改变路径权重排序。
\end{certificate}

\begin{proof}[证明（谓词因果链）]
由 Certificate~\ref{cert:rtsc-disturbance-spike}：
\[
\begin{aligned}
A:&\ \dot V=\langle\nabla V,f+Bu\rangle+\langle\nabla V,Ed\rangle,\\
B:&\ \langle\nabla V,Ed\rangle\ \text{主导且为正}\Rightarrow V\ \text{上升并可尖峰},\\
C:&\ \mathcal{S}_{\mathrm{ctrl}}[w]=\int_0^T(\cdots+V+\cdots)\,dt.
\end{aligned}
\]
由 $A\wedge B$ 得到 $V$ 的尖峰；由 $C$ 得到尖峰将放大对应路径的 Action 并改变 $\exp(-\beta\mathcal{S}_{\mathrm{ctrl}})$ 的相对权重排序，从而触发在线重规划。证毕。
\end{proof}

\subsection{定义：同伦扭曲（纤维规范变换/切面切换）作为快层动作}

\begin{definition}[同伦扭曲：时间依赖切面与纤维规范变换]
\label{def:rtsc-homotopy-twist}
在基准主丛 $\pi_\star:P_\star\to M_\ast$（Definition~\ref{def:first-order-benchmark-bundle}）与表示结构对象 $G_{\mathrm{rep}}
  $（Definition~\ref{def:g-rep}）下，
令 $s(t):U\subseteq M_\ast\to P_\star$ 为时间依赖截面（表示“当前采用的路线/管线/表示/控制协议族”）。
若允许选择时间依赖的纤维规范变换 $g(t)\in G_{\mathrm{rep}}$，并定义新截面
\[
s'(t):=s(t)\cdot g(t),
\]
则称 $(s(t)\to s'(t))$ 为一类\textbf{同伦扭曲}（快层动作）：在不必立刻大幅改变基底状态 $m(t)$ 的前提下，先在纤维方向快速重整表示、管线细节、容错参数与证书口径，以把系统拉回到可修复的切面/包络邻域内。
\end{definition}

\begin{remark}[与切面包络 $m_{ij}$ 的对应关系]
同伦扭曲可在第 10 章的“切面包络”口径下具体化为：在给定路线切面 $s_i^{(j)}$ 的包络邻域 $E_{ij}$ 内，
通过调节包络自由度 $m_{ij}\in \Xi_{ij}$（Definition~\ref{def:section-envelope}）实现快速鲁棒化（例如切换更保守的证书口径、增大容错裕度、启用更鲁棒的求解管线等）。
该快层动作的目的首先是压住 $V$ 中与 $\mathrm{Cert}/\mathrm{Res}/\mathrm{Unc}$ 相关的爆炸项，然后再由慢层控制把 $G_{\mathrm{tot}}$ 拉回低值区域。
\end{remark}

\subsection{定理（公理降级）：扰动下回归定理（ISS 型界）}

\begin{theorem}[扰动下回归（公理降级）：快层同伦扭曲 + 慢层基底回归蕴含阈值回归界]
\label{thm:rtsc-disturbance-regression}
在 Definition~\ref{def:rtsc-potential-V} 的势能 $V$ 下，考虑离散时间步 $n=0,1,2,\dots$ 的采样系统。
若存在一类可行控制策略（包含快层同伦扭曲与慢层基底回归）使得对某常数 $\kappa\in(0,1)$、$c\ge 0$ 与任意扰动序列 $\{d_n\}$ 都满足
\[
V_{n+1}\le (1-\kappa)\,V_n + c\,\|d_n\|,
\qquad V_n:=V(m_n),
\]
则对任意 $N\in\mathbb{N}$ 有
\[
V_N\le (1-\kappa)^N V_0 + c\sum_{j=0}^{N-1}(1-\kappa)^{N-1-j}\|d_j\|.
\]
因此：当扰动有界（$\|d_n\|\le d_{\max}$）时，$V_n$ 进入并保持在阈值邻域
\[
V_n\ \lesssim\ \frac{c}{\kappa}\,d_{\max}
\]
的尺度内（上式为几何级数上界口径），表现为“扰动下回归可控”的 ISS 型性质。
\end{theorem}

\begin{certificate}[数学归纳法证书：按时间步 $N$ 归纳（几何级数界）]
\label{cert:rtsc-disturbance-regression-induction}
定义谓词 $P(N)$ 为：
\[
P(N):\ V_N\le (1-\kappa)^N V_0 + c\sum_{j=0}^{N-1}(1-\kappa)^{N-1-j}\|d_j\|.
\]
证书化要求为：
\[
P(0)\ \text{成立},\qquad (\forall N\in\mathbb{N})\ \bigl(P(N)\Rightarrow P(N+1)\bigr).
\]
\end{certificate}

\begin{proof}[证明（数学归纳法证书；谓词因果链）]
由 Certificate~\ref{cert:rtsc-disturbance-regression-induction}：
\begin{itemize}
  \item 基步 $N=0$：右侧和式为空，$V_0\le V_0$ 成立，故 $P(0)$ 成立；
  \item 归纳步：假设 $P(N)$ 成立。由定理前提的递推不等式得
  \[
  V_{N+1}\le (1-\kappa)V_N + c\|d_N\|.
  \]
  将归纳假设中对 $V_N$ 的上界代入并整理，得到
  \[
  V_{N+1}\le (1-\kappa)^{N+1}V_0 + c\sum_{j=0}^{N}(1-\kappa)^{N-j}\|d_j\|,
  \]
  即 $P(N+1)$ 成立。
\end{itemize}
由数学归纳法得对任意 $N$ 成立，从而得到所述几何级数界。证毕。
\end{proof}

\subsection{推论：扰动消失时的零点逼近与“回到低总 JacGap 邻域”}

\begin{corollary}[扰动消失 $\Rightarrow$ 势能回落；证书可满足 $\Rightarrow$ 逼近 $G_{\mathrm{tot}}\to 0$]
\label{cor:rtsc-disturbance-vanish}
在 Theorem~\ref{thm:rtsc-disturbance-regression} 的条件下：
\begin{enumerate}
  \item 若扰动最终消失（存在 $N_0$ 使 $d_n=0$ 对 $n\ge N_0$），则 $V_n$ 以几何速率回落到 0 的邻域；
  \item 若进一步存在可行策略使 $\mathrm{Cert}(m_n)=1$ 在回归阶段保持（或最终稳定为 1），则 $V_n\to 0$ 蕴含 $G_{\mathrm{tot}}(m_n)\to 0$，从而逼近 \RTSC{}
    的“总 JacGap 零点 + 证书门槛”目标（Definition~\ref{def:rtsc-predicate-jacgap}）。
\end{enumerate}
\end{corollary}

\begin{certificate}[证书：$V$ 的分解给出 $V\to 0\Rightarrow G_{\mathrm{tot}}\to 0$]
\label{cert:rtsc-disturbance-vanish}
由 Definition~\ref{def:rtsc-potential-V}，$V$ 是由 $G_{\mathrm{tot}}$ 与非负项 $\lambda\,\mathrm{Res}$、$\eta\,\mathrm{Unc}$、
  $M(1-\mathrm{Cert})$ 的和构成。
因此在 $\mathrm{Cert}=1$ 可维持（或最终稳定）且 $V\to 0$ 时，各非负项均趋于 0，从而 $G_{\mathrm{tot}}\to 0$。
\end{certificate}

\begin{proof}[证明（谓词因果链）]
由 Certificate~\ref{cert:rtsc-disturbance-vanish}，取谓词链：
\[
\begin{aligned}
A:&\ d_n=0\ \text{最终成立}\Rightarrow V_{n+1}\le (1-\kappa)V_n,\\
B:&\ V_{n+1}\le (1-\kappa)V_n\Rightarrow V_n\to 0,\\
C:&\ \mathrm{Cert}=1\ \text{保持或最终稳定}\wedge V\to 0\Rightarrow G_{\mathrm{tot}}\to 0.
\end{aligned}
\]
由 $A\Rightarrow B$ 得到势能回落；由 $C$ 得到零点逼近结论。证毕。
\end{proof}

\begin{remark}[与“格点可变调节”条件的关系]
Theorem~\ref{thm:rtsc-disturbance-regression} 的递推不等式前提，在工程上对应“主导子节点的间隙均可被真实可作用的控制方向压下去”。
特别地，若某主导子节点的 JacGap 受晶格/结构障碍主导而缺少可作用的晶格控制自由度，则可能出现正下界锁死（Proposition~\ref{prop:rtsc-lattice-necessary}），从而无法满足回归不等式。
因此：第 11 章的“格点可变调节/晶格控制”不仅是静态可控性必要接口，也是扰动下回归可控的必要工程前提。
\end{remark}

\clearpage
%%%%%%%%%%%%%%%%%%%%%%%%%%%%%%%%%%%%%%%%%%%%%%%%%%%%%%%%%%%%
\section{纯数卷口径：法空间测地线与同伦扭曲回归}
\label{sec:puremath-lawspace-geodesic-homotopy-twist}
%%%%%%%%%%%%%%%%%%%%%%%%%%%%%%%%%%%%%%%%%%%%%%%%%%%%%%%%%%%%

\subsection{形式逻辑表达约定}

\begin{remark}[翻译目标：把“扰动尖峰 + 同伦扭曲回归”转写为法空间几何对象]
本章把第 12 章的“扰动尖峰 + 同伦扭曲回归”严格翻译到纯数卷口径的对象与术语中：
\begin{enumerate}
  \item \textbf{法空间对象：}用 $\GZLS=(\mathfrak{L};J,\mathrm{Loc},\Sigma)$ 表示法空间，其中 $J$ 给出度量（从而定义测地线），$\mathrm{Loc}$
    给出可局部化/可操作域，$\Sigma$ 给出连续–离散切换阈值；
  \item \textbf{联络与曲率对象：}用 $\GZLOC$ 表示法联络（law-connection），用 $\GZLCurv=\mathcal{F}_{\mathrm{law}}$ 表示法曲率（law-curvature），
    必要时引入三层联络 $\GZTLC=\mathcal{A}$ 以在对象层–参数层–法层之间实现一致的水平运输（horizontal law-flow）；
  \item \textbf{缺陷势能对象：}用 $\JacGap(L)$ 作为 Jacobiator-gap 证书，并将“最小电阻”翻译为“最小缺陷势能” $\Phi(L):=\|\JacGap(L)\|_{\mathrm{def}}$；

  \item \textbf{回归机制对象：}用法空间上的测地线 $\gamma$ 选取回归方向（方向优先），并在固定测地线基底之上进行同伦扭曲（纤维方向修复高阶括号数据），从而把系统拉回到低缺陷扇区。
\end{enumerate}
\end{remark}

\subsection{定义：法空间 $\GZLS=(\mathfrak{L};J,\mathrm{Loc},\Sigma)$}

\begin{definition}[法空间（law-space）]
\label{def:law-space-gzls}
定义法空间为结构元组
\[
\GZLS=(\mathfrak{L};J,\mathrm{Loc},\Sigma),
\]
其中：
\begin{enumerate}
  \item $J$ 为法空间的度量数据（用于定义测地线与最短回归方向）；
  \item $\mathrm{Loc}$ 为可局部化/可操作域（用于编码“控制只能在可操作域内发生”）；
  \item $\Sigma$ 为连续–离散切换阈值（用于编码“跨阈触发同伦修复/严格化机制接管”）。
\end{enumerate}
\end{definition}

\subsection{定义：法联络 $\GZLOC$、法曲率 $\GZLCurv$ 与三层联络 $\GZTLC$}

\begin{definition}[法联络与法曲率]
\label{def:law-connection-curvature}
在法空间 $\GZLS$ 上，称 $\GZLOC$ 为\textbf{法联络}（law-connection），其曲率记为
\[
\GZLCurv=\mathcal{F}_{\mathrm{law}}
\]
并称为\textbf{法曲率}（law-curvature）。
若需要在对象层–参数层–法层之间一致地进行水平运输，可引入三层联络
\[
\GZTLC=\mathcal{A},
\]
并用其提升版本控制“水平法流”（horizontal law-flow）的传输一致性。
\end{definition}

\subsection{定义：法层算子系统 $L$ 与 Jacobiator-gap 证书 $\JacGap(L)$}

\begin{definition}[法层算子系统：同伦扇区中的 $L_\infty$ 口径对象]
\label{def:law-level-operator-system}
在纯数卷 Part III 的口径下，将“系统”抽象为法空间同伦扇区中的一个法层算子系统
\[
L\in (\GZLS)_{\mathrm{homo}},
\qquad
L\ \text{携带}\ (V^{-1}\oplus V^{0},\ell_1,\ell_2,\ell_3^H).
\]
其中 $\ell_2$ 是二元括号（graded commutator），$\ell_3^H$ 用于编码 Jacobiator 的“可控失败”，例如以关系
\[
\mathrm{Jac}_{\ell_2}=\ell_1\circ \ell_3
\]
刻画“严格 Jacobi 失配的同伦可修复性”。
\end{definition}

\begin{definition}[Jacobiator-gap 证书与缺陷范数]
\label{def:jacgap-certificate}
称映射
\[
\JacGap:(\GZLS)_{\mathrm{homo}}\to \mathbb{R}_{\ge 0}^k,
\qquad
\JacGap(L)=(c_1(L),\dots,c_k(L))
\]
为 Jacobiator-gap 证书（\textbf{JacGap}），并给定一个缺陷聚合范数 $\|\cdot\|_{\mathrm{def}}$（例如加权 $\ell_1$ 或谱半径型聚合）。
定义缺陷势能（defect potential）为
\[
\Phi(L):=\|\JacGap(L)\|_{\mathrm{def}}.
\]
\end{definition}

\subsection{定理（公理降级）：$\JacGap(L)=0$ 当且仅当可严格化（低阶可观测量受保护）}

\begin{theorem}[严格化判据（公理降级）：JacGap 零点等价于可严格化]
\label{thm:jacgap-zero-iff-strictifiable}
设给定一个“低阶可观测量投影”谓词 $\mathrm{Obs}_{\le r}$（表示被声明为 protected 的低阶行为）。
若 $\JacGap$ 满足下列口径：
\begin{enumerate}
  \item 若 $\JacGap(L)=0$，则存在 $L_{\mathrm{strict}}$ 使 $L$ 在保持 $\mathrm{Obs}_{\le r}$ 不变的前提下可被提升/严格化到 $L_{\mathrm{strict}}
    $；
  \item 若 $L$ 可在保持 $\mathrm{Obs}_{\le r}$ 不变的前提下严格化到某个 $L_{\mathrm{strict}}$，则 $\JacGap(L)=0$；
\end{enumerate}
则有等价：
\[
\JacGap(L)=0
\quad\Longleftrightarrow\quad
\exists L_{\mathrm{strict}}\ \bigl(\mathrm{Obs}_{\le r}(L)=\mathrm{Obs}_{\le r}(L_{\mathrm{strict}})\bigr).
\]
\end{theorem}

\begin{certificate}[数学归纳法证书：按同伦修复步数 $n$ 归纳]
\label{cert:jacgap-strictification-induction}
将“严格化”视为由有限序列同伦动作（homotopy moves）实现的修复过程。
定义谓词 $P(n)$ 为：“存在 $n$ 步同伦动作把 $L$ 变形为 $L^{(n)}$，且 $\mathrm{Obs}_{\le r}$ 不变，并且 $\JacGap(L^{(n)})=0$”。
证书化要求为：
\[
P(0)\ \text{成立当且仅当 }\JacGap(L)=0,\qquad
(\forall n\in\mathbb{N})\ \bigl(P(n)\Rightarrow P(n+1)\bigr).
\]
\end{certificate}

\begin{proof}[证明（数学归纳法证书；谓词因果链）]
由 Certificate~\ref{cert:jacgap-strictification-induction}，取谓词链：
\[
\begin{aligned}
A:&\ \JacGap(L)=0\Rightarrow \exists L_{\mathrm{strict}}\ \bigl(\mathrm{Obs}_{\le r}\ \text{不变}\bigr),\\
B:&\ \exists L_{\mathrm{strict}}\ \bigl(\mathrm{Obs}_{\le r}\ \text{不变}\bigr)\Rightarrow \JacGap(L)=0.
\end{aligned}
\]
由 $A$ 得到“零点 $\Rightarrow$ 可严格化”，由 $B$ 得到“可严格化 $\Rightarrow$ 零点”，两者合取即得等价结论。证毕。
\end{proof}

\subsection{定理（公理降级）：法曲率尖峰导致 $\Phi(L(t))$ 的扰动尖峰}

\begin{theorem}[扰动尖峰（公理降级）：曲率–同伦控制不等式在尖峰输入下放大 JacGap]
\label{thm:law-curvature-spike-jacgap}
设外生扰动 $d(t)$ 通过对象层到法层的嵌入映射作用到 $L(t)\in(\GZLS)_{\mathrm{homo}}$ 上，并沿 $\GZLOC/\GZTLC$ 诱导的水平法流演化。
若存在可证书化的控制不等式使得
\[
\frac{d}{dt}\,\Phi(L(t))
\ \le\
\alpha\,\|\GZLCurv(L(t))\|
\,+\,
\beta\,\|d(t)\|,
\qquad \alpha,\beta\ge 0,
\]
则当 $\|\GZLCurv(L(t))\|$ 在短时窗口出现尖峰（spike）时，$\Phi(L(t))$ 的增长上界在该窗口被放大，从而表现为“缺陷势能尖峰”（对应第 12 章的 $V$ 尖峰口径）。
\end{theorem}

\begin{certificate}[证书：JacGap 沿水平法流的变化受曲率与同伦项控制]
\label{cert:jacgap-controlled-by-curvature}
纯数卷口径要求：沿任意水平法流 $t\mapsto L(t)$，$\JacGap(L(t))$ 的变化由法曲率与同伦项控制；
因此可将其汇总为对缺陷势能 $\Phi$ 的微分不等式上界（其中扰动 $d(t)$ 进入同伦项或外生项）。
\end{certificate}

\begin{proof}[证明（谓词因果链）]
由 Certificate~\ref{cert:jacgap-controlled-by-curvature}，取谓词链：
\[
\begin{aligned}
A:&\ \dot\Phi\ \text{受 }\|\GZLCurv\|\ \text{与 }\|d\|\ \text{控制},\\
B:&\ \|\GZLCurv\|\ \text{出现尖峰}\Rightarrow \dot\Phi\ \text{上界增大},\\
C:&\ \dot\Phi\ \text{上界增大}\Rightarrow \Phi\ \text{在该窗口更易出现尖峰}.
\end{aligned}
\]
由 $A\wedge B\Rightarrow C$ 得到结论。证毕。
\end{proof}

\subsection{定义：测地线回归 $\gamma$ 与“基底锁定”的同伦扭曲 $H$}

\begin{definition}[法空间测地线：由度量 $J$ 选取回归方向]
\label{def:lawspace-geodesic}
在度量 $J$ 下，称 $\gamma:[0,T]\to\GZLS$ 为法空间测地线，若其为 $J$ 的测地线解。
在扰动后回归问题中，可取 $\gamma(0)=L_{\mathrm{shock}}$，并要求 $\gamma(T)$ 落入某个低缺陷邻域（例如 $\Phi\le \varepsilon$ 的区域）。
\end{definition}

\begin{definition}[沿测地线的同伦扭曲：基底锁定 + 纤维内修复]
\label{def:geodesic-homotopy-twist}
给定测地线 $\gamma$，称映射
\[
H:[0,1]\times[0,T]\to (\GZLS)_{\mathrm{homo}}
\]
为“沿测地线的同伦扭曲”，若满足：
\begin{enumerate}
  \item \textbf{基底锁定（base locked）：}对每个 $s\in[0,1]$，$\mathrm{Obs}_{\le r}(H(s,t))$ 的投影与 $\gamma(t)$ 一致（低阶可观测量沿同一回归方向推进）；
  \item \textbf{端点条件：}$H(0,t)=L_\gamma(t)$ 为沿 $\GZLOC/\GZTLC$ 的水平提升，且 $H(1,t)=:L_\gamma^{\mathrm{tw}}(t)$；
  \item \textbf{缺陷单调性：}$\Phi(L_\gamma^{\mathrm{tw}}(t))\le \Phi(L_\gamma(t))$，必要时要求严格下降以实现回归。
\end{enumerate}
该定义把“扭曲”限定为纤维方向上的同伦修复：在不改变被保护的低阶可观测量前提下，允许调整高阶同伦数据（典型为 $\ell_3$ 及必要时的更高括号）。
\end{definition}

\subsection{命题：测地线选择（方向）与同伦扭曲（构造）在路径积分中可统一编码}

\begin{proposition}[法空间路径积分：用 $\mathcal{Z}_{\mathrm{law}}$ 封装“最优回归”]
\label{prop:lawspace-path-integral}
对允许的回归对 $(\gamma,H)$，定义法作用量
\[
\mathcal{S}_{\mathrm{law}}(\gamma,H)
:=
\int_0^T\Bigl(
\Phi(L(t))
\,+\,
\alpha\,\|\GZLCurv(\gamma(t))\|
\,+\,
\mathrm{Penalty}_{\mathrm{Loc},\Sigma}(\gamma(t))
\Bigr)\,dt,
\qquad \alpha\ge 0,
\]
其中 $\mathrm{Penalty}_{\mathrm{Loc},\Sigma}$ 用于编码局部化约束 $\mathrm{Loc}$ 与跨阈代价 $\Sigma$。
定义 law-space path integral（GZ-LSPI）分配函数为
\[
\mathcal{Z}_{\mathrm{law}}
:=
\sum_{(\gamma,H)}\exp\!\bigl(-\beta\,\mathcal{S}_{\mathrm{law}}(\gamma,H)\bigr),
\qquad \beta>0.
\]
则在大 $\beta$ 极限下，最小 $\mathcal{S}_{\mathrm{law}}$ 的“测地线回归方向 + 同伦扭曲修复”对将主导概率质量，从而把“扰动后如何回归”封装为法空间上的可计算最优化对象。
\end{proposition}

\begin{certificate}[证书：Boltzmann 权重在大 $\beta$ 极限由最小作用量主导]
\label{cert:lawspace-path-integral}
对任意有限候选集合，$\exp(-\beta\mathcal{S})$ 在 $\beta\to\infty$ 时由最小 $\mathcal{S}$ 项主导；
若候选为可测族，则相同结论以 Laplace 原理口径表达（此处作为工作笔记证书口径使用）。
\end{certificate}

\begin{proof}[证明（谓词因果链）]
由 Certificate~\ref{cert:lawspace-path-integral}，取谓词链：
\[
\begin{aligned}
A:&\ \mathcal{Z}_{\mathrm{law}}=\sum \exp(-\beta\mathcal{S}_{\mathrm{law}}),\\
B:&\ \beta\to\infty\Rightarrow \exp(-\beta\mathcal{S})\ \text{集中于最小 }\mathcal{S}\ \text{项},\\
C:&\ \text{集中}\Rightarrow \text{最优回归由最小 }\mathcal{S}_{\mathrm{law}}\ \text{的 }(\gamma,H)\ \text{主导}.
\end{aligned}
\]
由 $A\wedge B\Rightarrow C$ 得证。证毕。
\end{proof}

\subsection{推论：纯数卷口径的“方向优先于构造”与第 12 章快慢层机制同型}

\begin{corollary}[方向优先（纯数卷版）：先选测地线方向，再做同伦扭曲修复]
\label{cor:puremath-direction-before-construction}
在 Proposition~\ref{prop:lawspace-path-integral} 的口径下：
\begin{enumerate}
  \item 选择测地线 $\gamma$ 对应“方向/诊断”层：在度量 $J$ 的几何意义下选取最短/最省能的回归路径；
  \item 在固定 $\gamma$ 上做同伦扭曲 $H$ 对应“构造/治疗”层：在保持低阶可观测量受保护的前提下，修复高阶同伦数据以降低缺陷势能 $\Phi$。
\end{enumerate}
因此第 12 章的“快层同伦扭曲 + 慢层基底回归”与本章的“沿测地线水平推进 + 基底锁定的同伦扭曲”在形式系统口径上同型：回归首先是方向问题（测地线/诊断），其次才是类内构造（扭曲/修复）。
\end{corollary}

\begin{certificate}[证书：由 GZ-LSPI 的最小作用量主导分解为“方向 + 修复”两层]
\label{cert:puremath-direction-before-construction}
由 Proposition~\ref{prop:lawspace-path-integral}，在大 $\beta$ 极限下，概率质量由最小作用量 $\mathcal{S}_{\mathrm{law}}$ 的候选对 $(\gamma,H)$
  主导。
其中 $\gamma$ 作为法空间回归方向（测地线/诊断层），而 $H$ 作为在固定方向下的同伦扭曲（修复层），用于在相对约束下压低缺陷势能 $\Phi$。
因此“先选 $\gamma$ 再选 $H$”是该最小化结构的直接分解口径。
\end{certificate}

\begin{proof}[证明（谓词因果链）]
由 Certificate~\ref{cert:puremath-direction-before-construction}，取谓词链：
\[
\begin{aligned}
A:&\ \beta\to\infty\Rightarrow (\gamma,H)\ \text{由最小 }\mathcal{S}_{\mathrm{law}}\ \text{项主导（Proposition~\ref{prop:lawspace-path-integral})},\\
B:&\ \mathcal{S}_{\mathrm{law}}\ \text{的参数拆分为方向变量 }\gamma\ \text{与修复变量 }H,\\
C:&\ A\wedge B\Rightarrow \text{求解顺序可表述为先选 }\gamma\ \text{（诊断/方向）再选 }H\ \text{（修复/构造）}.
\end{aligned}
\]
由 $C$ 得证。证毕。
\end{proof}

\begin{remark}[与第 11 章总 JacGap 判据的接口]
第 11 章把 \RTSC{} 的可控性写为“总 JacGap 零点 + 证书门槛”；本章将其翻译为纯数卷的缺陷势能 $\Phi(L)=\|\JacGap(L)\|_{\mathrm{def}}$
  与严格化判据（Theorem~\ref{thm:jacgap-zero-iff-strictifiable}）。
当外部扰动引发法曲率尖峰时，缺陷势能尖峰可由曲率–同伦控制不等式解释（Theorem~\ref{thm:law-curvature-spike-jacgap}）。
回归则由“测地线方向 + 同伦扭曲修复”统一编码（Proposition~\ref{prop:lawspace-path-integral}），并可作为跨任务复用的诊断—修复资产沉淀到 $\mathcal{Z}_{\mathrm{law}}
  $ 的统计结构中。
\end{remark}

\clearpage
%%%%%%%%%%%%%%%%%%%%%%%%%%%%%%%%%%%%%%%%%%%%%%%%%%%%%%%%%%%%
\section{工业几何体工程设计基准：扩展扇区、包络丛与条件可行的逆向生成}
\label{sec:industrial-geometry-benchmark-inverse}
%%%%%%%%%%%%%%%%%%%%%%%%%%%%%%%%%%%%%%%%%%%%%%%%%%%%%%%%%%%%

\subsection{形式逻辑表达约定}

\begin{remark}[本章目标：在 $\GZLS$ 语言下引入工业工程基准并给出可行性判据]
本章在第 10 章的“切面包络/包络塔”（Definition~\ref{def:section-envelope}）与第 13 章的“法空间测地线 + 同伦扭曲 + GZ-LSPI（$\mathcal{Z}_{\mathrm{law}}$）
  ”（Proposition~\ref{prop:lawspace-path-integral}）口径之上，引入“大型工业几何体工程设计基准”，并用包络变量把“工业几何—仿真离散化—制造/装配—运行扰动—证书链”统一纳入同一套
  $\mathrm{Loc}/\Sigma$ 管理，从而实现\textbf{可证书化的逆向生成工程设计}。
本章按三条主线组织：
\begin{enumerate}
  \item \textbf{（A）扩展基底：}把 RTSC 基准底空间 $M_\ast$ 扩为工程‑RTSC 联合扇区 $M_\ast^{\mathrm{eng}}$；
  \item \textbf{（B）包络丛：}对每条工程路线切面 $s_i^{(j)}$ 引入二阶（必要时三阶）包络变量 $m_{ij}\in\Xi_{ij}$，使“工程细节”从隐变量变为可控自由度；
  \item \textbf{（C）条件可行判据：}用缺陷势能 $\Phi=\|\JacGap\|_{\mathrm{def}}$、阈值 $\Sigma$ 与沿 $J$‑测地线的同伦扭曲回归，给出逆向生成在形式系统层面可行、
    在计算层面条件可行的判据。
\end{enumerate}
\end{remark}

\subsection{定义：工程‑RTSC 联合扇区的扩展基底 $M_\ast^{\mathrm{eng}}$}

\begin{definition}[扩展基底：从 $M_\ast$ 到 $M_\ast^{\mathrm{eng}}$]
\label{def:eng-extended-base}
令 $M_\ast^{\mathrm{RTSC}}$ 表示已建立的 \RTSC{} 基准底空间（可取 Definition~\ref{def:rtsc-jacgap-base-subnodes} 的 $M_\ast$ 口径）。
引入三个工程相关法对象扇区：
\[
M_\ast^{\mathrm{Geom}}\ (\text{工业几何与边界条件}),
\qquad
M_\ast^{\mathrm{Grid}}\ (\text{运行与外部扰动}),
\qquad
M_\ast^{\mathrm{QC}}\ (\text{量子抗退相干}).
\]
定义工程‑RTSC 联合扇区（扩展基底）为
\[
\begin{aligned}
M_\ast^{\mathrm{eng}}
:=
&\,M_\ast^{\mathrm{RTSC}}
\times
M_\ast^{\mathrm{Geom}}
\\
&\times
M_\ast^{\mathrm{Grid}}
\times
M_\ast^{\mathrm{QC}}.
\end{aligned}
\]
并引入工程局部化域与阈值集合
\[
M_\ast^{\mathrm{eng}}\subseteq \mathrm{Loc}_{\mathrm{eng}}\subseteq \mathrm{Loc},
\qquad
\Sigma_{\mathrm{eng}}\subseteq \Sigma,
\]
其中 $\mathrm{Loc}_{\mathrm{eng}}$ 表示“工业几何—仿真离散化—制造/装配—运行扰动—证书链”在工程意义下相互一致、可操作、可审计的可局部化域，
$\Sigma_{\mathrm{eng}}$ 表示允许的连续–离散切换阈值（例如故障态/保护态/切换态等的触发集合）。
\end{definition}

\begin{definition}[扩展一阶主丛：$\pi_\ast^{\mathrm{eng}}:P_\ast^{\mathrm{eng}}\to M_\ast^{\mathrm{eng}}$]
\label{def:eng-extended-benchmark-bundle}
在扩展基底 $M_\ast^{\mathrm{eng}}$ 上，定义工程基准一阶主丛为
\[
\pi_\ast^{\mathrm{eng}}:P_\ast^{\mathrm{eng}}\to M_\ast^{\mathrm{eng}}.
\]
其表示冗余结构对象取为（半直积仅表示“组合并带有相容作用”，不要求固定某种具体群结构）
\[
\begin{aligned}
G_{\mathrm{rep}}^{\mathrm{eng}}
:=
&\,G_{\mathrm{rep}}^{\mathrm{RTSC}}
\rtimes
G_{\mathrm{rep}}^{\mathrm{Geom}}
\\
&\rtimes
G_{\mathrm{rep}}^{\mathrm{Sim}}
\rtimes
G_{\mathrm{rep}}^{\mathrm{Proc}}.
\end{aligned}
\]
其中 $G_{\mathrm{rep}}^{\mathrm{Geom}}$ 吸收 CAD/参数化的等价表示，
$G_{\mathrm{rep}}^{\mathrm{Sim}}$ 吸收网格/离散化/求解器/误差控制等仿真表示冗余，
$G_{\mathrm{rep}}^{\mathrm{Proc}}$ 吸收制造/装配路径的等价表述冗余。
在工作笔记的平凡化模型中可取 $P_\ast^{\mathrm{eng}}=M_\ast^{\mathrm{eng}}\times G_{\mathrm{rep}}^{\mathrm{eng}}$ 与
  $\pi_\ast^{\mathrm{eng}}(x,g)=x$，但本章后续判据并不依赖平凡化。
\end{definition}

\subsection{定理（公理降级）：扩展扇区闭合与水平法流可用性}

\begin{theorem}[扇区闭合（公理降级）：联合基底仍可局部化并支持水平法流]
\label{thm:eng-sector-closure}
在 Definition~\ref{def:eng-extended-base} 的设定下，若满足：
\begin{enumerate}
  \item 四个扇区分别可局部化：$M_\ast^{\mathrm{RTSC}},M_\ast^{\mathrm{Geom}},M_\ast^{\mathrm{Grid}},M_\ast^{\mathrm{QC}}\subseteq
    \mathrm{Loc}_{\mathrm{eng}}$；
  \item $\mathrm{Loc}_{\mathrm{eng}}$ 对扇区拼接闭合：若各分量在 $\mathrm{Loc}_{\mathrm{eng}}$ 内，则其直积点仍在 $\mathrm{Loc}_{\mathrm{eng}}
    $ 内；
  \item 法联络/三层联络的水平推进可按分量一致提升（Definition~\ref{def:law-connection-curvature} 的 $\GZLOC/\GZTLC$ 可对扇区拼接给出一致的水平法流）；
\end{enumerate}
则 $M_\ast^{\mathrm{eng}}\subseteq \mathrm{Loc}_{\mathrm{eng}}\subseteq \mathrm{Loc}$，并且工程基准主丛 $P_\ast^{\mathrm{eng}}$
  上存在与 $\GZLOC/\GZTLC$ 相容的水平法流，使得“工程‑RTSC 联合扇区”可以直接纳入第 13 章的测地线‑同伦扭曲‑路径积分（$\mathcal{Z}_{\mathrm{law}}$）统一语言中。
\end{theorem}

\begin{certificate}[数学归纳法证书：按扇区块数 $n\in\{1,2,3,4\}$ 归纳]
\label{cert:eng-sector-closure-induction}
将四个扇区按顺序记为 $M^{(1)},M^{(2)},M^{(3)},M^{(4)}$，并定义
\[
M^{(\le n)}:=\prod_{t=1}^{n} M^{(t)}.
\]
令谓词 $P(n)$ 为：“$M^{(\le n)}\subseteq \mathrm{Loc}_{\mathrm{eng}}$，且 $\GZLOC/\GZTLC$ 在 $M^{(\le n)}$ 上诱导一致的水平法流”。
证书化要求为：
\[
P(1)\ \text{成立},\qquad
(\forall n\in\{1,2,3\})\ \bigl(P(n)\Rightarrow P(n+1)\bigr).
\]
\end{certificate}

\begin{proof}[证明（数学归纳法证书；谓词因果链）]
由 Certificate~\ref{cert:eng-sector-closure-induction}：
\begin{itemize}
  \item 基步 $n=1$：由定理前提（1）得 $M^{(1)}\subseteq \mathrm{Loc}_{\mathrm{eng}}$，且水平法流在单扇区上可用，故 $P(1)$ 成立；
  \item 归纳步：假设 $P(n)$ 成立。由前提（1）得 $M^{(n+1)}\subseteq \mathrm{Loc}_{\mathrm{eng}}$；由前提（2）得
  $M^{(\le n+1)}=M^{(\le n)}\times M^{(n+1)}\subseteq \mathrm{Loc}_{\mathrm{eng}}$；由前提（3）得 $\GZLOC/\GZTLC$
    的水平推进可对拼接扇区一致提升，从而 $M^{(\le n+1)}$ 上水平法流仍可用，故 $P(n+1)$ 成立。
\end{itemize}
由数学归纳法得 $P(4)$ 成立，即 $M_\ast^{\mathrm{eng}}=M^{(\le 4)}\subseteq \mathrm{Loc}_{\mathrm{eng}}$ 且水平法流可用。证毕。
\end{proof}

\subsection{定义：工程路线切面与二阶包络变量 $m_{ij}\in \Xi_{ij}$}

\begin{definition}[工程基准的一阶切换与切面]
\label{def:eng-base-switch-section}
在工程基准主丛 $\pi_\ast^{\mathrm{eng}}:P_\ast^{\mathrm{eng}}\to M_\ast^{\mathrm{eng}}$（Definition~\ref{def:eng-extended-benchmark-bundle}）上，取一阶基底切换映射
\[
p_i: M_\ast^{\mathrm{eng}}\twoheadrightarrow M_i,
\]
并定义诱导投影 $\pi_i:=p_i\circ \pi_\ast^{\mathrm{eng}}:P_\ast^{\mathrm{eng}}\to M_i$（与 Definition~\ref{def:first-order-base-switching} 同型）。
称局部映射
\[
s_i^{(j)}:\ U_i^{(j)}\subseteq M_i\to P_\ast^{\mathrm{eng}}
\]
为工程路线切面（局部截面），若满足 $\pi_i\circ s_i^{(j)}=\mathrm{Id}_{U_i^{(j)}}$，其中索引 $j$ 标记“同一基底下的不同路线/管线/协议/表示选择”。
\end{definition}

\begin{definition}[工程二阶包络：切面邻域 $E_{ij}$ 与包络坐标分块]
\label{def:eng-envelope}
给定工程路线切面 $s_i^{(j)}$（Definition~\ref{def:eng-base-switch-section}）。
称 $E_{ij}\subseteq P_\ast^{\mathrm{eng}}$ 为其二阶包络邻域，若
\[
s_i^{(j)}\bigl(U_i^{(j)}\bigr)\subseteq E_{ij}.
\]
引入包络参数空间 $\Xi_{ij}$，并记包络变量 $m_{ij}\in\Xi_{ij}$。
若存在参数化映射
\[
\Phi_{ij}:\ U_i^{(j)}\times \Xi_{ij}\to E_{ij}
\quad\text{满足}\quad
\pi_i\circ \Phi_{ij}=\mathrm{pr}_{U_i^{(j)}},
\]
则可将“切面附近的纤维自由度”显式化为二阶细粒度化基底
\[
M_{ij}^{(2)}:=U_i^{(j)}\times \Xi_{ij}
\]
（与第 10 章的形态 2 写法同型）。

在工业几何体逆向生成场景中，将包络空间最小分块为：
\[
\Xi_{ij}
:=
\Xi_{ij}^{\mathrm{Geom}}
\times
\Xi_{ij}^{\mathrm{Sim}}
\times
\Xi_{ij}^{\mathrm{Proc}}.
\]
分别对应几何细粒度包络（形状/拓扑/间隙/连接方式等的可控微扰族）、仿真离散化包络（网格/基函数/误差控制/多物理耦合接口等）与制造装配包络（工艺序列、公差、批次变异等的等价类参数化）。
\end{definition}

\begin{lemma}[最小包络三块：几何/仿真/制造装配]
\label{lem:eng-envelope-minimal-three}
在 Definition~\ref{def:eng-envelope} 的口径下，若希望工程路线切面在扰动尖峰后仍能通过“同伦扭曲（纤维快层）+ 基底回归（慢层）”回到 $\mathrm{Loc}_{\mathrm{eng}}$
  且保持证书可审计，
则对每个关键路线切面 $s_i^{(j)}$，其包络参数空间 $\Xi_{ij}$ 至少应覆盖三类分块
$\Xi_{ij}^{\mathrm{Geom}},\Xi_{ij}^{\mathrm{Sim}},\Xi_{ij}^{\mathrm{Proc}}$。
\end{lemma}

\begin{certificate}[证书：三类包络对应三类“可控失败”来源]
\label{cert:eng-envelope-minimal-three}
三类包络分块分别吸收三类工程层面的可控失败来源：
\begin{enumerate}
  \item 几何细粒度：若缺少 $\Xi_{ij}^{\mathrm{Geom}}$，则“关键失效模式”可能落在切面不可达的几何邻域之外，导致缺陷势能无法被真实降低；
  \item 仿真离散化：若缺少 $\Xi_{ij}^{\mathrm{Sim}}$，则残差 $\mathrm{Res}$ 难以在跨网格/跨求解器意义下被证书化控制，容易出现“代价下降但残差/证书不改善”的自洽偏航；
  \item 制造装配：若缺少 $\Xi_{ij}^{\mathrm{Proc}}$，则生成的方案可能脱离可制造/可装配/可运行的 $\mathrm{Loc}_{\mathrm{eng}}$，从而证书门槛在工程落地时失效。
\end{enumerate}
\end{certificate}

\begin{proof}[证明（谓词因果链）]
由 Certificate~\ref{cert:eng-envelope-minimal-three}，取谓词链：
\[
\begin{aligned}
A:&\ \Xi_{ij}\ \text{缺少 }\Xi_{ij}^{\mathrm{Geom}}
\Rightarrow
\text{切面邻域无法覆盖关键几何失效模式},\\
B:&\ \Xi_{ij}\ \text{缺少 }\Xi_{ij}^{\mathrm{Sim}}
\Rightarrow
\mathrm{Res}/\mathrm{Cert}\ \text{难以跨离散化稳定},\\
C:&\ \Xi_{ij}\ \text{缺少 }\Xi_{ij}^{\mathrm{Proc}}
\Rightarrow
\text{方案易脱离 }\mathrm{Loc}_{\mathrm{eng}}\ \text{从而证书落空}.
\end{aligned}
\]
$A\vee B\vee C$ 任一成立都会破坏“扰动后可回归 + 可证书化”所需的可控自由度闭合性，因此三类分块均为最低要求。证毕。
\end{proof}

\subsection{命题：包络显式化抑制路线自洽偏航并支撑逆向生成}

\begin{proposition}[反偏航命题：包络显式化 $\Rightarrow$ 可证书化抑制自洽偏航]
\label{prop:eng-envelope-anti-drift}
对工程基准 $M_\ast^{\mathrm{eng}}$，若对关键路线切面均给出满足 Definition~\ref{def:eng-envelope} 的包络变量 $m_{ij}\in\Xi_{ij}$，
并将 $\mathrm{Res}/\mathrm{Unc}/\mathrm{Cert}$ 作为一等惩罚/门槛并入代价（参见 Proposition~\ref{prop:anti-drift-constraint} 与
  Corollary~\ref{cor:hallucination-certificate}），
则：
\begin{enumerate}
  \item 任何“仅在固定切面内自洽”的局部最优，都可通过对 $m_{ij}$ 的扫描与统计被暴露为残差/证书上的偏航信号；
  \item 一旦触发偏航预警，即可回退到诊断层（切换 $p_i$ 或扩张 $\Xi_{ij}$）进行再探索，而非继续在错误切面内加码治疗（与 Corollary~\ref{cor:envelope-anti-drift} 同型）。
\end{enumerate}
\end{proposition}

\begin{certificate}[证书：由包络塔与偏航预警口径直接推出]
\label{cert:eng-envelope-anti-drift}
\begin{enumerate}
  \item 由 Corollary~\ref{cor:envelope-anti-drift}，细粒度包络自由度的显式化提供“反偏航的可审计接口”；
  \item 由 Corollary~\ref{cor:hallucination-certificate}，固定路线内“代价下降但残差/证书不改善”构成可计算的偏航预警；
  \item 由 Definition~\ref{def:eng-envelope}，工程包络 $\Xi_{ij}^{\mathrm{Geom}}/\Xi_{ij}^{\mathrm{Sim}}/\Xi_{ij}
    ^{\mathrm{Proc}}$ 将几何/仿真/工艺三类最易隐藏的偏航自由度显式化，因此可直接纳入上述预警统计与再诊断触发。
\end{enumerate}
\end{certificate}

\begin{proof}[证明（谓词因果链）]
由 Certificate~\ref{cert:eng-envelope-anti-drift}，取谓词链：
\[
\begin{aligned}
A:&\ \text{包络显式化}\Rightarrow \text{偏航自由度可被审计与统计},\\
B:&\ \text{偏航统计满足预警条件}\Rightarrow \text{触发再诊断/再探索},\\
C:&\ \text{再诊断/再探索}\Rightarrow \text{避免在错误切面内自洽坍缩}.
\end{aligned}
\]
由 $A\wedge B\Rightarrow C$ 得到结论（1）（2）。证毕。
\end{proof}

\subsection{定理（公理降级）：逆向生成工程设计的条件可行性判据}

\begin{definition}[工程缺陷势能与可行性谓词]
\label{def:eng-feasibility-predicate}
在工程联合扇区上，仍以 $L\in(\GZLS)_{\mathrm{homo}}$ 表示法层算子系统（Definition~\ref{def:law-level-operator-system} 的同型口径），并以
\[
\Phi(L):=\|\JacGap(L)\|_{\mathrm{def}}
\]
表示缺陷势能（Definition~\ref{def:jacgap-certificate}）。
称工程逆向生成在阈值 $\varepsilon>0$ 下\textbf{可证书化可行}，若存在一条由 $\GZLOC/\GZTLC$ 允许的水平法流及其“沿测地线的同伦扭曲”对 $(\gamma,H)$（第 13 章口径），
  使得沿演化轨线满足：
\[
L(t)\in \mathrm{Loc}_{\mathrm{eng}},
\qquad
\Sigma_{\mathrm{eng}}\ \text{条件下允许切换},
\qquad
\Phi(L(T))\le \varepsilon,
\qquad
\mathrm{Cert}(L(T))=1.
\]
该谓词把“工程可落地”与“形式系统可严格化/可证书化”统一编码为同一条可审计判据。
\end{definition}

\begin{theorem}[条件可行（公理降级）：包络显式化 + 测地线同伦扭曲回归 $\Rightarrow$ 可证书化逆向生成]
\label{thm:eng-conditional-feasibility}
在 Definition~\ref{def:eng-extended-base}--Definition~\ref{def:eng-feasibility-predicate} 的设定下，若满足：
\begin{enumerate}
  \item 工程联合扇区闭合且水平法流可用（Theorem~\ref{thm:eng-sector-closure}）；
  \item 对关键路线切面均给出包络变量 $m_{ij}\in\Xi_{ij}$，且 $\Xi_{ij}$ 至少覆盖三类分块（Lemma~\ref{lem:eng-envelope-minimal-three}）：
  \[
  \Xi_{ij}^{\mathrm{Geom}},\qquad
  \Xi_{ij}^{\mathrm{Sim}},\qquad
  \Xi_{ij}^{\mathrm{Proc}}.
  \]
  \item 存在沿 $J$‑测地线的回归方向 $\gamma$ 与基底锁定的同伦扭曲 $H$，使缺陷势能在扰动尖峰后可回落并最终满足 Definition~\ref{def:eng-feasibility-predicate} 的门槛；
\end{enumerate}
则工程逆向生成在形式系统层面可行，并在计算意义上\textbf{条件可行}：只要把 $\Xi_{ij}$ 的包络分层（多保真/分辨率分层）并将 $\mathrm{Res}/\mathrm{Unc}/\mathrm{Cert}$ 纳入
  $\mathrm{Loc}_{\mathrm{eng}}$ 的可审计门槛，则第 13 章的 GZ‑LSPI（$\mathcal{Z}_{\mathrm{law}}$）可在 $\mathrm{Loc}_{\mathrm{eng}}$
  内给出“最小作用量路径族”的可计算近似，从而实现可复用的逆向生成与诊断资产沉淀。
\end{theorem}

\begin{certificate}[数学归纳法证书：按包络细化层级 $r$ 归纳（细化不劣化）]
\label{cert:eng-envelope-refine-induction}
设 $\Xi_{ij}^{(\le r)}:=\prod_{t=1}^{r}\Xi_{ij}^{[t]}$ 表示 $r$ 层包络细化后的参数空间（与第 10 章丛塔口径同型）。
令谓词 $P(r)$ 为：“在包络细化层级 $r$ 下存在一对 $(\gamma,H)$ 使可行性谓词成立（Definition~\ref{def:eng-feasibility-predicate}）”。
证书化要求为：
\[
P(r)\Rightarrow P(r+1),
\qquad
\text{并以某个 }r_0\ \text{作为可行层级的起点。}
\]
\end{certificate}

\begin{proof}[证明（数学归纳法证书；谓词因果链）]
由 Certificate~\ref{cert:eng-envelope-refine-induction}，取谓词链：
\[
\begin{aligned}
A:&\ P(r)\Rightarrow \exists (\gamma,H)\ \text{在 }\Xi_{ij}^{(\le r)}\text{ 上满足可行性谓词},\\
B:&\ \Xi_{ij}^{(\le r)}\subseteq \Xi_{ij}^{(\le r+1)}\Rightarrow \text{可复用同一 }(\gamma,H),\\
C:&\ \text{可复用}\Rightarrow P(r+1).
\end{aligned}
\]
由 $A\wedge B\Rightarrow C$ 得到 $P(r)\Rightarrow P(r+1)$（细化不劣化）。
结合定理前提（1）（2）（3），可在某个有限细化层级 $r_0$ 上构造满足可行性谓词的 $(\gamma,H)$；由细化不劣化，后续更高层级仍保持可行。
在此基础上，包络分层 + $\mathrm{Loc}_{\mathrm{eng}}$ 门槛化使得路径积分不会被“缺证书但自洽”的虚假极小值主导（Proposition~\ref{prop:eng-envelope-anti-drift}），
  从而得到“条件可行”的可计算实现口径。证毕。
\end{proof}

\subsection{推论：统一性——法拉第笼/旁路电阻/抗退相干的同一口径}

\begin{corollary}[统一性推论：不同工业任务是同一工程联合扇区的子扇区特化]
\label{cor:eng-unified-subsectors}
在 Definition~\ref{def:eng-extended-base} 的联合扇区 $M_\ast^{\mathrm{eng}}$ 上：
\begin{enumerate}
  \item “特高压节点的 \RTSC{} 法拉第笼/屏蔽体设计”可视为对 $M_\ast^{\mathrm{Geom}}\times M_\ast^{\mathrm{Grid}}$ 的子域限制，并选择相应几何‑边界条件目标谓词；
  \item “\RTSC{} 旁路电阻/故障态旁路设计”可视为对 $M_\ast^{\mathrm{Grid}}$ 的阈值切换子结构加权，并选择相应热‑电‑结构耦合目标谓词；
  \item “量子计算抗退相干设计”可视为对 $M_\ast^{\mathrm{QC}}$ 的子域限制，并选择相应噪声‑控制‑材料一致性目标谓词。
\end{enumerate}
因此三类任务共享同一条“缺陷势能 + 阈值 + 测地线同伦扭曲回归”的可行性判据（Theorem~\ref{thm:eng-conditional-feasibility}）；差异仅体现为子扇区限制与目标谓词不同，而非三套互不相干的形式系统。

\end{corollary}

\begin{certificate}[证书：子扇区限制对应子丛限制，判据在限制下保持]
\label{cert:eng-unified-subsectors}
由第 9 章的子丛限制口径（Lemma~\ref{lem:subbundle-restriction}），对子域 $M'\subset M_\ast^{\mathrm{eng}}$ 的限制对应主丛限制
  $P_\ast^{\mathrm{eng}}|_{M'}$。
可行性谓词仅依赖于 $\mathrm{Loc}_{\mathrm{eng}}/\Sigma_{\mathrm{eng}}$、缺陷势能 $\Phi$ 与证书门槛；在子域限制下，这些对象仍可按限制继承，因此判据保持成立。
\end{certificate}

\begin{proof}[证明（谓词因果链）]
由 Certificate~\ref{cert:eng-unified-subsectors}，取谓词链：
\[
\begin{aligned}
A:&\ M'\subset M_\ast^{\mathrm{eng}}\Rightarrow P_\ast^{\mathrm{eng}}|_{M'}\ \text{存在},\\
B:&\ \text{限制后 }\mathrm{Loc}_{\mathrm{eng}}/\Sigma_{\mathrm{eng}}/\Phi/\mathrm{Cert}\ \text{仍可用},\\
C:&\ A\wedge B\Rightarrow \text{可行性判据在子扇区上保持}.
\end{aligned}
\]
由 $C$ 得证。证毕。
\end{proof}

\begin{remark}[落地建议：最小基底块/最小包络块/受保护低阶可观测量]
若希望把本章内容进一步压缩为“可直接写入纯数卷附录的基准条目”，可按以下三项最小集合落地：
\begin{enumerate}
  \item 明确 $M_\ast^{\mathrm{Geom}}$ 的最小法对象坐标块：哪些几何自由度必须作为基底坐标被证书化（否则无法诊断与比较）；
  \item 明确 $\Xi_{ij}^{\mathrm{Geom}}/\Xi_{ij}^{\mathrm{Sim}}/\Xi_{ij}^{\mathrm{Proc}}$ 的最小包络坐标块：哪些必须显式化才能支持扰动尖峰后的同伦扭曲回归；
  \item 明确 $\mathrm{Obs}_{\le r}$（Theorem~\ref{thm:jacgap-zero-iff-strictifiable}）中受保护的低阶可观测量：哪些工程低阶行为在同伦扭曲中不允许改变，
    从而使“修复/严格化”可被审计与复现。
\end{enumerate}
\end{remark}

\clearpage
%%%%%%%%%%%%%%%%%%%%%%%%%%%%%%%%%%%%%%%%%%%%%%%%%%%%%%%%%%%%
\section{纯数卷口径：工业几何体包络与逆向生成（\textsf{IG-RTSC} 模板）}
\label{sec:puremath-industrial-envelope}
%%%%%%%%%%%%%%%%%%%%%%%%%%%%%%%%%%%%%%%%%%%%%%%%%%%%%%%%%%%%

\subsection{形式逻辑表达约定}

\begin{remark}[翻译目标：用纯数卷语言编码“工业逆向生成”的可行性与回归机制]
本章严格使用纯数卷（PureMath 卷）口径，对“在 \RTSC{} 基准上叠加大型工业几何体设计，并通过纤维丛包络实现逆向生成工程设计”的可能性给出一个可直接落入形式系统的评估与构造模板。
核心对象与接口固定为：
\begin{enumerate}
  \item \textbf{几何到法空间的嵌入：}对象层几何通过 $\Phi_{\mathrm{geom}}:\mathsf{GeomObj}\to\GZLS$ 升为法空间对象；
  \item \textbf{law-level operator system：}在 $(\GZLS)_{\mathrm{homo}}$ 或 $(\GZLS)_{\text{vf-op}}$ 扇区组织为 $L$；
  \item \textbf{缺陷证书：}用 Jacobiator-gap 证书 $\JacGap(L)$ 与缺陷势能 $\Phi(L)=\|\JacGap(L)\|_{\mathrm{def}}$ 刻画“可控性/可逆向生成性”；
  \item \textbf{包络显式化：}用“切面 $s$ + 包络邻域 $E$ + 包络坐标 $\Xi$”把工业级几何/仿真/工艺自由度从隐变量提升为可控纤维方向；
  \item \textbf{扰动回归：}外生扰动引发法曲率尖峰时，采用“沿 $J$‑测地线的同伦扭曲（geodesic homotopy twist）”回归到 $\mathrm{Loc}$ 的证书化工作区。
\end{enumerate}
本章与上一章（Section~\ref{sec:industrial-geometry-benchmark-inverse}）的关系是：上一章给出“工程联合扇区 + 包络 + 条件可行”的应用口径；本章给出同一思想在纯数卷的
  $\Phi_{\mathrm{geom}}$ 与 \textsf{GZ-OHU}（failure-free construction）语言下的可引用模板。
\end{remark}

\subsection{定义：对象层几何与 $\Phi_{\mathrm{geom}}$ 嵌入（统一为法空间扇区）}

\begin{definition}[对象层类别与工业‑RTSC 基准扇区]
\label{def:ig-geom-sector}
在纯数卷口径下，设存在对象层几何类别 $\mathsf{GeomObj}$，以及几何到法空间的嵌入函子
\[
\Phi_{\mathrm{geom}}:\mathsf{GeomObj}\longrightarrow \GZLS,
\]
其作用是把“主丛‑联络‑曲率‑辅助结构”等对象层几何数据映射为法空间 $\GZLS$ 中的法对象，并将 gauge 等价映射为同构法对象。

定义工业‑RTSC 基准的对象层复合类别为
\[
\mathsf{GeomObj}^{\textsf{IG-RTSC}}
:=
\mathsf{GeomObj}^{\mathrm{RTSC}}
\times
\mathsf{GeomObj}^{\mathrm{Cage}}
\times
\mathsf{GeomObj}^{\mathrm{Bypass}}
\times
\mathsf{GeomObj}^{\mathrm{QDec}},
\]
并令其在法空间中的像落入某个可局部化扇区
\[
M^{\textsf{IG-RTSC}}
\subseteq
\GZLS,
\qquad
M^{\textsf{IG-RTSC}}\subseteq \mathrm{Loc}.
\]
其中 “\textsf{Cage}/\textsf{Bypass}/\textsf{QDec}” 分别对应“法拉第笼（屏蔽体）/旁路电阻（故障旁路）/量子抗退相干装置”三类工业几何体任务的对象层几何与边界‑运行条件。
\end{definition}

\begin{certificate}[文本证书：$\Phi_{\mathrm{geom}}$、$\mathsf{GeomObj}$ 与 \textsf{GZ-OHU} 的纯数卷来源]
\label{cert:puremath-phi-geom-ohu}
在 \code{docs/PureMath/Pub_GFramework_PureMath.tex} 中给出：
\begin{enumerate}
  \item 法空间 $\GZLS=(\mathfrak{L};J,\mathrm{Loc},\Sigma)$ 以及对象层几何到法对象的嵌入函子 $\Phi_{\mathrm{geom}}:\mathsf{GeomObj}
    \to\GZLS$；
  \item $\Phi_{\mathrm{geom}}$ 的同伦/变分扩展口径（例如 $\Phi_{\mathrm{geom}}^{\mathrm{homo}}$ 与 $\Phi_{\mathrm{geom}}^{\text{vf-op}
    }$ 的扩展版本）；
  \item \textsf{GZ-OHU}（Operator--Homotopy Universality and Failure-Free Construction）框架：用 Jacobiator-gap 证书刻画“可严格化”的边界，
    并给出 failure-free constructive schemes 的抽象结论。
\end{enumerate}
本章仅抽取这些条款作为“逆向生成的形式系统接口”，不依赖纯数卷的具体实现细节。
\end{certificate}

\begin{remark}[与上一章工程联合扇区的对应]
Definition~\ref{def:ig-geom-sector} 的 $M^{\textsf{IG-RTSC}}$ 可视为上一章 $M_\ast^{\mathrm{eng}}$（Definition~\ref{def:eng-extended-base}）在纯数卷语言下的重命名：二者都表达“将 RTSC‑几何‑运行扰动‑量子装置统一纳入同一 $\mathrm{Loc}/\Sigma$ 管理的可局部化扇区”，差异仅在于采用的记号与引用侧重点不同。
\end{remark}

\subsection{定义：law-level operator system 与 JacGap/defect potential 接口}

\begin{definition}[\textsf{IG-RTSC} 的 law-level 编码：$L$ 的扇区与运输]
\label{def:ig-law-level-system}
在 $M^{\textsf{IG-RTSC}}\subseteq \GZLS$ 上，称映射
\[
L:\ M^{\textsf{IG-RTSC}}\to (\GZLS)_{\mathrm{homo}}
\quad\text{或}\quad
L:\ M^{\textsf{IG-RTSC}}\to (\GZLS)_{\text{vf-op}}
\]
为 \textsf{IG-RTSC} 的 law-level operator system，若对每个基底点 $x\in M^{\textsf{IG-RTSC}}$：
\begin{enumerate}
  \item $L(x)$ 至少携带 $(\ell_1,\ell_2,\ell_3)$ 等同伦括号数据（见 Definition~\ref{def:law-level-operator-system} 的同型口径），用于编码 Jacobi
    失配的同伦可修复性；
  \item $L(x)$ 的演化/运输受法联络 $\GZLOC$ 及（必要时）三层联络 $\GZTLC$ 的水平法流控制（Definition~\ref{def:law-connection-curvature}）。
\end{enumerate}
其中 $(\GZLS)_{\text{vf-op}}$ 仅表示“变分场‑算子（vf-op）”意义下的扩展扇区，用于把高维工程自由度、评价泛函与控制路径统一为可优化对象；本章不依赖其具体实现。
\end{definition}

\begin{definition}[Jacobiator-gap 证书与缺陷势能（纯数卷接口）]
\label{def:ig-jacgap-defect}
对 law-level 对象 $L(x)$，定义 Jacobiator-gap 证书与缺陷势能为
\[
\JacGap\bigl(L(x)\bigr)\in \mathbb{R}_{\ge 0}^k,
\qquad
\Phi\bigl(L(x)\bigr):=\bigl\|\JacGap\bigl(L(x)\bigr)\bigr\|_{\mathrm{def}},
\]
其中 $\|\cdot\|_{\mathrm{def}}$ 为缺陷聚合范数（Definition~\ref{def:jacgap-certificate}）。
本章把“可逆向生成/可严格化”统一编码为：“在保持受保护低阶可观测量不变的前提下，使 $\Phi$ 被严格降低并在可行时达到 $\JacGap=0$ 的严格扇区（Theorem~\ref{thm:jacgap-zero-iff-strictifiable}）”。
\end{definition}

\subsection{定义：受保护低阶可观测量（protected observables）的工程语义}

\begin{definition}[三类工业基准的受保护评价泛函族]
\label{def:ig-protected-observables}
在纯数卷口径下，将“工程指标”抽象为一组有限的低阶评价泛函（evaluation functionals），并声明其为受保护（protected）低阶可观测量：
\begin{enumerate}
  \item（法拉第笼/屏蔽体）取有限集合
  \[
  \mathcal{E}^{\mathrm{Cage}}
  :=
  \{\mathrm{Ev}^{\mathrm{shield}}_r,\mathrm{Ev}^{\mathrm{leak}}_s,\mathrm{Ev}^{\mathrm{thermal}}_t,\mathrm{Ev}
    ^{\mathrm{mech}}_u,\ldots\},
  \]
  例如屏蔽效能、泄漏/耦合、热稳定、机械完整性等的低阶评价；
  \item（旁路电阻/故障旁路）取有限集合
  \[
  \mathcal{E}^{\mathrm{Bypass}}
  :=
  \{\mathrm{Ev}^{\mathrm{IV}}_r,\mathrm{Ev}^{\mathrm{fault}}_s,\mathrm{Ev}^{\mathrm{energy}}_t,\mathrm{Ev}
    ^{\mathrm{switch}}_u,\ldots\},
  \]
  例如 I‑V 低阶特征、故障态响应、能量/热冲击低阶响应、切换一致性等；
  \item（量子抗退相干）取有限集合
  \[
  \mathcal{E}^{\mathrm{QDec}}
  :=
  \{\mathrm{Ev}^{T_1}_r,\mathrm{Ev}^{T_2}_s,\mathrm{Ev}^{\mathrm{noise}}_t,\mathrm{Ev}^{\mathrm{gate}}_u,\ldots\},
  \]
  例如退相干时间、噪声谱低阶特征、门操作低阶误差等。
\end{enumerate}
定义工业‑RTSC 受保护评价泛函族为
\[
\mathcal{E}^{\textsf{IG-RTSC}}
:=
\mathcal{E}^{\mathrm{Cage}}
\cup
\mathcal{E}^{\mathrm{Bypass}}
\cup
\mathcal{E}^{\mathrm{QDec}},
\]
并用谓词 $\mathrm{Obs}_{\le r}^{\textsf{IG-RTSC}}$ 表示 “$L$ 的受保护低阶可观测量取值”（与 Theorem~\ref{thm:jacgap-zero-iff-strictifiable} 的
  $\mathrm{Obs}_{\le r}$ 同型）。
\end{definition}

\begin{remark}[为何必须声明受保护量]
纯数卷的严格化判据强调：$\JacGap(L)=0$ 的“严格化”必须在不改变受保护低阶可观测量的前提下发生（见 Theorem~\ref{thm:jacgap-zero-iff-strictifiable}）。
因此在工业基准中，若不先声明 $\mathcal{E}^{\textsf{IG-RTSC}}$（即不固定“工程上必须保持的低阶行为”），则“$\JacGap=0$”将失去工程语义：同伦修复可能通过改变工程目标本身来使缺陷消失，
  导致不可审计的伪严格化。
\end{remark}

\subsection{定义：切面 + 包络邻域（把工业级自由度显式化为纤维方向）}

\begin{definition}[切面与包络：$s$、$E$ 与 $\Xi$]
\label{def:ig-section-envelope}
在局部化域 $U\subseteq M^{\textsf{IG-RTSC}}$ 上，取一阶切面（局部截面）
\[
s:\ U\to P,
\qquad
\pi\circ s=\mathrm{Id}_U,
\]
其语义为“固定一套工程路线/仿真‑评价‑制造口径”（即固定一套表示、求解管线与协议）。
取切面包络邻域
\[
E\subseteq P,
\qquad
s(U)\subseteq E,
\]
并引入包络坐标空间 $\Xi$ 以显式参数化切面附近的细粒度纤维自由度：
\[
\Psi:\ U\times \Xi \to E,
\qquad
\pi\circ \Psi=\mathrm{pr}_U.
\]
在工业几何体逆向生成场景中，将包络坐标最小分块为
\[
\Xi
:=
\Xi^{\mathrm{Geom}}
\times
\Xi^{\mathrm{Sim}}
\times
\Xi^{\mathrm{Proc}},
\]
分别对应几何细粒度自由度、仿真离散化/误差控制自由度、制造装配/公差自由度。
该结构与上一章的二阶包络（Definition~\ref{def:eng-envelope}）在形式上同型：都是把“最易导致自洽偏航的隐藏自由度”提升为可审计的纤维坐标。
\end{definition}

\begin{lemma}[最小包络分块（纯数卷版）：缺少任一分块都会破坏证书化下降]
\label{lem:ig-minimal-envelope}
在 Definition~\ref{def:ig-section-envelope} 的口径下，若希望在保持受保护低阶可观测量不变的前提下，通过有限同伦动作与包络移动使缺陷势能 $\Phi$ 被严格降低，
则包络空间 $\Xi$ 至少应包含 $\Xi^{\mathrm{Geom}},\Xi^{\mathrm{Sim}},\Xi^{\mathrm{Proc}}$ 三类分块。
\end{lemma}

\begin{certificate}[证书：三类分块分别对应“缺陷下降”的三类必要自由度来源]
\label{cert:ig-minimal-envelope}
\begin{enumerate}
  \item 若缺少 $\Xi^{\mathrm{Geom}}$，则切面邻域无法覆盖关键几何失效模式，缺陷势能的几何来源缺少可调方向；
  \item 若缺少 $\Xi^{\mathrm{Sim}}$，则残差/不确定度难以在跨离散化意义下证书化控制，容易出现“自洽但偏航”的虚假下降；
  \item 若缺少 $\Xi^{\mathrm{Proc}}$，则方案可能脱离 $\mathrm{Loc}$（可制造/可装配/可运行域），导致“形式上下降但工程上不可落地”的伪可行。
\end{enumerate}
\end{certificate}

\begin{proof}[证明（谓词因果链）]
由 Certificate~\ref{cert:ig-minimal-envelope}，取谓词链：
\[
\begin{aligned}
A:&\ \neg \Xi^{\mathrm{Geom}}\Rightarrow \text{缺陷势能缺少几何可调方向},\\
B:&\ \neg \Xi^{\mathrm{Sim}}\Rightarrow \mathrm{Res}/\mathrm{Unc}/\mathrm{Cert}\ \text{不稳定}\Rightarrow \text{偏航风险上升},\\
C:&\ \neg \Xi^{\mathrm{Proc}}\Rightarrow \text{易脱离 }\mathrm{Loc}\Rightarrow \text{工程不可落地}.
\end{aligned}
\]
$A\vee B\vee C$ 任一成立都会破坏“在 $\mathrm{Loc}$ 内保持受保护量并严格降低 $\Phi$”的必要自由度闭合性，因此三类分块均为最低要求。证毕。
\end{proof}

\subsection{定理（公理降级）：\textsf{GZ-OHU} 口径下的逆向生成判据}

\begin{theorem}[逆向生成判据（公理降级）：有限证书 + 受保护量 + 包络下降方向]
\label{thm:ig-invgen-ohu}
设 $L$ 为 \textsf{IG-RTSC} 的 law-level operator system（Definition~\ref{def:ig-law-level-system}），并固定受保护低阶可观测量谓词
  $\mathrm{Obs}_{\le r}^{\textsf{IG-RTSC}}$（Definition~\ref{def:ig-protected-observables}）。
若满足：
\begin{enumerate}
  \item（有限证书）$\JacGap(L)$ 为有限维证书，且缺陷势能 $\Phi=\|\JacGap\|_{\mathrm{def}}$ 可被用于比较与驱动下降；
  \item（保护量一致性）允许的同伦动作（homotopy moves）在不改变 $\mathrm{Obs}_{\le r}^{\textsf{IG-RTSC}}$ 的前提下调整高阶同伦数据；
  \item（包络下降方向）存在切面‑包络结构 $(s,E,\Xi)$（Definition~\ref{def:ig-section-envelope}），使得在 $\Xi$ 上存在可计算的包络移动（几何/仿真/工艺分块联合移动），
    从而与同伦动作共同提供“严格降低 $\Phi$ 且保持 $\mathrm{Loc}$”的下降机制；
\end{enumerate}
则在纯数卷意义下，工业几何体逆向生成具有可证书化的 failure-free 构造口径：存在一条有限步的“同伦动作 + 包络移动”序列，使缺陷势能单调下降，并在可行时进入 $\JacGap=0$ 的严格扇区（从而可严格化，见
  Theorem~\ref{thm:jacgap-zero-iff-strictifiable}）。
\end{theorem}

\begin{certificate}[数学归纳法证书：按同伦动作步数 $n$ 归纳（单调下降证书）]
\label{cert:ig-invgen-induction}
令 $n\in\mathbb{N}$ 表示“同伦动作步数”（每一步允许伴随一次包络移动）。
定义谓词 $P(n)$ 为：“存在 $n$ 步动作序列，使得（i）受保护量 $\mathrm{Obs}_{\le r}^{\textsf{IG-RTSC}}$ 不变，（ii）轨线始终留在 $\mathrm{Loc}$，（iii）
  缺陷势能形成单调链 $\Phi_0\ge \Phi_1\ge\cdots\ge \Phi_n$，且在必要时出现严格下降 $\Phi_{t+1}<\Phi_t$”。
证书化要求为：
\[
P(0)\ \text{成立},\qquad
(\forall n\in\mathbb{N})\ \bigl(P(n)\Rightarrow P(n+1)\bigr).
\]
\end{certificate}

\begin{proof}[证明（数学归纳法证书；谓词因果链）]
由 Certificate~\ref{cert:ig-invgen-induction}，取谓词链：
\[
\begin{aligned}
A:&\ \text{（包络下降方向 + 同伦动作）}\Rightarrow \exists\ \text{一步动作}\ \text{使 }\Phi\ \text{不增且可严格降},\\
B:&\ \text{该一步动作}\Rightarrow \mathrm{Obs}_{\le r}^{\textsf{IG-RTSC}}\ \text{不变}\wedge \text{仍在 }\mathrm{Loc},\\
C:&\ P(n)\wedge (A\wedge B)\Rightarrow P(n+1).
\end{aligned}
\]
基步 $P(0)$ 由“零步动作序列”成立。归纳步由 $C$ 给出。
因此存在有限步动作序列得到单调下降链；当下降把 $\Phi$ 推入 0 的邻域且证书允许时，由 Theorem~\ref{thm:jacgap-zero-iff-strictifiable} 得到进入严格扇区/可严格化的结论。证毕。
\end{proof}

\subsection{定理（公理降级）：外生扰动下的测地线同伦扭曲回归条款}

\begin{theorem}[扰动回归条款（公理降级）：曲率尖峰诱发缺陷尖峰；测地线 + 同伦扭曲回归]
\label{thm:ig-disturbance-regression}
设外生扰动（如强外场/浪涌/太阳风暴类冲击）通过对象层到法空间的映射作用于 $L(t)$ 的水平法流，导致法曲率 $\GZLCurv$ 在短时窗口出现尖峰。
若满足：
\begin{enumerate}
  \item（尖峰放大）缺陷势能的变化受曲率–同伦控制不等式约束（Theorem~\ref{thm:law-curvature-spike-jacgap} 的同型口径），从而尖峰输入会引发 $\Phi(L(t))$ 的尖峰；
  \item（回归机制）存在 $J$‑测地线方向 $\gamma$ 与基底锁定的同伦扭曲 $H$，并允许在包络坐标 $\Xi$ 上做快层修复（Definition~\ref{def:ig-section-envelope}），
    使得扰动后能把 $L(t)$ 拉回 $\mathrm{Loc}$ 且使 $\Phi$ 回落；
\end{enumerate}
则工业逆向生成过程在扰动下具备“回到证书化工作区”的动态条款：扰动尖峰并不必然导致路线崩溃，而会触发“先选测地线方向（诊断/方向层）+ 再做同伦扭曲与包络修复（治疗/构造层）”的回归机制（与 Corollary~\ref{cor:puremath-direction-before-construction} 同型）。
\end{theorem}

\begin{certificate}[证书：由曲率尖峰定理与方向优先推论合取得到]
\label{cert:ig-disturbance-regression}
\begin{enumerate}
  \item 由 Theorem~\ref{thm:law-curvature-spike-jacgap}，曲率尖峰输入会放大缺陷势能的变化并导致尖峰；
  \item 由 Corollary~\ref{cor:puremath-direction-before-construction}，回归优先是“选测地线方向”，其次才是在固定方向上做同伦扭曲修复；
  \item 由 Definition~\ref{def:ig-section-envelope}，包络坐标 $\Xi$ 提供快层修复所需的显式纤维自由度，使回归过程可被证书化与审计。
\end{enumerate}
\end{certificate}

\begin{proof}[证明（谓词因果链）]
由 Certificate~\ref{cert:ig-disturbance-regression}，取谓词链：
\[
\begin{aligned}
A:&\ \GZLCurv\ \text{尖峰}\Rightarrow \Phi(L(t))\ \text{尖峰},\\
B:&\ \text{存在 }\gamma\Rightarrow \text{先在方向层选取回归路径},\\
C:&\ \text{存在 }H\ \text{且 }\Xi\ \text{可用}\Rightarrow \text{在固定 }\gamma\text{ 上做同伦扭曲/包络修复使 }\Phi\ \text{回落并回到 }
  \mathrm{Loc}.
\end{aligned}
\]
由 $A\wedge B\wedge C$ 得到“扰动尖峰可被回归机制吸收”的动态条款。证毕。
\end{proof}

\subsection{推论：结论性评估（肯定但条件化）}

\begin{corollary}[结论：工业逆向生成的可能性肯定，但成立条件可形式化]
\label{cor:ig-conditional-yes}
在纯数卷口径下，“在 \RTSC{} 基准上叠加工业几何体设计并通过包络实现逆向生成”具有肯定的形式系统表达，但其成立依赖于以下三条必要条件：
\begin{enumerate}
  \item（可局部化嵌入）对象层工业几何与运行条件可通过 $\Phi_{\mathrm{geom}}$ 嵌入到 $\GZLS$ 的可局部化扇区并受 $\mathrm{Loc}/\Sigma$ 管理（Definition~\ref{def:ig-geom-sector}）；
  \item（有限缺陷证书 + 保护量）扩展系统 admit 有限的 $\JacGap(L)$，并明确受保护低阶可观测量 $\mathcal{E}^{\textsf{IG-RTSC}}$，使 $\JacGap=0$
    的严格化具有工程语义（Definition~\ref{def:ig-jacgap-defect} 与 Definition~\ref{def:ig-protected-observables}）；
  \item（包络显式化 + 下降机制）每条关键路线切面都具备包含 $\Xi^{\mathrm{Geom}}/\Xi^{\mathrm{Sim}}/\Xi^{\mathrm{Proc}}$ 的包络，
    并提供与同伦动作耦合的缺陷势能下降方向（Definition~\ref{def:ig-section-envelope} 与 Theorem~\ref{thm:ig-invgen-ohu}）。
\end{enumerate}
当三条条件满足时，逆向生成可被表述为纯数卷意义下的 failure-free constructive scheme；在外生扰动下，还可用“测地线方向 + 同伦扭曲回归”条款维持其可操作性（Theorem~\ref{thm:ig-disturbance-regression}）。
\end{corollary}

\begin{proof}[证明（谓词因果链）]
取谓词链：
\[
\begin{aligned}
A:&\ \text{条件（1）}\Rightarrow \text{工业对象可进入 }\mathrm{Loc}\ \text{并受阈值管理},\\
B:&\ \text{条件（2）}\Rightarrow \JacGap\ \text{与严格化判据具备工程语义},\\
C:&\ \text{条件（3）}\Rightarrow \text{存在证书化下降机制并得到逆向生成构造}(\text{Theorem~\ref{thm:ig-invgen-ohu}}),\\
D:&\ A\wedge B\wedge C\Rightarrow \text{逆向生成的可能性在形式系统内成立}.
\end{aligned}
\]
由 $D$ 得证；扰动下的可操作性由 Theorem~\ref{thm:ig-disturbance-regression} 追加。证毕。
\end{proof}

\clearpage
%%%%%%%%%%%%%%%%%%%%%%%%%%%%%%%%%%%%%%%%%%%%%%%%%%%%%%%%%%%%
\section{纯数卷口径：相对同伦最小包络类（受保护评价泛函）}
\label{sec:puremath-relative-homotopy-minimal-envelope}
%%%%%%%%%%%%%%%%%%%%%%%%%%%%%%%%%%%%%%%%%%%%%%%%%%%%%%%%%%%%

\subsection{形式逻辑表达约定}

\begin{remark}[核心主张：包络“本质”由相对同伦类决定，而非层数]
本章把如下主张严格化为纯数卷口径的定义链与可引用命题：
\begin{enumerate}
  \item \textbf{相对同伦视角：}形式化包络（envelope / iterative extension / $L_\infty$-envelope）的“最小同伦类”应理解为
  \emph{相对一组受保护评价泛函 $\mathcal{E}_{\mathrm{prot}}$ 的最小同伦类}，而不是“绝对最小”；
  \item \textbf{层数非本质：}包络层数（两层或更深递归解耦，甚至包含非因果解耦）本身不决定方案本质；若新增层是满足相对约束的纤维可回缩扩展，则不会改变相对同伦类；
  \item \textbf{受保护量必要：}若不指定 $\mathcal{E}_{\mathrm{prot}}$（pure 卷的 observable low-order behaviour），则“最小同伦类”不具备可判定语义，且
    $\JacGap$ 与缺陷势能 $\Phi=\|\JacGap\|_{\mathrm{def}}$ 会失去工程含义（同伦扭曲可通过改变观测口径做“假严格化/假回归”）。
\end{enumerate}
\end{remark}

\subsection{定义：受保护评价泛函族与相对约束}

\begin{definition}[受保护评价泛函族（protected evaluation functionals）]
\label{def:protected-evaluation-functionals}
在纯数卷口径下，取一组有限的受保护评价泛函
\[
\mathcal{E}_{\mathrm{prot}}
:=
\{\mathrm{Ev}_a\}_{a=1}^{q},
\qquad
\mathrm{Ev}_a:\ \mathrm{Loc}\subseteq \GZLS\to \mathbb{R},
\]
并记向量化映射
\[
\mathrm{Ev}:=(\mathrm{Ev}_1,\ldots,\mathrm{Ev}_q):\mathrm{Loc}\to \mathbb{R}^q.
\]
称某条同伦过程（同伦扭曲/同伦动作序列的连续极限）$H:[0,1]\times[0,T]\to (\GZLS)_{\mathrm{homo}}$ 满足\textbf{相对约束}，若对任意 $t\in[0,T]$ 都有
\[
\mathrm{Ev}\bigl(H(s,t)\bigr)=\mathrm{Ev}\bigl(H(0,t)\bigr)
\qquad (\forall s\in[0,1]).
\]
该条件即 pure 卷“保持 observable low-order behaviour 不变”的严格表达。
\end{definition}

\begin{certificate}[证书：$\JacGap=0$ 的语义依赖于受保护低阶行为]
\label{cert:jacgap-needs-protected-observables}
纯数卷的 Jacobiator-gap 证书口径要求：$\JacGap(L)=0$ 当且仅当 $L$ 可被严格化到某个 $L_{\mathrm{strict}}$，并且该严格化过程不改变被声明为受保护的低阶可观测行为。

在本工作笔记中，这一口径已被编码为 Theorem~\ref{thm:jacgap-zero-iff-strictifiable} 中的“低阶可观测量投影谓词” $\mathrm{Obs}_{\le r}$；
本章以 Definition~\ref{def:protected-evaluation-functionals} 的 $\mathcal{E}_{\mathrm{prot}}$ 作为一个可计算的具体实现：用有限评价泛函族替代抽象谓词，
  从而使“相对同伦”与“严格化”具有可审计语义。
\end{certificate}

\begin{proof}[证明（谓词因果链）]
由 Certificate~\ref{cert:jacgap-needs-protected-observables}，取谓词链：
\[
\begin{aligned}
A:&\ \JacGap(L)=0\ \text{的纯数卷语义}\Rightarrow \text{存在“保持受保护量不变”的严格化},\\
B:&\ \text{存在保持受保护量不变的严格化}\Rightarrow \text{需先固定受保护量的取值口径},\\
C:&\ \text{固定受保护量口径}\Rightarrow \text{可用有限评价泛函族 }\mathcal{E}_{\mathrm{prot}}\ \text{表达相对约束}.
\end{aligned}
\]
由 $A\wedge B\Rightarrow C$ 得证。证毕。
\end{proof}

\begin{remark}[与 \textsf{IG-RTSC} 章的接口]
本章的 $\mathcal{E}_{\mathrm{prot}}$ 是对上一章中 $\mathcal{E}^{\textsf{IG-RTSC}}$（Definition~\ref{def:ig-protected-observables}）
  的抽象化：后者给出三类工业任务的受保护评价泛函族示例；本章则把“受保护量”抽象为任意有限族 $\mathcal{E}_{\mathrm{prot}}$，用于定义相对同伦类与最小包络类。
\end{remark}

\subsection{定义：包络丛塔（iterative extensions）与非因果解耦}

\begin{definition}[包络丛塔：由参数块递归细化得到的 tower]
\label{def:envelope-tower-iterative}
令 $U$ 为某个局部化基底域（例如 $U\subseteq \mathrm{Loc}$ 的可操作域）。
称序列 $\{E^{(r)}\}_{r\ge 0}$ 为一条\textbf{包络丛塔}，若存在一组参数块 $\{\Xi^{(r)}\}_{r\ge 1}$ 使得
\[
E^{(0)}:=U,
\qquad
E^{(r)}:=E^{(r-1)}\times \Xi^{(r)}\quad (r\ge 1),
\]
并对每层定义投影
\[
q^{(r)}:\ E^{(r)}\to E^{(r-1)},
\qquad
q^{(r)}(x,\xi)=x.
\]
其中每层参数块 $\Xi^{(r)}$ 允许为“非因果解耦”的：它不要求对应时间顺序，仅要求其作为纤维方向自由度受到 $\mathrm{Loc}$ 与阈值 $\Sigma$ 的约束，并可被 $\GZLOC/\widetilde{\GZTLC}
  $ 的水平结构一致管理。
\end{definition}

\begin{remark}[与“切面包络/丛塔”章节的同型关系]
Definition~\ref{def:envelope-tower-iterative} 与第 10 章的包络塔（Theorem~\ref{thm:bundle-tower-closure}）同型：
  都是通过递归引入参数块把“隐变量/细节自由度”显式化为可审计坐标。
本章的差异是：进一步引入 $\mathcal{E}_{\mathrm{prot}}$，把“层数可变”转写为“相对同伦等价类”的稳定性问题。
\end{remark}

\subsection{定义：相对同伦与相对同伦等价（relative homotopy equivalence）}

\begin{definition}[相对同伦：保持受保护评价泛函不变的同伦]
\label{def:relative-homotopy}
设 $F_0,F_1:E\to E'$ 为两映射，且存在（法对象）映射 $\Phi_{\mathrm{law}}:E'\to \mathrm{Loc}\subseteq\GZLS$（例如由 $\Phi_{\mathrm{geom}}$
  与表示选择复合得到；见 Certificate~\ref{cert:puremath-phi-geom-ohu} 的口径）。
称 $F_0$ 与 $F_1$ 在 $\mathcal{E}_{\mathrm{prot}}$ 下\textbf{相对同伦}，若存在连续映射
\[
H_F:\ [0,1]\times E\to E'
\]
使得
\[
H_F(0,x)=F_0(x),\qquad H_F(1,x)=F_1(x),
\]
并且对任意 $x\in E$，曲线 $s\mapsto \Phi_{\mathrm{law}}(H_F(s,x))$ 满足 Definition~\ref{def:protected-evaluation-functionals}
  的相对约束：
\[
\mathrm{Ev}\bigl(\Phi_{\mathrm{law}}(H_F(s,x))\bigr)
=
\mathrm{Ev}\bigl(\Phi_{\mathrm{law}}(H_F(0,x))\bigr).
\]
记作 $F_0\simeq_{\mathcal{E}_{\mathrm{prot}}}F_1$。
\end{definition}

\begin{definition}[$\mathcal{E}_{\mathrm{prot}}$-相对同伦等价：包络塔之间的等价]
\label{def:relative-homotopy-equivalence}
给定两条包络丛塔 $\{E^{(r)}\}$ 与 $\{(E')^{(r)}\}$。
称二者在 $\mathcal{E}_{\mathrm{prot}}$ 下\textbf{相对同伦等价}，若存在某些层级 $r,r'$ 以及映射
\[
F:E^{(r)}\to (E')^{(r')},
\qquad
G:(E')^{(r')}\to E^{(r)},
\]
满足：
\begin{enumerate}
  \item（投影交换）$F,G$ 与各自丛塔投影交换（保持“同一基准下”的语义）；
  \item（保护量保持）$\Phi_{\mathrm{law}}\circ F$ 与 $\Phi_{\mathrm{law}}\circ G$ 均保持 $\mathcal{E}_{\mathrm{prot}}$ 的评价向量
    $\mathrm{Ev}$；
  \item（相对同伦到恒等）$G\circ F\simeq_{\mathcal{E}_{\mathrm{prot}}}\mathrm{Id}$ 且 $F\circ G\simeq_{\mathcal{E}_{\mathrm{prot}}}
    \mathrm{Id}$（Definition~\ref{def:relative-homotopy}）。
\end{enumerate}
并将其等价类称为“$\mathcal{E}_{\mathrm{prot}}$-相对同伦类”。
\end{definition}

\subsection{命题：层数不变量——可回缩纤维扩展不改变相对同伦类}

\begin{definition}[相对可回缩扩展：纤维可回缩且保持 $\mathcal{E}_{\mathrm{prot}}$]
\label{def:relative-retractable-extension}
在包络丛塔中，称第 $r$ 层扩展 $E^{(r)}=E^{(r-1)}\times \Xi^{(r)}$ 为\textbf{$\mathcal{E}_{\mathrm{prot}}$-相对可回缩}，若存在截面
\[
\iota^{(r)}:\ E^{(r-1)}\to E^{(r)}
\]
与投影 $q^{(r)}$ 满足 $q^{(r)}\circ \iota^{(r)}=\mathrm{Id}_{E^{(r-1)}}$，
并且复合 $\iota^{(r)}\circ q^{(r)}$ 与 $\mathrm{Id}_{E^{(r)}}$ 在 Definition~\ref{def:relative-homotopy} 意义下相对同伦：
\[
\iota^{(r)}\circ q^{(r)}\simeq_{\mathcal{E}_{\mathrm{prot}}}\mathrm{Id}_{E^{(r)}}.
\]
\end{definition}

\begin{proposition}[层数不变量：相对可回缩扩展不改变最小同伦类]
\label{prop:relative-homotopy-depth-invariant}
在 Definition~\ref{def:relative-retractable-extension} 的设定下，若第 $r$ 层扩展是 $\mathcal{E}_{\mathrm{prot}}$-相对可回缩的，
则 $E^{(r)}$ 与 $E^{(r-1)}$ 在 $\mathcal{E}_{\mathrm{prot}}$ 下相对同伦等价（Definition~\ref{def:relative-homotopy-equivalence}）。
因此：在 $\mathcal{E}_{\mathrm{prot}}$ 固定时，包络层数的增加（若为相对可回缩扩展）不会改变方案落入的“最小同伦类”。
\end{proposition}

\begin{certificate}[证书：相对可回缩给出互为逆的相对同伦]
\label{cert:relative-homotopy-depth-invariant}
取 $F:=q^{(r)}:E^{(r)}\to E^{(r-1)}$ 与 $G:=\iota^{(r)}:E^{(r-1)}\to E^{(r)}$。
则：
\[
F\circ G=\mathrm{Id}_{E^{(r-1)}},
\qquad
G\circ F=\iota^{(r)}\circ q^{(r)}\simeq_{\mathcal{E}_{\mathrm{prot}}}\mathrm{Id}_{E^{(r)}},
\]
从而满足 Definition~\ref{def:relative-homotopy-equivalence} 的等价条件。
\end{certificate}

\begin{proof}[证明（谓词因果链）]
由 Certificate~\ref{cert:relative-homotopy-depth-invariant}，取谓词链：
\[
\begin{aligned}
A:&\ q^{(r)}\circ \iota^{(r)}=\mathrm{Id}\Rightarrow F\circ G=\mathrm{Id}_{E^{(r-1)}},\\
B:&\ \iota^{(r)}\circ q^{(r)}\simeq_{\mathcal{E}_{\mathrm{prot}}}\mathrm{Id}\Rightarrow G\circ F\simeq_{\mathcal{E}
  _{\mathrm{prot}}}\mathrm{Id}_{E^{(r)}},\\
C:&\ A\wedge B\Rightarrow \text{$E^{(r)}$ 与 $E^{(r-1)}$ 相对同伦等价}.
\end{aligned}
\]
由 $C$ 得证。证毕。
\end{proof}

\subsection{定理（公理降级）：相对最小包络同伦类（相对脊柱）}

\begin{definition}[相对最小包络类（core/spine envelope）]
\label{def:relative-minimal-envelope-class}
在固定 $\mathcal{E}_{\mathrm{prot}}$ 的前提下，称某个包络层 $E^{(r_\ast)}$ 为\textbf{相对脊柱（core/spine envelope）}，
若对任意 $r\ge r_\ast$，从 $E^{(r_\ast)}$ 到 $E^{(r)}$ 的层级扩展可分解为有限个 $\mathcal{E}_{\mathrm{prot}}$-相对可回缩扩展（Definition~\ref{def:relative-retractable-extension}）。
称其等价类为\textbf{$\mathcal{E}_{\mathrm{prot}}$-相对最小包络同伦类}。
\end{definition}

\begin{theorem}[相对最小类（公理降级）：层数变化不影响本质类]
\label{thm:relative-minimal-envelope-class}
固定 $\mathcal{E}_{\mathrm{prot}}$ 与法对象映射 $\Phi_{\mathrm{law}}$（Definition~\ref{def:relative-homotopy}）。
若两条包络丛塔 $\{E^{(r)}\}$ 与 $\{(E')^{(r)}\}$ 都满足：
\begin{enumerate}
  \item 它们在某个层级上承载同一组受保护评价向量 $\mathrm{Ev}$ 的语义（Definition~\ref{def:protected-evaluation-functionals}）；
  \item 它们的层级差异仅来自有限步的 $\mathcal{E}_{\mathrm{prot}}$-相对可回缩扩展与其逆（Definition~\ref{def:relative-retractable-extension}
    的“加一层/去一层”）；
\end{enumerate}
则两塔落入同一 $\mathcal{E}_{\mathrm{prot}}$-相对最小包络同伦类（Definition~\ref{def:relative-minimal-envelope-class}）。
换言之：包络层数（包括更深递归解耦、非因果解耦）并不决定方案本质；本质由其在 $\mathcal{E}_{\mathrm{prot}}$ 下的相对同伦类决定。
\end{theorem}

\begin{certificate}[数学归纳法证书：按“可回缩扩展步数 $N$”归纳]
\label{cert:relative-minimal-envelope-induction}
设两条丛塔之间存在一个长度为 $N$ 的“加一层/去一层”序列，每一步要么应用 Proposition~\ref{prop:relative-homotopy-depth-invariant}，
要么应用其逆向（去层）。
令谓词 $P(N)$ 为：“经过 $N$ 步相对可回缩变换后，两塔在 $\mathcal{E}_{\mathrm{prot}}$ 下相对同伦等价，并落入同一相对最小包络类”。
证书化要求为：
\[
P(0)\ \text{成立},\qquad
(\forall N\in\mathbb{N})\ \bigl(P(N)\Rightarrow P(N+1)\bigr).
\]
\end{certificate}

\begin{proof}[证明（数学归纳法证书；谓词因果链）]
由 Certificate~\ref{cert:relative-minimal-envelope-induction}，取谓词链：
\[
\begin{aligned}
A:&\ P(0)\Rightarrow \text{零步时两塔相对同伦等价（平凡）},\\
B:&\ P(N)\wedge \text{第 $N{+}1$ 步为相对可回缩扩展/其逆}\Rightarrow P(N+1),
\end{aligned}
\]
其中 $B$ 由 Proposition~\ref{prop:relative-homotopy-depth-invariant}（及其逆向）给出。
由数学归纳法得对任意有限 $N$ 成立，从而两塔在有限可回缩变换下落入同一相对最小包络类。证毕。
\end{proof}

\subsection{定理（公理降级）：没有受保护泛函则“最小同伦类”不可判定}

\begin{theorem}[必要性（公理降级）：缺少 $\mathcal{E}_{\mathrm{prot}}$ 则 $\JacGap$ 与最小类失去工程语义]
\label{thm:protected-evaluations-necessary}
若不固定受保护评价泛函族 $\mathcal{E}_{\mathrm{prot}}$（或等价地，不固定“observable low-order behaviour”的取值口径），
则：
\begin{enumerate}
  \item “相对同伦”等价关系退化为“无约束同伦”，从而包络层数变化可通过改变观测口径被错误地视为同一类；
  \item $\JacGap(L)=0$ 的严格化语义失去约束锚点：同伦扭曲可以通过改变低阶可观测量来制造 $\JacGap$ 的虚假消失，
  使 defect potential 的下降不再对应“同一工程语义下的修复/严格化”；
  \item 因此“最小同伦类”不具备可判定语义：缺少 $\mathcal{E}_{\mathrm{prot}}$ 时，所谓“最小类”无法区分
  “真实降低缺陷且保持工程指标”与“改变工程指标使缺陷消失”。
\end{enumerate}
\end{theorem}

\begin{certificate}[证书：严格化判据必须相对受保护量]
\label{cert:protected-evaluations-necessary}
由 Theorem~\ref{thm:jacgap-zero-iff-strictifiable}，$\JacGap(L)=0$ 的等价陈述显式依赖于“低阶可观测量谓词”。
若该谓词不固定，则“严格化不改变什么”不固定，从而 $\JacGap$ 的消失可以通过改变观测口径而获得，导致证书语义漂移。
这等价于在工程上允许“通过改变评价指标把缺陷势能人为压低”，从而形成不可审计的路线自洽偏航。
\end{certificate}

\begin{proof}[证明（谓词因果链）]
由 Certificate~\ref{cert:protected-evaluations-necessary}，取谓词链：
\[
\begin{aligned}
A:&\ \text{不固定 }\mathcal{E}_{\mathrm{prot}}\Rightarrow \text{不固定 }\mathrm{Obs}_{\le r},\\
B:&\ \text{不固定 }\mathrm{Obs}_{\le r}\Rightarrow \JacGap(L)=0\ \text{的“保持不变”语义漂移},\\
C:&\ \text{语义漂移}\Rightarrow \text{缺陷势能下降可由改变工程指标造成},\\
D:&\ C\Rightarrow \text{最小同伦类不可判定（无法区分真修复与假修复）}.
\end{aligned}
\]
由 $A\wedge B\Rightarrow C$，再由 $C\Rightarrow D$ 得证。证毕。
\end{proof}

\subsection{推论：测地线同伦扭曲必须是相对同伦过程（防止假回归）}

\begin{corollary}[相对回归推论：沿 $J$‑测地线的同伦扭曲必须保持 $\mathcal{E}_{\mathrm{prot}}$]
\label{cor:relative-geodesic-twist}
在第 13 章与 \textsf{IG-RTSC} 章的口径下（Theorem~\ref{thm:ig-disturbance-regression}），扰动后回归由“测地线方向 $\gamma$ + 同伦扭曲 $H$”组成。
若希望回归具有可审计工程语义，则该同伦扭曲必须满足 Definition~\ref{def:protected-evaluation-functionals} 的相对约束：
\[
\mathrm{Ev}\bigl(H(s,t)\bigr)=\mathrm{Ev}\bigl(H(0,t)\bigr).
\]
否则回归可能通过改变受保护评价泛函制造“假回归”，与 Theorem~\ref{thm:protected-evaluations-necessary} 相矛盾。
\end{corollary}

\begin{certificate}[证书：缺少相对约束将触发“假回归/假严格化”]
\label{cert:relative-geodesic-twist}
由 Theorem~\ref{thm:protected-evaluations-necessary}，若不固定并保持 $\mathcal{E}_{\mathrm{prot}}$，则 $\JacGap$
  与缺陷势能下降可被“改变评价口径”伪造。
因此在采用“测地线方向 $\gamma$ + 同伦扭曲 $H$”编码回归时（Theorem~\ref{thm:ig-disturbance-regression}），
$H$ 必须是相对同伦：在扭曲过程中保持 $\mathcal{E}_{\mathrm{prot}}$ 不变。
\end{certificate}

\begin{proof}[证明（谓词因果链）]
取谓词链：
\[
\begin{aligned}
A:&\ \text{希望回归可审计}\Rightarrow \text{回归必须保持工程语义不变},\\
B:&\ \text{保持工程语义不变}\Rightarrow \text{需保持 }\mathcal{E}_{\mathrm{prot}}\ \text{不变},\\
C:&\ \text{保持 }\mathcal{E}_{\mathrm{prot}}\ \text{不变}\Rightarrow \text{同伦扭曲是相对同伦过程}.
\end{aligned}
\]
由 $A\Rightarrow B\Rightarrow C$ 得证。证毕。
\end{proof}

\begin{remark}[一句话收束]
包络层数并非本质；本质是 $\mathcal{E}_{\mathrm{prot}}$ 给定时的相对同伦类。
层数增加若为相对可回缩扩展，则不改变“最小同伦类”；而缺少 $\mathcal{E}_{\mathrm{prot}}$ 会使 $\JacGap$ 与回归机制失去可判定语义锚点，从而导致不可审计的“假严格化/假回归”。
\end{remark}

\clearpage
%%%%%%%%%%%%%%%%%%%%%%%%%%%%%%%%%%%%%%%%%%%%%%%%%%%%%%%%%%%%
\section{纯数卷口径：开放系统对抗博弈与防御策略同伦类（防御塔示例）}
\label{sec:puremath-open-system-adversarial-game}
%%%%%%%%%%%%%%%%%%%%%%%%%%%%%%%%%%%%%%%%%%%%%%%%%%%%%%%%%%%%

\subsection{形式逻辑表达约定}

\begin{remark}[目标：把“防御塔—基地—敌方攻击”写成相对同伦上的开放系统对抗博弈]
本章用纯数卷语言把“防御塔—基地—敌方攻击”的直觉严格化为一个开放系统对抗博弈（open-system adversarial game）结构：
\begin{enumerate}
  \item \textbf{被保护评价：}基地健康评价 $\mathrm{Ev}_{\mathrm{HP}}$ 作为受保护评价泛函族 $\mathcal{E}_{\mathrm{prot}}$
    的锚点（Definition~\ref{def:protected-evaluation-functionals}）；
  \item \textbf{工程方案：}防御工事/塔布局与升级由切面 $s$ 与包络 $(E,\Xi)$ 生成（Definition~\ref{def:ig-section-envelope}）；
  \item \textbf{对抗输入：}敌方攻击序列作为外生/对抗输入 $d(t)$，通过法联络/三层联络诱导水平法流并可触发曲率尖峰；
  \item \textbf{策略等价：}工程设计在相对 $\mathcal{E}_{\mathrm{prot}}$ 的同伦等价下，对应同一防御策略等价类（即同一相对同伦类）。
\end{enumerate}
因此“工程设计同伦等价于防御策略”的严格含义是：在固定 $\mathcal{E}_{\mathrm{prot}}$（尤其 $\mathrm{Ev}_{\mathrm{HP}}$）的前提下，对包络化设计取相对同伦商，
  与对抗博弈中的策略等价类取商，得到同一个“策略类对象”。
\end{remark}

\subsection{定义：基地健康评价 $\mathrm{Ev}_{\mathrm{HP}}$ 与相对锚点}

\begin{definition}[基地健康评价：$\mathrm{Ev}_{\mathrm{HP}}$ 作为 $\mathcal{E}_{\mathrm{prot}}$ 的锚点]
\label{def:base-health-evhp}
在 $\mathrm{Loc}\subseteq \GZLS$ 上，取受保护评价泛函族 $\mathcal{E}_{\mathrm{prot}}$（Definition~\ref{def:protected-evaluation-functionals}）。
称 $\mathrm{Ev}_{\mathrm{HP}}:\mathrm{Loc}\to\mathbb{R}_{\ge 0}$ 为基地健康评价泛函，若它被纳入受保护族：
\[
\mathrm{Ev}_{\mathrm{HP}}\in \mathcal{E}_{\mathrm{prot}}.
\]
给定阈值 $\underline{h}\ge 0$，称“基地不崩溃”约束为
\[
\mathrm{Ev}_{\mathrm{HP}}(L(t))\ge \underline{h}\qquad (\forall t\in[0,T]).
\]
并将其视为相对同伦约束的一部分：允许的同伦扭曲/包络动作不得通过“改变评价口径”来伪造基地健康（参见 Theorem~\ref{thm:protected-evaluations-necessary}）。
\end{definition}

\subsection{定义：开放系统对抗法流（双输入水平法流）}

\begin{definition}[开放系统对抗法流：$\dot L=\mathcal{H}(L,u,d;\GZLOC,\GZTLC)$]
\label{def:open-system-adversarial-law-flow}
称系统在法空间上形成开放系统对抗结构，若存在演化算子
\[
\mathcal{H}:\ \mathrm{Loc}\times \mathcal{U}\times \mathcal{D}\to T(\GZLS),
\]
使得对任意防御输入 $u(t)\in\mathcal{U}$ 与对抗输入 $d(t)\in\mathcal{D}$，法对象轨线 $L(t)\in\mathrm{Loc}\subseteq\GZLS$ 满足（schematic）
\[
\dot L(t)=\mathcal{H}\bigl(L(t),u(t),d(t);\GZLOC,\GZTLC\bigr),
\qquad
L(t)\in \mathrm{Loc}.
\]
其中：
\begin{enumerate}
  \item $u(t)$ 表示我方动作（建设/调控/包络移动/同伦修复等的统称）；
  \item $d(t)$ 表示外部扰动/对抗输入（敌方攻击序列、外场冲击、浪涌等的统称）；
  \item $\mathrm{Loc}$ 与阈值集合 $\Sigma$ 约束系统只能在可操作域内演化，并允许在跨阈时进行连续–离散切换（Definition~\ref{def:law-space-gzls}）。
\end{enumerate}
\end{definition}

\begin{remark}[曲率尖峰与缺陷势能尖峰]
在 Definition~\ref{def:open-system-adversarial-law-flow} 下，外生对抗输入往往通过对象层嵌入诱发法曲率 $\GZLCurv$ 的短时尖峰；
由 Theorem~\ref{thm:law-curvature-spike-jacgap}，曲率尖峰会放大缺陷势能 $\Phi=\|\JacGap\|_{\mathrm{def}}$ 的变化并产生尖峰，
这对应“敌方攻击会使总雅可比间隙/缺陷势能骤升”的法空间表达。
\end{remark}

\subsection{定义：对抗目标（基地健康硬约束 + 缺陷势能软目标）}

\begin{definition}[对抗目标：$\min_u\max_d$ 的值函数口径]
\label{def:adversarial-objective}
给定缺陷势能 $\Phi(L)=\|\JacGap(L)\|_{\mathrm{def}}$（Definition~\ref{def:ig-jacgap-defect}）与控制代价 $\mathrm{Cost}(u)\ge 0$。
定义在时间窗 $[0,T]$ 上的对抗代价为
\[
\mathcal{J}(u(\cdot),d(\cdot))
:=
\int_0^T
\Bigl(\Phi(L(t))+\mathrm{Cost}\bigl(u(t)\bigr)\Bigr)\,dt,
\]
并将基地健康作为硬约束（Definition~\ref{def:base-health-evhp}）：
\[
\mathrm{Ev}_{\mathrm{HP}}(L(t))\ge \underline{h}\qquad (\forall t\in[0,T]).
\]
定义开放系统对抗博弈的值函数（schematic）为
\[
V^\star
:=
\inf_{u(\cdot)\in\mathcal{U}_{\mathrm{adm}}}
\ \sup_{d(\cdot)\in\mathcal{D}_{\mathrm{adm}}}
\ \mathcal{J}\bigl(u(\cdot),d(\cdot)\bigr),
\]
其中 $\mathcal{U}_{\mathrm{adm}},\mathcal{D}_{\mathrm{adm}}$ 表示在 $\mathrm{Loc}/\Sigma$ 约束下允许的输入族，并要求演化轨线满足基地健康硬约束。
\end{definition}

\subsection{定义：工程设计（切面+包络）作为策略实现}

\begin{definition}[工程设计作为策略实现：切面 $s$ 与包络 $\Xi$]
\label{def:design-as-defense-strategy}
在防御塔语境中，将一个“具体防御工事方案”抽象为切面与包络结构（Definition~\ref{def:ig-section-envelope}）：
\[
s:\ U\to P,\qquad s(U)\subseteq E,\qquad \Xi=\Xi^{\mathrm{Geom}}\times \Xi^{\mathrm{Sim}}\times \Xi^{\mathrm{Proc}}.
\]
其中：
\begin{enumerate}
  \item 切面 $s$ 表示固定的路线/表示/协议/管线选择（对应固定的一套“塔布置与规则口径”）；
  \item 包络坐标 $\Xi$ 表示可升级/可重构/可调整的细粒度自由度（塔参数、位置微扰、阈值策略、联动规则、鲁棒参数等），但其允许变化必须满足受保护评价泛函的相对约束（Definition~\ref{def:protected-evaluation-functionals}）。
\end{enumerate}
将“策略”理解为：在包络坐标与同伦修复允许的动作集合内，随时间生成 $u(t)$ 的规则；从而工程设计（含包络）给出对抗博弈中的一个可执行策略族。
\end{definition}

\subsection{定义：策略的相对同伦等价（相对 $\mathcal{E}_{\mathrm{prot}}$）}

\begin{definition}[策略等价：相对同伦下的防御策略类]
\label{def:relative-strategy-equivalence}
给定两套工程设计/策略实现 $\mathcal{S}_1,\mathcal{S}_2$（Definition~\ref{def:design-as-defense-strategy}）。
称它们属于同一防御策略等价类，若存在相对同伦等价（Definition~\ref{def:relative-homotopy-equivalence}）使得：
\begin{enumerate}
  \item 相对约束以 $\mathcal{E}_{\mathrm{prot}}$ 为锚点，且包含基地健康评价 $\mathrm{Ev}_{\mathrm{HP}}$（Definition~\ref{def:base-health-evhp}）；
  \item 等价映射保持“基地健康硬约束 + 缺陷势能下降机制”的语义：即在相同的 $\mathrm{Loc}/\Sigma$ 管理口径下，二者对 $\Phi$ 的可下降性与对 $\mathrm{Ev}_{\mathrm{HP}}$
    的可维持性属于同一相对同伦类。
\end{enumerate}
记作 $\mathcal{S}_1\simeq_{\mathcal{E}_{\mathrm{prot}}}\mathcal{S}_2$，并将商空间称为“防御策略同伦类”。
\end{definition}

\subsection{命题：设计—策略对应（在相对同伦商意义下）}

\begin{proposition}[对应命题：工程设计（含包络）在相对同伦商下等价于防御策略类]
\label{prop:design-strategy-correspondence}
在 Definition~\ref{def:adversarial-objective}--Definition~\ref{def:relative-strategy-equivalence} 的设定下，
对“切面+包络”工程设计取相对同伦商（以 $\mathcal{E}_{\mathrm{prot}}$ 为锚点，包含 $\mathrm{Ev}_{\mathrm{HP}}$），
与对抗博弈中“可执行策略”取相对同伦商（同一锚点）在语义上同型：二者都刻画同一类“在基地健康语义固定时可实现的防御策略等价类”。
\end{proposition}

\begin{certificate}[证书：由相对同伦锚点与包络/策略实现关系直接推出]
\label{cert:design-strategy-correspondence}
\begin{enumerate}
  \item 由 Definition~\ref{def:design-as-defense-strategy}，工程设计给出策略实现的动作空间（切面固定路线，包络提供可调自由度）；
  \item 由 Definition~\ref{def:protected-evaluation-functionals} 与 Definition~\ref{def:base-health-evhp}，$\mathrm{Ev}
    _{\mathrm{HP}}$ 作为受保护量固定了“防御成功”的语义锚点；
  \item 由 Definition~\ref{def:relative-strategy-equivalence}，策略等价正是相对 $\mathcal{E}_{\mathrm{prot}}$ 的同伦等价；
    因此对设计/对策略取商得到同一类对象：相对同伦类本身。
\end{enumerate}
\end{certificate}

\begin{proof}[证明（谓词因果链）]
由 Certificate~\ref{cert:design-strategy-correspondence}，取谓词链：
\[
\begin{aligned}
A:&\ \text{设计}\Rightarrow \text{策略实现（动作空间）},\\
B:&\ \mathrm{Ev}_{\mathrm{HP}}\in\mathcal{E}_{\mathrm{prot}}\Rightarrow \text{防御语义锚点固定},\\
C:&\ \text{等价关系}= \simeq_{\mathcal{E}_{\mathrm{prot}}}\Rightarrow \text{商对象为相对同伦类}.
\end{aligned}
\]
由 $A\wedge B\wedge C$ 得到“设计商对象与策略商对象同型”的结论。证毕。
\end{proof}

\subsection{定理（公理降级）：相对同伦策略类的不变量（基地健康可行性）}

\begin{theorem}[不变量（公理降级）：相对同伦等价保持基地健康可行性]
\label{thm:relative-strategy-feasibility-invariant}
设 $\mathcal{S}_1\simeq_{\mathcal{E}_{\mathrm{prot}}}\mathcal{S}_2$（Definition~\ref{def:relative-strategy-equivalence}），
  且等价映射与离散时间步的演化算子相容：
存在一步演化映射 $\mathcal{T}$ 使 $L_{n+1}=\mathcal{T}(L_n,u_n,d_n)$，并且 $\mathcal{S}_1,\mathcal{S}_2$ 的等价映射在相对约束下与 $\mathcal{T}$
  交换（conjugacy on-step，schematic）。
则在同一攻击输入族 $\mathcal{D}_{\mathrm{adm}}$ 下，“存在防御策略使 $\mathrm{Ev}_{\mathrm{HP}}(L_n)\ge \underline{h}$ 对所有 $n\le N$
  成立”的可行性在 $\mathcal{S}_1,\mathcal{S}_2$ 之间保持不变。
\end{theorem}

\begin{certificate}[数学归纳法证书：按离散步数 $N$ 归纳]
\label{cert:relative-strategy-feasibility-induction}
令谓词 $P(N)$ 为：“对任意允许对抗输入序列 $\{d_n\}_{n=0}^{N-1}$，$\mathcal{S}_1$ 可行（存在 $\{u_n\}$ 使 $\mathrm{Ev}_{\mathrm{HP}}(L_n)\ge
  \underline{h}$）当且仅当 $\mathcal{S}_2$ 可行”。
证书化要求为：
\[
P(0)\ \text{成立},\qquad
(\forall N\in\mathbb{N})\ \bigl(P(N)\Rightarrow P(N+1)\bigr).
\]
\end{certificate}

\begin{proof}[证明（数学归纳法证书；谓词因果链）]
由 Certificate~\ref{cert:relative-strategy-feasibility-induction}，取谓词链：
\[
\begin{aligned}
A:&\ \mathcal{S}_1\simeq_{\mathcal{E}_{\mathrm{prot}}}\mathcal{S}_2\Rightarrow \mathrm{Ev}_{\mathrm{HP}}\
  \text{评价在等价映射下保持},\\
B:&\ \text{等价映射与 }\mathcal{T}\ \text{交换}\Rightarrow \text{可将 }\mathcal{S}_1\text{ 的可行控制序列映射为 }\mathcal{S}_2\text{
  的可行控制序列},\\
C:&\ A\wedge B\Rightarrow P(N)\Rightarrow P(N+1).
\end{aligned}
\]
基步 $P(0)$ 平凡成立。由 $C$ 得到归纳步，从而对任意 $N$ 成立。证毕。
\end{proof}

\subsection{备注：沿测地线的同伦扭曲在防御语境中的含义}

\begin{remark}[“先稳住基地健康，再降缺陷势能”的纯数卷翻译]
当对抗输入导致法曲率尖峰、使缺陷势能 $\Phi$ 突增时，纯数卷的回归机制可读作：
\begin{enumerate}
  \item \textbf{方向层：}在度量 $J$ 下选取回归测地线 $\gamma$（先选方向/诊断）；
  \item \textbf{构造层：}在保持受保护量（尤其 $\mathrm{Ev}_{\mathrm{HP}}$）不变的前提下，在固定 $\gamma$ 上做同伦扭曲 $H$ 并配合包络变量的快层修复（再做构造/治疗），使
    $\Phi$ 严格下降并回到 $\mathrm{Loc}$ 的证书化工作区。
\end{enumerate}
这与 Corollary~\ref{cor:relative-geodesic-twist}、Theorem~\ref{thm:ig-disturbance-regression} 的口径一致：防御塔直觉里的“先稳住基地健康，
  再升级重构回到优势策略域”恰好对应“相对同伦约束下的测地线同伦扭曲回归”。
\end{remark}

\clearpage
%%%%%%%%%%%%%%%%%%%%%%%%%%%%%%%%%%%%%%%%%%%%%%%%%%%%%%%%%%%%
\section{纯数卷口径：非交换算子词的设计生成与开放对抗策略生成等价（更接近 \textsf{OpenRA}）}
\label{sec:puremath-noncommutative-design-strategy}
%%%%%%%%%%%%%%%%%%%%%%%%%%%%%%%%%%%%%%%%%%%%%%%%%%%%%%%%%%%%

\subsection{形式逻辑表达约定}

\begin{remark}[等价链条：生成工程设计 $\simeq_{\mathcal{E}_{\mathrm{prot}}}$ 生成开放对抗策略]
本章把如下等价链条写成可直接引用的形式系统结论：
\begin{enumerate}
  \item “生成工程设计方案（含非交换算子）”可理解为：生成一套可执行的构造/调控方案，使法对象 $L(t)\in\mathrm{Loc}\subseteq\GZLS$ 在外生扰动/对抗输入下仍满足受保护评价泛函族
    $\mathcal{E}_{\mathrm{prot}}$ 的语义锚点；
  \item “生成开放博弈系统中的防御/对抗策略（含非交换算子）”可理解为：在开放系统对抗法流（Definition~\ref{def:open-system-adversarial-law-flow}）下生成策略 $\pi$
    与可执行算子词（operator words），并在相同的 $\mathrm{Loc}/\Sigma$ 与 $\mathcal{E}_{\mathrm{prot}}$ 锚点下比较/取商；
  \item 二者的等价必须是\textbf{相对} $\mathcal{E}_{\mathrm{prot}}$ 的等价：若不固定受保护评价泛函族，则会退化为“改变评价口径即可宣称胜利”的伪等价（Theorem~\ref{thm:protected-evaluations-necessary}）。
\end{enumerate}
本章用“非交换算子词 + 相对同伦策略类”的口径，把上述三点组织为定义—命题—定理—推论结构，并给出为何其时间/信息结构更接近 \textsf{OpenRA} 这类实时开放对抗，而非 \textsf{AlphaGo} 式回合制封闭对弈。
\end{remark}

\subsection{定义：非交换算子体系、算子词与在线策略}

\begin{definition}[非交换算子体系：我方算子与对抗算子]
\label{def:noncommutative-operator-systems}
定义我方可执行算子幺半群（或幺半群胚）与对抗算子幺半群为
\[
\mathcal{M}_{\mathrm{self}},
\qquad
\mathcal{M}_{\mathrm{adv}},
\]
并允许一般情况下的非交换性：
\[
\exists\,u_1,u_2\in\mathcal{M}_{\mathrm{self}}\ \text{使}\ u_1u_2\neq u_2u_1,
\qquad
\exists\,d_1,d_2\in\mathcal{M}_{\mathrm{adv}}\ \text{使}\ d_1d_2\neq d_2d_1.
\]
记对应算子词空间为
\[
\mathcal{W}_{\mathrm{self}}:=\mathcal{M}_{\mathrm{self}}^\ast,
\qquad
\mathcal{W}_{\mathrm{adv}}:=\mathcal{M}_{\mathrm{adv}}^\ast,
\]
其中词的乘法为串接；非交换性表达“动作顺序改变结果”的工程事实（工艺顺序、装配顺序、布置/布线顺序、控制协议顺序等）。
\end{definition}

\begin{definition}[观测与策略：从低阶可观测信息到非交换动作]
\label{def:noncommutative-policy}
在 $\mathrm{Loc}\subseteq\GZLS$ 上，取观测映射
\[
\mathrm{Obs}:\ \mathrm{Loc}\to \mathcal{O},
\]
其中 $\mathcal{O}$ 表示可观测低阶信息/局部化信息（可包含证书字段与阈值状态）。
称策略 $\pi$ 为一条（确定性或随机的）在线策略，若
\[
\pi:\ \mathcal{O}\to \mathcal{M}_{\mathrm{self}}
\quad\text{或更一般}\quad
\pi:\ \mathcal{O}\to \Delta(\mathcal{M}_{\mathrm{self}}),
\]
并由此生成非交换算子词（在线串接）：
\[
w_{\mathrm{self}}=\pi(o_0)\,\pi(o_1)\,\cdots\,\pi(o_{N-1})\in \mathcal{W}_{\mathrm{self}},
\qquad
o_n=\mathrm{Obs}(L_n).
\]
对抗侧同理可由 $d(\cdot)$ 生成 $w_{\mathrm{adv}}\in\mathcal{W}_{\mathrm{adv}}$。
\end{definition}

\subsection{命题：生成设计等价于生成策略（相对 $\mathcal{E}_{\mathrm{prot}}$）}

\begin{proposition}[对应命题：开放系统下的“设计生成”必须显式给出策略生成器]
\label{prop:design-generation-equals-strategy-generation}
在开放系统对抗法流（Definition~\ref{def:open-system-adversarial-law-flow}）下，若“工程设计方案”要求在对抗输入 $d(\cdot)$ 存在时仍保持受保护评价泛函族 $\mathcal{E}
  _{\mathrm{prot}}$ 的语义锚点（Definition~\ref{def:protected-evaluation-functionals}，并可包含基地健康评价 $\mathrm{Ev}_{\mathrm{HP}}$，
  Definition~\ref{def:base-health-evhp}），
则“生成工程设计方案”在形式系统层面等价于“生成在线策略 $\pi$（以及其诱导的非交换算子词）”。
\end{proposition}

\begin{certificate}[证书：开放性 + 非交换性 $\Rightarrow$ 静态对象不足以定义方案]
\label{cert:design-generation-equals-strategy-generation}
\begin{enumerate}
  \item（开放性）在开放对抗结构中，$d(\cdot)$ 持续注入并可触发曲率尖峰与阈值切换；仅给出一个静态对象无法指定扰动到来后的可执行响应；
  \item（非交换性）动作顺序影响结果；因此必须显式给出“如何串接动作”的生成规则（算子词/策略）；
  \item（受保护锚点）若不相对 $\mathcal{E}_{\mathrm{prot}}$ 取等价，则可通过改变评价口径伪造“成功”，与 Theorem~\ref{thm:protected-evaluations-necessary}
    相冲突。
\end{enumerate}
\end{certificate}

\begin{proof}[证明（谓词因果链）]
由 Certificate~\ref{cert:design-generation-equals-strategy-generation}，取谓词链：
\[
\begin{aligned}
A:&\ \text{开放性}\Rightarrow \text{需要在线响应规则},\\
B:&\ \text{非交换性}\Rightarrow \text{响应规则必须生成算子词},\\
C:&\ \text{受保护锚点}\Rightarrow \text{等价必须相对 }\mathcal{E}_{\mathrm{prot}}.
\end{aligned}
\]
由 $A\wedge B\Rightarrow$ “设计生成 $\Rightarrow$ 策略生成”，并由 $C$ 给出“相对锚点”要求。证毕。
\end{proof}

\subsection{定理（公理降级）：含非交换算子的等价链在相对同伦商下成立}

\begin{theorem}[等价链（公理降级）：非交换工程设计生成与开放对抗策略生成同型]
\label{thm:noncommutative-design-strategy-equivalence}
固定受保护评价泛函族 $\mathcal{E}_{\mathrm{prot}}$（Definition~\ref{def:protected-evaluation-functionals}）。
在开放系统对抗法流（Definition~\ref{def:open-system-adversarial-law-flow}）与对抗目标（Definition~\ref{def:adversarial-objective}）下，
对“切面+包络”的工程设计（Definition~\ref{def:design-as-defense-strategy}）取相对同伦商（Definition~\ref{def:relative-strategy-equivalence}），

与对抗博弈中的在线策略取相对同伦商在语义上同型；因此成立相对等价：
\[
\text{生成（含非交换算子）的工程设计}
\ \simeq_{\mathcal{E}_{\mathrm{prot}}}\
\text{生成（含非交换算子）的开放对抗策略}.
\]
\end{theorem}

\begin{certificate}[数学归纳法证书：按离散步数 $N$ 的策略可行性归纳]
\label{cert:noncommutative-design-strategy-equivalence-induction}
令 $N$ 为离散时间步数。定义谓词 $P(N)$ 为：“对任意长度为 $N$ 的对抗输入词 $w_{\mathrm{adv}}\in\mathcal{W}_{\mathrm{adv}}$，
工程设计（切面+包络）所诱导的策略可在相对同伦商意义下被识别为同一策略类，且保持受保护评价泛函族的语义锚点”。
证书化要求为：
\[
P(0)\ \text{成立},\qquad
(\forall N\in\mathbb{N})\ \bigl(P(N)\Rightarrow P(N+1)\bigr).
\]
\end{certificate}

\begin{proof}[证明（数学归纳法证书；谓词因果链）]
由 Certificate~\ref{cert:noncommutative-design-strategy-equivalence-induction}，取谓词链：
\[
\begin{aligned}
A:&\ \text{设计—策略对应在相对同伦商下同型}\Rightarrow \text{一步等价可被保持},\\
B:&\ \text{一步等价可被保持}\Rightarrow P(N)\Rightarrow P(N+1),\\
C:&\ \text{相对同伦锚点固定}\Rightarrow \text{等价不允许通过改变评价口径实现}.
\end{aligned}
\]
基步 $P(0)$ 平凡成立。归纳步由 $A\wedge B$ 给出。$C$ 排除伪等价（Theorem~\ref{thm:protected-evaluations-necessary}）。
因此对任意 $N$ 成立，得到所述相对等价链。证毕。
\end{proof}

\subsection{推论：时间/信息结构更接近 \textsf{OpenRA} 而非 \textsf{AlphaGo}}

\begin{corollary}[结构对比推论：实时开放对抗更贴近 \textsf{OpenRA} 语义]
\label{cor:openra-vs-alphago}
在 Theorem~\ref{thm:noncommutative-design-strategy-equivalence} 的设定下，
“工程设计生成/策略生成”对应的开放系统对抗结构更接近 \textsf{OpenRA} 这类实时开放对抗（开放、实时、非平稳、部分可观测、强非交换），
而非 \textsf{AlphaGo} 式回合制封闭对弈（规则固定、状态转移离散且环境封闭）。
\end{corollary}

\begin{certificate}[证书：三类差异——开放性、非交换性、信息结构]
\label{cert:openra-vs-alphago}
\begin{enumerate}
  \item（开放性）存在持续注入的对抗输入 $d(\cdot)$ 与阈值切换 $\Sigma$，需要在线回归与同伦扭曲修复（Theorem~\ref{thm:ig-disturbance-regression}）；
  \item（非交换性）动作顺序敏感，策略天然以非交换算子词编码（Definition~\ref{def:noncommutative-policy}）；
  \item（信息结构）策略依赖局部化可观测信息 $\mathrm{Obs}(L)$，并在不完整信息与模型失配下运行，需以 $\mathcal{E}_{\mathrm{prot}}$ 锚定语义避免“换口径即胜利”。
\end{enumerate}
上述三点与实时开放对抗范式一致，而与回合制封闭对弈范式不同。
\end{certificate}

\begin{proof}[证明（谓词因果链）]
由 Certificate~\ref{cert:openra-vs-alphago}，取谓词链：
\[
\begin{aligned}
A:&\ \text{开放性}\wedge \text{阈值切换}\Rightarrow \text{需要在线回归（非封闭）},\\
B:&\ \text{非交换性}\Rightarrow \text{策略以算子词/时序编码（实时性强）},\\
C:&\ \text{部分可观测}\wedge \mathcal{E}_{\mathrm{prot}}\Rightarrow \text{需证书化语义锚点（避免伪胜利）}.
\end{aligned}
\]
由 $A\wedge B\wedge C$ 得到“更贴近 \textsf{OpenRA} 语义”的结论。证毕。
\end{proof}

\begin{remark}[推广：发动机、芯片等一般工程系统]
在 pure 卷口径下，上述等价并不依赖“防御塔”语义本身，而依赖三项结构事实：开放性（外生扰动/变异持续注入）、非交换性（顺序敏感的工艺/控制链）、以及受保护评价泛函族 $\mathcal{E}_{\mathrm{prot}}
  $（工程语义锚点）。
因此对发动机、芯片等工程系统，只要把性能/安全/可制造性等评价写成 $\mathcal{E}_{\mathrm{prot}}$ 的有限族，并把可执行动作组织为非交换算子体系，就可沿同一模板把“设计生成”转写为“开放对抗策略生成”，
  并用同伦扭曲/包络修复在扰动下维持可控性与可审计性。
\end{remark}

\clearpage
%%%%%%%%%%%%%%%%%%%%%%%%%%%%%%%%%%%%%%%%%%%%%%%%%%%%%%%%%%%%
\section{纯数卷口径：开放对抗 \OPB{} 的基准模板（芯片/发动机最小基准包）}
\label{sec:puremath-open-adversarial-opb-benchmark-template}
%%%%%%%%%%%%%%%%%%%%%%%%%%%%%%%%%%%%%%%%%%%%%%%%%%%%%%%%%%%%

\subsection{形式逻辑表达约定}

\begin{remark}[从等价链到“可运行基准包”]
Theorem~\ref{thm:noncommutative-design-strategy-equivalence} 已把“生成工程设计（非交换算子词）”与“生成开放对抗策略（非交换算子词）”在相对同伦商下严格对应。
本章不再重复证明，而把该等价链压成一个可直接复用的\textbf{基准模板}：
\begin{enumerate}
  \item \textbf{对象：}把开放系统对抗法流（Definition~\ref{def:open-system-adversarial-law-flow}）与主丛化表示冗余（$G_{\mathrm{rep}}$）统一为“开放对抗
    \OPB{}”对象；
  \item \textbf{锚点：}以受保护评价泛函族 $\mathcal{E}_{\mathrm{prot}}$ 锁定工程语义（Definition~\ref{def:protected-evaluation-functionals}）
    ，避免“换口径即胜利”的伪等价（Theorem~\ref{thm:protected-evaluations-necessary}）；
  \item \textbf{最小包：}给出两个工程实例（芯片与发动机）的“最小基准包”
  \[
  \bigl(
    \mathcal{E}_{\mathrm{prot}},
    \mathcal{G}_{\mathrm{self}},
    \mathcal{G}_{\mathrm{adv}},
    \Xi,
    \JacGap_{\min}
  \bigr),
  \]
  以说明它们并非两套体系，而是同一 open adversarial \OPB{} 模板下的不同扇区/切面/包络选择。
\end{enumerate}
\end{remark}

\subsection{定义：开放对抗 \OPB{}（open adversarial OPB）}

\begin{definition}[开放对抗 \OPB{}：主丛化的开放博弈基准对象]
\label{def:open-adversarial-opb-object}
在法空间 $\GZLS=(\mathfrak{L};J,\mathrm{Loc},\Sigma)$（Definition~\ref{def:law-space-gzls}）上，
称一个工程域的“开放对抗 \OPB{} 基准对象”由如下数据给出：
\[
\OPB_{\mathrm{adv}}
:=
\bigl(
  M,\ P,\ \pi,\ G_{\mathrm{rep}},\ \Conn;\
  \mathcal{E}_{\mathrm{prot}},\ \mathcal{M}_{\mathrm{self}},\ \mathcal{M}_{\mathrm{adv}},\ \Xi;\
  \mathcal{H},\ \JacGap,\ \Phi
\bigr),
\]
其中：
\begin{enumerate}
  \item $M$ 为“可判定基底”（工程基准空间），并假定存在嵌入 $\iota:M\hookrightarrow \mathrm{Loc}$ 以将基底点解释为法对象代表；
  \item $\pi:P\to M$ 为主 $G_{\mathrm{rep}}$‑丛，$G_{\mathrm{rep}}$ 吸收表示冗余（几何参数化、离散化、求解器设置、工艺叙述等）；
  \item $\Conn$ 为允许的联络/三层联络数据（与 $\GZLOC/\GZTLC$ 兼容），用于刻画水平法流与曲率效应；
  \item $\mathcal{E}_{\mathrm{prot}}$ 为受保护评价泛函族（Definition~\ref{def:protected-evaluation-functionals}），作为相对同伦约束锚点；
  \item $\mathcal{M}_{\mathrm{self}},\mathcal{M}_{\mathrm{adv}}$ 为我方与对抗侧的算子幺半群（Definition~\ref{def:noncommutative-operator-systems}），允许一般非交换；
  \item $\Xi=\Xi^{\mathrm{Geom}}\times\Xi^{\mathrm{Sim}}\times\Xi^{\mathrm{Proc}}$ 为包络自由度（Definition~\ref{def:ig-section-envelope}），作为纤维方向的可调细粒度变量；
  \item $\mathcal{H}$ 为开放系统对抗法流的演化算子（Definition~\ref{def:open-system-adversarial-law-flow}），以 $u(t)\in\mathcal{M}
    _{\mathrm{self}}$、$d(t)\in\mathcal{M}_{\mathrm{adv}}$ 作为双输入；
  \item $\JacGap$ 与缺陷势能 $\Phi$ 为证书化缺陷对象（Definition~\ref{def:ig-jacgap-defect}），用于度量“可严格化/可回归”的失败程度。
\end{enumerate}
\end{definition}

\begin{remark}[解释：基底负责“判定”，纤维负责“吸收与修复”]
Definition~\ref{def:open-adversarial-opb-object} 的直观含义是：开放对抗发生在 $M$（或 $\iota(M)\subset\mathrm{Loc}$）上的法流演化；
而一切“等价表示/细节自由度”（含非因果重参数化、递归解耦等）进入纤维，由 $G_{\mathrm{rep}}$ 与包络 $\Xi$ 管理。
受保护评价泛函族 $\mathcal{E}_{\mathrm{prot}}$ 则规定了纤维方向允许怎样变化：允许修复/扭曲，但不允许偷换评价口径。
\end{remark}

\subsection{定义：最小缺陷三分量与缺陷势能（chart--lift--threshold）}

\begin{definition}[最小 \JacGap{} 三分量证书：chart--lift--threshold]
\label{def:opb-jacgap-minimal}
在 open adversarial \OPB{} 中，为了得到可迁移、可复用且可证书化的“最小缺陷势能”，可取最小三分量缺陷证书（schematic）：
\[
\JacGap_{\min}(L)
:=
\bigl(
  c_{\mathrm{chart}}(L),\ c_{\mathrm{lift}}(L),\ c_{\mathrm{thres}}(L)
\bigr),
\qquad
\Phi_{\min}(L):=\mathcal{V}\bigl(\JacGap_{\min}(L)\bigr),
\]
其中 $\mathcal{V}$ 为固定的聚合（范数/谱聚合等），并约定：
\begin{enumerate}
  \item $c_{\mathrm{chart}}$：跨表示/跨评估器（chart）一致性的闭合失配；
  \item $c_{\mathrm{lift}}$：在 $\GZTLC/\widetilde{\GZTLC}$ 的水平提升中，“保持受保护低阶可观测量”与“高阶同伦修复”之间的失配；
  \item $c_{\mathrm{thres}}$：跨阈值集合 $\Sigma$ 的一致性失配（连续–离散切换的证书口径）。
\end{enumerate}
当该三分量可计算且满足 pure 卷的“可严格化/可下降”要求时，$\Phi_{\min}$ 可直接作为开放对抗中的缺陷势能项进入值函数（Definition~\ref{def:adversarial-objective}）。
\end{definition}

\subsection{定理（公理降级）：基准模板的等价链可直接调用}

\begin{theorem}[基准模板定理（公理降级）：设计生成/策略生成/切面包络选择三者等价]
\label{thm:opb-benchmark-template-equivalence}
在 open adversarial \OPB{} 对象（Definition~\ref{def:open-adversarial-opb-object}）下，固定受保护评价泛函族 $\mathcal{E}_{\mathrm{prot}}$，

并允许非交换算子体系（Definition~\ref{def:noncommutative-operator-systems}）与在线策略（Definition~\ref{def:noncommutative-policy}）。
则以下三件事在相对 $\mathcal{E}_{\mathrm{prot}}$ 的意义下等价：
\[
\begin{aligned}
&\text{(i) 生成工程设计方案（允许非交换算子词）}\\
&\Longleftrightarrow\ \text{(ii) 生成开放系统对抗中的在线策略（允许非交换算子词）}\\
&\Longleftrightarrow\ \text{(iii) 在 }\OPB_{\mathrm{adv}}\text{ 上选取切面/包络并在相对同伦类内做法流与同伦扭曲 }H.
\end{aligned}
\]
\end{theorem}

\begin{certificate}[数学归纳法证书：复用离散步数归纳框架]
\label{cert:opb-benchmark-template-induction}
将定理的等价链表述为随离散步数 $N$ 的谓词 $P(N)$（“长度为 $N$ 的对抗输入词下，等价链在相对同伦商意义下成立并保持 $\mathcal{E}_{\mathrm{prot}}$”）。
其数学归纳法证书可直接复用 Certificate~\ref{cert:noncommutative-design-strategy-equivalence-induction} 与
Certificate~\ref{cert:relative-strategy-feasibility-induction} 的归纳结构：基步 $P(0)$ 平凡成立，归纳步由“一步演化交换性 + 相对约束保持”给出。
\end{certificate}

\begin{proof}[证明（数学归纳法证书；谓词因果链）]
由 Proposition~\ref{prop:design-generation-equals-strategy-generation} 得到 (i)$\Longleftrightarrow$(ii) 的“开放性 + 非交换性
  $\Rightarrow$ 方案必须给出策略生成器”结论；
由 Theorem~\ref{thm:noncommutative-design-strategy-equivalence} 得到 (ii)$\Longleftrightarrow$(iii) 的“相对同伦商下同型”结论；
由 Definition~\ref{def:open-adversarial-opb-object} 知 (iii) 的对象即为“切面+包络”在 $\OPB_{\mathrm{adv}}$
  上的选择与允许的同伦扭曲动作（Definition~\ref{def:design-as-defense-strategy}）。

形式化地，取谓词链：
\[
\begin{aligned}
A:&\ \text{开放系统对抗}\Rightarrow \text{静态对象不足以定义方案}\Rightarrow \text{需要在线策略生成器},\\
B:&\ \text{切面+包络}\Rightarrow \text{给出策略实现空间}\Rightarrow \text{(i) 与 (ii) 互释},\\
C:&\ \text{相对同伦商}\Rightarrow \text{(ii) 与 (iii) 的商对象同型}.
\end{aligned}
\]
由 $A\wedge B\wedge C$ 得到三者等价。归纳细节由 Certificate~\ref{cert:opb-benchmark-template-induction} 保证。证毕。
\end{proof}

\subsection{实例：芯片设计基准（Chip-\OPB{} Minimal Pack）}

\begin{definition}[芯片基准包：$\bigl(\mathcal{E}_{\mathrm{prot}},\mathcal{G}_{\mathrm{self}},\mathcal{G}_{\mathrm{adv}},\Xi,
  \JacGap_{\min}\bigr)$]
\label{def:chip-opb-minimal-pack}
定义芯片工程的最小基准包（benchmark pack）为
\[
\mathfrak{B}_{\mathrm{chip}}
:=
\bigl(
  \mathcal{E}_{\mathrm{prot}}^{\mathrm{chip}},
  \mathcal{G}_{\mathrm{self}}^{\mathrm{chip}},
  \mathcal{G}_{\mathrm{adv}}^{\mathrm{chip}},
  \Xi^{\mathrm{chip}},
  \JacGap_{\min}^{\mathrm{chip}}
\bigr),
\]
其中：
\begin{enumerate}
  \item（受保护评价）
  \[
  \mathcal{E}_{\mathrm{prot}}^{\mathrm{chip}}
  :=
  \bigl\{
    \mathrm{Ev}^{P},\ \mathrm{Ev}^{A},\ \mathrm{Ev}^{\mathrm{tim}},\ \mathrm{Ev}^{\mathrm{rel}},\ \mathrm{Ev}
      ^{\mathrm{mfg}}
  \bigr\},
  \]
  分别锚定功耗、面积、时序可行性、可靠性/裕度、可制造/可验证门槛；
  \item（我方生成元）取顺序敏感的动作生成元集合（示意）：
  \[
  \begin{aligned}
  \mathcal{G}_{\mathrm{self}}^{\mathrm{chip}}
  :=
  \{\, &g_{\mathrm{arch}},\ g_{\mathrm{logic}},\ g_{\mathrm{place}},\ g_{\mathrm{route}},\\
  &g_{\mathrm{clk}},\ g_{\mathrm{power}},\ g_{\mathrm{eco}}\,\},
  \end{aligned}
  \]
  并令 $\mathcal{M}_{\mathrm{self}}^{\mathrm{chip}}:=\langle \mathcal{G}_{\mathrm{self}}^{\mathrm{chip}}\rangle$；
  \item（对抗生成元）取开放扰动/变异生成元集合（示意）：
  \[
  \begin{aligned}
  \mathcal{G}_{\mathrm{adv}}^{\mathrm{chip}}
  :=
  \{\, &h_{\mathrm{proc}},\ h_{\mathrm{temp}},\ h_{\mathrm{vdd}},\ h_{\mathrm{noise}},\ h_{\mathrm{aging}}\,\},
  \end{aligned}
  \]
  并令 $\mathcal{M}_{\mathrm{adv}}^{\mathrm{chip}}:=\langle \mathcal{G}_{\mathrm{adv}}^{\mathrm{chip}}\rangle$；
  \item（包络自由度）
  \[
  \Xi^{\mathrm{chip}}:=\Xi^{\mathrm{Geom}}\times \Xi^{\mathrm{Sim}}\times \Xi^{\mathrm{Proc}},
  \]
  其中 $\Xi^{\mathrm{Geom}}$ 表示布局/几何细粒度可调自由度，$\Xi^{\mathrm{Sim}}$ 表示模型/离散化/验证口径自由度，$\Xi^{\mathrm{Proc}}$ 表示签核/制造约束口径自由度；
  \item（最小缺陷证书）取 Definition~\ref{def:opb-jacgap-minimal} 的三分量：
  \[
  \JacGap_{\min}^{\mathrm{chip}}(L)
  :=
  \bigl(c_{\mathrm{chart}}(L),c_{\mathrm{lift}}(L),c_{\mathrm{thres}}(L)\bigr).
  \]
\end{enumerate}
\end{definition}

\begin{remark}[芯片：非交换性来自“顺序敏感的闭环收敛”]
芯片设计流程的非交换性并非修辞：布局/布线/时钟/功耗/签核的顺序交换会改变可行域与证书字段；
因此其“方案”天然是非交换算子词（Definition~\ref{def:noncommutative-policy}），并且必须显式给出 $\Xi^{\mathrm{Sim}}$ 与 $\Xi^{\mathrm{Proc}}$，
  否则容易在单一评估口径下形成自洽偏航。
\end{remark}

\subsection{实例：发动机设计基准（Engine-\OPB{} Minimal Pack）}

\begin{definition}[发动机基准包：$\bigl(\mathcal{E}_{\mathrm{prot}},\mathcal{G}_{\mathrm{self}},\mathcal{G}_{\mathrm{adv}},\Xi,
  \JacGap_{\min}\bigr)$]
\label{def:engine-opb-minimal-pack}
定义发动机工程的最小基准包为
\[
\mathfrak{B}_{\mathrm{eng}}
:=
\bigl(
  \mathcal{E}_{\mathrm{prot}}^{\mathrm{eng}},
  \mathcal{G}_{\mathrm{self}}^{\mathrm{eng}},
  \mathcal{G}_{\mathrm{adv}}^{\mathrm{eng}},
  \Xi^{\mathrm{eng}},
  \JacGap_{\min}^{\mathrm{eng}}
\bigr),
\]
其中：
\begin{enumerate}
  \item（受保护评价）
  \[
  \mathcal{E}_{\mathrm{prot}}^{\mathrm{eng}}
  :=
  \bigl\{
    \mathrm{Ev}^{\eta},\ \mathrm{Ev}^{\mathrm{perf}},\ \mathrm{Ev}^{\mathrm{safe}},\ \mathrm{Ev}^{\mathrm{life}},\
      \mathrm{Ev}^{\mathrm{mfg}}
  \bigr\},
  \]
  分别锚定效率、性能输出、安全阈值、寿命/可靠性、制造装配可行性；
  \item（我方生成元）取顺序敏感的设计/控制生成元集合（示意）：
  \[
  \begin{aligned}
  \mathcal{G}_{\mathrm{self}}^{\mathrm{eng}}
  :=
  \{\, g_{\mathrm{shape}},\ g_{\mathrm{mat}},\ g_{\mathrm{cool}},\ g_{\mathrm{ctrl}},\ g_{\mathrm{proc}}\,\},
  \end{aligned}
  \]
  并令 $\mathcal{M}_{\mathrm{self}}^{\mathrm{eng}}:=\langle \mathcal{G}_{\mathrm{self}}^{\mathrm{eng}}\rangle$；
  \item（对抗生成元）取开放工况/扰动生成元集合（示意）：
  \[
  \begin{aligned}
  \mathcal{G}_{\mathrm{adv}}^{\mathrm{eng}}
  :=
  \{\, h_{\mathrm{env}},\ h_{\mathrm{load}},\ h_{\mathrm{wear}},\ h_{\mathrm{fault}},\ h_{\mathrm{unc}}\,\},
  \end{aligned}
  \]
  并令 $\mathcal{M}_{\mathrm{adv}}^{\mathrm{eng}}:=\langle \mathcal{G}_{\mathrm{adv}}^{\mathrm{eng}}\rangle$；
  \item（包络自由度）
  \[
  \Xi^{\mathrm{eng}}:=\Xi^{\mathrm{Geom}}\times \Xi^{\mathrm{Sim}}\times \Xi^{\mathrm{Proc}},
  \]
  其中 $\Xi^{\mathrm{Geom}}$ 为几何/结构细粒度自由度，$\Xi^{\mathrm{Sim}}$ 为多物理离散化与误差控制自由度，$\Xi^{\mathrm{Proc}}$ 为制造装配与公差口径自由度；
  \item（最小缺陷证书）同样取三分量：
  \[
  \JacGap_{\min}^{\mathrm{eng}}(L)
  :=
  \bigl(c_{\mathrm{chart}}(L),c_{\mathrm{lift}}(L),c_{\mathrm{thres}}(L)\bigr).
  \]
\end{enumerate}
\end{definition}

\begin{remark}[发动机：阈值与外生注入使其天然“开放对抗”]
发动机系统的安全阈值（$\mathrm{Ev}^{\mathrm{safe}}$）、故障态切换、磨损演化与外部工况注入，使其天然满足 Definition~\ref{def:open-system-adversarial-law-flow}
  的开放结构；
因此“设计”若要在真实工况下成立，必须以策略/算子词形式给出，并在 $\mathcal{E}_{\mathrm{prot}}^{\mathrm{eng}}$ 约束下通过同伦扭曲与包络修复维持可行域的闭合。
\end{remark}

\subsection{命题：两个实例落入同一模板（而非两套体系）}

\begin{proposition}[实例命题：Chip-\OPB{} 与 Engine-\OPB{} 共享同一 open adversarial \OPB{} 模板]
\label{prop:chip-engine-share-opb-template}
由 Definition~\ref{def:chip-opb-minimal-pack} 与 Definition~\ref{def:engine-opb-minimal-pack} 给出的两个最小基准包，
在 open adversarial \OPB{} 的语义下只是对 $\mathcal{E}_{\mathrm{prot}},\mathcal{M}_{\mathrm{self}},\mathcal{M}_{\mathrm{adv}},
  \Xi,\JacGap_{\min}$ 的不同实例化；
因此二者可被同一 Theorem~\ref{thm:opb-benchmark-template-equivalence} 的等价链统一：它们都等价于“在 $\OPB_{\mathrm{adv}}$
  上选切面/包络并在相对同伦类内做法流与同伦扭曲”。
\end{proposition}

\begin{certificate}[证书：实例到模板的代入映射]
\label{cert:chip-engine-share-opb-template}
对每个实例 $X\in\{\mathrm{chip},\mathrm{eng}\}$，由 Definition~\ref{def:open-adversarial-opb-object} 取代入：
\[
\mathcal{E}_{\mathrm{prot}}\leftarrow \mathcal{E}_{\mathrm{prot}}^{X},\quad
\mathcal{M}_{\mathrm{self}}\leftarrow \mathcal{M}_{\mathrm{self}}^{X},\quad
\mathcal{M}_{\mathrm{adv}}\leftarrow \mathcal{M}_{\mathrm{adv}}^{X},\quad
\Xi\leftarrow \Xi^{X},\quad
\JacGap_{\min}\leftarrow \JacGap_{\min}^{X}.
\]
代入后得到的 $\OPB_{\mathrm{adv}}^{X}$ 仍满足 open adversarial \OPB{} 的全部结构项与相对约束（受保护评价泛函族不变口径），从而可直接调用 Theorem~\ref{thm:opb-benchmark-template-equivalence}。
\end{certificate}

\begin{proof}[证明（谓词因果链）]
由 Certificate~\ref{cert:chip-engine-share-opb-template}，取谓词链：
\[
\begin{aligned}
A:&\ \text{给定实例化 }(\mathcal{E}_{\mathrm{prot}}^{X},\mathcal{M}_{\mathrm{self}}^{X},\mathcal{M}_{\mathrm{adv}}^{X},\Xi^{X}
  ,\JacGap_{\min}^{X})\\
&\Rightarrow \text{得到 }\OPB_{\mathrm{adv}}^{X}\text{ 的一组数据},\\
B:&\ \mathcal{E}_{\mathrm{prot}}^{X}\ \text{作为相对锚点}\Rightarrow \text{两实例仅改变扇区/切面/包络的取值，而不改变模板语义},\\
C:&\ A\wedge B\Rightarrow \text{二者均满足 Theorem~\ref{thm:opb-benchmark-template-equivalence} 的前提}.
\end{aligned}
\]
由 $C$ 得到命题结论。证毕。
\end{proof}

\begin{remark}[下一步：把某一实例落到“可跑基准”]
若希望把上述模板落到“可跑”的工程基准（含可复现数据与证书字段），通常需要进一步指定：
\begin{enumerate}
  \item $\mathcal{E}_{\mathrm{prot}}$ 的具体可计算接口（每个 $\mathrm{Ev}$ 的数据来源与单位/口径）；
  \item $\Xi^{\mathrm{Sim}}$ 与 $\Xi^{\mathrm{Proc}}$ 的最小字段（使 $\mathrm{Res}/\mathrm{Unc}/\mathrm{Cert}$ 可审计）；
  \item $\JacGap_{\min}$ 三分量的具体实现（如何由跨评估器不一致、水平提升失配、阈值切换不一致产生可计算证书）。
\end{enumerate}
在该信息明确后，即可按 Theorem~\ref{thm:ig-invgen-ohu} 的 “failure-free + defect decrease” 口径，把实例写成可直接调用的基准段落与检查清单。
\end{remark}

\clearpage
%%%%%%%%%%%%%%%%%%%%%%%%%%%%%%%%%%%%%%%%%%%%%%%%%%%%%%%%%%%%
\section{纯数卷口径：\RTSC{}-工业几何体 Benchmark（开放博弈系统对抗主丛结构实例）}
\label{sec:puremath-rtsc-industrial-geometry-benchmark}
%%%%%%%%%%%%%%%%%%%%%%%%%%%%%%%%%%%%%%%%%%%%%%%%%%%%%%%%%%%%

\subsection{形式逻辑表达约定}

\begin{remark}[\RTSC{}-工业几何体 Benchmark：可逆向生成实例化目标]
本章在 Pure 卷符号体系（$\GZLS,\GZLOC,\GZTLC,\GZLCurv,\mathrm{Loc},\Sigma,\JacGap$ 等）下，
给出一个“\RTSC{}-工业几何体 Benchmark”的\textbf{可逆向生成}形式化实例。
工业几何体可涵盖（但不限于）特高压变电节点的 \RTSC{} 法拉第笼、\RTSC{} 旁路电阻、以及量子计算抗退相干装置；
它们在 Benchmark 中统一为同一类 law-objects 的不同受保护评价约束与不同切面/包络选择。
\end{remark}

\subsection{定义：在 $(\mathrm{Loc},\Sigma)$ 下的基底 $M$（受保护评价与分层标签）}

\begin{definition}[基底 $M$：由受保护评价与阈值分层诱导的“可判定状态域”]
\label{def:rtsc-ig-benchmark-base}
固定 Pure 卷 law-space
\[
\GZLS=(\mathfrak{L};J,\mathrm{Loc},\Sigma),
\]
并固定 $\GZLOC,\GZTLC$ 与 $\GZLCurv$（Definition~\ref{def:law-space-gzls}）。
在 $\mathrm{Loc}$ 上取一组受保护评价泛函族
\[
\mathcal{E}_{\mathrm{prot}}
=
\{\mathrm{Ev}_a\}_{a=1}^{q},
\qquad
\mathrm{Ev}_a:\mathrm{Loc}\to\mathbb{R},
\]
并要求该族至少包含下列语义锚点（以抽象评价泛函记号表示）：
\[
\mathrm{Ev}_{\mathrm{RTSC}},\ \mathrm{Ev}_{\mathrm{Cert}},\
\mathrm{Ev}_{\mathrm{Cage}},\ \mathrm{Ev}_{\mathrm{Safety}},\
\mathrm{Ev}_{\mathrm{Bypass}},\
\mathrm{Ev}_{\mathrm{QDec}}.
\]
记向量化评价映射为
\[
\mathrm{Ev}:=(\mathrm{Ev}_1,\dots,\mathrm{Ev}_q):\mathrm{Loc}\to\mathbb{R}^q.
\]
由阈值集合 $\Sigma$ 诱导一个分层/工况标签映射（regime map）
\[
\rho_{\Sigma}:\mathrm{Loc}\to \mathcal{R}_{\Sigma},
\]
其中 $\mathcal{R}_{\Sigma}$ 为有限或可数的分层标签集合（例如正常态/故障态/保护切换态/测量口径切换态等）。

定义 \RTSC{}-工业几何体 Benchmark 的基底为集合
\[
\begin{aligned}
M
:=
\Bigl\{
(\mathbf{y},r)\in \mathbb{R}^q\times \mathcal{R}_{\Sigma}\ \Big|\
&\exists\,L\in \mathrm{Loc}\subset\GZLS,\\
&\mathrm{Ev}(L)=\mathbf{y},\ \rho_{\Sigma}(L)=r
\Bigr\}.
\end{aligned}
\]
\end{definition}

\begin{remark}[解释：基底 $M$ 只承载“可判定语义”]
Definition~\ref{def:rtsc-ig-benchmark-base} 的 $M$ 不是“全状态空间”，而是由受保护评价向量 $\mathbf{y}$ 与阈值分层标签 $r$ 组成的可判定基底；
所有表示冗余、实现冗余与包络细节将被并入纤维方向，以避免在基底层出现不可审计的“路线自洽偏航”。
\end{remark}

\subsection{构造：主丛 $\pi:P\to M$、切面与包络（$E\simeq U\times\Xi$）}

\begin{definition}[表示冗余结构：$G_{\mathrm{rep}}^{\mathrm{RTSC\text{-}IG}}$]
\label{def:rtsc-ig-grep}
取一个表示冗余结构对象（群或广义群胚）
\[
G_{\mathrm{rep}}^{\mathrm{RTSC\text{-}IG}},
\]
其语义至少覆盖：工业几何体的等价参数化（CAD/几何表述等价）、多物理仿真的离散化与求解器等价（网格、容差、耦合接口等）、制造/装配路径的等价表述（工艺序列与公差参数化），以及测量/证书口径的允许变换（但不得改变受保护评价的口径锚点，见
  Definition~\ref{def:protected-evaluation-functionals}）。
\end{definition}

\begin{definition}[开放对抗主丛：$P=\{(L,\chi)\}$ 与投影 $\pi:P\to M$]
\label{def:rtsc-ig-principal-bundle}
定义总空间为带有法对象代表与实现数据的集合（schematic）：
\[
\begin{aligned}
P
:=
\Bigl\{
(L,\chi)\ \Big|\
&L\in \mathrm{Loc}\subset\GZLS,\ (\mathrm{Ev}(L),\rho_{\Sigma}(L))\in M,\\
&\chi\in \mathcal{X}(L)
\Bigr\},
\end{aligned}
\]
其中 $\chi$ 表示实现/表示/求解/工艺等附加数据，并由 $G_{\mathrm{rep}}^{\mathrm{RTSC\text{-}IG}}$ 作用。
定义投影
\[
\pi:P\to M,\qquad
\pi(L,\chi)=(\mathrm{Ev}(L),\rho_{\Sigma}(L)).
\]
定义右作用
\[
(L,\chi)\cdot g=(L,\chi\cdot g),\qquad g\in G_{\mathrm{rep}}^{\mathrm{RTSC\text{-}IG}},
\]
并要求其保持基底不变：
\[
\pi\bigl((L,\chi)\cdot g\bigr)=\pi(L,\chi).
\]
\end{definition}

\begin{definition}[切面与包络：路线/管线选择与细粒度纤维自由度]
\label{def:rtsc-ig-section-envelope}
对任一可操作开集 $U\subset M$，称映射
\[
s:U\to P,\qquad \pi\circ s=\mathrm{Id}_U
\]
为一条切面（section），其语义是：对每个 $(\mathbf{y},r)\in U$ 选取一个具体法对象代表 $L$ 与实现口径 $\chi$（路线/管线/协议的选择）。

给定切面 $s$，取一个包络邻域
\[
E\subset P,\qquad s(U)\subset E,
\]
并假设局部存在包络坐标化同构
\[
\varepsilon:\ U\times\Xi \xrightarrow{\ \cong\ } E.
\]
其中包络自由度取最小分解
\[
\Xi
=
\Xi^{\mathrm{Geom}}
\times
\Xi^{\mathrm{Sim}}
\times
\Xi^{\mathrm{Proc}}
\times
\Xi^{\mathrm{Ctrl}}.
\]
其语义分别为：工业几何体细粒度自由度（$\Xi^{\mathrm{Geom}}$）、仿真表示与误差控制自由度（$\Xi^{\mathrm{Sim}}$）、制造/装配路径与公差自由度（$\Xi^{\mathrm{Proc}}$），
  以及运行期控制协议细粒度自由度（$\Xi^{\mathrm{Ctrl}}$；其顺序结构一般非交换）。
\end{definition}

\begin{remark}[主丛语义：对抗/策略/设计统一为同一载体]
在上述构造下，“开放博弈系统对抗主丛结构”的语义可读作：
对抗与策略作用于基底 $M$ 上的可判定语义（受保护评价与阈值分层），
而一切表示冗余与可调细节进入纤维（$G_{\mathrm{rep}}^{\mathrm{RTSC\text{-}IG}}$ 与 $\Xi$）；
因此“工程设计生成”与“策略生成”可在同一 $P$ 上被解释为“切面选择 + 包络内动作（含同伦修复）”。
\end{remark}

\subsection{定义：最小 $\JacGap(L)$ 三分量证书与缺陷势能 $\Phi$}

\begin{definition}[最小三分量证书：$c_{\mathrm{chart}},c_{\mathrm{lift}},c_{\mathrm{thres}}$]
\label{def:rtsc-ig-jacgap-minimal}
对 \RTSC{}-工业几何体 Benchmark 的法层算子系统 $L$（Pure 卷允许 $L\in(\GZLS)_{\mathrm{homo}}$ 或 $L\in(\GZLS)_{\mathrm{vf\text{-}op}}$），
定义最小三分量 Jacobiator-gap 证书为
\[
\JacGap(L)=\bigl(c_{\mathrm{chart}}(L),c_{\mathrm{lift}}(L),c_{\mathrm{thres}}(L)\bigr),
\]
并作如下语义约定（与 Definition~\ref{def:opb-jacgap-minimal} 一致）：
\begin{enumerate}
  \item $c_{\mathrm{chart}}(L)$：跨图册/多口径一致性缺陷（不同建模层、不同仿真--验证口径、不同 charts 在重叠域上的对齐失败）；
  \item $c_{\mathrm{lift}}(L)$：水平提升/同伦修复一致性缺陷（在保持受保护评价口径的前提下，同伦修复与 $\GZLOC/\GZTLC$ 的运输是否相容）；
  \item $c_{\mathrm{thres}}(L)$：阈值分层/切换一致性缺陷（跨 $\Sigma$ 的切换图式是否保持必要同伦相容性与 $\mathrm{Loc}$ 闭合）。
\end{enumerate}
\end{definition}

\begin{definition}[缺陷势能：三分量证书的单调聚合]
\label{def:rtsc-ig-defect-potential}
取权重 $w_{\mathrm{chart}},w_{\mathrm{lift}},w_{\mathrm{thres}}>0$。
定义缺陷势能（defect potential）为
\[
\Phi(L)
=
w_{\mathrm{chart}}\|c_{\mathrm{chart}}(L)\|
w_{\mathrm{lift}}\|c_{\mathrm{lift}}(L)\|
w_{\mathrm{thres}}\|c_{\mathrm{thres}}(L)\|,
\]
其中 $\|\cdot\|$ 为与证书载体空间相容的范数（或谱聚合）。
Pure 卷语义要求：若 $\JacGap(L)=0$，则 $L$ 可在保持 observable low-order behaviour（即 $\mathcal{E}_{\mathrm{prot}}$ 的评价口径锚点）下被
  promoted 到严格对象；
并且可允许有限同伦动作严格下降某个 defect potential（此处即 $\Phi$）。
\end{definition}

\subsection{命题：开放对抗下的相对同伦等价（设计生成 $\simeq$ 策略生成）}

\begin{proposition}[相对同伦等价命题：defect potential 下降链刻画“设计生成=策略生成”]
\label{prop:rtsc-ig-design-strategy-equivalence}
在开放系统对抗法流（Definition~\ref{def:open-system-adversarial-law-flow}）下，
若对任意给定对抗输入 $d(\cdot)$，存在：
\begin{enumerate}
  \item 一条切面 $s:U\to P$ 及其包络 $E\simeq U\times\Xi$（Definition~\ref{def:rtsc-ig-section-envelope}）；
  \item 一列有限同伦动作（Pure 卷 homotopy moves，schematic）与包络有限步调整，使得在保持受保护评价口径不变的前提下，缺陷势能严格下降：
  \[
  \begin{aligned}
  \mathrm{Ev}\bigl(L^{(k)}(t)\bigr) &= \mathrm{Ev}\bigl(L^{(0)}(t)\bigr)\quad(\forall k,t),\\
  \Phi\bigl(L^{(k+1)}(t)\bigr) &< \Phi\bigl(L^{(k)}(t)\bigr).
  \end{aligned}
  \]
  并且第二行不等式在一段非零测度的 $t$-区间成立；
  其中 $L^{(k)}(t)\in\mathrm{Loc}$ 为同伦修复与包络动作后的法对象轨线；
\end{enumerate}
则在 $\mathcal{E}_{\mathrm{prot}}$-相对同伦意义下，
“生成工程设计方案（切面+包络+同伦动作）”与“生成开放对抗策略（非交换算子词）”等价。
\end{proposition}

\begin{certificate}[证书：相对约束 + 可下降缺陷势能 $\Rightarrow$ 等价链可调用]
\label{cert:rtsc-ig-design-strategy-equivalence}
证书化断言由三条可核验前提组成：
\begin{enumerate}
  \item（相对锚点）受保护评价泛函族 $\mathcal{E}_{\mathrm{prot}}$ 固定评价口径，排除“换口径即胜利”（Theorem~\ref{thm:protected-evaluations-necessary}）；

  \item（包络显式化）存在 $\Xi^{\mathrm{Geom}}\times\Xi^{\mathrm{Sim}}\times\Xi^{\mathrm{Proc}}\times\Xi^{\mathrm{Ctrl}}$ 使修复动作在
    $\mathrm{Loc}$ 内可实现（Definition~\ref{def:rtsc-ig-section-envelope}）；
  \item（可下降性）存在有限同伦动作使 defect potential 严格下降（与 Theorem~\ref{thm:ig-invgen-ohu} 的 “finite homotopy moves strictly
    decrease defect potential” 口径一致）。
\end{enumerate}
上述三点成立时，可直接调用 Theorem~\ref{thm:opb-benchmark-template-equivalence} 与 Theorem~\ref{thm:noncommutative-design-strategy-equivalence} 的等价链结论。
\end{certificate}

\begin{proof}[证明（谓词因果链）]
由 Certificate~\ref{cert:rtsc-ig-design-strategy-equivalence}，取谓词链：
\[
\begin{aligned}
A:&\ \mathcal{E}_{\mathrm{prot}}\ \text{固定}\Rightarrow \text{相对同伦语义锚点成立},\\
B:&\ \Xi\ \text{显式化}\Rightarrow \text{工程设计可表为切面+包络上的可执行动作序列},\\
C:&\ \Phi\ \text{可严格下降}\Rightarrow \text{同伦修复/包络动作给出可证书化回归链}.
\end{aligned}
\]
由 $A\wedge B$ 得到“设计生成具有策略（算子词）语义”，并由 $A\wedge C$ 得到“相对同伦类内的 defect potential 下降链可判定”。
结合 Theorem~\ref{thm:opb-benchmark-template-equivalence} 的等价链，即得所述相对等价结论。证毕。
\end{proof}

\subsection{备注：与回合制封闭对抗的结构差异（OpenRA 型而非 AlphaGo 型）}

\begin{remark}[非自治法流 + 阈值切换 + 曲率尖峰 $\Rightarrow$ 必须在线同伦回归]
\RTSC{}-工业几何体 Benchmark 的关键差异点在于：法流是开放且非自治的（有 $d(\cdot)$ 外生注入），并受阈值集合 $\Sigma$ 管理；
外生冲击可诱发法曲率尖峰并放大 $\JacGap$ 的变化，从而使缺陷势能 $\Phi$ 突增（参见 Theorem~\ref{thm:law-curvature-spike-jacgap}）。
因此系统必须在线执行“先选测地线方向，再做相对同伦扭曲”的回归机制（Corollary~\ref{cor:relative-geodesic-twist}），
其时间/信息结构与实时开放对抗更贴近（Corollary~\ref{cor:openra-vs-alphago}）。
\end{remark}

\clearpage
%%%%%%%%%%%%%%%%%%%%%%%%%%%%%%%%%%%%%%%%%%%%%%%%%%%%%%%%%%%%
\section{纯数卷口径：\textsf{RTSC-approx} 工业低阻法拉第笼与低阻旁路熔断（实时对抗控制系统）}
\label{sec:puremath-rtsc-approx-lowres-cage-bypass}
%%%%%%%%%%%%%%%%%%%%%%%%%%%%%%%%%%%%%%%%%%%%%%%%%%%%%%%%%%%%

\subsection{形式逻辑表达约定}

\begin{remark}[定位：不预设“已实现 RTSC”，而把“低阻逼近+证书化+阈值切换+回归”写成可计算对象]
本章把“工业低阻法拉第笼 + 旁路电压低阻/电流熔断（触发式限流/隔离/恢复）”在 Pure 卷符号体系下形式化为一个
``\textsf{RTSC-approx} 实时开放对抗控制系统''：
不在物理层面硬编码 ``$\rho(300K)=0$'' 的强断言，而在法空间 $\GZLS$ 内把
``低阻逼近 + 证书口径 + 阈值分层切换 + 曲率尖峰下回归'' 组织为可证书、可比较、可迁移的主丛化对象。
\end{remark}

\subsection{定义：\textsf{RTSC-approx} 的工业低阻目标（受保护评价阈值）}

\begin{definition}[\textsf{RTSC-approx}：室温低阻评价的阈值约束]
\label{def:rtsc-approx-goal}
在 $\mathrm{Loc}\subset\GZLS$ 上，取一条“室温低阻”低阶可观测评价泛函
\[
\mathrm{Ev}_{\rho,300}:\ \mathrm{Loc}\to \mathbb{R}_{\ge 0},
\]
并给定工程阈值 $\varepsilon_{\rho}>0$。
称法对象 $L\in\mathrm{Loc}$ 满足 \textsf{RTSC-approx} 的低阻目标，当且仅当
\[
\mathrm{Ev}_{\rho,300}(L)\le \varepsilon_{\rho}.
\]
该目标不等同于 ``$\rho(300K)=0$'' 的物理断言；其工程语义由受保护评价泛函族 $\mathcal{E}_{\mathrm{prot}}$ 的口径锚点与证书字段共同固定（见 Definition~\ref{def:rtsc-approx-eprot-base}）。
\end{definition}

\subsection{定义：合成 law-object（法拉第笼 + 旁路熔断 + 网络耦合 + 测量证书）}

\begin{definition}[合成系统：$L=L_{\mathrm{cage}}\sqcup L_{\mathrm{bypass}}\sqcup L_{\mathrm{grid}}\sqcup L_{\mathrm{meas}}$]
\label{def:rtsc-approx-composite-law-object}
固定 Pure 卷 law-space $\GZLS=(\mathfrak{L};J,\mathrm{Loc},\Sigma)$ 及 $\GZLOC,\GZTLC,\GZLCurv$（Definition~\ref{def:law-space-gzls}）。
称“工业低阻法拉第笼 + 低阻旁路熔断/限流/隔离”的合成系统对应一个 law-object
\[
L
:=
L_{\mathrm{cage}}
\sqcup
L_{\mathrm{bypass}}
\sqcup
L_{\mathrm{grid}}
\sqcup
L_{\mathrm{meas}},
\qquad
L\in\mathrm{Loc}\subset\GZLS,
\]
其中：
\begin{enumerate}
  \item $L_{\mathrm{cage}}$ 表示工业低阻法拉第笼（几何--边界--接地拓扑--材料状态在 law-space 的像）；
  \item $L_{\mathrm{bypass}}$ 表示低阻旁路与触发式熔断/限流/隔离装置（含阈值触发逻辑与状态机）；
  \item $L_{\mathrm{grid}}$ 表示 UHV 节点外部网络耦合与运行工况（外生输入主要来源）；
  \item $L_{\mathrm{meas}}$ 表示测量与证书口径（用于 $\mathrm{Res}/\mathrm{Unc}/\mathrm{Cert}$ 与 protected observables 的固定）。
\end{enumerate}
\end{definition}

\subsection{定义：$\mathcal{E}_{\mathrm{prot}}$ 与基底 $M$（可判定语义层）}

\begin{definition}[受保护评价泛函族与基底 $M$：低阻--安全--阈值--证书的语义压缩]
\label{def:rtsc-approx-eprot-base}
在 $\mathrm{Loc}\subset\GZLS$ 上，取有限族受保护评价泛函
\[
\mathcal{E}_{\mathrm{prot}}=\{\mathrm{Ev}_a\}_{a=1}^{q},
\qquad
\mathrm{Ev}_a:\mathrm{Loc}\to\mathbb{R},
\]
并要求其最小闭包至少包含四类语义锚点（写作抽象评价泛函，不绑定具体物理公式）：
\begin{enumerate}
  \item（低阻逼近与载流能力）
  \[
  \mathrm{Ev}_{\rho,300}(L),\qquad \mathrm{Ev}_{I}(L);
  \]
  \item（屏蔽/泄漏与接地安全）
  \[
  \mathrm{Ev}_{\mathrm{shield}}(L),\qquad
  \mathrm{Ev}_{\mathrm{leak}}(L),\qquad
  \mathrm{Ev}_{\mathrm{gnd}}(L);
  \]
  \item（旁路熔断/限流的阈值语义）
  \[
  \mathrm{Ev}_{\mathrm{trip}}(L),\qquad
  \mathrm{Ev}_{\mathrm{arc}}(L),\qquad
  \mathrm{Ev}_{\mathrm{recover}}(L);
  \]
  \item（证书/残差/不确定度口径）
  \[
  \mathrm{Ev}_{\mathrm{Cert}}(L),\qquad
  \mathrm{Ev}_{\mathrm{Res}}(L),\qquad
  \mathrm{Ev}_{\mathrm{Unc}}(L).
  \]
\end{enumerate}
记向量化评价映射为
\[
\mathrm{Ev}:=(\mathrm{Ev}_1,\dots,\mathrm{Ev}_q):\mathrm{Loc}\to\mathbb{R}^q.
\]

由阈值集合 $\Sigma$ 诱导分层/工况标签映射
\[
\rho_{\Sigma}:\mathrm{Loc}\to \mathcal{R}_{\Sigma},
\]
其中 $\mathcal{R}_{\Sigma}$ 表示正常态/故障态/触发态/恢复态等分层标签集合。

定义可判定基底为
\[
\begin{aligned}
M
:=
\Bigl\{
(\mathbf{y},r)\in \mathbb{R}^q\times \mathcal{R}_{\Sigma}\ \Big|\
&\exists\,L\in \mathrm{Loc},\\
&\mathrm{Ev}(L)=\mathbf{y},\ \rho_{\Sigma}(L)=r
\Bigr\}.
\end{aligned}
\]
并将 “\textsf{RTSC-approx} 目标”理解为 $M$ 上的阈值约束（由 $\mathrm{Ev}_{\rho,300}$ 分量给出），而非直接在对象层宣称“已实现 RTSC”。
\end{definition}

\begin{remark}[相对同伦约束：禁止通过改变评价口径获得伪下降]
任意允许的同伦扭曲/包络延拓必须保持 $\mathcal{E}_{\mathrm{prot}}$ 的评价语义（observable low-order behaviour）不变：
否则可通过改变 $\mathrm{Ev}_{\rho,300}$ 等评价口径制造“低阻伪像”，与 Theorem~\ref{thm:protected-evaluations-necessary} 相矛盾。
\end{remark}

\subsection{构造：主丛、切面与包络（显式化实时对抗控制自由度）}

\begin{definition}[主丛 $\pi:P\to M$：对抗/策略/设计的同一载体]
\label{def:rtsc-approx-principal-bundle}
取表示冗余结构对象（群或广义群胚）
\[
G_{\mathrm{rep}}^{\mathrm{IG\text{-}RTSC}},
\]
吸收 CAD 参数化、仿真离散化、求解器、工艺表述、控制协议参数化等实现冗余。
定义总空间
\[
P
:=
\Bigl\{
(L,\chi)\ \Big|\
L\in\mathrm{Loc},\ (\mathrm{Ev}(L),\rho_{\Sigma}(L))\in M,\ \chi\in\mathcal{X}(L)
\Bigr\},
\]
并定义投影
\[
\pi:P\to M,\qquad \pi(L,\chi)=(\mathrm{Ev}(L),\rho_{\Sigma}(L)).
\]
定义右作用 $(L,\chi)\cdot g=(L,\chi\cdot g)$（$g\in G_{\mathrm{rep}}^{\mathrm{IG\text{-}RTSC}}$），并要求保持基底不变：
\[
\pi\bigl((L,\chi)\cdot g\bigr)=\pi(L,\chi).
\]
\end{definition}

\begin{definition}[切面与包络：$E\simeq U\times\Xi$，其中 $\Xi^{\mathrm{Ctrl}}$ 参数化非交换算子词]
\label{def:rtsc-approx-envelope}
对任一开集 $U\subset M$，切面为
\[
s:U\to P,\qquad \pi\circ s=\mathrm{Id}_U,
\]
表示固定的一条工程路线/管线/协议/测量证书口径。
取包络邻域 $E\subset P$ 满足 $s(U)\subset E$，并假设局部存在坐标化同构
\[
\varepsilon:\ U\times\Xi \xrightarrow{\ \cong\ } E.
\]
对实时对抗控制系统，包络自由度至少取分解
\[
\Xi
=
\Xi^{\mathrm{Geom}}
\times
\Xi^{\mathrm{Sim}}
\times
\Xi^{\mathrm{Proc}}
\times
\Xi^{\mathrm{Ctrl}}.
\]
其中 $\Xi^{\mathrm{Ctrl}}$ 可视为我方控制算子词空间的一种参数化：存在算子幺半群 $\mathcal{M}_{\mathrm{self}}$ 使
\[
\Xi^{\mathrm{Ctrl}}\ \sim\ \mathcal{M}_{\mathrm{self}}^\ast,
\qquad
u=u_1u_2\cdots u_n,\ \ u_i\in\mathcal{M}_{\mathrm{self}},
\]
并允许一般非交换性（Definition~\ref{def:noncommutative-operator-systems}）。
对抗侧同理由 $\mathcal{M}_{\mathrm{adv}}^\ast$ 生成对抗输入词 $d$。
\end{definition}

\subsection{定义：最小三分量 $\JacGap(L)$ 与缺陷势能 $\Phi$}

\begin{definition}[最小三分量证书：chart--lift--threshold]
\label{def:rtsc-approx-jacgap}
对合成 law-object $L$（Definition~\ref{def:rtsc-approx-composite-law-object}），取最小三分量 Jacobiator-gap 证书
\[
\JacGap(L)=\bigl(c_{\mathrm{chart}}(L),c_{\mathrm{lift}}(L),c_{\mathrm{thres}}(L)\bigr),
\]
其语义分别为：多口径/多模型层一致性缺陷（$c_{\mathrm{chart}}$）、水平提升与同伦修复相容性缺陷（$c_{\mathrm{lift}}$）、以及跨阈值分层切换的一致性缺陷（$c_{\mathrm{thres}}$）（参见
  Definition~\ref{def:opb-jacgap-minimal} 与 Definition~\ref{def:rtsc-ig-jacgap-minimal}）。
\end{definition}

\begin{definition}[缺陷势能：$\Phi=w_{\mathrm{chart}}\|c_{\mathrm{chart}}\|+w_{\mathrm{lift}}\|c_{\mathrm{lift}}
  \|+w_{\mathrm{thres}}\|c_{\mathrm{thres}}\|$]
\label{def:rtsc-approx-defect-potential}
取权重 $w_{\mathrm{chart}},w_{\mathrm{lift}},w_{\mathrm{thres}}>0$，定义 defect potential 为
\[
\Phi(L)
:=
w_{\mathrm{chart}}\|c_{\mathrm{chart}}(L)\|
w_{\mathrm{lift}}\|c_{\mathrm{lift}}(L)\|
w_{\mathrm{thres}}\|c_{\mathrm{thres}}(L)\|.
\]
直观上，$\Phi$ 小意味着“低阻逼近 + 证书口径 + 阈值切换逻辑 + 多模型一致性”更接近可严格化/可复用的稳定区域；
外生冲击下 $\Phi$ 的尖峰往往与 $\GZLCurv$ 尖峰共现（Theorem~\ref{thm:law-curvature-spike-jacgap}）。
\end{definition}

\subsection{命题：开放对抗下的相对同伦等价（设计生成 $\simeq$ 策略生成）}

\begin{proposition}[\textsf{RTSC-approx} 实时对抗：设计生成与策略生成在相对同伦意义下等价]
\label{prop:rtsc-approx-design-strategy-equivalence}
在开放系统对抗法流（Definition~\ref{def:open-system-adversarial-law-flow}）下，
若对任意允许的对抗输入词 $d(\cdot)$，存在：
\begin{enumerate}
  \item 某切面 $s:U\to P$ 及其包络 $E\simeq U\times\Xi$（Definition~\ref{def:rtsc-approx-envelope}）；
  \item 一列有限同伦动作（Pure 卷 homotopy moves；schematic，可理解为对高阶同伦数据的有限修正）与包络有限步移动，使得：
  \[
  \mathrm{Ev}\bigl(L^{(k)}(t)\bigr)=\mathrm{Ev}\bigl(L^{(0)}(t)\bigr)\quad(\forall k,t),
  \]
  且存在一段非零测度的时间区间使
  \[
  \Phi\bigl(L^{(k+1)}(t)\bigr)<\Phi\bigl(L^{(k)}(t)\bigr);
  \]
\end{enumerate}
则在 $\mathcal{E}_{\mathrm{prot}}$-相对同伦意义下，
“生成工业低阻方案（法拉第笼+旁路熔断）”与“生成实时开放对抗控制策略（非交换算子词）”等价。
\end{proposition}

\begin{certificate}[证书：相对锚点 + 包络显式化 + 缺陷势能可下降]
\label{cert:rtsc-approx-design-strategy-equivalence}
证书化前提为：
\begin{enumerate}
  \item（相对锚点）$\mathcal{E}_{\mathrm{prot}}$ 固定评价口径，排除伪下降（Theorem~\ref{thm:protected-evaluations-necessary}）；
  \item（包络显式化）$\Xi^{\mathrm{Geom}}\times\Xi^{\mathrm{Sim}}\times\Xi^{\mathrm{Proc}}\times\Xi^{\mathrm{Ctrl}}$
    提供足够自由度承载在线控制与修复（Definition~\ref{def:rtsc-approx-envelope}）；
  \item（可下降性）存在有限同伦动作使 $\Phi$ 严格下降（与 Theorem~\ref{thm:ig-invgen-ohu} 的 defect decrease 口径一致）。
\end{enumerate}
\end{certificate}

\begin{proof}[证明（谓词因果链）]
由 Certificate~\ref{cert:rtsc-approx-design-strategy-equivalence}，取谓词链：
\[
\begin{aligned}
A:&\ \mathcal{E}_{\mathrm{prot}}\ \text{固定}\Rightarrow \text{相对同伦语义锚点成立},\\
B:&\ \Xi^{\mathrm{Ctrl}}\sim\mathcal{M}_{\mathrm{self}}^\ast\Rightarrow \text{工程方案自然对应非交换策略词},\\
C:&\ \Phi\ \text{可严格下降}\Rightarrow \text{在同一锚点下存在可证书化回归/修复链}.
\end{aligned}
\]
由 $A\wedge B\wedge C$ 得到“设计生成”与“策略生成”在相对同伦意义下可互释的结论。证毕。
\end{proof}

\subsection{备注：OpenRA 型即时开放对抗，而非 AlphaGo 型回合制封闭对抗}

\begin{remark}[非自治法流 + 阈值切换 + 曲率尖峰 $\Rightarrow$ 必须在线回归]
该系统的法流为非自治开放法流：
\[
\dot L(t)=\mathcal{H}\bigl(L(t),u(t),d(t);\GZLOC,\GZTLC\bigr),
\]
其中 $d(t)$ 持续外生注入并可触发阈值切换（熔断/隔离/恢复等）。
跨阈值的一致性由 $c_{\mathrm{thres}}$ 捕捉，而强扰动可诱发 $\GZLCurv$ 尖峰并使 $\Phi$ 突增；
因此必须在线执行“先选测地线方向，再做相对同伦扭曲”的回归机制（Corollary~\ref{cor:relative-geodesic-twist}）。
这正对应实时开放对抗（OpenRA 型）的时间/信息结构，而非回合制封闭对弈（AlphaGo 型）（Corollary~\ref{cor:openra-vs-alphago}）。
\end{remark}

\subsection{备注：一句话定义（作为实时对抗控制系统的可调用口径）}

\begin{remark}[一句话定义：\textsf{RTSC-approx} 低阻工程系统的主丛控制问题]
\[
\boxed{
\begin{aligned}
&\text{\textsf{RTSC-approx} 工业低阻法拉第笼 + 低阻旁路熔断}\\
&\equiv\ \text{在 }\GZLS\text{ 上相对 }\mathcal{E}_{\mathrm{prot}}\text{ 的开放对抗主丛控制问题：}\\
&\hspace{2em}\text{通过包络 }\Xi^{\mathrm{Ctrl}}\text{ 的非交换算子词在线选择，使 }\Phi\downarrow\text{ 且跨 }\Sigma\text{ 可回归。}
\end{aligned}
}
\]
\end{remark}

\clearpage
%%%%%%%%%%%%%%%%%%%%%%%%%%%%%%%%%%%%%%%%%%%%%%%%%%%%%%%%%%%%
\section{纯数卷口径：边缘算子（Edge Operators）与有限元包络（让阈值过渡更平滑）}
\label{sec:puremath-edge-operators-fem-envelope}
%%%%%%%%%%%%%%%%%%%%%%%%%%%%%%%%%%%%%%%%%%%%%%%%%%%%%%%%%%%%

\subsection{形式逻辑表达约定}

\begin{remark}[定位：边缘算子不是补丁，而是包络与水平法流机制天然预留的细粒度接口]
在我们已建立的 \RTSC{}-工业几何体开放对抗主丛结构中（Section~\ref{sec:puremath-rtsc-industrial-geometry-benchmark} 与 Section~\ref{sec:puremath-rtsc-approx-lowres-cage-bypass}），
引入“边缘算子（edge operators）”不是额外补丁，而是由包络 $E\simeq U\times\Xi$（Definition~\ref{def:rtsc-approx-envelope}）与 $\GZLOC/\GZTLC$
  的水平法流机制天然预留的细粒度自由度：
它们把原本会表现为“阈值跃迁/曲率尖峰/非交换词导致的突变”的变化，分解为可控的局部微步，
从而使过渡在 law-metric $J$ 与 defect potential $\Phi$ 意义下更平滑（更易严格下降且保持 $\mathcal{E}_{\mathrm{prot}}$）。
\end{remark}

\subsection{定义：有限元/格点离散作为包络载体（网格复形与边集合）}

\begin{definition}[有限元网格复形：$\mathcal{T}_h=(V_h,E_h,K_h)$]
\label{def:fem-mesh-complex}
令工业几何域（法拉第笼/旁路结构/节点区域）对应对象层几何为 $\Omega$。
取一套有限元离散（可自适应）为复形
\[
\mathcal{T}_h=(V_h,E_h,K_h),
\]
其中 $V_h$ 为节点集合（vertices），$E_h$ 为边集合（edges），$K_h$ 为单元集合（elements）。
用边--节点关联（incidence）定义离散微分结构：例如离散外微分 $d_h$ 或边--节点关联矩阵 $B_h$，
从而可定义离散 Dirichlet 能量与离散 Laplacian。
\end{definition}

\begin{remark}[Pure 卷解释：$\mathcal{T}_h$ 属于 $\Xi^{\mathrm{Geom}}/\Xi^{\mathrm{Sim}}$ 的可操作坐标化]
从 Pure 卷视角看，$\mathcal{T}_h$ 及其离散微分结构属于包络中的仿真表示/几何坐标化数据：
它们可被纳入 $\Xi^{\mathrm{Sim}}$（离散化、求解器、容差等）与 $\Xi^{\mathrm{Geom}}$（几何参数化与局部自由度），
其等价离散化/等价参数化由 $G_{\mathrm{rep}}^{\mathrm{Sim}}$ 与 $G_{\mathrm{rep}}^{\mathrm{Geom}}$ 吸收（参见 Definition~\ref{def:rtsc-approx-principal-bundle} 的表示冗余口径）。
\end{remark}

\subsection{定义：四类边缘算子与非交换算子字母表 $\mathcal{M}_{\mathrm{edge}}$}

\begin{definition}[边缘算子生成元：四类 $\mathsf{E}^{\bullet}_e$（$e\in E_h$）]
\label{def:edge-operator-generators}
在复形 $\mathcal{T}_h=(V_h,E_h,K_h)$（Definition~\ref{def:fem-mesh-complex}）上，
定义四类边缘算子生成元集合为
\[
\mathcal{G}_{\mathrm{edge}}
:=
\Bigl\{
\mathsf{E}^{\mathrm{mat}}_{e},\
\mathsf{E}^{\mathrm{env}}_{e},\
\mathsf{E}^{\mathrm{prec}}_{e},\
\mathsf{E}^{\mathrm{comp}}_{e}
\ \Big|\ e\in E_h
\Bigr\}.
\]
其语义分别为：
\begin{enumerate}
  \item（材料格点边缘算子）$\mathsf{E}^{\mathrm{mat}}_{e}$：对材料局部格点/界面/缺陷等价数据做边方向局部更新（“格点可变调节”的可执行入口）；
  \item（环境控制边缘算子）$\mathsf{E}^{\mathrm{env}}_{e}$：对环境场/边界条件的离散自由度做边方向局部更新（局部屏蔽/接地/边界加载等）；
  \item（控制精度边缘算子）$\mathsf{E}^{\mathrm{prec}}_{e}$：对控制/测量/阈值触发的离散精度参数做局部更新（量化步长、滞回带、触发抖动抑制参数等）；
  \item（算力支持边缘算子）$\mathsf{E}^{\mathrm{comp}}_{e}$：对 $\Xi^{\mathrm{Sim}}$ 的局部算力分配与自适应参数做局部更新（例如局部 $h/p$ 加密与容差调整），
    用于抑制由数值口径诱导的突变。
\end{enumerate}
定义由其生成的（一般非交换）幺半群/算子词空间为
\[
\mathcal{M}_{\mathrm{edge}}:=\langle \mathcal{G}_{\mathrm{edge}}\rangle\subseteq \mathrm{End}(\mathcal{X}),
\qquad
\mathcal{W}_{\mathrm{edge}}:=\mathcal{M}_{\mathrm{edge}}^\ast,
\]
其中 $\mathcal{X}$ 表示实现/仿真/工艺等“law-realisation data”的承载空间（与 $\chi\in\mathcal{X}(L)$ 一致）。
\end{definition}

\subsection{定理（公理降级）：局部耦合边缘算子导致非交换性（算子词是必要载体）}

\begin{theorem}[非交换性（公理降级）：共享单元/共享端点的边更新一般不可交换]
\label{thm:edge-operator-noncommutative}
设 $e_1,e_2\in E_h$ 共享端点或共享单元（即在复形上局部耦合）。
则在一般情形下，存在边缘算子 $A\in\{\mathsf{E}^{\mathrm{mat}}_{e_1},\mathsf{E}^{\mathrm{env}}_{e_1},\mathsf{E}^{\mathrm{prec}}_{e_1},
  \mathsf{E}^{\mathrm{comp}}_{e_1}\}$ 与
$B\in\{\mathsf{E}^{\mathrm{mat}}_{e_2},\mathsf{E}^{\mathrm{env}}_{e_2},\mathsf{E}^{\mathrm{prec}}_{e_2},\mathsf{E}
  ^{\mathrm{comp}}_{e_2}\}$，
使得
\[
AB\neq BA,
\]
从而 $\mathcal{M}_{\mathrm{edge}}$ 一般非交换；相应地，控制/对抗动作的可执行对象天然是算子词（operator words）。
\end{theorem}

\begin{certificate}[数学归纳法证书：按算子词长度 $N$ 归纳]
\label{cert:edge-operator-noncommutative-induction}
令谓词 $P(N)$ 为：“存在长度不超过 $N$ 的两个边缘算子词 $w_1,w_2\in \mathcal{W}_{\mathrm{edge}}$ 使 $w_1w_2\neq w_2w_1$”。
证书化要求为
\[
P(1)\ \text{成立},\qquad
(\forall N\in\mathbb{N})\ \bigl(P(N)\Rightarrow P(N+1)\bigr).
\]
\end{certificate}

\begin{proof}[证明（数学归纳法证书；谓词因果链）]
由 Theorem~\ref{thm:edge-operator-noncommutative}，存在 $A,B\in\mathcal{M}_{\mathrm{edge}}$ 使 $AB\neq BA$，从而 $P(1)$ 成立（取
  $w_1=A,w_2=B$）。
归纳步：若 $P(N)$ 成立，则同一对 $w_1,w_2$ 亦满足长度不超过 $N+1$，因此 $P(N+1)$ 成立。
故对任意 $N$ 成立。由 $P(1)$ 推出 $\mathcal{M}_{\mathrm{edge}}$ 非交换，因而“算子词”是必要载体。证毕。
\end{proof}

\subsection{构造：将边缘算子自由度并入包络（$\Xi\to\Xi^{+}$）}

\begin{definition}[边包络扩展：$\Xi^{+}:=\Xi\times \Xi^{\mathrm{edge}}$]
\label{def:edge-envelope-extension}
在既有包络分解
\[
\Xi=\Xi^{\mathrm{Geom}}\times\Xi^{\mathrm{Sim}}\times\Xi^{\mathrm{Proc}}\times\Xi^{\mathrm{Ctrl}}
\quad\text{（Definition~\ref{def:rtsc-approx-envelope}）}
\]
之上，引入边缘算子自由度作为天然预留的细粒度接口，定义扩展包络
\[
\begin{aligned}
\Xi^{+}
&:=\Xi\times \Xi^{\mathrm{edge}},\\
\Xi^{\mathrm{edge}}
&:=\Xi^{\mathrm{mat}}_{\mathrm{edge}}
\times
\Xi^{\mathrm{env}}_{\mathrm{edge}}
\times
\Xi^{\mathrm{prec}}_{\mathrm{edge}}
\times
\Xi^{\mathrm{comp}}_{\mathrm{edge}}.
\end{aligned}
\]
其中每一块可用“边集合上的截面”建模（schematic）：
\[
\Xi^{\mathrm{mat}}_{\mathrm{edge}}\cong \Gamma(E_h,\mathcal{B}^{\mathrm{mat}}),\quad
\Xi^{\mathrm{env}}_{\mathrm{edge}}\cong \Gamma(E_h,\mathcal{B}^{\mathrm{env}}),\quad
\Xi^{\mathrm{prec}}_{\mathrm{edge}}\cong \Gamma(E_h,\mathcal{B}^{\mathrm{prec}}),\quad
\Xi^{\mathrm{comp}}_{\mathrm{edge}}\cong \Gamma(E_h,\mathcal{B}^{\mathrm{comp}}).
\]
相应地，包络坐标化提升为
\[
\varepsilon:\ U\times \Xi^{+}\xrightarrow{\ \cong\ } E^{+}\subset P,
\qquad s(U)\subset E^{+},
\]
表示在同一条切面（同一工程路线/口径）附近，允许沿 $E_h$ 做局部微调与可计算的微步串联。
\end{definition}

\subsection{定义：边-Dirichlet 正则与边正则化法作用量 $\mathcal{S}_{\mathrm{law}}^{\mathrm{edge}}$}

\begin{definition}[边-Dirichlet 正则：$\mathcal{R}_{\mathrm{edge}}(\xi)$]
\label{def:edge-dirichlet-regularizer}
令 $\xi$ 为边包络变量（可取 $\xi(t)\in\Xi^{\mathrm{edge}}$ 或其某一分量）。
定义离散边平滑正则（示意）为：
\[
\mathcal{R}_{\mathrm{edge}}(\xi)
:=
\sum_{e=(v,w)\in E_h}\|\xi(v)-\xi(w)\|^2,
\]
其中 $\xi(v)$ 为将边变量投影到端点/邻域的局部表示；或等价地，可用离散外微分写作
\[
\mathcal{R}_{\mathrm{edge}}(\xi):=\|d_h\xi\|^2.
\]
\end{definition}

\begin{definition}[边正则化法作用量：$\mathcal{S}_{\mathrm{law}}^{\mathrm{edge}}$]
\label{def:edge-regularised-law-action}
在缺陷势能 $\Phi(L)$（Definition~\ref{def:rtsc-approx-defect-potential}）基础上，
定义边正则化法作用量为
\[
\mathcal{S}_{\mathrm{law}}^{\mathrm{edge}}[L,\xi]
:=
\int_0^T
\Bigl(
\Phi(L(t))
\;+\alpha\,\|\GZLCurv(L(t))\|
\;+\lambda\,\mathcal{R}_{\mathrm{edge}}\bigl(\xi(t)\bigr)
\;+\mathrm{Pen}_{\Sigma}\bigl(L(t)\bigr)
\Bigr)\,dt,
\]
其中 $\alpha,\lambda>0$，$\mathrm{Pen}_{\Sigma}$ 为跨阈值集合 $\Sigma$ 的罚项（用于约束触发/熔断/隔离/恢复图式的一致性），$\|\cdot\|$ 为与证书载体相容的范数。
\end{definition}

\subsection{命题：边正则化提供局部下降方向并抑制阈值突变}

\begin{proposition}[Edge-Regularised Descent：局部微步 $\Rightarrow$ 更平滑的回归路径族]
\label{prop:edge-regularised-descent}
在扩展包络 $\Xi^{+}$（Definition~\ref{def:edge-envelope-extension}）下，
若存在一列有限步边缘算子词与同伦修复动作，使得：
\begin{enumerate}
  \item（相对保持）受保护评价口径保持不变：
  \[
  \mathrm{Ev}\bigl(L^{(k)}(t)\bigr)=\mathrm{Ev}\bigl(L^{(0)}(t)\bigr)\quad(\forall k,t),
  \]
  \item（严格下降）边正则化作用量严格下降：存在 $k$ 使
  \[
  \mathcal{S}_{\mathrm{law}}^{\mathrm{edge}}[L^{(k+1)},\xi^{(k+1)}]
  <
  \mathcal{S}_{\mathrm{law}}^{\mathrm{edge}}[L^{(k)},\xi^{(k)}],
  \]
\end{enumerate}
则阈值切换一致性缺陷 $c_{\mathrm{thres}}$ 可被压制到“更平滑的回归路径族”中：即切换过程可由边微步串联近似为在 law-metric $J$ 意义下分段更平滑的路径，并在 defect potential $\Phi$
  意义下更易严格下降。
\end{proposition}

\begin{certificate}[证书：局部平滑能量将突变转化为可控代价]
\label{cert:edge-regularised-descent}
证书化断言由三条结构前提组成：
\begin{enumerate}
  \item（局部微步）边缘算子提供沿 $E_h$ 的可执行局部更新，使“硬切换”可分解为有限微步串联（Definition~\ref{def:edge-operator-generators}、Definition~\ref{def:edge-envelope-extension}）；
  \item（平滑正则）$\mathcal{R}_{\mathrm{edge}}$ 把空间突变压成能量项，迫使最优路径在网格上局部一致、逐步过渡（Definition~\ref{def:edge-dirichlet-regularizer}
    ）；
  \item（相对约束）$\mathcal{E}_{\mathrm{prot}}$ 固定评价口径，排除通过改口径获得伪下降（Theorem~\ref{thm:protected-evaluations-necessary}）。
\end{enumerate}
\end{certificate}

\begin{proof}[证明（谓词因果链）]
由 Certificate~\ref{cert:edge-regularised-descent}，取谓词链：
\[
\begin{aligned}
A:&\ \text{边微步存在}\Rightarrow \text{阈值切换可由有限局部更新串联近似},\\
B:&\ \mathcal{R}_{\mathrm{edge}}\ \text{入作用量}\Rightarrow \text{空间突变被惩罚}\Rightarrow \text{路径更平滑},\\
C:&\ \mathcal{E}_{\mathrm{prot}}\ \text{固定}\Rightarrow \text{下降不可由换口径伪造}.
\end{aligned}
\]
由 $A\wedge B$ 得到“更平滑的回归路径族”结论，由 $C$ 保证该平滑性与缺陷下降具有可审计语义。
结合 Proposition~\ref{prop:edge-regularised-descent} 的严格下降前提，即得命题结论。证毕。
\end{proof}

\begin{remark}[与测地线同伦扭曲一致：方向（测地线）+微步（边算子）+能量（正则）]
Pure 卷回归控制的宏观结构是“先选测地线方向，再做相对同伦扭曲修复”（Corollary~\ref{cor:relative-geodesic-twist}）。
边缘算子机制提供了该结构的可计算微步实现：测地线给出方向层，边缘算子给出可执行局部微步，
而 $\mathcal{R}_{\mathrm{edge}}$ 把空间突变压成能量项，使得在 $\Phi$ 与 $J$ 意义下更易获得严格下降与平滑过渡。
算力支持算子 $\mathsf{E}^{\mathrm{comp}}_{e}$ 进一步提供“算力越多 $\Rightarrow$ 网格越细 $\Rightarrow$ 过渡越平滑”的单调通道。
\end{remark}

\clearpage
%%%%%%%%%%%%%%%%%%%%%%%%%%%%%%%%%%%%%%%%%%%%%%%%%%%%%%%%%%%%
\section{阶段性总结：七组增益（可写入 Pure 卷、可扩展到工业工程、可算法化）}
\label{sec:discussion-gains-summary}
%%%%%%%%%%%%%%%%%%%%%%%%%%%%%%%%%%%%%%%%%%%%%%%%%%%%%%%%%%%%

\subsection{总览：从直觉到统一形式系统}

\begin{remark}[总览：将“虚拟实验 + 逆向构造 + 防路径偏航”收敛为可引用结构]
本次讨论的增益，可以归纳为把“\RTSC{} 虚拟实验 + 逆向构造 + 防路径偏航”的直觉，收敛成一套
\textbf{可写入 Pure 卷、可扩展到工业工程、可被算法化}的统一形式系统。
核心收获分为七组，分别对应：基底切换与纤维包络、总雅可比间隙可控性、法空间动态回归、相对同伦最小类、
设计生成与开放对抗策略生成等价、可打印 Benchmark 体例、以及边缘算子与 FEM 正则化的平滑下降机制。
\end{remark}

\subsection{增益 1：\PTOPB{}/\LBOPB{} 的统一 —— 基底切换 + 纤维包络}

\begin{remark}[基底切换、子丛限制与切面包络]
将 \PTOPB{}/\LBOPB{} 的“多套路线/多套组合”统一为同一基准下的两类操作：
\begin{enumerate}
  \item \textbf{基底切换：}通过 $p_i:M_\star\twoheadrightarrow M_i$ 把部分自由度并入纤维（Definition~\ref{def:base-switch-map}），并诱导包络投影
    $\pi_i=p_i\circ\pi_\star$（Definition~\ref{def:bundle-envelope}）；
  \item \textbf{子丛限制：}在诊断确认后对 $M'\subseteq M_\star$ 做限制进入类内深挖（Section~\ref{sec:ptopb-rtsc-total-jacgap}
    的“诊断--治疗”结构与相关子域论证）；
  \item \textbf{切面与包络：}将路线/管线选择严格化为切面 $s_i^{(j)}$，并将细粒度可调空间严格化为切面包络 $E_{ij}\simeq U\times\Xi_{ij}$（Definition~\ref{def:ig-section-envelope}）。
\end{enumerate}
这一步把“方向比构造更重要”变成：先在基底（诊断空间）上选同伦类/机制类，再在包络（治疗/构造空间）里做可控下降（Theorem~\ref{thm:bundle-tower-closure}）。
\end{remark}

\subsection{增益 2：\RTSC{} 可控性 —— 总雅可比间隙（Total JacGap）与格点可调必要条件}

\begin{remark}[可证书化谓词与必要可控自由度]
把“室温电阻满足 \RTSC{}”提升为可证书化谓词（参见 Section~\ref{sec:ptopb-rtsc-total-jacgap}）：
\[
\mathrm{RTSC}(m)
\Longleftrightarrow
\bigl(G_{\mathrm{tot}}(m)=0\bigr)\ \wedge\ \mathrm{Cert}(m)=1,
\]
其中 $G_{\mathrm{tot}}$ 为子节点间隙向量聚合得到的总雅可比间隙/缺陷势能主项。
可控性被严密化为：在允许算子词集合下能否把作用量压到阈值内（Proposition~\ref{prop:rtsc-jacgap-controllability-main}）。
关键必要条件也被硬化为：子节点必须包含格点/结构可变调节等可作用自由度，否则会出现不可消除下界并锁死 $G_{\mathrm{tot}}$。
\end{remark}

\subsection{增益 3：动态回归 —— 法空间、测地线与相对同伦扭曲}

\begin{remark}[扰动尖峰与回归机制的 Pure 卷化]
在 Pure 卷口径下，固定
\[
\GZLS=(\mathfrak{L};J,\mathrm{Loc},\Sigma),\quad \GZLOC,\quad \GZTLC,\quad \GZLCurv,
\]
将“扰动导致 JacGap/defect potential 急升”严格化为曲率尖峰放大缺陷变化上界（Theorem~\ref{thm:law-curvature-spike-jacgap}）。
将“动态调整回归”严格化为：在度量 $J$ 下选取回归测地线 $\gamma$（方向），
并在保持受保护评价口径的前提下做同伦扭曲 $H$（纤维方向修复），使 defect potential 单调下降并回到 $\mathrm{Loc}$ 的证书化域（Theorem~\ref{thm:ig-disturbance-regression}，Corollary~\ref{cor:relative-geodesic-twist}）。
\end{remark}

\subsection{增益 4：最小同伦类 —— 相对 $\mathcal{E}_{\mathrm{prot}}$ 的最小类}

\begin{remark}[相对同伦锚点与“最小包络类”]
形式化证明了：无论包络是两层还是更深递归（包括非因果解耦），只要额外层在 $\mathcal{E}_{\mathrm{prot}}$ 下可回缩，
它们都落在同一个\textbf{相对同伦等价类}（Section~\ref{sec:puremath-relative-homotopy-minimal-envelope}，Proposition~\ref{prop:relative-homotopy-depth-invariant}）。
因此 $\mathcal{E}_{\mathrm{prot}}$（Definition~\ref{def:protected-evaluation-functionals}）是语义锚点；
缺少该锚点会使“同伦扭曲/间隙下降”退化为“更换评价口径”的伪下降（Theorem~\ref{thm:protected-evaluations-necessary}）。
\end{remark}

\subsection{增益 5：工程生成 $\simeq$ 开放对抗策略生成 —— 非交换算子词范式}

\begin{remark}[从工程逆向生成到开放对抗策略生成]
把工程生成过程写成非交换算子词生成，并严格嵌入开放博弈系统对抗主丛结构：
\[
\text{生成工程设计（含非交换算子词）}
\ \simeq_{\mathcal{E}_{\mathrm{prot}}}\
\text{生成开放对抗策略（含非交换算子词）}.
\]
其时间/信息结构由“非自治法流 + 阈值 $\Sigma$ + 曲率尖峰 + 在线回归”决定，更接近实时开放对抗而非回合制封闭对弈（Theorem~\ref{thm:noncommutative-design-strategy-equivalence}，Corollary~\ref{cor:openra-vs-alphago}）。
并可进一步压成 open adversarial \OPB{} 的基准模板（Theorem~\ref{thm:opb-benchmark-template-equivalence}）。
\end{remark}

\subsection{增益 6：可直接复用的 \RTSC{}-工业几何体 Benchmark（Pure 卷体例）}

\begin{remark}[Benchmark 作为“可打印基准件”]
已给出可直接复用的 \RTSC{}-工业几何体 Benchmark（Definition/Construction/Proposition/Remark 体例）：
由受保护评价与分层标签诱导基底 $M$，由主丛 $\pi:P\to M$ 与切面/包络组织路线与细节，
并以最小三分量证书 $\JacGap=(c_{\mathrm{chart}},c_{\mathrm{lift}},c_{\mathrm{thres}})$ 与 defect potential $\Phi$ 作为统一判据，
从而把“设计生成 $\simeq$ 策略生成”写成可调用命题（Proposition~\ref{prop:rtsc-ig-design-strategy-equivalence}，Proposition~\ref{prop:rtsc-approx-design-strategy-equivalence}）。
对应段落见 Section~\ref{sec:puremath-rtsc-industrial-geometry-benchmark} 与 Section~\ref{sec:puremath-rtsc-approx-lowres-cage-bypass}。
\end{remark}

\subsection{增益 7：边缘算子 + FEM 优化空间 —— 把突变正则化为平滑下降}

\begin{remark}[边-局部微步 + 边正则 action]
引入边缘算子与有限元网格复形，将原本的阈值突变与数值口径诱导突变分解为可控局部微步，并以边-Dirichlet 正则化压制突变：
\[
\mathcal{S}_{\mathrm{law}}^{\mathrm{edge}}[L,\xi]
=
\int_0^T
\Bigl(
\Phi(L(t))
\;+\alpha\|\GZLCurv(L(t))\|
\;+\lambda\mathcal{R}_{\mathrm{edge}}(\xi(t))
\;+\mathrm{Pen}_{\Sigma}(L(t))
\Bigr)\,dt
\]
（Definition~\ref{def:edge-regularised-law-action}）。
从结构上解释了“过渡变平滑”：$c_{\mathrm{thres}}$ 与 $c_{\mathrm{chart}}$ 更容易获得局部下降方向，且下降过程更可控、可复现（Proposition~\ref{prop:edge-regularised-descent}）。
\end{remark}

\subsection{形式化收束：引理/命题/定理/推论（带证书与证明）}

\begin{lemma}[索引化引理：七组增益均有文内可引用落点]
\label{lem:discussion-gains-indexed}
本节所述七组增益中的每一组，都可由本文既有的“定义--命题/定理--证书--证明”条目在形式上支撑与引用；
因此“讨论增益”可被视为一个可复用的 Pure 卷形式系统段落，而非仅叙事性总结。
\end{lemma}

\begin{certificate}[证书：七组增益与文内标签映射]
\label{cert:discussion-gains-map}
七组增益的形式化落点分别为：
\begin{enumerate}
  \item（基底切换/包络）Definition~\ref{def:base-switch-map}、Definition~\ref{def:bundle-envelope}、Theorem~\ref{thm:bundle-tower-closure}；
  \item（总 JacGap 可控性）Section~\ref{sec:ptopb-rtsc-total-jacgap}、Proposition~\ref{prop:rtsc-jacgap-controllability-main}；
  \item（测地线同伦扭曲回归）Theorem~\ref{thm:law-curvature-spike-jacgap}、Theorem~\ref{thm:ig-disturbance-regression}、
    Corollary~\ref{cor:relative-geodesic-twist}；
  \item（相对最小同伦类）Section~\ref{sec:puremath-relative-homotopy-minimal-envelope}、Proposition~\ref{prop:relative-homotopy-depth-invariant}、Theorem~\ref{thm:protected-evaluations-necessary}；
  \item（设计生成 $\simeq$ 策略生成）Section~\ref{sec:puremath-open-system-adversarial-game}、Theorem~\ref{thm:noncommutative-design-strategy-equivalence}、Theorem~\ref{thm:opb-benchmark-template-equivalence}；
  \item（\RTSC{}-工业几何体 Benchmark）Section~\ref{sec:puremath-rtsc-industrial-geometry-benchmark}、Section~\ref{sec:puremath-rtsc-approx-lowres-cage-bypass}；
  \item（边缘算子与 FEM 正则）Section~\ref{sec:puremath-edge-operators-fem-envelope}、Definition~\ref{def:edge-regularised-law-action}、Proposition~\ref{prop:edge-regularised-descent}。
\end{enumerate}
\end{certificate}

\begin{proof}[证明（谓词因果链）]
由 Certificate~\ref{cert:discussion-gains-map}，对每一组增益取谓词 $A_i$ 为“存在一组可引用的定义/定理/命题支撑该增益的可证书化表述”。
证书给出了 $A_1,\dots,A_7$ 的具体引用，因此得到
\[
A_1\wedge A_2\wedge \cdots \wedge A_7.
\]
由合取推出“七组增益均有文内可引用落点”，从而 Lemma~\ref{lem:discussion-gains-indexed} 成立。证毕。
\end{proof}

\begin{proposition}[统一模板命题：七组增益构成可算法化的诊断--治疗--回归--正则框架]
\label{prop:discussion-gains-template}
在固定受保护评价泛函族 $\mathcal{E}_{\mathrm{prot}}$ 的前提下，
本次讨论收敛得到的结构可被统一理解为：
\textbf{在开放对抗主丛与包络上，先做诊断（基底切换/同伦类选择），再做治疗（包络内构造优化），并在扰动下通过测地线同伦扭曲回归，同时用边缘算子与 FEM 正则把突变压成可控能量项}。
\end{proposition}

\begin{certificate}[证书：模板四要素对应的形式化条目]
\label{cert:discussion-gains-template}
\begin{enumerate}
  \item（诊断）Definition~\ref{def:diagnosis-and-treatment} 与基底切换结构 Definition~\ref{def:base-switch-map}；
  \item（治疗）切面包络与可逆向生成条目 Definition~\ref{def:ig-section-envelope}、Theorem~\ref{thm:ig-invgen-ohu}；
  \item（回归）Theorem~\ref{thm:ig-disturbance-regression} 与 Corollary~\ref{cor:relative-geodesic-twist}；
  \item（正则）Definition~\ref{def:edge-regularised-law-action} 与 Proposition~\ref{prop:edge-regularised-descent}。
\end{enumerate}
\end{certificate}

\begin{proof}[证明（谓词因果链）]
由 Certificate~\ref{cert:discussion-gains-template}，取谓词链：
\[
\begin{aligned}
A:&\ \text{诊断层}= \text{基底切换/同伦类选择},\\
B:&\ \text{治疗层}= \text{切面包络内的构造优化},\\
C:&\ \text{回归层}= \text{测地线方向 + 相对同伦扭曲修复},\\
D:&\ \text{正则层}= \text{边缘算子微步 + FEM 正则能量}.
\end{aligned}
\]
由 $A\wedge B\wedge C\wedge D$ 得到“七组增益可压缩为诊断--治疗--回归--正则的统一模板”。证毕。
\end{proof}

\begin{theorem}[收敛定理（公理降级）：七组增益形成可写入 Pure 卷且可算法化的统一形式系统]
\label{thm:discussion-convergence}
在相对受保护评价泛函族 $\mathcal{E}_{\mathrm{prot}}$ 的约束下，
本次讨论得到的七组增益可被一致地写入 Pure 卷语言，并可扩展到工业工程与算法实现：
其核心对象为开放对抗主丛及其包络（含边缘算子扩展），核心判据为最小三分量 $\JacGap$ 与 defect potential $\Phi$，
核心机制为“基底切换诊断 + 包络治疗 + 测地线同伦扭曲回归 + 边正则平滑下降”。
\end{theorem}

\begin{certificate}[数学归纳法证书：按“纳入的增益组数” $n$ 归纳]
\label{cert:discussion-convergence-induction}
令谓词 $P(n)$ 表示：“当纳入前 $n$ 组增益时，所得结构仍可被表述为相对 $\mathcal{E}_{\mathrm{prot}}$ 的开放对抗主丛问题，并存在相应的 defect potential 下降机制”。
证书化要求为：
\[
P(1)\ \text{成立},\qquad
(\forall n\in\{1,\dots,6\})\ \bigl(P(n)\Rightarrow P(n+1)\bigr).
\]
\end{certificate}

\begin{proof}[证明（数学归纳法证书；谓词因果链）]
由 Lemma~\ref{lem:discussion-gains-indexed} 与 Proposition~\ref{prop:discussion-gains-template}，每一组增益都对应一个可调用的形式化条目，
  并且可被压入同一模板。
取谓词链：
\[
\begin{aligned}
A:&\ P(1)\ \text{由基底切换/包络结构的封闭性给出（Theorem~\ref{thm:bundle-tower-closure})},\\
B:&\ P(n)\Rightarrow P(n+1)\ \text{由“相对锚点保持 + 包络可扩展”给出（Proposition~\ref{prop:relative-homotopy-depth-invariant})},\\
C:&\ \text{每一步扩展均保持 defect potential 可下降机制（Theorem~\ref{thm:ig-invgen-ohu}，Proposition~\ref{prop:edge-regularised-descent})}.
\end{aligned}
\]
由 $A$ 得到基步，由 $B\wedge C$ 得到归纳步，从而对 $n=1,\dots,7$ 成立。
因此七组增益整体收敛为可写入 Pure 卷且可算法化的统一形式系统。证毕。
\end{proof}

\begin{corollary}[一句话推论：\RTSC{}-工业几何体逆向工程设计 $\equiv$ 相对同伦下的开放对抗主丛问题]
\label{cor:discussion-one-sentence}
本次讨论的增益可压成一句话：
\[
\boxed{
\begin{aligned}
&\text{\RTSC{}-工业几何体逆向工程设计}\ \equiv\ \text{开放对抗主丛问题；}\\
&\text{域：在 }\GZLS\text{ 上；相对 }\mathcal{E}_{\mathrm{prot}}\text{；}\\
&\text{判据：}\JacGap/\Phi,\quad \text{机制：测地线同伦扭曲 + 边正则；}\\
&\text{目标：实现可控下降。}
\end{aligned}
}
\]
\end{corollary}

\begin{certificate}[证书：由收敛定理直接推出]
\label{cert:discussion-one-sentence}
由 Theorem~\ref{thm:discussion-convergence}，统一对象（开放对抗主丛及包络）、统一判据（$\JacGap/\Phi$）与统一机制（诊断--治疗--回归--正则）均已确立；
因此将其压缩为一句话结论得到推论所述等价口径。
\end{certificate}

\begin{proof}[证明（谓词因果链）]
由 Certificate~\ref{cert:discussion-one-sentence}，取谓词链：
\[
\begin{aligned}
A:&\ \text{统一对象已确立},\\
B:&\ \text{统一判据已确立},\\
C:&\ \text{统一机制已确立}.
\end{aligned}
\]
由 $A\wedge B\wedge C$ 直接得到推论陈述。证毕。
\end{proof}

\clearpage
%%%%%%%%%%%%%%%%%%%%%%%%%%%%%%%%%%%%%%%%%%%%%%%%%%%%%%%%%%%%
\section{障碍消去与执行降级定理（\RTSC{}-工业几何体 Benchmark）}
\label{sec:puremath-obstruction-to-pi0}
%%%%%%%%%%%%%%%%%%%%%%%%%%%%%%%%%%%%%%%%%%%%%%%%%%%%%%%%%%%%

\subsection{设定：法空间、局部化、阈值与被保护评价}

\begin{remark}[定位与重要说明：讨论的是形式系统内部的“障碍 $\Rightarrow$ 连通”降级]
您指出得对：在本次 Pure 卷形式系统语境里，应避免用外部叙事句式替代内部数学对象。
在我们已建立的 \RTSC{}-工业几何体 Benchmark（Section~\ref{sec:puremath-rtsc-industrial-geometry-benchmark}）与
\textsf{RTSC-approx} 系统（Section~\ref{sec:puremath-rtsc-approx-lowres-cage-bypass}）中，关键是：
\begin{enumerate}
  \item 形式系统内部的“不可行”对应为上同调/相对上同调意义下的障碍（obstruction），并由证书接口可计算地表征；
  \item 本次讨论构造出一个\textbf{质变的相对同伦类}，使这些障碍在该类内可被有限同伦动作消去（$\Phi\downarrow 0$）；
  \item 从而把执行问题从“障碍类计算”\textbf{降级（decategorify）}为“连通类/可达性”（$\pi_0$ 与 path-connectivity）。
\end{enumerate}
\textbf{重要说明：}本节证明完全在 Pure 卷形式系统内部进行：证明的是“在该形式系统定义下，障碍类 $\Rightarrow$ 连通类”的降级；
它不等同于对自然界是否存在某个室温超导材料的实验性证明。
\end{remark}

沿用既有 Benchmark 设定（Definition~\ref{def:rtsc-ig-benchmark-base}、Definition~\ref{def:rtsc-ig-principal-bundle}）：
固定
\[
\GZLS=(\mathfrak{L};J,\mathrm{Loc},\Sigma),\qquad \GZLOC,\ \GZTLC,\ \GZLCurv,
\]
固定受保护评价泛函族
\[
\mathcal{E}_{\mathrm{prot}}=\{\mathrm{Ev}_a\}_{a=1}^{q},\qquad
\mathrm{Ev}_a:\mathrm{Loc}\to\mathbb{R},\qquad
\mathrm{Ev}=(\mathrm{Ev}_1,\dots,\mathrm{Ev}_q):\mathrm{Loc}\to\mathbb{R}^q,
\]
并由 $\Sigma$ 诱导阈值分层标签 $\rho_{\Sigma}:\mathrm{Loc}\to\mathcal{R}_{\Sigma}$。
由此得到可判定语义基底（Definition~\ref{def:rtsc-ig-benchmark-base}）：
\[
M=\Bigl\{
(\mathbf{y},r)\in\mathbb{R}^q\times \mathcal{R}_{\Sigma}\ \Big|\ \exists\,L\in\mathrm{Loc},\ \mathrm{Ev}(L)=\mathbf{y},\
  \rho_{\Sigma}(L)=r
\Bigr\}.
\]

\subsection{下降数据：开覆盖、局部切面与过渡函数}

\begin{definition}[下降数据：开覆盖、局部切面与过渡函数]
\label{def:appendix-descent-data}
取基底 $M$ 的一个开覆盖 $\mathcal{U}=\{U_\alpha\}_{\alpha\in I}$。
在每个 $U_\alpha$ 上选取局部切面（局部实现）
\[
s_\alpha:U_\alpha\to P,\qquad \pi\circ s_\alpha=\mathrm{Id}_{U_\alpha},
\]
其中 $\pi:P\to M$ 是表示冗余入纤维的开放对抗主丛载体（Definition~\ref{def:rtsc-ig-principal-bundle}），结构对象记为 $G_{\mathrm{rep}}$。
在重叠域 $U_{\alpha\beta}:=U_\alpha\cap U_\beta$ 上，由主丛结构得到过渡函数（或过渡 1-箭头）
\[
g_{\alpha\beta}:U_{\alpha\beta}\to G_{\mathrm{rep}}
\quad\text{使得}\quad
s_\alpha=s_\beta\cdot g_{\beta\alpha}\ \text{于 }U_{\alpha\beta}.
\]
由于 $\pi:P\to M$ 为主 $G_{\mathrm{rep}}$-丛，过渡函数满足 Čech 1-cocycle 条件：
\[
g_{\alpha\beta}\,g_{\beta\gamma}\,g_{\gamma\alpha}=e\qquad \text{于 }U_{\alpha\beta\gamma}:=U_\alpha\cap U_\beta\cap
  U_\gamma.
\]
\end{definition}

\subsection{三类“上同调障碍类”的严格化：中心扩张与 2-上同调类}

\begin{definition}[chart 障碍类：中心扩张抽取与 $\check H^2$ 表述]
\label{def:chart-obstruction-class}
为把“非交换过渡数据”抽取为可计算的上同调障碍，我们采用标准的中心扩张/可交换系抽取做法：
设存在可交换层（Abelian sheaf）$\mathcal{A}$ 与短正合列（群丛或层的短正合列）
\[
1\longrightarrow \mathcal{A}\longrightarrow \widetilde{G}\longrightarrow G_{\mathrm{rep}}\longrightarrow 1.
\]
对 Definition~\ref{def:appendix-descent-data} 中的 $g_{\alpha\beta}$ 选择提升
$\tilde g_{\alpha\beta}:U_{\alpha\beta}\to \widetilde{G}$。
在三重交叠域上定义 Čech 2-cochain：
\[
\omega^{\mathrm{chart}}_{\alpha\beta\gamma}
:=
\tilde g_{\alpha\beta}\,\tilde g_{\beta\gamma}\,\tilde g_{\gamma\alpha}
\in \mathcal{A}(U_{\alpha\beta\gamma}).
\]
\end{definition}

\begin{lemma}[引理：$\omega^{\mathrm{chart}}$ 为 2-cocycle 且同调类不依赖提升]
\label{lem:chart-omega-cocycle}
在 Definition~\ref{def:appendix-descent-data} 与 Definition~\ref{def:chart-obstruction-class} 的设定下：
\begin{enumerate}
  \item $\omega^{\mathrm{chart}}$ 满足 Čech 2-cocycle 条件 $\delta\omega^{\mathrm{chart}}=1$；
  \item 不同提升 $\{\tilde g_{\alpha\beta}\}$ 得到的 $\omega^{\mathrm{chart}}$ 只相差一个 Čech 2-coboundary；
\end{enumerate}
从而定义了良定的上同调障碍类
\[
\mathfrak{o}_{\mathrm{chart}}
:=
[\omega^{\mathrm{chart}}]\in \check{H}^2(M;\mathcal{A}).
\]
\end{lemma}

\begin{certificate}[证书：中心扩张与 Čech 上同调的标准构造事实]
\label{cert:chart-omega-standard}
在中心扩张 $1\to\mathcal{A}\to\widetilde{G}\to G_{\mathrm{rep}}\to 1$ 下，
对 Čech 1-cocycle $\{g_{\alpha\beta}\}$ 的任意提升 $\{\tilde g_{\alpha\beta}\}$，
由
$g_{\alpha\beta}g_{\beta\gamma}g_{\gamma\alpha}=e$
可推出 $\tilde g_{\alpha\beta}\tilde g_{\beta\gamma}\tilde g_{\gamma\alpha}\in \ker(\widetilde{G}\to G_{\mathrm{rep}})
  =\mathcal{A}$，
并且更换提升仅改变一个 2-coboundary，因此同调类良定。
\end{certificate}

\begin{proof}[证明（谓词因果链）]
取谓词链：
\[
\begin{aligned}
A:&\ g_{\alpha\beta}g_{\beta\gamma}g_{\gamma\alpha}=e
\Rightarrow
\tilde g_{\alpha\beta}\tilde g_{\beta\gamma}\tilde g_{\gamma\alpha}\in \mathcal{A},
\\
B:&\ A\Rightarrow \delta\omega^{\mathrm{chart}}=1\ \text{（由结合律/中心性与 Čech 余边界定义）},
\\
C:&\ \tilde g'_{\alpha\beta}=a_{\alpha\beta}\tilde g_{\alpha\beta},\ a_{\alpha\beta}\in\mathcal{A}
\Rightarrow
\omega'=\omega\cdot \delta a,
\\
D:&\ C\Rightarrow [\omega']=[\omega]\in \check H^2(M;\mathcal{A}).
\end{aligned}
\]
由 $B$ 得到 $\omega^{\mathrm{chart}}$ 为 2-cocycle；由 $D$ 得到同调类不依赖提升。
因此 $\mathfrak{o}_{\mathrm{chart}}:=[\omega^{\mathrm{chart}}]$ 良定。证毕。
\end{proof}

\begin{definition}[lift/threshold 障碍类（schematic）]
\label{def:lift-thres-obstructions}
同理，将“水平提升一致性缺陷”（$\GZLOC/\GZTLC$ 的运输与同伦修复的相容性失配）
与“阈值切换一致性缺陷”（跨 $\Sigma$ 的切换图式失配）抽取为可交换障碍：
设存在可交换层 $\mathcal{B},\mathcal{C}$，使得对应障碍类分别落在
\[
\mathfrak{o}_{\mathrm{lift}}\in \check{H}^2(M;\mathcal{B}),
\qquad
\mathfrak{o}_{\mathrm{thres}}\in \check{H}^2(M,\Sigma;\mathcal{C}).
\]
\end{definition}

\subsection{证书化投影：$\JacGap$ 作为障碍代表元}

\begin{definition}[障碍类三元组与三分量 $\JacGap$ 证书：$\mathfrak{o}_{\mathrm{chart}},\mathfrak{o}_{\mathrm{lift}},\mathfrak{o}
  _{\mathrm{thres}}$]
\label{def:obstruction-triple}
在本 Benchmark 中，取最小三分量证书
\[
\JacGap(L)=\bigl(c_{\mathrm{chart}}(L),c_{\mathrm{lift}}(L),c_{\mathrm{thres}}(L)\bigr),
\qquad
\Phi(L)=\sum_{\bullet} w_{\bullet}\,\|c_{\bullet}(L)\|,\quad w_{\bullet}>0.
\]
其中 $c_{\mathrm{chart}}(L),c_{\mathrm{lift}}(L),c_{\mathrm{thres}}(L)$ 被视为
Definition~\ref{def:chart-obstruction-class}--Lemma~\ref{lem:chart-omega-cocycle} 与 Definition~\ref{def:lift-thres-obstructions}
所给三类障碍类
$(\mathfrak{o}_{\mathrm{chart}},\mathfrak{o}_{\mathrm{lift}},\mathfrak{o}_{\mathrm{thres}})$
的\textbf{可计算代表/证书化投影}（可审计接口：当代表为零时应能见证障碍类为零）。
\end{definition}

\begin{certificate}[证书：$\JacGap$ 分量为障碍类的证书化代表]
\label{cert:jacgap-as-obstruction}
在本 Benchmark 的证书口径中，约定以下蕴含成立（作为“可计算代表/投影”的可审计定义边界）：
\[
c_{\mathrm{chart}}(L)=0\Rightarrow \mathfrak{o}_{\mathrm{chart}}=0,\qquad
c_{\mathrm{lift}}(L)=0\Rightarrow \mathfrak{o}_{\mathrm{lift}}=0,\qquad
c_{\mathrm{thres}}(L)=0\Rightarrow \mathfrak{o}_{\mathrm{thres}}=0.
\]
其语义是：当证书化代表元在所选口径下为零时，对应障碍类在相应（相对）Čech 上同调中消失。
\end{certificate}

\subsection{相对同伦与有限同伦动作}

\begin{definition}[\textrm{$\mathcal{E}_{\mathrm{prot}}$}-相对同伦与有限同伦动作]
\label{def:relative-homotopy-finite-moves}
称 $L_0,L_1\in\mathrm{Loc}$ 在 $\mathcal{E}_{\mathrm{prot}}$ 下相对同伦，若存在连续路径
\[
H:[0,1]\to \mathrm{Loc},\qquad H(0)=L_0,\quad H(1)=L_1,
\]
并满足对所有 $t\in[0,1]$：
\[
\mathrm{Ev}\bigl(H(t)\bigr)=\mathrm{Ev}(L_0),
\]
且阈值分层标签 $\rho_{\Sigma}(H(t))$ 仅按 $\Sigma$ 的允许图式进行（允许有限次切换，但评价口径不被偷换）。

称“有限同伦动作”指存在有限序列
\[
L^{(0)}\to L^{(1)}\to\cdots\to L^{(N)},
\]
其中每一步由 Pure 卷允许的 homotopy moves（例如高阶同伦数据的有限调整）与包络可执行更新（$\xi\mapsto\xi'$）生成，
并始终满足上述相对约束（$\mathrm{Ev}$ 口径固定）。
\end{definition}

\begin{definition}[质变相对同伦类：$\mathcal{H}_{\mathrm{Q}}$]
\label{def:qualitative-relative-homotopy-class}
在相对约束下，称一条同伦类为“质变同伦类”，若它包含一条有限同伦动作序列使 defect potential 可严格下降并达到零：
存在 $L^{(0)}$ 及有限序列 $L^{(0)}\to\cdots\to L^{(N)}$ 使
\[
\Phi\bigl(L^{(k+1)}\bigr)<\Phi\bigl(L^{(k)}\bigr)\quad(\forall k),\qquad
\Phi\bigl(L^{(N)}\bigr)=0.
\]
记该类为 $\mathcal{H}_{\mathrm{Q}}$（relative qualitative homotopy class）。
\end{definition}

\subsection{定理：障碍消去与执行降级（由上同调障碍到连通类 $\pi_0$）}

\subsubsection{定理 A：障碍类消去（$\Phi\to 0$ $\Rightarrow$ 障碍类为零）}

\begin{theorem}[障碍消去定理：有限同伦下降到 $\Phi=0$ 蕴含障碍类消失]
\label{thm:obstruction-vanishing}
若存在有限同伦动作序列 $L^{(0)}\to\cdots\to L^{(N)}$ 满足：
\begin{enumerate}
  \item（相对保护）$\mathrm{Ev}(L^{(k)})=\mathrm{Ev}(L^{(0)})$ 对所有 $k$ 成立，且阈值切换仅按 $\Sigma$ 的允许图式进行；
  \item（严格下降并归零）$\Phi(L^{(k+1)})<\Phi(L^{(k)})$ 对所有 $k$ 成立，且 $\Phi(L^{(N)})=0$；
\end{enumerate}
则三类障碍类同时消失：
\[
\mathfrak{o}_{\mathrm{chart}}=0,\qquad
\mathfrak{o}_{\mathrm{lift}}=0,\qquad
\mathfrak{o}_{\mathrm{thres}}=0.
\]
\end{theorem}

\begin{certificate}[证书：$\Phi=0$ $\Rightarrow$ 三分量证书全为零]
\label{cert:phi0-implies-components0}
在定义 $\Phi(L)=\sum_{\bullet} w_{\bullet}\|c_{\bullet}(L)\|$ 且 $w_{\bullet}>0$、$\|\cdot\|\ge 0$ 的前提下，
\[
\Phi(L)=0\Rightarrow c_{\mathrm{chart}}(L)=0,\ c_{\mathrm{lift}}(L)=0,\ c_{\mathrm{thres}}(L)=0.
\]
\end{certificate}

\begin{proof}[证明（谓词因果链）]
由条件 $\Phi(L^{(N)})=0$ 与 Certificate~\ref{cert:phi0-implies-components0}，得到
\[
c_{\mathrm{chart}}(L^{(N)})=0,\quad c_{\mathrm{lift}}(L^{(N)})=0,\quad c_{\mathrm{thres}}(L^{(N)})=0.
\]
再由 Certificate~\ref{cert:jacgap-as-obstruction}，取谓词链：
\[
\begin{aligned}
A:&\ c_{\mathrm{chart}}(L^{(N)})=0\Rightarrow \mathfrak{o}_{\mathrm{chart}}=0,\\
B:&\ c_{\mathrm{lift}}(L^{(N)})=0\Rightarrow \mathfrak{o}_{\mathrm{lift}}=0,\\
C:&\ c_{\mathrm{thres}}(L^{(N)})=0\Rightarrow \mathfrak{o}_{\mathrm{thres}}=0.
\end{aligned}
\]
由 $A\wedge B\wedge C$ 得到三类障碍类同时消失。证毕。
\end{proof}

\subsubsection{补充定理（公理降级）：零障碍类可编译为有限同伦动作（$\Phi\downarrow 0$）}

\begin{theorem}[可执行完备性（公理降级）：障碍类为零蕴含存在 $\Phi\to 0$ 的有限下降链]
\label{thm:obstruction-zero-implies-finite-descent}
设在某一相对约束类内（Definition~\ref{def:relative-homotopy-finite-moves}），三类障碍类为零：
\[
\mathfrak{o}_{\mathrm{chart}}=0,\qquad
\mathfrak{o}_{\mathrm{lift}}=0,\qquad
\mathfrak{o}_{\mathrm{thres}}=0.
\]
并且存在一个有限口径的包络坐标化 $E\simeq U\times\Xi$ 使得允许的 homotopy moves 与包络更新可以在该包络中被实现。
则存在有限同伦动作序列 $L^{(0)}\to\cdots\to L^{(N)}$ 满足
\[
\Phi(L^{(N)})=0,
\]
从而存在质变相对同伦类 $\mathcal{H}_{\mathrm{Q}}$（Definition~\ref{def:qualitative-relative-homotopy-class}）。
\end{theorem}

\begin{certificate}[证书：Čech 余边界见证可编译为包络内有限修复动作]
\label{cert:obstruction-zero-compilable}
当 $\mathfrak{o}_{\mathrm{chart}}=0$ 时，存在 1-cochain 见证 $b^{\mathrm{chart}}$ 使 $\omega^{\mathrm{chart}}=\delta
  b^{\mathrm{chart}}$；
当 $\mathfrak{o}_{\mathrm{lift}}=0$ 时，存在相应的 1-cochain/同伦修复见证 $b^{\mathrm{lift}}$；
当 $\mathfrak{o}_{\mathrm{thres}}=0$ 时，存在相对 1-cochain/切换图式见证 $b^{\mathrm{thres}}$。
并假定这些见证均可被编译为有限步的“局部重参数化 + 高阶同伦数据修复 + 阈值图式修复”动作，
从而在所选包络坐标化内实现
\[
c_{\mathrm{chart}}=c_{\mathrm{lift}}=c_{\mathrm{thres}}=0.
\]
\end{certificate}

\begin{proof}[证明（谓词因果链）]
由 Certificate~\ref{cert:obstruction-zero-compilable}，取谓词链：
\[
\begin{aligned}
A:&\ \mathfrak{o}_{\mathrm{chart}}=\mathfrak{o}_{\mathrm{lift}}=\mathfrak{o}_{\mathrm{thres}}=0
\Rightarrow \exists\ \text{余边界/修复见证可编译为有限动作},\\
B:&\ A \Rightarrow c_{\mathrm{chart}}=c_{\mathrm{lift}}=c_{\mathrm{thres}}=0,\\
C:&\ B \Rightarrow \Phi=0.
\end{aligned}
\]
由 $A\Rightarrow B\Rightarrow C$ 得到存在有限同伦动作序列使 $\Phi(L^{(N)})=0$。证毕。
\end{proof}

\subsubsection{定理 B：执行降级为连通类（$\pi_0$）}

\begin{definition}[严格可执行子集：$\Xi_{\mathrm{exec}}(\mathbf{y},r)$]
\label{def:xi-exec}
在一个包络坐标化 $E\simeq U\times\Xi$ 中，固定基底点 $(\mathbf{y},r)\in U\subset M$。
定义严格可执行子集为
\[
\Xi_{\mathrm{exec}}(\mathbf{y},r)
:=
\Bigl\{
\xi\in\Xi\ \Big|\ ((\mathbf{y},r),\xi)\in U\times\Xi\ \text{且对应的法对象 }L\ \text{满足 }\Phi(L)=0
\Bigr\}.
\]
其语义是：在固定“可判定语义”（$\mathbf{y},r$）下，筛出能达到 $\Phi=0$（障碍被消去、strict sector）的一类可执行实现细节。
\end{definition}

\begin{theorem}[执行降级定理：障碍消去后分类仅剩 $\pi_0(\Xi_{\mathrm{exec}})$]
\label{thm:decategorify-to-pi0}
在 Theorem~\ref{thm:obstruction-vanishing} 的结论成立且 $E\simeq U\times\Xi$ 给定的前提下，
若对每个固定基底点 $(\mathbf{y},r)\in U$，严格可执行集合 $\Xi_{\mathrm{exec}}(\mathbf{y},r)$ 路径连通，
则在该基底点下，“严格执行”的分类降级为连通类：
任意两种严格可执行实现 $\xi_0,\xi_1\in\Xi_{\mathrm{exec}}(\mathbf{y},r)$ 在相对同伦意义下等价；
更一般地，严格执行的分类同构于
\[
\pi_0\!\bigl(\Xi_{\mathrm{exec}}(\mathbf{y},r)\bigr).
\]
\end{theorem}

\begin{certificate}[证书：固定基底点后纤维与 $\Xi$ 同胚，路径连通 $\Rightarrow$ 存在纤维路径]
\label{cert:xi-exec-path}
由 $E\simeq U\times\Xi$，固定 $(\mathbf{y},r)$ 后纤维同胚于 $\Xi$。
若 $\Xi_{\mathrm{exec}}(\mathbf{y},r)$ 路径连通，则对任意 $\xi_0,\xi_1\in\Xi_{\mathrm{exec}}(\mathbf{y},r)$ 存在连续路径
\[
\gamma:[0,1]\to \Xi_{\mathrm{exec}}(\mathbf{y},r),\qquad \gamma(0)=\xi_0,\ \gamma(1)=\xi_1.
\]
\end{certificate}

\begin{proof}[证明（谓词因果链）]
由 Certificate~\ref{cert:xi-exec-path} 得到 $\gamma$。
将路径推送到包络总空间：令
\[
\widetilde{\gamma}(t):=\varepsilon\bigl((\mathbf{y},r),\gamma(t)\bigr)\in E\subset P.
\]
由于基底点固定，$\pi(\widetilde{\gamma}(t))=(\mathbf{y},r)$ 不变，故受保护评价口径保持不变；
又因为 $\gamma(t)\in \Xi_{\mathrm{exec}}(\mathbf{y},r)$，对应法对象满足 $\Phi=0$，故在 strict sector 内沿纤维运动不触发障碍复活。
因此 $\xi_0,\xi_1$ 在相对同伦意义下等价；分类仅剩连通分支，即 $\pi_0(\Xi_{\mathrm{exec}}(\mathbf{y},r))$。证毕。
\end{proof}

\subsection{推论：质变同伦类 $\mathcal{H}_{\mathrm{Q}}$ 与 strict sector 可达性}

\begin{corollary}[strict sector 可达性（形式系统内）$\Longleftrightarrow$ 障碍可消去，执行降级到 $\pi_0$]
\label{cor:qualitative-breakthrough}
定义 strict sector（以 \textsf{RTSC-approx} 口径示意）为
\[
\mathrm{Strict}_{\mathrm{RTSC}}
:=
\Bigl\{
L\in\mathrm{Loc}\ \Big|\ \Phi(L)=0,\ \mathrm{Ev}_{\rho,300}(L)\le \varepsilon_{\rho},\ \mathrm{Ev}_{\mathrm{Cert}}(L)=1
\Bigr\}.
\]
并固定一个基底点 $(\mathbf{y},r)\in M$，要求其受保护评价分量满足上述 \textsf{RTSC-approx} 门槛（该门槛作为 $\mathbf{y}$ 的分量语义被固定）。
则在 Pure 卷形式系统内部，以下断言等价：
\begin{enumerate}
  \item 存在与 $(\mathbf{y},r)$ 相容的质变同伦类 $\mathcal{H}_{\mathrm{Q}}$，使得从该基底点的某个实现出发，存在有限同伦动作把 $\Phi$ 严格下降到 $0$，并落入
    $\mathrm{Strict}_{\mathrm{RTSC}}$；
  \item 三类障碍类可在相对同伦中被消去：$\mathfrak{o}_{\mathrm{chart}}=\mathfrak{o}_{\mathrm{lift}}=\mathfrak{o}_{\mathrm{thres}}=0$；
  \item 障碍消去后，现实执行降级为拓扑连通/可达性问题：分类由 $\pi_0(\Xi_{\mathrm{exec}}(\mathbf{y},r))$ 给出。
\end{enumerate}
\end{corollary}

\begin{certificate}[证书：由定理 A/B 与质变同伦类定义合成]
\label{cert:qualitative-breakthrough}
\begin{enumerate}
  \item 由 Definition~\ref{def:qualitative-relative-homotopy-class}，(1) 等价于“存在有限序列使 $\Phi\downarrow 0$ 且落入 strict sector”；

  \item 由 Theorem~\ref{thm:obstruction-vanishing}，$\Phi\to 0$ 蕴含障碍类为零；
  \item 由 Theorem~\ref{thm:obstruction-zero-implies-finite-descent}，障碍类为零在本形式系统口径下可编译为 $\Phi\to 0$ 的有限下降链；
  \item 由 Theorem~\ref{thm:decategorify-to-pi0}，在 strict sector 内执行分类由 $\pi_0(\Xi_{\mathrm{exec}})$ 给出。
\end{enumerate}
\end{certificate}

\begin{proof}[证明（谓词因果链）]
由 Certificate~\ref{cert:qualitative-breakthrough}，取谓词链：
\[
\begin{aligned}
A:&\ \mathcal{H}_{\mathrm{Q}}\ \text{存在}\Rightarrow \exists\,L\ \text{使 }\Phi(L)=0,\\
B:&\ \Phi(L)=0\Rightarrow \mathfrak{o}_{\mathrm{chart}}=\mathfrak{o}_{\mathrm{lift}}=\mathfrak{o}_{\mathrm{thres}}=0,\\
C:&\ \Phi=0\ \text{的严格可执行域内}\Rightarrow \text{分类降级为 }\pi_0(\Xi_{\mathrm{exec}}).
\end{aligned}
\]
由 $A\Rightarrow B$ 与 $B\Rightarrow C$ 得到 (1)$\Rightarrow$(2)$\Rightarrow$(3)。
反向方向：由 Theorem~\ref{thm:obstruction-zero-implies-finite-descent} 得到 (2)$\Rightarrow$(1)，从而三者等价。证毕。
\end{proof}

\subsection{备注：为什么这不是“把没有与上同调等同”}

\begin{remark}[区分外部叙事与内部对象：避免“现实尚未普遍实现 $\Leftrightarrow$ 障碍类仍在”的误读]
您指出的要点是：不能把“现实尚未普遍实现/尚未普遍体现”当作“上同调障碍仍不为零”的同义句。
上同调障碍类 $\mathfrak{o}_{\mathrm{chart}},\mathfrak{o}_{\mathrm{lift}},\mathfrak{o}_{\mathrm{thres}}$ 是本形式系统内部的对象
（位于 $\check H^2(\cdot)$ 或相对 $\check H^2(\cdot,\Sigma)$ 中），其是否为零只由本节的 descent/证书/同伦动作语义判定；
外部叙事不参与定义。

因此在 Pure 卷语义里，更精确的表述应当是：
\begin{enumerate}
  \item \textbf{理论层面：}本次讨论的增益是构造出一个质变的相对同伦类 $\mathcal{H}_{\mathrm{Q}}$，
  将现实成立的上同调/相对上同调障碍显式化为 $\JacGap$ 并在相对同伦中把它们消去（$\Phi\to 0$）；
  \item \textbf{执行层面：}一旦障碍类消去，现实执行不再是上同调问题，而是可执行空间的拓扑连通/可达性问题，
  分类由 $\pi_0(\Xi_{\mathrm{exec}}(\mathbf{y},r))$ 表达（Theorem~\ref{thm:decategorify-to-pi0}）。
\end{enumerate}
这正是“由障碍类降级为连通类”的严格含义。
\end{remark}

\clearpage
%%%%%%%%%%%%%%%%%%%%%%%%%%%%%%%%%%%%%%%%%%%%%%%%%%%%%%%%%%%%
\section{强结论——工业级 \RTSC{} strict 可达性的非材料决定性与开放对抗治理等价}
\label{sec:puremath-nonmaterial-governance}
%%%%%%%%%%%%%%%%%%%%%%%%%%%%%%%%%%%%%%%%%%%%%%%%%%%%%%%%%%%%

\subsection{重要限定：形式系统内部证明}

\begin{remark}[限定与边界：内部可证 $\neq$ 外部实证]
本节在我们已建立的 Pure 卷 \RTSC{}-工业几何体 Benchmark 内给出一个“强结论”的形式化版本：
\begin{enumerate}
  \item “工业级 \RTSC{}（或极限低阻）strict 可达性”一般不是由材料坐标投影单独决定的性质；
  \item 更不可约的口径是：它等价于一个开放、抗扰动、非交换动态对抗系统的治理问题（存在保持 $\mathcal{E}_{\mathrm{prot}}$ 语义锚点的策略，使 $\Phi\downarrow 0$ 并进入 strict
    sector）；
  \item 一旦 strict sector 可达，则关键难点从“上同调障碍类”降级为“可执行集合的连通类/可达性”（$\pi_0$）。
\end{enumerate}
以上证明完全在形式系统内部进行：证明的是“在该形式系统的定义与证书口径下，结论如何成立”；
它不等同于对自然界是否存在某种室温超导材料的实验性证明。
\end{remark}

\subsection{定义：材料投影与材料决定性}

\begin{definition}[材料投影（遗忘投影）：$p_{\mathrm{mat}}$]
\label{def:p-mat-projection}
在 $\mathrm{Loc}\subseteq\GZLS$ 上，取一个遗忘投影（forgetful projection）
\[
p_{\mathrm{mat}}:\ \mathrm{Loc}\to M_{\mathrm{mat}},
\]
其中 $M_{\mathrm{mat}}$ 仅保留“材料本体”坐标（例如组分、格点/晶格、相、微结构等），并\emph{不}包含环境/几何/控制协议、
仿真离散化/求解器口径、算力配置、测量与证书链等信息。
\end{definition}

\begin{definition}[材料决定性（strict 扇区归属由材料投影决定）]
\label{def:material-determinacy}
称“工业级 \RTSC{} strict 扇区归属”是材料决定的，若存在子集 $S_{\mathrm{mat}}\subseteq M_{\mathrm{mat}}$，
使得对任意 $L\in\mathrm{Loc}$：
\[
L\in \mathrm{Strict}_{\mathrm{RTSC}}
\quad\Longleftrightarrow\quad
p_{\mathrm{mat}}(L)\in S_{\mathrm{mat}}.
\]
\end{definition}

\begin{remark}[“不是材料问题”的形式含义]
Definition~\ref{def:material-determinacy} 把“是否为材料问题”翻译为一个可证伪命题：是否存在 $S_{\mathrm{mat}}$ 使 strict 扇区判定完全由材料投影决定。
其否定即意味着：在同一材料投影下，非材料自由度（环境/几何/阈值治理/数值口径/算力等）可以改变 strict 扇区归属。
\end{remark}

\subsection{定理：工业级 strict 扇区归属一般非材料决定}

\begin{theorem}[非材料决定性定理：存在材料保持但可改变 strict 归属的动作 $\Rightarrow$ 非材料决定]
\label{thm:nonmaterial-determinacy}
若存在一个“非材料动作” $U\in\mathcal{M}_{\mathrm{self}}$ 满足：
\begin{enumerate}
  \item（材料投影保持）$p_{\mathrm{mat}}(U\cdot L)=p_{\mathrm{mat}}(L)$ 对所有 $L\in\mathrm{Loc}$ 成立；
  \item（strict 归属可变）存在某个 $L\in\mathrm{Loc}$ 使
  \[
  L\in \mathrm{Strict}_{\mathrm{RTSC}}
  \quad\text{且}\quad
  U\cdot L\notin \mathrm{Strict}_{\mathrm{RTSC}}
  \qquad\text{或相反。}
  \]
\end{enumerate}
则不存在 $S_{\mathrm{mat}}\subseteq M_{\mathrm{mat}}$ 使 Definition~\ref{def:material-determinacy} 成立。
\end{theorem}

\begin{certificate}[证书：存在材料保持但可改变 strict 判定的纤维动作]
\label{cert:nonmaterial-action-exists}
在 Section~\ref{sec:puremath-governance-generators-counterexample} 的显式反例中，
给出了仅由 $(p_{\mathrm{mat}})$-保持治理生成元生成的治理词 $u$，以及初始法对象 $L_0$，
使得 $p_{\mathrm{mat}}(u\cdot L_0)=p_{\mathrm{mat}}(L_0)$ 且 $\Phi(u\cdot L_0)=0$ 而 $\Phi(L_0)>0$
（见 Theorem~\ref{thm:explicit-noncommutative-counterexample} 的构造与计算）。
取 $U:=u$、$L:=L_0$ 即得到 Theorem~\ref{thm:nonmaterial-determinacy} 所需的存在性见证。
\end{certificate}

\begin{proof}[证明（谓词因果链）]
反证。假设存在 $S_{\mathrm{mat}}\subseteq M_{\mathrm{mat}}$ 使 Definition~\ref{def:material-determinacy} 成立。
由 Theorem~\ref{thm:nonmaterial-determinacy}(1) 得到
\[
p_{\mathrm{mat}}(U\cdot L)=p_{\mathrm{mat}}(L).
\]
于是由材料决定性可得
\[
U\cdot L\in \mathrm{Strict}_{\mathrm{RTSC}}
\Longleftrightarrow
p_{\mathrm{mat}}(U\cdot L)\in S_{\mathrm{mat}}
\Longleftrightarrow
p_{\mathrm{mat}}(L)\in S_{\mathrm{mat}}
\Longleftrightarrow
L\in \mathrm{Strict}_{\mathrm{RTSC}}.
\]
这与 Theorem~\ref{thm:nonmaterial-determinacy}(2)（strict 归属可变）矛盾。
故不存在这样的 $S_{\mathrm{mat}}$。证毕。
\end{proof}

\begin{remark}[为何证书在本 Benchmark 中“天然成立”（对齐边缘算子与三分量证书）]
在 \RTSC{}-工业几何体 Benchmark 中，我们已显式引入包络与边缘算子（Section~\ref{sec:puremath-edge-operators-fem-envelope}）：
\begin{enumerate}
  \item $\Xi^{\mathrm{env}},\Xi^{\mathrm{prec}},\Xi^{\mathrm{comp}}$ 等自由度可在\emph{不改变}材料本体坐标的前提下改变
  $c_{\mathrm{chart}},c_{\mathrm{lift}},c_{\mathrm{thres}}$ 的可下降性与可消去性；
  \item strict sector 的定义包含 $\Phi=0$ 以及证书/阈值口径（Corollary~\ref{cor:qualitative-breakthrough}），因此“保持材料投影但改变非材料口径/阈值治理”可以改变
    strict 归属并不违背形式系统设定。
\end{enumerate}
因此将 Theorem~\ref{thm:nonmaterial-determinacy} 的前提作为证书约定与本 Benchmark 的结构是一致的。
\end{remark}

\subsection{定义：抗扰动治理策略（非交换算子词）}

\begin{definition}[治理策略：从观测到非交换动作词]
\label{def:governance-policy}
取观测映射 $\mathrm{Obs}:\mathrm{Loc}\to\mathcal{O}$（Definition~\ref{def:noncommutative-policy}），并取在线策略
\[
\pi:\ \mathcal{O}\to \mathcal{M}_{\mathrm{self}}.
\]
在开放对抗法流（Definition~\ref{def:open-system-adversarial-law-flow}）下，给定允许的对抗输入族 $\mathcal{D}_{\mathrm{adm}}$，
称 $\pi$ 为一条“抗扰动治理策略”，若对任意对抗输入序列 $d(\cdot)\in\mathcal{D}_{\mathrm{adm}}$，
策略诱导的自行动作序列 $u_n=\pi(\mathrm{Obs}(L_n))$ 能在相对 $\mathcal{E}_{\mathrm{prot}}$ 约束下推动系统进入 strict sector：
存在 $N$ 使 $L_N\in \mathrm{Strict}_{\mathrm{RTSC}}$。
\end{definition}

\subsection{定理：strict 可达性 $\Longleftrightarrow$ 存在治理策略}

\begin{theorem}[治理等价定理：strict 可达 $\Longleftrightarrow$ 存在抗扰动治理策略（相对 $\mathcal{E}_{\mathrm{prot}}$）]
\label{thm:governance-equivalence}
在本 Benchmark 中，以下两种表述等价：
\begin{enumerate}
  \item（strict 可达的动作链表述）存在初始 $L_0\in\mathrm{Loc}$ 与一条在相对约束下的有限动作链，
  使得在允许对抗输入族 $\mathcal{D}_{\mathrm{adm}}$ 下可达 $L_\ast\in \mathrm{Strict}_{\mathrm{RTSC}}$；
  \item（治理策略表述）存在抗扰动治理策略 $\pi$（Definition~\ref{def:governance-policy}），使得对任意 $d(\cdot)\in\mathcal{D}_{\mathrm{adm}}$，
  都存在 $N$ 使 $L_N\in \mathrm{Strict}_{\mathrm{RTSC}}$。
\end{enumerate}
\end{theorem}

\begin{certificate}[证书：动作链 $\leftrightarrow$ 策略的编译与展开]
\label{cert:policy-compilation}
在证书口径下，约定：
\begin{enumerate}
  \item 任意有限动作链都可视为一个（可能忽略观测的）策略实例：在 $n=0,\dots,N-1$ 输出预定 $u_n$；
  \item 任意策略 $\pi$ 沿任意给定对抗输入序列 $d(\cdot)$ 的执行，都可展开为一条有限动作链（取首次进入 $\mathrm{Strict}_{\mathrm{RTSC}}$ 的步数作为终止时刻）。
\end{enumerate}
该证书只涉及“策略生成非交换算子词”的语义展开（Definition~\ref{def:noncommutative-policy}），不增加新的物理假设。
\end{certificate}

\begin{proof}[证明（谓词因果链）]
由 Certificate~\ref{cert:policy-compilation}，取谓词链：
\[
\begin{aligned}
A:&\ \text{(1) 的有限动作链}\Rightarrow \exists\,\pi\ \text{可逐步输出该链（忽略观测亦可）},\\
B:&\ A\Rightarrow \text{(2) 的治理策略存在},\\
C:&\ \text{(2) 的治理策略}\Rightarrow \forall d(\cdot)\ \text{其执行可展开为到达 strict sector 的有限链}.
\end{aligned}
\]
由 $A\Rightarrow B$ 得到 (1)$\Rightarrow$(2)；由 $C$ 得到 (2)$\Rightarrow$(1)。证毕。
\end{proof}

\subsection{推论：治理内容不可约（至少依赖非材料自由度）}

\begin{corollary}[治理不可约推论：严格可达性不能压缩为纯材料谓词]
\label{cor:governance-irreducible}
若 Theorem~\ref{thm:nonmaterial-determinacy} 的前提成立，则任意对 strict 可达性的完备描述必须包含至少一类非材料自由度（环境/几何/阈值治理/数值口径/算力等）所生成的治理算子；
因此把问题描述为“纯材料问题”在形式系统内是不完备的，更不可约的口径是“开放、抗扰动、非交换动态对抗系统的治理问题”。
\end{corollary}

\begin{certificate}[证书：非材料决定性迫使治理自由度进入接口]
\label{cert:governance-irreducible}
由 Theorem~\ref{thm:nonmaterial-determinacy}，strict 判定不可能由材料投影 $p_{\mathrm{mat}}$ 单独决定；
因此若要对 strict 可达性给出\emph{完备}的形式系统描述，必须显式保留至少一类非材料自由度及其生成的治理算子族
（环境/几何/阈值治理/数值口径/算力等），否则描述将缺失必要的判定信息。
\end{certificate}

\begin{proof}[证明（谓词因果链）]
取谓词链：
\[
\begin{aligned}
A:&\ \text{Theorem~\ref{thm:nonmaterial-determinacy}}\Rightarrow \text{strict 判定不由 }p_{\mathrm{mat}}\text{ 单独决定},\\
B:&\ A\Rightarrow \text{若忽略所有非材料治理算子，则 strict 判定无法完备描述},\\
C:&\ B\Rightarrow \text{必须在 }\mathcal{M}_{\mathrm{self}}\text{ 中保留非材料自由度对应的生成元（治理内容不可约）}.
\end{aligned}
\]
由 $A\wedge B\wedge C$ 得到推论。证毕。
\end{proof}

\subsection{定理：里程碑意义——从障碍论到连通论的可执行降级}

\begin{theorem}[里程碑意义定理：strict 可达 $\Rightarrow$ 障碍消去且执行降级为 $\pi_0$]
\label{thm:milestone-significance}
若存在抗扰动治理策略使系统进入 strict sector（Theorem~\ref{thm:governance-equivalence}），则：
\begin{enumerate}
  \item 三类上同调障碍类在形式系统内同时消失：
  \[
  \mathfrak{o}_{\mathrm{chart}}=\mathfrak{o}_{\mathrm{lift}}=\mathfrak{o}_{\mathrm{thres}}=0;
  \]
  \item 在 strict sector 内，执行分类降级为可执行包络集合的连通类：
  \[
  \text{执行分类}\ \cong\ \pi_0\bigl(\Xi_{\mathrm{exec}}(\mathbf{y},r)\bigr)
  \quad(\text{见 Definition~\ref{def:xi-exec} 与 Theorem~\ref{thm:decategorify-to-pi0}}).
  \]
\end{enumerate}
\end{theorem}

\begin{certificate}[证书：由障碍消去与执行降级定理链直接推出]
\label{cert:milestone-significance}
\begin{enumerate}
  \item 由 Theorem~\ref{thm:governance-equivalence}，strict 可达意味着存在有限同伦动作/包络更新使 $\Phi\downarrow 0$ 并进入 strict sector；
  \item 由 Theorem~\ref{thm:obstruction-vanishing}，$\Phi\to 0$ 蕴含 $\mathfrak{o}_{\mathrm{chart}}=\mathfrak{o}
    _{\mathrm{lift}}=\mathfrak{o}_{\mathrm{thres}}=0$；
  \item 由 Theorem~\ref{thm:decategorify-to-pi0}，在 strict sector 内执行分类由 $\pi_0(\Xi_{\mathrm{exec}})$ 给出。
\end{enumerate}
\end{certificate}

\begin{proof}[证明（谓词因果链）]
由 Certificate~\ref{cert:milestone-significance}，取谓词链：
\[
\begin{aligned}
A:&\ \text{strict 可达}\Rightarrow \Phi\to 0,\\
B:&\ \Phi\to 0\Rightarrow \mathfrak{o}_{\mathrm{chart}}=\mathfrak{o}_{\mathrm{lift}}=\mathfrak{o}_{\mathrm{thres}}=0,\\
C:&\ \Phi=0\ \text{的 strict 扇区内}\Rightarrow \text{执行分类降级为 }\pi_0(\Xi_{\mathrm{exec}}).
\end{aligned}
\]
由 $A\Rightarrow B$ 与 $B\Rightarrow C$ 得到结论两点。证毕。
\end{proof}

\begin{remark}[“里程碑意义”的形式化内核：问题类降级]
Theorem~\ref{thm:milestone-significance} 把“里程碑意义”严格化为一个\emph{问题类的降级}：
\begin{enumerate}
  \item strict 不可达时，关键困难由上同调/相对上同调障碍类支配（$\check H^2$ 层面）；
  \item 一旦 strict 可达并使 $\Phi=0$，障碍论层面的困难消失，剩余执行变为“在 $\Xi_{\mathrm{exec}}$ 中选连通分支并沿路径实现”，即 $\pi_0$ 层面问题。
\end{enumerate}
这解释了为何该结论与“是否已发现某个具体材料”无关：它讨论的是一旦 strict sector 在任何候选体系上可达，系统工程难点将如何结构性降级。
\end{remark}

\subsection{结论：一句 Pure 卷式三句强陈述}

\begin{corollary}[三句强陈述（Pure 卷式）]
\label{cor:strong-conclusion}
在 \RTSC{}-工业几何体 Benchmark 的形式系统内部，可将强结论压成三句：
\[
\boxed{
\begin{aligned}
&\text{(i) strict 可达性一般不是材料投影 }p_{\mathrm{mat}}\text{ 可决定的性质；}\\
&\text{(ii) strict 可达 }\equiv\ \text{存在相对 }\mathcal{E}_{\mathrm{prot}}\text{ 的抗扰动治理策略，使 }\Phi\downarrow 0;\\
&\text{(iii) 一旦 }\Phi=0,\ \text{障碍类消去且执行分类降级为 }\pi_0(\Xi_{\mathrm{exec}}).
\end{aligned}}
\]
\end{corollary}

\begin{certificate}[证书：三句分别由三条主定理链压缩得到]
\label{cert:strong-conclusion}
三句强陈述分别是三条主结论链的“接口级压缩”：
\begin{enumerate}
  \item (i) 由 Theorem~\ref{thm:nonmaterial-determinacy} 给出非材料决定性；
  \item (ii) 由 Theorem~\ref{thm:governance-equivalence} 给出 strict 可达与抗扰动治理策略存在的等价；
  \item (iii) 由 Theorem~\ref{thm:milestone-significance} 给出 strict 进入后的障碍消去与执行降级（$\pi_0$）。
\end{enumerate}
\end{certificate}

\begin{proof}[证明（谓词因果链）]
取谓词链：
\[
\begin{aligned}
A:&\ \text{Theorem~\ref{thm:nonmaterial-determinacy}}\Rightarrow \text{(i)},\\
B:&\ \text{Theorem~\ref{thm:governance-equivalence}}\Rightarrow \text{(ii)},\\
C:&\ \text{Theorem~\ref{thm:milestone-significance}}\Rightarrow \text{(iii)}.
\end{aligned}
\]
由 $A\wedge B\wedge C$ 得到三句强陈述。证毕。
\end{proof}

\clearpage
%%%%%%%%%%%%%%%%%%%%%%%%%%%%%%%%%%%%%%%%%%%%%%%%%%%%%%%%%%%%
\section{治理算子最小生成元表与显式非交换反例（同一材料投影下的 $\Phi$ 分离）}
\label{sec:puremath-governance-generators-counterexample}
%%%%%%%%%%%%%%%%%%%%%%%%%%%%%%%%%%%%%%%%%%%%%%%%%%%%%%%%%%%%

\subsection{补齐项：最小生成元表与显式反例}

\begin{remark}[补齐两件事：生成元表 + 构造性反例]
本节补齐两件可在 Pure 卷正文中直接引用的内容：
\begin{enumerate}
  \item 给出治理算子最小生成元表，并严格区分哪些生成元保持材料投影 $p_{\mathrm{mat}}$ 不变、哪些会改变 $p_{\mathrm{mat}}$；
  \item 给出一个完全显式的非交换反例：在同一材料投影下，存在两条由 $(p_{\mathrm{mat}})$-保持治理生成元生成的非交换策略词，使 $\Phi$ 分别取 $0$ 与 $>0$。
\end{enumerate}
重要限定：本节是对已固定 Benchmark（$\GZLS,\mathrm{Loc},\Sigma,\mathcal{E}_{\mathrm{prot}},\JacGap,\Phi$）的\emph{形式系统内部}证明；
不替代任何实验性论证。
\end{remark}

\subsection{定义：治理算子字母表与 $(p_{\mathrm{mat}})$-保持/改变划分}

\begin{definition}[治理算子字母表与算子词]
\label{def:governance-operator-alphabet}
在我方可执行算子幺半群 $\mathcal{M}_{\mathrm{self}}$（Definition~\ref{def:noncommutative-operator-systems}）中，
取一组生成元（字母表）$\mathcal{G}_{\mathrm{self}}$ 使
\[
\mathcal{M}_{\mathrm{self}}=\langle \mathcal{G}_{\mathrm{self}}\rangle,
\qquad
u=u_1u_2\cdots u_n\in \mathcal{M}_{\mathrm{self}}^\ast,\ \ u_i\in\mathcal{G}_{\mathrm{self}}.
\]
这里算子词的乘法为串接；一般情况下不交换，反映“顺序改变结果”的工程事实。
\end{definition}

\begin{definition}[$(p_{\mathrm{mat}})$-保持/改变生成元]
\label{def:p-mat-preserve-change-generators}
给定材料投影 $p_{\mathrm{mat}}$（Definition~\ref{def:p-mat-projection}），
称生成元 $U\in\mathcal{G}_{\mathrm{self}}$ 为：
\begin{enumerate}
  \item \textbf{$(p_{\mathrm{mat}})$-保持}，若对任意 $L\in\mathrm{Loc}$ 有 $p_{\mathrm{mat}}(U\cdot L)=p_{\mathrm{mat}}(L)$；
  \item \textbf{$(p_{\mathrm{mat}})$-改变}，若存在 $L\in\mathrm{Loc}$ 使 $p_{\mathrm{mat}}(U\cdot L)\neq p_{\mathrm{mat}}(L)$。
\end{enumerate}
记两类生成元集合分别为
\[
\mathcal{G}_{\mathrm{self}}^{\mathrm{pres}}
:=
\{U\in\mathcal{G}_{\mathrm{self}}\mid p_{\mathrm{mat}}(U\cdot L)=p_{\mathrm{mat}}(L)\ \forall L\},
\qquad
\mathcal{G}_{\mathrm{self}}^{\Delta}
:=
\mathcal{G}_{\mathrm{self}}\setminus \mathcal{G}_{\mathrm{self}}^{\mathrm{pres}}.
\]
\end{definition}

\subsection{定义：治理算子最小生成元表（可直接引用）}

\begin{definition}[最小生成元表：按 $p_{\mathrm{mat}}$ 是否保持划分]
\label{def:governance-generators-minimal-table}
在 \RTSC{}-工业几何体 Benchmark 的包络语义（Definition~\ref{def:rtsc-ig-section-envelope}、Section~\ref{sec:puremath-edge-operators-fem-envelope}）下，
可取如下“最小治理生成元表”（示意，但足以承载三分量证书的可控下降机制）。

\textbf{（A）治理生成元：保持 $p_{\mathrm{mat}}$ 不变（非材料自由度）}
\begin{center}
\footnotesize
\renewcommand{\arraystretch}{1.15}
\begin{tabular}{|p{0.18\linewidth}|p{0.36\linewidth}|p{0.10\linewidth}|p{0.26\linewidth}|}
\hline
生成元 & 作用域（包络坐标/对象） & $p_{\mathrm{mat}}$ & 主要影响分量（示意）\\
\hline
$\mathsf{E}^{\mathrm{env}}_{e}$ & $\Xi^{\mathrm{env}}_{\mathrm{edge}}$（边界/外场/接地/屏蔽等） & 保持 & $c_{\mathrm{thres}},
  c_{\mathrm{lift}}$\\
\hline
$\mathsf{E}^{\mathrm{geom}}$ & $\Xi^{\mathrm{Geom}}$（几何形状/拓扑/连接方式等） & 保持 & $c_{\mathrm{chart}},c_{\mathrm{thres}}$\\
\hline
$\mathsf{E}^{\mathrm{prec}}_{e}$ & $\Xi^{\mathrm{prec}}_{\mathrm{edge}}$（精度/滞回/采样/抖动抑制等） & 保持 & $c_{\mathrm{lift}},
  c_{\mathrm{thres}}$\\
\hline
$\mathsf{E}^{\mathrm{comp}}_{e}$ & $\Xi^{\mathrm{comp}}_{\mathrm{edge}}$（网格/容差/多保真/算力分配等） & 保持 & $c_{\mathrm{chart}},
  c_{\mathrm{lift}}$\\
\hline
$\mathsf{E}^{\mathrm{ctrl}}$ & $\Xi^{\mathrm{Ctrl}}$（阈值调度/触发图式/恢复图式） & 保持 & $c_{\mathrm{thres}}$\\
\hline
$\mathsf{E}^{\mathrm{cert}}$ & $\Xi^{\mathrm{Sim}}\times\Xi^{\mathrm{meas}}$（证书口径/残差对齐/测量口径对齐） & 保持 & $c_{\mathrm{chart}}
  $\\
\hline
\end{tabular}
\end{center}

\textbf{（B）材料生成元：改变 $p_{\mathrm{mat}}$（材料身份/长期等价类）}
\begin{center}
\footnotesize
\renewcommand{\arraystretch}{1.15}
\begin{tabular}{|p{0.22\linewidth}|p{0.56\linewidth}|p{0.12\linewidth}|}
\hline
生成元 & 作用域（示意） & $p_{\mathrm{mat}}$\\
\hline
$\mathsf{E}^{\Delta\mathrm{comp}}$ & 组分/掺杂/配比改变（进入不同材料身份） & 改变\\
\hline
$\mathsf{E}^{\Delta\mathrm{phase}}$ & 相结构/晶系族切换（进入不同材料身份） & 改变\\
\hline
$\mathsf{E}^{\Delta\mathrm{defect}}$ & 不可逆缺陷谱迁移/长期等价类改变 & 改变\\
\hline
\end{tabular}
\end{center}

上述划分用于将“材料范畴内的搜索”（B 类）与“治理范畴内的策略合成”（A 类）在形式系统内严格分离。
\end{definition}

\begin{remark}[最小性口径（结构，而非叙事）]
在本 Benchmark 中，$\Phi=\sum_{\bullet}w_{\bullet}\|c_{\bullet}\|$ 的三分量分别刻画 chart/lift/threshold 三类缺陷；
因此“最小治理生成元表”的口径是：\emph{至少提供能够影响每个缺陷分量的 $(p_{\mathrm{mat}})$-保持生成元}，
以保证 $\Phi$ 在治理词空间内存在下降方向（否则将出现某分量对所有治理词不变的刚性障碍）。
\end{remark}

\subsection{定理（公理降级）：缺陷分量的必要覆盖与词不变性}

\begin{definition}[影响谓词：$\mathrm{Affect}_{\bullet}(U)$]
\label{def:affect-predicate}
对 $\bullet\in\{\mathrm{chart},\mathrm{lift},\mathrm{thres}\}$ 与生成元 $U\in\mathcal{G}_{\mathrm{self}}$，
定义影响谓词：
\[
\mathrm{Affect}_{\bullet}(U)
\quad:\Longleftrightarrow\quad
\exists\,L\in\mathrm{Loc}\ \text{使}\ c_{\bullet}(U\cdot L)\neq c_{\bullet}(L).
\]
\end{definition}

\begin{theorem}[必要覆盖定理（公理降级）：若对某分量无影响生成元，则该分量对任意治理词不变]
\label{thm:component-invariance-under-words}
固定 $\bullet\in\{\mathrm{chart},\mathrm{lift},\mathrm{thres}\}$。
若对所有生成元 $U\in\mathcal{G}_{\mathrm{self}}$ 都有 $\neg\mathrm{Affect}_{\bullet}(U)$，
则对任意治理词 $u\in\langle \mathcal{G}_{\mathrm{self}}\rangle^\ast$ 与任意 $L\in\mathrm{Loc}$，
\[
c_{\bullet}(u\cdot L)=c_{\bullet}(L).
\]
特别地，若存在 $L_0$ 使 $c_{\bullet}(L_0)>0$，则对任意治理词 $u$ 都不可能使 $\Phi(u\cdot L_0)=0$。
\end{theorem}

\begin{certificate}[数学归纳法证书：按词长 $n$ 归纳]
\label{cert:component-invariance-induction}
令谓词 $P(n)$ 为：对任意长度为 $n$ 的治理词 $u=u_1\cdots u_n$（$u_i\in\mathcal{G}_{\mathrm{self}}$）与任意 $L\in\mathrm{Loc}$，
有 $c_{\bullet}(u\cdot L)=c_{\bullet}(L)$。
证书化要求：
\[
P(0)\ \text{成立},\qquad
(\forall n\in\mathbb{N})\ \bigl(P(n)\Rightarrow P(n+1)\bigr).
\]
\end{certificate}

\begin{proof}[证明（数学归纳法证书；谓词因果链）]
由 Certificate~\ref{cert:component-invariance-induction}，取谓词链：
\[
\begin{aligned}
A:&\ \neg\mathrm{Affect}_{\bullet}(U)\ \forall U\in\mathcal{G}_{\mathrm{self}}
\Rightarrow c_{\bullet}(U\cdot L)=c_{\bullet}(L)\ \forall U,L,\\
B:&\ P(n)\wedge A \Rightarrow P(n+1)\quad(\text{对 }u=u'\,U\text{ 分解并应用 }A),\\
C:&\ P(0)\ \text{平凡成立}.
\end{aligned}
\]
由 $C$ 得基步，由 $B$ 得归纳步，从而对任意词长 $n$ 成立；即对任意治理词 $u$ 有 $c_{\bullet}(u\cdot L)=c_{\bullet}(L)$。
若 $c_{\bullet}(L_0)>0$，则 $\Phi(u\cdot L_0)\ge w_{\bullet}\,c_{\bullet}(L_0)>0$，不可能为零。证毕。
\end{proof}

\subsection{显式反例：同一 $p_{\mathrm{mat}}$ 下的非交换治理词导致 $\Phi$ 分离}

\begin{definition}[toy 局部包络坐标（仅保留构造所需最小变量）]
\label{def:toy-envelope-variables}
取一个局部包络坐标化 $E\simeq U\times\Xi$，并在纤维 $\Xi$ 中取最小坐标
\[
\xi=(b,g,\tau,h,\theta)\in\mathbb{R}^5,
\]
其语义（抽象）为：
\[
\begin{aligned}
b:&\ \text{环境/边界综合量（外场--屏蔽--接地等低阶等效量）},\\
g:&\ \text{几何结构综合量（几何--拓扑的低阶等效量）},\\
\tau:&\ \text{仿真/计算口径综合量（网格/容差/保真等）},\\
h:&\ \text{控制精度/滞回带综合量},\\
\theta:&\ \text{阈值调度综合量（触发/恢复门槛等）}.
\end{aligned}
\]
并约定该局部构造固定材料身份：在该局部域内 $p_{\mathrm{mat}}$ 为常量，
即材料投影与 $(b,g,\tau,h,\theta)$ 无关。
\end{definition}

\begin{definition}[toy 三分量证书与缺陷势能 $\Phi$]
\label{def:toy-defect-potential}
在 Definition~\ref{def:toy-envelope-variables} 的局部坐标下，定义最小三分量证书为
\[
\begin{aligned}
c_{\mathrm{chart}}&:=|\tau-gb|,\\
c_{\mathrm{lift}}&:=|h-b|,\\
c_{\mathrm{thres}}&:=|\theta-b|,
\end{aligned}
\]
并取权重为 $1$ 的 defect potential：
\[
\Phi:=c_{\mathrm{chart}}+c_{\mathrm{lift}}+c_{\mathrm{thres}}.
\]
\end{definition}

\begin{definition}[toy 治理生成元（均保持 $p_{\mathrm{mat}}$）]
\label{def:toy-governance-operators}
定义 5 个治理算子（全为 $(p_{\mathrm{mat}})$-保持动作）在 $\xi$ 上的作用：
\[
\begin{aligned}
U_{\mathrm{env}}(b,g,\tau,h,\theta)&:=(1,g,\tau,h,\theta),\\
U_{\mathrm{prec}}(b,g,\tau,h,\theta)&:=(b,g,\tau,b,\theta),\\
U_{\mathrm{ctrl}}(b,g,\tau,h,\theta)&:=(b,g,\tau,h,h),\\
U_{\mathrm{geom}}(b,g,\tau,h,\theta)&:=\left(b,\frac{\tau}{b},\tau,h,\theta\right)\quad(b\neq 0),\\
U_{\mathrm{comp}}(b,g,\tau,h,\theta)&:=(b,g,gb,h,\theta).
\end{aligned}
\]
这些算子可分别视为 $\mathsf{E}^{\mathrm{env}},\mathsf{E}^{\mathrm{prec}},\mathsf{E}^{\mathrm{ctrl}},\mathsf{E}^{\mathrm{geom}},
  \mathsf{E}^{\mathrm{comp}}$
在局部坐标上的一个最小抽象实现。
\end{definition}

\begin{lemma}[最小非交换见证：同一生成元集合的不同顺序可导致不同结果]
\label{lem:toy-noncommutative-witness}
在上述 toy 模型中，存在初值 $\xi_0$ 使得
\[
U_{\mathrm{env}}\circ U_{\mathrm{geom}}(\xi_0)\neq U_{\mathrm{geom}}\circ U_{\mathrm{env}}(\xi_0),
\]
从而治理生成元一般不交换。
\end{lemma}

\begin{certificate}[证书：给出一组初值并直接计算得到不交换]
\label{cert:toy-noncommutative-witness}
取 $\xi_0=(2,1,3,0,0)$，分别计算 $U_{\mathrm{env}}\circ U_{\mathrm{geom}}(\xi_0)$ 与 $U_{\mathrm{geom}}\circ U_{\mathrm{env}}
  (\xi_0)$，
两者给出不同结果，从而提供最小非交换见证。
\end{certificate}

\begin{proof}[证明（构造性计算）]
取 $\xi_0=(2,1,3,0,0)$。
则
\[
U_{\mathrm{env}}\circ U_{\mathrm{geom}}(\xi_0)
=
U_{\mathrm{env}}\left(2,\frac{3}{2},3,0,0\right)
=
\left(1,\frac{3}{2},3,0,0\right),
\]
而
\[
U_{\mathrm{geom}}\circ U_{\mathrm{env}}(\xi_0)
=
U_{\mathrm{geom}}(1,1,3,0,0)
=
(1,3,3,0,0),
\]
两者不同。证毕。
\end{proof}

\begin{theorem}[显式反例定理：同一材料投影下存在两条非交换治理词使 $\Phi$ 一为零一为正]
\label{thm:explicit-noncommutative-counterexample}
在 Definition~\ref{def:toy-envelope-variables}--Definition~\ref{def:toy-governance-operators} 的局部 toy 模型中，
取初始状态 $\xi_0=(2,1,3,0,0)$，并定义两条治理词（同一生成元集合，不同顺序）：
\[
u:=U_{\mathrm{comp}}\circ U_{\mathrm{geom}}\circ U_{\mathrm{ctrl}}\circ U_{\mathrm{prec}}\circ U_{\mathrm{env}},
\qquad
u':=U_{\mathrm{env}}\circ U_{\mathrm{ctrl}}\circ U_{\mathrm{prec}}\circ U_{\mathrm{comp}}\circ U_{\mathrm{geom}}.
\]
则：
\begin{enumerate}
  \item $p_{\mathrm{mat}}(u\cdot L_0)=p_{\mathrm{mat}}(u'\cdot L_0)=p_{\mathrm{mat}}(L_0)$（同一材料投影）；
  \item $\Phi(u\cdot L_0)=0$ 而 $\Phi(u'\cdot L_0)>0$。
\end{enumerate}
\end{theorem}

\begin{certificate}[证书：逐步计算得到 $\Phi$ 分离]
\label{cert:explicit-counterexample}
在该 toy 模型中：
\[
u\cdot \xi_0=(1,3,3,1,1)\Rightarrow \Phi(u\cdot L_0)=0,
\qquad
u'\cdot \xi_0=\left(1,\frac{3}{2},3,2,2\right)\Rightarrow \Phi(u'\cdot L_0)=\frac{7}{2}.
\]
其中 $L_0$ 表示与 $\xi_0$ 对应的法对象。
\end{certificate}

\begin{proof}[证明（谓词因果链；显式计算）]
由 Definition~\ref{def:toy-envelope-variables}，材料投影与 $\xi$ 无关，故 (1) 成立。
对 (2)，按 Certificate~\ref{cert:explicit-counterexample} 逐步计算：
\[
\begin{aligned}
U_{\mathrm{env}}(2,1,3,0,0)&=(1,1,3,0,0),\\
U_{\mathrm{prec}}(1,1,3,0,0)&=(1,1,3,1,0),\\
U_{\mathrm{ctrl}}(1,1,3,1,0)&=(1,1,3,1,1),\\
U_{\mathrm{geom}}(1,1,3,1,1)&=(1,3,3,1,1),\\
U_{\mathrm{comp}}(1,3,3,1,1)&=(1,3,3,1,1),
\end{aligned}
\]
从而 $u\cdot \xi_0=(1,3,3,1,1)$，代入 Definition~\ref{def:toy-defect-potential} 得 $\Phi(u\cdot L_0)=0$。
另一方面：
\[
\begin{aligned}
U_{\mathrm{geom}}(2,1,3,0,0)&=\left(2,\frac{3}{2},3,0,0\right),\\
U_{\mathrm{comp}}\left(2,\frac{3}{2},3,0,0\right)&=\left(2,\frac{3}{2},3,0,0\right),\\
U_{\mathrm{prec}}\left(2,\frac{3}{2},3,0,0\right)&=\left(2,\frac{3}{2},3,2,0\right),\\
U_{\mathrm{ctrl}}\left(2,\frac{3}{2},3,2,0\right)&=\left(2,\frac{3}{2},3,2,2\right),\\
U_{\mathrm{env}}\left(2,\frac{3}{2},3,2,2\right)&=\left(1,\frac{3}{2},3,2,2\right),
\end{aligned}
\]
从而 $u'\cdot \xi_0=\left(1,\frac{3}{2},3,2,2\right)$，代入 Definition~\ref{def:toy-defect-potential} 得
\[
\Phi(u'\cdot L_0)=\left|3-\frac{3}{2}\cdot 1\right|+|2-1|+|2-1|=\frac{3}{2}+1+1=\frac{7}{2}>0.
\]
故 $\Phi$ 在同一材料投影下被两条非交换治理词分离。证毕。
\end{proof}

\begin{corollary}[构造性否证：strict 可达性非材料决定]
\label{cor:explicit-nonmaterial-determinacy}
在上述 toy 构造中，同一材料投影下存在 $L_1,L_2\in\mathrm{Loc}$ 使 $p_{\mathrm{mat}}(L_1)=p_{\mathrm{mat}}(L_2)$ 且
$\Phi(L_1)=0$、$\Phi(L_2)>0$。
因此不存在 $S_{\mathrm{mat}}\subseteq M_{\mathrm{mat}}$ 使 strict 扇区归属由材料投影决定（Definition~\ref{def:material-determinacy}）。
\end{corollary}

\begin{certificate}[证书：由显式非交换反例给出同一材料投影下的 $\Phi$ 分离]
\label{cert:explicit-nonmaterial-determinacy}
由 Theorem~\ref{thm:explicit-noncommutative-counterexample}，同一材料投影下可由两条非交换治理词得到 $\Phi=0$ 与 $\Phi>0$ 的分离样本；
由于 strict 扇区以 $\Phi=0$ 为门槛（Definition~\ref{def:markdown-strict-sector}），strict 判定因此不可能仅由材料投影决定。
\end{certificate}

\begin{proof}[证明（谓词因果链）]
取 $L_1:=u\cdot L_0$、$L_2:=u'\cdot L_0$。
由 Theorem~\ref{thm:explicit-noncommutative-counterexample} 得到
\[
p_{\mathrm{mat}}(L_1)=p_{\mathrm{mat}}(L_2),\qquad \Phi(L_1)=0,\qquad \Phi(L_2)>0.
\]
若存在 $S_{\mathrm{mat}}$ 使 Definition~\ref{def:material-determinacy} 成立，则同一 $p_{\mathrm{mat}}$ 下 strict 判定必须一致；
但 $\Phi$ 的分离与 strict 扇区的 $\Phi=0$ 门槛相矛盾。故不存在该 $S_{\mathrm{mat}}$。证毕。
\end{proof}

\begin{remark}[从 toy 反例到真实管线接口：一条可复用的证明模板]
toy 构造的作用不是替代真实物理，而是提供一条可复用的“形式系统证明模板”：
\begin{enumerate}
  \item 把 $(b,g,\tau,h,\theta)$ 映射到真实系统中的受保护评价与阈值图式接口（环境等效量、几何参数族、FEM 网格/容差、多保真切换、滞回与触发阈值等）；
  \item 把 toy 的 $\Phi=|\tau-gb|+|h-b|+|\theta-b|$ 替换为真实三分量证书范数 $\|c_{\mathrm{chart}}\|+\|c_{\mathrm{lift}}
    \|+\|c_{\mathrm{thres}}\|$；
  \item 保持同一证明结构：在同一 $p_{\mathrm{mat}}$ 下，通过两条非交换治理词使 $\Phi$ 取不同值，从而将“非材料决定性”落到可审计、可复现的证书计算上。
\end{enumerate}
\end{remark}

\clearpage
%%%%%%%%%%%%%%%%%%%%%%%%%%%%%%%%%%%%%%%%%%%%%%%%%%%%%%%%%%%%
\section{\RTSC{}-工业几何体 Benchmark 的“非材料决定性”与“治理问题”等价（Markdown 口径）}
\label{sec:puremath-markdown-nonmaterial-governance}
%%%%%%%%%%%%%%%%%%%%%%%%%%%%%%%%%%%%%%%%%%%%%%%%%%%%%%%%%%%%

\subsection{范围声明（形式系统内）}

\begin{remark}[范围声明（形式系统内）]
以下所有结论均在我们已建立的 Pure 卷形式系统
$(\GZLS,\mathrm{Loc},\Sigma,\GZLOC,\GZTLC,\GZLCurv,\mathcal{E}_{\mathrm{prot}},\JacGap,\Phi)$
内部严格成立。
它证明的是：\textbf{工业级 \RTSC{}/极限低阻 strict 可达性}在该形式系统内不可能由“材料投影”单独决定，
而等价于\textbf{开放、抗扰动、非交换动态对抗治理}的策略合成问题；
并且一旦 strict 进入，现实执行可降级为可执行集合的连通类（$\pi_0$）问题。
\end{remark}

\subsection{预备：法空间与可判定语义}

固定 Pure 卷法空间与数据：
\[
\GZLS=(\mathfrak{L};J,\mathrm{Loc},\Sigma),\qquad \GZLOC,\ \GZTLC,\ \GZLCurv.
\]
固定被保护评价泛函（protected observables）：
\[
\mathcal{E}_{\mathrm{prot}}=\{\mathrm{Ev}_a\}_{a=1}^{q},\qquad \mathrm{Ev}_a:\mathrm{Loc}\to\mathbb{R},
\qquad
\mathrm{Ev}:=(\mathrm{Ev}_1,\dots,\mathrm{Ev}_q):\mathrm{Loc}\to\mathbb{R}^q.
\]
由阈值集合 $\Sigma$ 诱导分层标签：
\[
\rho_\Sigma:\mathrm{Loc}\to \mathcal{R}_\Sigma.
\]

\subsection{定义：工业对象的 law-object 分解与材料投影 $p_{\mathrm{mat}}$}

\begin{definition}[工业 \RTSC{}-对象的 law-object 分解（schematic）]
\label{def:markdown-law-object-decomposition}
对“工业低阻法拉第笼 + 旁路低阻/电流熔断（含触发/隔离/恢复）”类对象，取一个可操作分解：
\[
L=\bigl(L_{\mathrm{mat}},L_{\mathrm{geom}},L_{\mathrm{env}},L_{\mathrm{sim}},L_{\mathrm{ctrl}},L_{\mathrm{cert}}\bigr)
  \in \mathrm{Loc}\subset\GZLS.
\]
其中：
\begin{enumerate}
  \item $L_{\mathrm{mat}}$：材料身份/材料等价类（组分/相族/晶格族/缺陷谱等）；
  \item $L_{\mathrm{geom}}$：几何结构（法拉第笼/旁路路径的几何--拓扑）；
  \item $L_{\mathrm{env}}$：环境与边界（外场、接地/回流网络、工况、扰动注入）；
  \item $L_{\mathrm{sim}}$：仿真表示（FEM 网格、离散化、容差、多保真耦合口径）；
  \item $L_{\mathrm{ctrl}}$：实时控制与阈值调度（触发、熔断/限流、隔离、恢复、滞回）；
  \item $L_{\mathrm{cert}}$：测量与证书口径（确保 $\mathcal{E}_{\mathrm{prot}}$ 语义不被偷换）。
\end{enumerate}
\end{definition}

\begin{definition}[材料投影（forgetful projection）]
\label{def:markdown-p-mat}
定义材料投影：
\[
p_{\mathrm{mat}}:\mathrm{Loc}\to M_{\mathrm{mat}},
\qquad
p_{\mathrm{mat}}(L):=L_{\mathrm{mat}}.
\]
\end{definition}

\subsection{定义：strict 扇区、$\JacGap$ 与 defect potential}

\begin{definition}[最小三分量证书与缺陷势能（回顾）]
\label{def:markdown-jacgap-phi}
取 Pure 卷最小三分量证书（chart / lift / threshold）：
\[
\JacGap(L)=\bigl(c_{\mathrm{chart}}(L),c_{\mathrm{lift}}(L),c_{\mathrm{thres}}(L)\bigr).
\]
取权重 $w_\bullet>0$，定义 defect potential：
\[
\Phi(L)=w_{\mathrm{chart}}\|c_{\mathrm{chart}}(L)\|
w_{\mathrm{lift}}\|c_{\mathrm{lift}}(L)\|
w_{\mathrm{thres}}\|c_{\mathrm{thres}}(L)\|.
\]
\end{definition}

\begin{definition}[\RTSC{}/极限低阻 strict 扇区（回顾）]
\label{def:markdown-strict-sector}
以“室温低阻逼近 + 证书化 + 缺陷归零”为 strict 目标（项目中可替换具体评价项），定义：
\[
\mathrm{Strict}_{\mathrm{RTSC}}
:=
\Bigl\{
L\in\mathrm{Loc}\ \Big|\ \Phi(L)=0,\ \mathrm{Ev}_{\rho,300}(L)\le \varepsilon_\rho,\ \mathrm{Ev}_{\mathrm{Cert}}(L)=1
\Bigr\}.
\]
该口径与 Corollary~\ref{cor:qualitative-breakthrough} 一致。
\end{definition}

\subsection{定义：治理算子字母表（严格区分是否保持 $p_{\mathrm{mat}}$）}

\begin{definition}[治理算子词：非交换动作字母表]
\label{def:markdown-governance-words}
将“工程设计/控制/验证”统一为非交换算子词生成。令
\[
\mathcal{M}_{\mathrm{self}}=\langle\mathcal{G}_{\mathrm{self}}\rangle,\qquad
u=u_1u_2\cdots u_n,\ \ u_i\in\mathcal{G}_{\mathrm{self}}.
\]
并按 $p_{\mathrm{mat}}$ 是否保持将生成元划分为
$\mathcal{G}_{\mathrm{self}}^{\mathrm{pres}}$ 与 $\mathcal{G}_{\mathrm{self}}^{\Delta}$
（参见 Definition~\ref{def:p-mat-preserve-change-generators}）。
\end{definition}

\begin{definition}[$(p_{\mathrm{mat}})$-保持的最小治理生成元集（Benchmark 动作字典）]
\label{def:markdown-minimal-governance-generators}
取最小治理字母表（仅列出 $(p_{\mathrm{mat}})$-保持生成元）：
\[
\mathcal{G}_{\mathrm{self}}^{(0)}
:=
\bigl\{
U_{\mathrm{env}},U_{\mathrm{geom}},U_{\mathrm{comp}},U_{\mathrm{prec}},U_{\mathrm{ctrl}},U_{\mathrm{cert}}
\bigr\}\subseteq \mathcal{G}_{\mathrm{self}}^{\mathrm{pres}}.
\]
其语义分别为：环境/边界控制、几何迭代设计、算力/FEM 自适应、控制精度治理、阈值调度治理、证书口径对齐；
更完整的最小生成元表见 Definition~\ref{def:governance-generators-minimal-table}。
\end{definition}

\begin{definition}[改变 $p_{\mathrm{mat}}$ 的材料生成元（示意）]
\label{def:markdown-material-generators}
改变材料身份/长期等价类的生成元可抽象为
\[
\mathcal{G}_{\mathrm{mat}}
:=
\{U_{\Delta\mathrm{comp}},U_{\Delta\mathrm{phase}},U_{\Delta\mathrm{defect}}\}
\subseteq \mathcal{G}_{\mathrm{self}}^{\Delta}.
\]
它们一般满足 $p_{\mathrm{mat}}(U\cdot L)\neq p_{\mathrm{mat}}(L)$。
\end{definition}

\subsection{引理：治理算子一般非交换（最小见证）}

\begin{lemma}[非交换性（最小见证）]
\label{lem:markdown-noncommutative}
在开放系统对抗/实时治理语义下，通常有：
\[
U_{\mathrm{env}}U_{\mathrm{ctrl}}\neq U_{\mathrm{ctrl}}U_{\mathrm{env}},
\qquad
U_{\mathrm{prec}}U_{\mathrm{ctrl}}\neq U_{\mathrm{ctrl}}U_{\mathrm{prec}},
\qquad
U_{\mathrm{comp}}U_{\mathrm{geom}}\neq U_{\mathrm{geom}}U_{\mathrm{comp}}.
\]
\end{lemma}

\begin{certificate}[证书：状态依赖输入导致顺序敏感]
\label{cert:markdown-noncommutative}
证书化理由是：$U_{\mathrm{ctrl}}$ 的阈值调度与 $U_{\mathrm{prec}}$ 的精度治理通常以“当前环境估计/当前精度参数”为输入；
因此先执行 $U_{\mathrm{env}}$ 或 $U_{\mathrm{prec}}$ 会改变 $U_{\mathrm{ctrl}}$ 的输入，进而改变输出，导致一般不交换。
同理，$U_{\mathrm{geom}}$ 与 $U_{\mathrm{comp}}$ 对 FEM 口径/几何参数的耦合更新为顺序敏感。
\end{certificate}

\begin{proof}[证明（谓词因果链）]
由 Certificate~\ref{cert:markdown-noncommutative}，取谓词链：
\[
\begin{aligned}
A:&\ U_{\mathrm{ctrl}}\ \text{的更新规则依赖当前环境/精度输入},\\
B:&\ A\Rightarrow U_{\mathrm{env}}\ \text{或 }U_{\mathrm{prec}}\ \text{的先后次序会改变 }U_{\mathrm{ctrl}}\ \text{的输入},\\
C:&\ B\Rightarrow U_{\mathrm{env}}U_{\mathrm{ctrl}}\neq U_{\mathrm{ctrl}}U_{\mathrm{env}},\ \ U_{\mathrm{prec}}
  U_{\mathrm{ctrl}}\neq U_{\mathrm{ctrl}}U_{\mathrm{prec}},\\
D:&\ \text{FEM 口径/几何耦合为顺序敏感}\Rightarrow U_{\mathrm{comp}}U_{\mathrm{geom}}\neq U_{\mathrm{geom}}U_{\mathrm{comp}}.
\end{aligned}
\]
由 $C\wedge D$ 得到结论。证毕。
\end{proof}

\subsection{定理：显式反例（同一材料投影下的 $\Phi$ 分离）}

\begin{theorem}[显式反例：同一 $p_{\mathrm{mat}}$ 下存在非交换治理词使 $\Phi$ 分离]
\label{thm:markdown-explicit-counterexample}
在同一材料投影 $p_{\mathrm{mat}}$ 下，存在两条仅由 $(p_{\mathrm{mat}})$-保持治理生成元构成的非交换治理词 $u,u'$ 与初始对象 $L_0$，
使得
\[
p_{\mathrm{mat}}(u\cdot L_0)=p_{\mathrm{mat}}(u'\cdot L_0)=p_{\mathrm{mat}}(L_0),
\qquad
\Phi(u\cdot L_0)=0,\quad \Phi(u'\cdot L_0)>0.
\]
\end{theorem}

\begin{certificate}[证书：由显式反例定理直接给出（含充分公式）]
\label{cert:markdown-explicit-counterexample}
证书可直接取 Theorem~\ref{thm:explicit-noncommutative-counterexample} 与 Certificate~\ref{cert:explicit-counterexample}：
在局部坐标 $\xi=(b,g,\tau,h,\theta)$、初值 $\xi_0=(2,1,3,0,0)$ 下，取
\[
u:=U_{\mathrm{comp}}\circ U_{\mathrm{geom}}\circ U_{\mathrm{ctrl}}\circ U_{\mathrm{prec}}\circ U_{\mathrm{env}},
\qquad
u':=U_{\mathrm{env}}\circ U_{\mathrm{ctrl}}\circ U_{\mathrm{prec}}\circ U_{\mathrm{comp}}\circ U_{\mathrm{geom}},
\]
可计算得到
\[
u\cdot \xi_0=(1,3,3,1,1)\Rightarrow \Phi(u\cdot L_0)=0,
\qquad
u'\cdot \xi_0=\left(1,\frac{3}{2},3,2,2\right)\Rightarrow \Phi(u'\cdot L_0)=\frac{7}{2}.
\]
并且这些操作均为 $(p_{\mathrm{mat}})$-保持动作。
\end{certificate}

\begin{proof}[证明（谓词因果链）]
由 Certificate~\ref{cert:markdown-explicit-counterexample}，取谓词链：
\[
\begin{aligned}
A:&\ u,u'\ \text{仅由 }(p_{\mathrm{mat}})\text{-保持治理生成元构成}\Rightarrow p_{\mathrm{mat}}(u\cdot L_0)=p_{\mathrm{mat}}
  (u'\cdot L_0)=p_{\mathrm{mat}}(L_0),\\
B:&\ \text{显式计算}\Rightarrow \Phi(u\cdot L_0)=0\ \wedge\ \Phi(u'\cdot L_0)=\frac{7}{2}>0.
\end{aligned}
\]
由 $A\wedge B$ 得到结论。证毕。
\end{proof}

\subsection{推论：strict 可达性不是材料问题，而是治理问题（严格版本）}

\begin{corollary}[非材料决定性：strict 判定不可能仅由 $p_{\mathrm{mat}}$ 决定]
\label{cor:markdown-nonmaterial-determinacy}
不存在集合 $S_{\mathrm{mat}}\subseteq M_{\mathrm{mat}}$ 使得
\[
L\in\mathrm{Strict}_{\mathrm{RTSC}}
\quad\Longleftrightarrow\quad
p_{\mathrm{mat}}(L)\in S_{\mathrm{mat}}.
\]
\end{corollary}

\begin{certificate}[证书：由显式反例与 strict 判据直接否证材料决定性]
\label{cert:markdown-nonmaterial-determinacy}
由 Theorem~\ref{thm:markdown-explicit-counterexample}，存在同一材料投影 $p_{\mathrm{mat}}$ 下的两对象使 $\Phi$ 一为零一为正；
而 Definition~\ref{def:markdown-strict-sector} 以 $\Phi=0$ 作为 strict 门槛，故 strict 判定在同一 $p_{\mathrm{mat}}$ 下不一致，
  从而材料决定性不可能成立。
\end{certificate}

\begin{proof}[证明（谓词因果链）]
由 Theorem~\ref{thm:markdown-explicit-counterexample}，存在同一材料投影下的两对象 $u\cdot L_0$ 与 $u'\cdot L_0$，满足 $\Phi$ 分离。
而 strict 扇区要求 $\Phi=0$（Definition~\ref{def:markdown-strict-sector}），故 strict 判定在同一 $p_{\mathrm{mat}}$ 下不一致。
因此不可能存在仅由 $p_{\mathrm{mat}}$ 决定 strict 判定的集合 $S_{\mathrm{mat}}$。证毕。
\end{proof}

\subsection{定理：治理等价（strict 可达 $\Longleftrightarrow$ 存在策略使 $\Phi\downarrow 0$ 且保持 $\mathcal{E}_{\mathrm{prot}}$）}

\begin{theorem}[治理等价（引用版）：strict 可达 $\Longleftrightarrow$ 抗扰动治理策略存在]
\label{thm:markdown-governance-equivalence}
在开放对抗非自治法流
\[
\frac{d}{dt}L(t)=X_{\mathrm{law}}\bigl(L(t);u(t),d(t)\bigr),\qquad L(t)\in\mathrm{Loc},
\]
且相对 $\mathcal{E}_{\mathrm{prot}}$ 约束固定的前提下，
strict 可达性等价于“存在从观测到非交换治理词的策略，使 $\Phi$ 在允许扰动族内可严格下降到 $0$ 并进入 strict 扇区”。
\end{theorem}

\begin{certificate}[证书：由治理等价定理直接调用]
\label{cert:markdown-governance-equivalence}
该定理可直接调用 Theorem~\ref{thm:governance-equivalence} 的严格表述与其证书链。
\end{certificate}

\begin{proof}[证明（谓词因果链）]
由 Certificate~\ref{cert:markdown-governance-equivalence}，取谓词链：
\[
\begin{aligned}
A:&\ \text{Theorem~\ref{thm:governance-equivalence}}\Rightarrow \text{strict 可达与治理策略存在互相蕴含},\\
B:&\ A\Rightarrow \text{本定理陈述成立}.
\end{aligned}
\]
证毕。
\end{proof}

\subsection{备注：里程碑意义（由障碍论降级到连通论 $\pi_0$）}

\begin{remark}[严格内核：strict 进入后执行降级为 $\pi_0$]
当存在治理策略使 $\Phi\to 0$ 并进入 strict 扇区时，
由 Theorem~\ref{thm:milestone-significance} 与 Theorem~\ref{thm:decategorify-to-pi0}，
关键难点从三类（相对）上同调障碍的消去，降级为 strict 可执行集合的连通类分类：
\[
\text{执行分类}\ \cong\ \pi_0\bigl(\Xi_{\mathrm{exec}}(\mathbf{y},r)\bigr).
\]
因此“里程碑意义”在 Pure 卷内部可严格化为一个\emph{问题类变换/降级}，而非外部叙事口号。
\end{remark}

\subsection{一句话总结（正文可用）}

\begin{corollary}[一句话总结（可作为正文结论句）]
\label{cor:markdown-one-line-summary}
在 \RTSC{}-工业几何体 Benchmark（Pure 卷）中，工业级 \RTSC{}/极限低阻 strict 可达性不可由材料投影 $p_{\mathrm{mat}}$ 单独决定；
它等价于开放、抗扰动、非交换动态对抗治理：通过环境/几何/FEM/精度/阈值调度/证书口径等治理词使 defect potential $\Phi\downarrow 0$
且保持 $\mathcal{E}_{\mathrm{prot}}$。
\end{corollary}

\begin{certificate}[证书：由非材料决定性与治理等价两条接口合成]
\label{cert:markdown-one-line-summary}
由 Corollary~\ref{cor:markdown-nonmaterial-determinacy} 得到“strict 不可由材料投影单独决定”；
由 Theorem~\ref{thm:markdown-governance-equivalence} 得到“strict 可达 $\Longleftrightarrow$ 存在抗扰动治理策略使 $\Phi\downarrow 0$ 且保持
  $\mathcal{E}_{\mathrm{prot}}$”。
将两者合取压缩为一句话即得本推论。
\end{certificate}

\begin{proof}[证明（谓词因果链）]
取谓词链：
\[
\begin{aligned}
A:&\ \text{Corollary~\ref{cor:markdown-nonmaterial-determinacy}}\Rightarrow \text{非材料决定性},\\
B:&\ \text{Theorem~\ref{thm:markdown-governance-equivalence}}\Rightarrow \text{strict 可达}\equiv \text{治理策略存在},\\
C:&\ A\wedge B\Rightarrow \text{一句话总结成立}.
\end{aligned}
\]
证毕。
\end{proof}

\clearpage
\appendix

%%%%%%%%%%%%%%%%%%%%%%%%%%%%%%%%%%%%%%%%%%%%%%%%%%%%%%%%%%%%
\section{符号体系统一注释}
\label{app:symbol-glossary}
%%%%%%%%%%%%%%%%%%%%%%%%%%%%%%%%%%%%%%%%%%%%%%%%%%%%%%%%%%%%

\subsection{总览（用于跨章节复用）}
\begin{remark}[符号说明的目的]
本附录用于把全文反复出现的“对象/算子/评价/证书”记号统一为可复用接口，避免在不同章节中出现同名异义或口径漂移。
\end{remark}

\subsection{对象、空间与主丛结构}
\begin{description}
  \item[$\LBOPB$] 生命总算子主纤维丛（Definition~\ref{def:lbopb-skeleton}）。
  \item[$\PTOPB,\ETOPB$] 物理/工程总算子主纤维丛（Definition~\ref{def:ptopb}）；本笔记中常按“可控构造算子 + 外生扰动算子”语义拆分讨论。
  \item[$M,P,G,\Conn,\Curv$] 底空间、主纤维丛、结构群、联络与曲率；曲率用于记录允许变化的不交换性与路径依赖（Theorem~\ref{thm:connection-curvature-semantics}）。
  \item[$\mathcal{D},\mathcal{M}^D$] 域集合与域算子幺半群；用以表示可组合的“操作/干预/治理”序列。
  \item[$\EDOM$] 刚性干预层/统一执行层的工程设计算子幺半群（Definition~\ref{def:edom}）。
  \item[$\GZLS$] Pure 卷法空间对象的总记号；其子记号 $\GZLOC,\GZTLC,\GZLCurv$ 分别承载局部化、阈值分层与法曲率相关数据。
\end{description}

\subsection{评价、缺陷势能与可达性接口}
\begin{description}
  \item[$\JacGap$] 总雅可比间隙（可控性判据接口；见 Section~\ref{sec:ptopb-rtsc-total-jacgap}）。
  \item[$\Phi$] defect potential（缺陷势能）；常作为 $\JacGap$ 的范数化/加权化表示，用于表达“严格化/回归/障碍消去”的可下降目标。
  \item[$\mathcal{E}_{\mathrm{prot}}$] 受保护评价泛函族（protected observables），用于固定工程语义锚点并排除“换口径伪下降”（Definition~\ref{def:protected-evaluation-functionals}）。
\end{description}

\clearpage
%%%%%%%%%%%%%%%%%%%%%%%%%%%%%%%%%%%%%%%%%%%%%%%%%%%%%%%%%%%%
\section{证书调用接口（标签命名与索引）}
\label{app:certificate-api}
%%%%%%%%%%%%%%%%%%%%%%%%%%%%%%%%%%%%%%%%%%%%%%%%%%%%%%%%%%%%

\subsection{标签命名约定（稳定接口名）}
本文默认使用如下前缀作为“稳定接口名”：
\begin{center}
\renewcommand{\arraystretch}{1.05}
\begin{tabular}{>{\raggedright\arraybackslash}p{0.22\textwidth} >{\raggedright\arraybackslash}p{0.70\textwidth}}
  \texttt{thm:}  & 定理 \\
  \texttt{lem:}  & 引理 \\
  \texttt{prop:} & 命题 \\
  \texttt{cor:}  & 推论 \\
  \texttt{def:}  & 定义 \\
  \texttt{cert:} & 证书（调用接口） \\
  \texttt{sec:}  & 章节 \\
\end{tabular}
\end{center}

\subsection{调用建议（断言 + 证书成对引用）}
\begin{remark}[最小调用对]
跨文档复用时，建议成对引用“断言 + 证书”，以确保调用具备可审计边界：
\[
\texttt{(Theorem/Lemma/Proposition/Corollary)~\textbackslash ref\{...\}}
\quad+\quad
\texttt{Certificate~\textbackslash ref\{...\}}.
\]
\end{remark}

\subsection{证书索引（证书标签作为调用接口）}
{\small
\setlength{\tabcolsep}{4pt}
\renewcommand{\arraystretch}{1.15}
\begin{longtable}{>{\raggedright\arraybackslash}p{0.18\textwidth} >{\raggedright\arraybackslash}p{0.32\textwidth}
  >{\raggedright\arraybackslash}p{0.44\textwidth}}
\hline
\textbf{证书编号} & \textbf{接口名} & \textbf{摘要} \\
\hline
\endfirsthead
\hline
\textbf{证书编号} & \textbf{接口名} & \textbf{摘要} \\
\hline
\endhead
% BEGIN_CERT_INDEX
% 证书索引（证书标签作为调用接口）
% 说明：原独立文件已并入正文；如需更新请根据正文证书环境重建。
证书~\ref{cert:lbopb-definition-source} & \code{cert:lbopb-definition-source} & 文本证书：\LBOPB{} 的结构定义位置 \\
证书~\ref{cert:lbopb-template-induction} & \code{cert:lbopb-template-induction} & 数学归纳法证书：按“结构元组替换步数”验证同型性 \\
证书~\ref{cert:connection-curvature-source} & \code{cert:connection-curvature-source} & 文本证书：联络/曲率与非交换性的语义来源 \\
证书~\ref{cert:rigid-layer-word-induction} & \code{cert:rigid-layer-word-induction} & 数学归纳法证书：按词长度 $T$ 归纳的“可执行计划=词”口径 \\
证书~\ref{cert:grl-sac-source} & \code{cert:grl-sac-source} & 文本证书：离散 GRL-SAC 的动作空间与目标形式 \\
证书~\ref{cert:isomorphic-migration} & \code{cert:isomorphic-migration} & 证书：同型迁移的结构对应表 \\
证书~\ref{cert:outer-loop-pruning} & \code{cert:outer-loop-pruning} & 证书：外环约束作为动作空间与词空间的可计算裁剪 \\
证书~\ref{cert:feasibility-criteria} & \code{cert:feasibility-criteria} & 证书：三项判据的可检验形式 \\
证书~\ref{cert:ptopb-sc-word-completeness-induction} & \code{cert:ptopb-sc-word-completeness-induction} & 数学归纳法证书：按词长 $T$
  的表示归纳 \\
证书~\ref{cert:natcomm-2025-max-tc} & \code{cert:natcomm-2025-max-tc} & 证书：常压常规（electron--phonon）机制下室温 $T_c$ 的证据性上限倾向 \\
证书~\ref{cert:nsr-superhydrides-review} & \code{cert:nsr-superhydrides-review} & 证书：高压氢化物（superhydrides）高 $T_c$
  的理论--实验一致性样本 \\
证书~\ref{cert:pubmed-retraction} & \code{cert:pubmed-retraction} & 证书：近常压/室温个别报道的可重复性风险样本 \\
证书~\ref{cert:rbph3-ambient-100k} & \code{cert:rbph3-ambient-100k} & 证书：常压附近 $\sim 100\,\mathrm{K}$ 量级理论候选与非简谐/量子核效应 \\
证书~\ref{cert:htsc-2025-dataset} & \code{cert:htsc-2025-dataset} & 证书：常压高温超导数据集与基准回归口径 \\
证书~\ref{cert:thresholded-cost} & \code{cert:thresholded-cost} & 证书：代价泛函的阈值化构造口径 \\
证书~\ref{cert:rtsc-tier-evidence} & \code{cert:rtsc-tier-evidence} & 证书：三类目标的外部证据对应 \\
证书~\ref{cert:uncertainty-necessity} & \code{cert:uncertainty-necessity} & 证书：前向模型成熟度与不确定度管理的必要性 \\
证书~\ref{cert:reproducibility-risk} & \code{cert:reproducibility-risk} & 证书：不可复现风险样本与门槛化必要性 \\
证书~\ref{cert:inverse-design-sufficient} & \code{cert:inverse-design-sufficient} & 证书：两条件合取的逻辑充分性 \\
证书~\ref{cert:plugin-structure} & \code{cert:plugin-structure} & 证书：路线插件化的结构约束口径 \\
证书~\ref{cert:engine-invariance} & \code{cert:engine-invariance} & 证书：引擎不变性与插件替换点 \\
证书~\ref{cert:k-as-theory-operator} & \code{cert:k-as-theory-operator} & 证书：机制标签与理论算子链的同构口径 \\
证书~\ref{cert:three-advantages-deps} & \code{cert:three-advantages-deps} & 证书：三项优势的依赖关系 \\
证书~\ref{cert:vol1-plugin} & \code{cert:vol1-plugin} & 证书：Vol I 路线作为插件的可核验输入项 \\
证书~\ref{cert:word-compilation-induction} & \code{cert:word-compilation-induction} & 数学归纳法证书：按词长 $T$ 的编译一致性 \\
证书~\ref{cert:cert-monotone-transfer} & \code{cert:cert-monotone-transfer} & 证书：证书谓词的单调迁移口径 \\
证书~\ref{cert:partition-sandwich} & \code{cert:partition-sandwich} & 证书：分配函数夹逼不等式 \\
证书~\ref{cert:boltzmann-ratio} & \code{cert:boltzmann-ratio} & 证书：Boltzmann 权重比的上界 \\
证书~\ref{cert:triple-reuse-deps} & \code{cert:triple-reuse-deps} & 证书：复用三元组的依赖链 \\
证书~\ref{cert:ctrl-first-class-completeness-induction} & \code{cert:ctrl-first-class-completeness-induction} & 数学归纳法证书：
  按词长 $T$ 的总词表示归纳 \\
证书~\ref{cert:route-hallucination} & \code{cert:route-hallucination} & 证书：幻觉成立的最小条件（失配 + 缺少校准项 + 坍缩） \\
证书~\ref{cert:entropy-calibration} & \code{cert:entropy-calibration} & 证书：熵底座与校准项对坍缩的抑制口径 \\
证书~\ref{cert:codesign-upper-envelope} & \code{cert:codesign-upper-envelope} & 证书：以 $w_c^\ast(w_m)$ 表示材料候选的可达上限 \\
证书~\ref{cert:hallucination-trigger} & \code{cert:hallucination-trigger} & 证书：触发条件的可观测统计量 \\
证书~\ref{cert:computable-diagnosis-induction} & \code{cert:computable-diagnosis-induction} & 数学归纳法证书：按词长 $T$ 的下界单调性 \\
证书~\ref{cert:diagnosis-priority} & \code{cert:diagnosis-priority} & 证书：全局最小值的嵌套下确界展开 \\
证书~\ref{cert:inf-lower-bound} & \code{cert:inf-lower-bound} & 证书：下确界是全体取值的下界 \\
证书~\ref{cert:inf-set-monotone} & \code{cert:inf-set-monotone} & 证书：下确界对候选集的集合单调性 \\
证书~\ref{cert:diagnosis-before-treatment} & \code{cert:diagnosis-before-treatment} & 证书：诊断优先的可计算收益来自下界排序与集合扩张单调性 \\
证书~\ref{cert:tex14-metasol-source} & \code{cert:tex14-metasol-source} & 文本证书：\textsf{MetaSol}/绝对锚点口径的定义位置 \\
证书~\ref{cert:m6-homotopy-reuse-induction} & \code{cert:m6-homotopy-reuse-induction} & 数学归纳法证书：按复用步数 $N$ 的链式闭包 \\
证书~\ref{cert:near-optimal-pushforward} & \code{cert:near-optimal-pushforward} & 证书：由自由能稳定性界推得最优值与近最优差界 \\
证书~\ref{cert:ctrl-first-class-necessary} & \code{cert:ctrl-first-class-necessary} & 证书：缩小动作集只会抬高诊断下界 \\
证书~\ref{cert:anti-drift-constraint} & \code{cert:anti-drift-constraint} & 证书：反偏航代价由残差一等化与证书门槛化拼接得到 \\
证书~\ref{cert:tail-risk-rd-chain} & \code{cert:tail-risk-rd-chain} & 证书：两阶段联动链由门槛化、诊断优先与反偏航机制合取得到 \\
证书~\ref{cert:block-lift-induction} & \code{cert:block-lift-induction} & 数学归纳法证书：按“已提升到基底的块数 $k$”归纳 \\
证书~\ref{cert:section-function-bijection} & \code{cert:section-function-bijection} & 证书：在平凡束 $P_A=M_A\times
  G_{\mathrm{rep}}$ 上截面与函数一一对应 \\
证书~\ref{cert:bundle-B-bilevel} & \code{cert:bundle-B-bilevel} & 证书：纤维下确界与最优截面的等价展开 \\
证书~\ref{cert:ctrl-as-base-diagnosis} & \code{cert:ctrl-as-base-diagnosis} & 证书：把控制放入基底等价于强制在控制维度做诊断排序 \\
证书~\ref{cert:diagnosis-on-invariants} & \code{cert:diagnosis-on-invariants} & 证书：下确界是函数取值的下界 \\
证书~\ref{cert:inverse-design-as-section} & \code{cert:inverse-design-as-section} & 证书：非空纤维 $\Rightarrow$ 可选代表元形成截面 \\
证书~\ref{cert:certificate-bundle-anti-hallucination} & \code{cert:certificate-bundle-anti-hallucination} & 证书：
  证书门槛化会显式排除不可复现区域 \\
证书~\ref{cert:rtsc-bundle-selection} & \code{cert:rtsc-bundle-selection} & 证书：选型依据来自“基底坐标决定诊断粒度”与“纤维内化实现细节” \\
证书~\ref{cert:subbundle-restriction} & \code{cert:subbundle-restriction} & 证书：主丛在子基底上的限制保持主丛结构 \\
证书~\ref{cert:envelope-fiber-formula} & \code{cert:envelope-fiber-formula} & 证书：复合映射的原像展开 \\
证书~\ref{cert:envelope-closure-induction} & \code{cert:envelope-closure-induction} & 数学归纳法证书：按压缩步数 $N$ 归纳 \\
证书~\ref{cert:commuting-square} & \code{cert:commuting-square} & 证书：商映射下的良定义判据 \\
证书~\ref{cert:principalization} & \code{cert:principalization} & 证书：自由传递作用 + 兼容提升 $\Rightarrow$ 主丛化 \\
证书~\ref{cert:af-unified} & \code{cert:af-unified} & 证书：包络投影与交换图给出“同一总空间的不同投影” \\
证书~\ref{cert:first-order-fiber-union} & \code{cert:first-order-fiber-union} & 证书：包络纤维公式由复合原像展开得到 \\
证书~\ref{cert:second-order-envelope} & \code{cert:second-order-envelope} & 证书：限制投影与坐标化嵌入的等价口径 \\
证书~\ref{cert:bundle-tower-closure-induction} & \code{cert:bundle-tower-closure-induction} & 数学归纳法证书：按包络细化层数 $\ell$ 归纳 \\
证书~\ref{cert:envelope-anti-drift} & \code{cert:envelope-anti-drift} & 证书：反偏航代价与偏航预警口径已在前文给出 \\
证书~\ref{cert:rtsc-jacgap-unactuated-induction} & \code{cert:rtsc-jacgap-unactuated-induction} & 数学归纳法证书：按子节点数 $N$
  归纳（加权和下界） \\
证书~\ref{cert:rtsc-jacgap-rank} & \code{cert:rtsc-jacgap-rank} & 证书：满秩 $\Rightarrow$ 局部可达（隐函数/正则值口径） \\
证书~\ref{cert:rtsc-lattice-necessary} & \code{cert:rtsc-lattice-necessary} & 证书：不可作用自由度导致子节点间隙有界下不去 \\
证书~\ref{cert:rtsc-envelope-lattice-qualified} & \code{cert:rtsc-envelope-lattice-qualified} & 证书：包络自由度显式化 + 格点可作用
  $\Rightarrow$ 避免固定晶格子空间偏航 \\
证书~\ref{cert:rtsc-jacgap-controllability-main} & \code{cert:rtsc-jacgap-controllability-main} & 证书：主命题由“定义等价 + 不可作用下界 +
  局部可达条件”拼接得到 \\
证书~\ref{cert:rtsc-disturbance-spike} & \code{cert:rtsc-disturbance-spike} & 证书：链式法则给出 $\dot V$ 的扰动项 \\
证书~\ref{cert:rtsc-disturbance-regression-induction} & \code{cert:rtsc-disturbance-regression-induction} & 数学归纳法证书：按时间步
  $N$ 归纳（几何级数界） \\
证书~\ref{cert:rtsc-disturbance-vanish} & \code{cert:rtsc-disturbance-vanish} & 证书：$V$ 的分解给出 $V\to 0\Rightarrow
  G_{\mathrm{tot}}\to 0$ \\
证书~\ref{cert:jacgap-strictification-induction} & \code{cert:jacgap-strictification-induction} & 数学归纳法证书：按同伦修复步数 $n$ 归纳
  \\
证书~\ref{cert:jacgap-controlled-by-curvature} & \code{cert:jacgap-controlled-by-curvature} & 证书：JacGap 沿水平法流的变化受曲率与同伦项控制
  \\
证书~\ref{cert:lawspace-path-integral} & \code{cert:lawspace-path-integral} & 证书：Boltzmann 权重在大 $\beta$ 极限由最小作用量主导 \\
证书~\ref{cert:puremath-direction-before-construction} & \code{cert:puremath-direction-before-construction} & 证书：由 GZ-LSPI
  的最小作用量主导分解为“方向 + 修复”两层 \\
证书~\ref{cert:eng-sector-closure-induction} & \code{cert:eng-sector-closure-induction} & 数学归纳法证书：按扇区块数 $n\in\{1,2,3,4\}$
  归纳 \\
证书~\ref{cert:eng-envelope-minimal-three} & \code{cert:eng-envelope-minimal-three} & 证书：三类包络对应三类“可控失败”来源 \\
证书~\ref{cert:eng-envelope-anti-drift} & \code{cert:eng-envelope-anti-drift} & 证书：由包络塔与偏航预警口径直接推出 \\
证书~\ref{cert:eng-envelope-refine-induction} & \code{cert:eng-envelope-refine-induction} & 数学归纳法证书：按包络细化层级 $r$ 归纳（细化不劣化）
  \\
证书~\ref{cert:eng-unified-subsectors} & \code{cert:eng-unified-subsectors} & 证书：子扇区限制对应子丛限制，判据在限制下保持 \\
证书~\ref{cert:puremath-phi-geom-ohu} & \code{cert:puremath-phi-geom-ohu} & 文本证书：$\Phi_{\mathrm{geom}}$、$\mathsf{GeomObj}$
  与 \textsf{GZ-OHU} 的纯数卷来源 \\
证书~\ref{cert:ig-minimal-envelope} & \code{cert:ig-minimal-envelope} & 证书：三类分块分别对应“缺陷下降”的三类必要自由度来源 \\
证书~\ref{cert:ig-invgen-induction} & \code{cert:ig-invgen-induction} & 数学归纳法证书：按同伦动作步数 $n$ 归纳（单调下降证书） \\
证书~\ref{cert:ig-disturbance-regression} & \code{cert:ig-disturbance-regression} & 证书：由曲率尖峰定理与方向优先推论合取得到 \\
证书~\ref{cert:jacgap-needs-protected-observables} & \code{cert:jacgap-needs-protected-observables} & 证书：$\JacGap=0$
  的语义依赖于受保护低阶行为 \\
证书~\ref{cert:relative-homotopy-depth-invariant} & \code{cert:relative-homotopy-depth-invariant} & 证书：相对可回缩给出互为逆的相对同伦 \\
证书~\ref{cert:relative-minimal-envelope-induction} & \code{cert:relative-minimal-envelope-induction} & 数学归纳法证书：按“可回缩扩展步数
  $N$”归纳 \\
证书~\ref{cert:protected-evaluations-necessary} & \code{cert:protected-evaluations-necessary} & 证书：严格化判据必须相对受保护量 \\
证书~\ref{cert:relative-geodesic-twist} & \code{cert:relative-geodesic-twist} & 证书：缺少相对约束将触发“假回归/假严格化” \\
证书~\ref{cert:design-strategy-correspondence} & \code{cert:design-strategy-correspondence} & 证书：由相对同伦锚点与包络/策略实现关系直接推出 \\
证书~\ref{cert:relative-strategy-feasibility-induction} & \code{cert:relative-strategy-feasibility-induction} & 数学归纳法证书：
  按离散步数 $N$ 归纳 \\
证书~\ref{cert:design-generation-equals-strategy-generation} & \code{cert:design-generation-equals-strategy-generation} &
  证书：开放性 + 非交换性 $\Rightarrow$ 静态对象不足以定义方案 \\
证书~\ref{cert:noncommutative-design-strategy-equivalence-induction} & \code{cert:
  noncommutative-design-strategy-equivalence-induction} & 数学归纳法证书：按离散步数 $N$ 的策略可行性归纳 \\
证书~\ref{cert:openra-vs-alphago} & \code{cert:openra-vs-alphago} & 证书：三类差异——开放性、非交换性、信息结构 \\
证书~\ref{cert:opb-benchmark-template-induction} & \code{cert:opb-benchmark-template-induction} & 数学归纳法证书：复用离散步数归纳框架 \\
证书~\ref{cert:chip-engine-share-opb-template} & \code{cert:chip-engine-share-opb-template} & 证书：实例到模板的代入映射 \\
证书~\ref{cert:rtsc-ig-design-strategy-equivalence} & \code{cert:rtsc-ig-design-strategy-equivalence} & 证书：相对约束 + 可下降缺陷势能
  $\Rightarrow$ 等价链可调用 \\
证书~\ref{cert:rtsc-approx-design-strategy-equivalence} & \code{cert:rtsc-approx-design-strategy-equivalence} & 证书：相对锚点 +
  包络显式化 + 缺陷势能可下降 \\
证书~\ref{cert:edge-operator-noncommutative-induction} & \code{cert:edge-operator-noncommutative-induction} & 数学归纳法证书：
  按算子词长度 $N$ 归纳 \\
证书~\ref{cert:edge-regularised-descent} & \code{cert:edge-regularised-descent} & 证书：局部平滑能量将突变转化为可控代价 \\
证书~\ref{cert:discussion-gains-map} & \code{cert:discussion-gains-map} & 证书：七组增益与文内标签映射 \\
证书~\ref{cert:discussion-gains-template} & \code{cert:discussion-gains-template} & 证书：模板四要素对应的形式化条目 \\
证书~\ref{cert:discussion-convergence-induction} & \code{cert:discussion-convergence-induction} & 数学归纳法证书：按“纳入的增益组数” $n$
  归纳 \\
证书~\ref{cert:discussion-one-sentence} & \code{cert:discussion-one-sentence} & 证书：由收敛定理直接推出 \\
证书~\ref{cert:chart-omega-standard} & \code{cert:chart-omega-standard} & 证书：中心扩张与 Čech 上同调的标准构造事实 \\
证书~\ref{cert:jacgap-as-obstruction} & \code{cert:jacgap-as-obstruction} & 证书：$\JacGap$ 分量为障碍类的证书化代表 \\
证书~\ref{cert:phi0-implies-components0} & \code{cert:phi0-implies-components0} & 证书：$\Phi=0$ $\Rightarrow$ 三分量证书全为零 \\
证书~\ref{cert:obstruction-zero-compilable} & \code{cert:obstruction-zero-compilable} & 证书：Čech 余边界见证可编译为包络内有限修复动作 \\
证书~\ref{cert:xi-exec-path} & \code{cert:xi-exec-path} & 证书：固定基底点后纤维与 $\Xi$ 同胚，路径连通 $\Rightarrow$ 存在纤维路径 \\
证书~\ref{cert:qualitative-breakthrough} & \code{cert:qualitative-breakthrough} & 证书：由定理 A/B 与质变同伦类定义合成 \\
证书~\ref{cert:nonmaterial-action-exists} & \code{cert:nonmaterial-action-exists} & 证书：存在材料保持但可改变 strict 判定的纤维动作 \\
证书~\ref{cert:policy-compilation} & \code{cert:policy-compilation} & 证书：动作链 $\leftrightarrow$ 策略的编译与展开 \\
证书~\ref{cert:governance-irreducible} & \code{cert:governance-irreducible} & 证书：非材料决定性迫使治理自由度进入接口 \\
证书~\ref{cert:milestone-significance} & \code{cert:milestone-significance} & 证书：由障碍消去与执行降级定理链直接推出 \\
证书~\ref{cert:strong-conclusion} & \code{cert:strong-conclusion} & 证书：三句分别由三条主定理链压缩得到 \\
证书~\ref{cert:component-invariance-induction} & \code{cert:component-invariance-induction} & 数学归纳法证书：按词长 $n$ 归纳 \\
证书~\ref{cert:toy-noncommutative-witness} & \code{cert:toy-noncommutative-witness} & 证书：给出一组初值并直接计算得到不交换 \\
证书~\ref{cert:explicit-counterexample} & \code{cert:explicit-counterexample} & 证书：逐步计算得到 $\Phi$ 分离 \\
证书~\ref{cert:explicit-nonmaterial-determinacy} & \code{cert:explicit-nonmaterial-determinacy} & 证书：由显式非交换反例给出同一材料投影下的
  $\Phi$ 分离 \\
证书~\ref{cert:markdown-noncommutative} & \code{cert:markdown-noncommutative} & 证书：状态依赖输入导致顺序敏感 \\
证书~\ref{cert:markdown-explicit-counterexample} & \code{cert:markdown-explicit-counterexample} & 证书：由显式反例定理直接给出（含充分公式） \\
证书~\ref{cert:markdown-nonmaterial-determinacy} & \code{cert:markdown-nonmaterial-determinacy} & 证书：由显式反例与 strict
  判据直接否证材料决定性 \\
证书~\ref{cert:markdown-governance-equivalence} & \code{cert:markdown-governance-equivalence} & 证书：由治理等价定理直接调用 \\
证书~\ref{cert:markdown-one-line-summary} & \code{cert:markdown-one-line-summary} & 证书：由非材料决定性与治理等价两条接口合成 \\
% END_CERT_INDEX
\end{longtable}
}

\clearpage
\begin{thebibliography}{20}

\bibitem{GaoZhengAppVol2}
Gao, Z. (2025).
\newblock GaoZheng G-Framework and GaoZheng G-Algebra: Applied Mathematics Volume II — LBOPB, Life-Science Monoids, and
  Generative Precision Medicine.
\newblock In \emph{GaoZheng G-Framework and GaoZheng G-Algebra: Applied Mathematics Volume II — LBOPB, Life-Science
  Monoids, and Generative Precision Medicine (v1.1)}.
\newblock Zenodo.
\newblock \url{https://doi.org/10.5281/zenodo.17744672}.


\bibitem{NatComm2025MaxTcConvAmbient}
【期刊论文】
Nature Communications.
\newblock \emph{The maximum $T_c$ of conventional superconductors at ambient pressure}.
\newblock 2025.
\newblock \url{https://www.nature.com/articles/s41467-025-63702-w}.

\bibitem{NSR2024Superhydrides}
【期刊论文】
National Science Review (Oxford Academic).
\newblock \emph{Pressure-induced hydrogen-dominant high-temperature superconductors}.
\newblock 2024.
\newblock \url{https://academic.oup.com/nsr/article/11/7/nwae004/7511108}.

\bibitem{PubMedRetractionLutetiumHydride}
【撤稿记录】
PubMed.
\newblock \emph{Retraction Note: Evidence of near-ambient superconductivity in a N-doped lutetium hydride}.
\newblock 2023.
\newblock \url{https://pubmed.ncbi.nlm.nih.gov/37935926/}.

\bibitem{ScienceDirectRbPH3}
【期刊论文】
ScienceDirect.
\newblock \emph{Ambient pressure high temperature superconductivity in RbPH$_3$ facilitated by ionic anharmonicity}.
\newblock 2025.
\newblock \url{https://www.sciencedirect.com/science/article/pii/S2950463525000195}.

\bibitem{arXivHTSC2025}
【预印本/数据集】
arXiv.
\newblock \emph{HTSC-2025: A Benchmark Dataset of Ambient-Pressure High-Temperature Superconductors for AI-Driven
  Critical Temperature Prediction}.
\newblock 2025.
\newblock \url{https://arxiv.org/html/2506.03837v1}.

\end{thebibliography}

\end{CJK*}
\end{document}
